% !TeX program = xelatex
% Standalone English volume. Use XeLaTeX until the references stabilize.
\documentclass[a4paper,11pt]{article}
\usepackage[T1]{fontenc}
\usepackage{lmodern}
\usepackage[english]{babel}
\usepackage[left=26mm,right=26mm,top=25mm,bottom=25mm,headheight=15pt,headsep=7mm,footskip=11mm]{geometry}
\usepackage{amsmath,amssymb,mathtools,amsthm}
\newtheorem{proposition}{Proposition}[section]
\newtheorem{lemma}[proposition]{Lemma}
\newtheorem{theorem}[proposition]{Theorem}
\newtheorem{corollary}[proposition]{Corollary}
\newtheorem{maintheorem}[proposition]{Main Theorem}
\theoremstyle{definition}
\newtheorem{definition}[proposition]{Definition}
\theoremstyle{remark}
\newtheorem{remark}[proposition]{Remark}
\makeatletter
\renewcommand*\l@section{\@dottedtocline{1}{0em}{2.4em}}
\renewcommand*\l@subsection{\@dottedtocline{2}{1.45em}{3.45em}}
\renewcommand*\l@subsubsection{\@dottedtocline{3}{3.8em}{4.55em}}
\makeatother
\usepackage{xcolor,fancyhdr,titlesec,needspace,xurl}
\usepackage{longtable,multicol}
\usepackage{hyperref}
\definecolor{ink}{HTML}{1A3542}
\definecolor{muted}{HTML}{4F606A}
\hypersetup{unicode=true,colorlinks=true,linkcolor=ink,filecolor=ink,urlcolor=ink,citecolor=ink,bookmarksnumbered=true,
 pdfdisplaydoctitle=true,pdflang={en},
 pdftitle={The generalized Riemann hypothesis for Dirichlet L-functions. Edition 930, Volume XI},
 pdfsubject={A natural-language reconstruction of the proposed proof},
 pdfauthor={A mathematical reconstruction}}
\titleformat{\section}{\Large\bfseries\color{ink}}{\thesection}{.65em}{}
\titlespacing*{\section}{0pt}{.8em}{.65em}
\setlength{\parindent}{1.2em}
\linespread{1.045}
\setlength{\parskip}{.22em plus .05em minus .03em}
\setlength{\emergencystretch}{2em}
\setcounter{topnumber}{2}
\setcounter{tocdepth}{3}\setcounter{secnumdepth}{3}
\widowpenalty=4500
\clubpenalty=4500
\pagestyle{fancy}\fancyhf{}
\fancyhead[L]{\small\color{muted}Dirichlet GRH \textperiodcentered\ Volume I}
\fancyhead[R]{\small\color{muted}Dirichlet zeros and power intervals}
\fancyfoot[C]{\small\thepage}
\renewcommand{\headrulewidth}{.3pt}
\newcommand{\lead}[1]{\par\smallskip\noindent{\bfseries\color{ink}#1.}\enspace}
\newcommand{\sour}[2]{\unskip\penalty150\hspace{.35em}{\small\hyperlink{ref#1}{[#2]}}}
\titleformat{\subsection}{\large\bfseries}{\thesubsection}{.6em}{}
\titlespacing*{\subsection}{0pt}{1.15em plus .15em}{.4em}
\setlength{\abovedisplayskip}{9pt plus 2pt minus 3pt}
\setlength{\belowdisplayskip}{9pt plus 2pt minus 3pt}
\newcommand{\RR}{\mathbb R}\newcommand{\CC}{\mathbb C}\newcommand{\QQ}{\mathbb Q}
\newcommand{\ZZ}{\mathbb Z}\newcommand{\ord}{\operatorname{ord}}
\newcommand{\Frob}{\operatorname{Frob}}\newcommand{\Gal}{\operatorname{Gal}}
\newcommand{\len}{\operatorname{length}}\newcommand{\Res}{\operatorname*{Res}}
\newcommand{\wmean}{\operatorname{av}}\newcommand{\Ker}{\operatorname{ker}}
\newcommand{\Lh}{\mathcal L}\newcommand{\Gg}{\mathcal G}
\newcommand{\real}{\operatorname{Re}}\newcommand{\imag}{\operatorname{Im}}
\newcommand{\abs}[1]{\left|#1\right|}\newcommand{\eps}{\varepsilon}
\newcommand{\anchoritem}[2]{\par\Needspace{2.8\baselineskip}\hypertarget{ref#1}{}\noindent{\bfseries\color{ink}#2.}\enspace}
\newcommand{\srcloc}[2]{{\small\nolinkurl{#1.lean}}, line~#2}
\numberwithin{equation}{section}
\raggedbottom

\newcommand{\Sel}{\mathcal S}
\newcommand{\Cl}{\operatorname{Cl}}
\newcommand{\Aut}{\operatorname{Aut}}
\newcommand{\GL}{\operatorname{GL}}
\newcommand{\OO}{\mathcal O}
\hypersetup{bookmarksdepth=3}


\addto\captionsenglish{\renewcommand{\partname}{Volume}}


\makeatletter
\expandafter\def\csname r@cor:detectingsplit650\endcsname{{22.54}{105}{}{proposition.22.54}{Dirichlet-GRH-v930-Volume-IV.pdf}}
\expandafter\def\csname r@cor:dual68\endcsname{{32.9}{33}{}{proposition.32.9}{Dirichlet-GRH-v930-Volume-VI.pdf}}
\expandafter\def\csname r@cor:firsttorus469\endcsname{{47.6}{81}{}{proposition.47.6}{Dirichlet-GRH-v930-Volume-IX.pdf}}
\expandafter\def\csname r@cor:originaletale424\endcsname{{44.43}{35}{}{proposition.44.43}{Dirichlet-GRH-v930-Volume-IX.pdf}}
\expandafter\def\csname r@cor:present116\endcsname{{32.20}{45}{}{proposition.32.20}{Dirichlet-GRH-v930-Volume-VI.pdf}}
\expandafter\def\csname r@cor:projectsymbol483\endcsname{{68.17}{97}{}{proposition.68.17}{Dirichlet-GRH-v930-Volume-XIII.pdf}}
\expandafter\def\csname r@cor:sameprimes517\endcsname{{9.34}{134}{}{proposition.9.34}{Dirichlet-GRH-v930-Volume-I.pdf}}
\expandafter\def\csname r@cor:shortcenter460\endcsname{{9.27}{130}{}{proposition.9.27}{Dirichlet-GRH-v930-Volume-I.pdf}}
\expandafter\def\csname r@cor:supportmass465\endcsname{{9.30}{132}{}{proposition.9.30}{Dirichlet-GRH-v930-Volume-I.pdf}}
\expandafter\def\csname r@cor:whittakerterminal527\endcsname{{96.8}{78}{}{proposition.96.8}{Dirichlet-GRH-v930-Volume-XVIII.pdf}}
\expandafter\def\csname r@eq:length12\endcsname{{{L}}{30}{}{AMS.3}{Dirichlet-GRH-v930-Volume-I.pdf}}
\expandafter\def\csname r@lem:algebra116\endcsname{{32.17}{43}{}{proposition.32.17}{Dirichlet-GRH-v930-Volume-VI.pdf}}
\expandafter\def\csname r@lem:commondyadictwist429\endcsname{{46.1}{63}{}{proposition.46.1}{Dirichlet-GRH-v930-Volume-IX.pdf}}
\expandafter\def\csname r@lem:connectedthree418\endcsname{{44.13}{13}{}{proposition.44.13}{Dirichlet-GRH-v930-Volume-IX.pdf}}
\expandafter\def\csname r@lem:cusp361\endcsname{{36.16}{16}{}{proposition.36.16}{Dirichlet-GRH-v930-Volume-VII.pdf}}
\expandafter\def\csname r@lem:cyclointegral377\endcsname{{40.1}{2}{}{proposition.40.1}{Dirichlet-GRH-v930-Volume-VIII.pdf}}
\expandafter\def\csname r@lem:des28\endcsname{{24.8}{8}{}{proposition.24.8}{Dirichlet-GRH-v930-Volume-V.pdf}}
\expandafter\def\csname r@lem:detangent372\endcsname{{37.35}{68}{}{proposition.37.35}{Dirichlet-GRH-v930-Volume-VII.pdf}}
\expandafter\def\csname r@lem:flag69\endcsname{{32.10}{34}{}{proposition.32.10}{Dirichlet-GRH-v930-Volume-VI.pdf}}
\expandafter\def\csname r@lem:fullinteger517\endcsname{{9.31}{132}{}{proposition.9.31}{Dirichlet-GRH-v930-Volume-I.pdf}}
\expandafter\def\csname r@lem:inertaux364\endcsname{{37.3}{35}{}{proposition.37.3}{Dirichlet-GRH-v930-Volume-VII.pdf}}
\expandafter\def\csname r@lem:kummer72\endcsname{{32.14}{39}{}{proposition.32.14}{Dirichlet-GRH-v930-Volume-VI.pdf}}
\expandafter\def\csname r@lem:locallevel345\endcsname{{22.16}{75}{}{proposition.22.16}{Dirichlet-GRH-v930-Volume-IV.pdf}}
\expandafter\def\csname r@lem:localweight362\endcsname{{36.19}{21}{}{proposition.36.19}{Dirichlet-GRH-v930-Volume-VII.pdf}}
\expandafter\def\csname r@lem:oldidempotent405\endcsname{{41.21}{49}{}{proposition.41.21}{Dirichlet-GRH-v930-Volume-VIII.pdf}}
\expandafter\def\csname r@lem:oppositeroot469\endcsname{{47.4}{79}{}{proposition.47.4}{Dirichlet-GRH-v930-Volume-IX.pdf}}
\expandafter\def\csname r@lem:planecurve417\endcsname{{44.9}{9}{}{proposition.44.9}{Dirichlet-GRH-v930-Volume-IX.pdf}}
\expandafter\def\csname r@lem:planeunits379\endcsname{{40.9}{7}{}{proposition.40.9}{Dirichlet-GRH-v930-Volume-VIII.pdf}}
\expandafter\def\csname r@lem:rel70\endcsname{{32.12}{36}{}{proposition.32.12}{Dirichlet-GRH-v930-Volume-VI.pdf}}
\expandafter\def\csname r@lem:shiftdivisors465\endcsname{{9.28}{131}{}{proposition.9.28}{Dirichlet-GRH-v930-Volume-I.pdf}}
\expandafter\def\csname r@lem:slope138\endcsname{{24.27}{22}{}{proposition.24.27}{Dirichlet-GRH-v930-Volume-V.pdf}}
\expandafter\def\csname r@lem:torsioncrys375\endcsname{{38.9}{76}{}{proposition.38.9}{Dirichlet-GRH-v930-Volume-VII.pdf}}
\expandafter\def\csname r@nav:III-74\endcsname{{22.5}{63}{}{subsection.22.5}{Dirichlet-GRH-v930-Volume-IV.pdf}}
\expandafter\def\csname r@prop:alloriginalpoints424\endcsname{{44.42}{34}{}{proposition.44.42}{Dirichlet-GRH-v930-Volume-IX.pdf}}
\expandafter\def\csname r@prop:auxfield361\endcsname{{36.15}{15}{}{proposition.36.15}{Dirichlet-GRH-v930-Volume-VII.pdf}}
\expandafter\def\csname r@prop:classrank365\endcsname{{37.6}{39}{}{proposition.37.6}{Dirichlet-GRH-v930-Volume-VII.pdf}}
\expandafter\def\csname r@prop:dualselmer362\endcsname{{36.20}{22}{}{proposition.36.20}{Dirichlet-GRH-v930-Volume-VII.pdf}}
\expandafter\def\csname r@prop:fixed116\endcsname{{32.18}{43}{}{proposition.32.18}{Dirichlet-GRH-v930-Volume-VI.pdf}}
\expandafter\def\csname r@prop:jointpatch417\endcsname{{44.10}{10}{}{proposition.44.10}{Dirichlet-GRH-v930-Volume-IX.pdf}}
\expandafter\def\csname r@prop:localmult399\endcsname{{41.4}{32}{}{proposition.41.4}{Dirichlet-GRH-v930-Volume-VIII.pdf}}
\expandafter\def\csname r@prop:odddescent376\endcsname{{38.14}{79}{}{proposition.38.14}{Dirichlet-GRH-v930-Volume-VII.pdf}}
\expandafter\def\csname r@prop:ordinaryclass405\endcsname{{41.22}{50}{}{proposition.41.22}{Dirichlet-GRH-v930-Volume-VIII.pdf}}
\expandafter\def\csname r@prop:ordinaryfamily363\endcsname{{36.23}{26}{}{proposition.36.23}{Dirichlet-GRH-v930-Volume-VII.pdf}}
\expandafter\def\csname r@prop:orighecke374\endcsname{{38.7}{74}{}{proposition.38.7}{Dirichlet-GRH-v930-Volume-VII.pdf}}
\expandafter\def\csname r@prop:origideal373\endcsname{{38.2}{71}{}{proposition.38.2}{Dirichlet-GRH-v930-Volume-VII.pdf}}
\expandafter\def\csname r@prop:planelabels379\endcsname{{40.11}{8}{}{proposition.40.11}{Dirichlet-GRH-v930-Volume-VIII.pdf}}
\expandafter\def\csname r@prop:qbound102\endcsname{{20.6}{27}{}{proposition.20.6}{Dirichlet-GRH-v930-Volume-IV.pdf}}
\expandafter\def\csname r@prop:quadratic345\endcsname{{36.3}{3}{}{proposition.36.3}{Dirichlet-GRH-v930-Volume-VII.pdf}}
\expandafter\def\csname r@prop:quarticconductor539\endcsname{{22.27}{84}{}{proposition.22.27}{Dirichlet-GRH-v930-Volume-IV.pdf}}
\expandafter\def\csname r@prop:res23\endcsname{{6.11}{32}{}{proposition.6.11}{Dirichlet-GRH-v930-Volume-I.pdf}}
\expandafter\def\csname r@prop:strictclasses365\endcsname{{37.5}{38}{}{proposition.37.5}{Dirichlet-GRH-v930-Volume-VII.pdf}}
\expandafter\def\csname r@prop:tracecrystal375\endcsname{{38.10}{76}{}{proposition.38.10}{Dirichlet-GRH-v930-Volume-VII.pdf}}
\expandafter\def\csname r@sec:adjcomponent372\endcsname{{37.9}{67}{}{subsection.37.9}{Dirichlet-GRH-v930-Volume-VII.pdf}}
\expandafter\def\csname r@sec:irrcurvesupport409\endcsname{{42.2}{66}{}{subsection.42.2}{Dirichlet-GRH-v930-Volume-VIII.pdf}}
\expandafter\def\csname r@sec:levelstrict345\endcsname{{22.13}{74}{}{subsection.22.13}{Dirichlet-GRH-v930-Volume-IV.pdf}}
\expandafter\def\csname r@sec:markedintersections420\endcsname{{44.6}{20}{}{subsection.44.6}{Dirichlet-GRH-v930-Volume-IX.pdf}}
\expandafter\def\csname r@sec:missing\endcsname{{6.6}{32}{}{subsection.6.6}{Dirichlet-GRH-v930-Volume-I.pdf}}
\expandafter\def\csname r@sec:oddquartic539\endcsname{{22.17}{84}{}{subsection.22.17}{Dirichlet-GRH-v930-Volume-IV.pdf}}
\expandafter\def\csname r@sec:origin116\endcsname{{32.8}{42}{}{subsection.32.8}{Dirichlet-GRH-v930-Volume-VI.pdf}}
\expandafter\def\csname r@sec:original68\endcsname{{32.4}{31}{}{subsection.32.4}{Dirichlet-GRH-v930-Volume-VI.pdf}}
\expandafter\def\csname r@sec:originalfixedhecke424\endcsname{{44.10}{34}{}{subsection.44.10}{Dirichlet-GRH-v930-Volume-IX.pdf}}
\expandafter\def\csname r@sec:shiftedsupport465\endcsname{{9.16}{130}{}{subsection.9.16}{Dirichlet-GRH-v930-Volume-I.pdf}}
\expandafter\def\csname r@sec:sympglobal469\endcsname{{47.2}{79}{}{subsection.47.2}{Dirichlet-GRH-v930-Volume-IX.pdf}}
\expandafter\def\csname r@sec:symplecticstrict429\endcsname{{46.1}{63}{}{subsection.46.1}{Dirichlet-GRH-v930-Volume-IX.pdf}}
\expandafter\def\csname r@sec:weightchart362\endcsname{{36.8}{21}{}{subsection.36.8}{Dirichlet-GRH-v930-Volume-VII.pdf}}
\expandafter\def\csname r@sec:weightintersections419\endcsname{{44.5}{16}{}{subsection.44.5}{Dirichlet-GRH-v930-Volume-IX.pdf}}
\expandafter\def\csname r@seriesguide\endcsname{{1.1}{4}{}{subsection.1.1}{Dirichlet-GRH-v930-Volume-I.pdf}}
\expandafter\def\csname r@thm:adjcomponent372\endcsname{{37.36}{69}{}{proposition.37.36}{Dirichlet-GRH-v930-Volume-VII.pdf}}
\expandafter\def\csname r@thm:adjointseed361\endcsname{{36.17}{17}{}{proposition.36.17}{Dirichlet-GRH-v930-Volume-VII.pdf}}
\expandafter\def\csname r@thm:alllabels375\endcsname{{38.11}{77}{}{proposition.38.11}{Dirichlet-GRH-v930-Volume-VII.pdf}}
\expandafter\def\csname r@thm:auxdescent361\endcsname{{36.18}{18}{}{proposition.36.18}{Dirichlet-GRH-v930-Volume-VII.pdf}}
\expandafter\def\csname r@thm:commonframe365\endcsname{{37.8}{41}{}{proposition.37.8}{Dirichlet-GRH-v930-Volume-VII.pdf}}
\expandafter\def\csname r@thm:commonsplit483\endcsname{{68.15}{96}{}{proposition.68.15}{Dirichlet-GRH-v930-Volume-XIII.pdf}}
\expandafter\def\csname r@thm:conductor232\endcsname{{19.15}{12}{}{proposition.19.15}{Dirichlet-GRH-v930-Volume-IV.pdf}}
\expandafter\def\csname r@thm:cot116\endcsname{{32.19}{44}{}{proposition.32.19}{Dirichlet-GRH-v930-Volume-VI.pdf}}
\expandafter\def\csname r@thm:crystalline121\endcsname{{33.11}{57}{}{proposition.33.11}{Dirichlet-GRH-v930-Volume-VI.pdf}}
\expandafter\def\csname r@thm:exactstricttrace470\endcsname{{47.8}{82}{}{proposition.47.8}{Dirichlet-GRH-v930-Volume-IX.pdf}}
\expandafter\def\csname r@thm:first72\endcsname{{32.15}{40}{}{proposition.32.15}{Dirichlet-GRH-v930-Volume-VI.pdf}}
\expandafter\def\csname r@thm:globalhecke423\endcsname{{44.39}{32}{}{proposition.44.39}{Dirichlet-GRH-v930-Volume-IX.pdf}}
\expandafter\def\csname r@thm:levelstrict345\endcsname{{22.17}{75}{}{proposition.22.17}{Dirichlet-GRH-v930-Volume-IV.pdf}}
\expandafter\def\csname r@thm:log147\endcsname{{20.4}{26}{}{proposition.20.4}{Dirichlet-GRH-v930-Volume-IV.pdf}}
\expandafter\def\csname r@thm:oddquarticfield539\endcsname{{22.28}{85}{}{proposition.22.28}{Dirichlet-GRH-v930-Volume-IV.pdf}}
\expandafter\def\csname r@thm:originalfinite424\endcsname{{44.41}{34}{}{proposition.44.41}{Dirichlet-GRH-v930-Volume-IX.pdf}}
\expandafter\def\csname r@thm:originalorder424\endcsname{{44.44}{35}{}{proposition.44.44}{Dirichlet-GRH-v930-Volume-IX.pdf}}
\expandafter\def\csname r@thm:pairedtrace651\endcsname{{62.23}{89}{}{proposition.62.23}{Dirichlet-GRH-v930-Volume-XII.pdf}}
\expandafter\def\csname r@thm:pres70\endcsname{{32.13}{37}{}{proposition.32.13}{Dirichlet-GRH-v930-Volume-VI.pdf}}
\expandafter\def\csname r@thm:primitivegram463\endcsname{{141.27}{26}{}{proposition.141.27}{Dirichlet-GRH-v930-Volume-XXVII.pdf}}
\expandafter\def\csname r@thm:realtwist404\endcsname{{41.17}{46}{}{proposition.41.17}{Dirichlet-GRH-v930-Volume-VIII.pdf}}
\expandafter\def\csname r@thm:reduciblefiber376\endcsname{{38.15}{80}{}{proposition.38.15}{Dirichlet-GRH-v930-Volume-VII.pdf}}
\expandafter\def\csname r@thm:reducibleminors378\endcsname{{40.7}{6}{}{proposition.40.7}{Dirichlet-GRH-v930-Volume-VIII.pdf}}
\expandafter\def\csname r@thm:restrictedtrace374\endcsname{{38.6}{74}{}{proposition.38.6}{Dirichlet-GRH-v930-Volume-VII.pdf}}
\expandafter\def\csname r@thm:slthree363\endcsname{{36.26}{30}{}{proposition.36.26}{Dirichlet-GRH-v930-Volume-VII.pdf}}
\expandafter\def\csname r@thm:smooth69\endcsname{{32.11}{35}{}{proposition.32.11}{Dirichlet-GRH-v930-Volume-VI.pdf}}
\expandafter\def\csname r@thm:sympordinary468\endcsname{{47.2}{77}{}{proposition.47.2}{Dirichlet-GRH-v930-Volume-IX.pdf}}
\expandafter\def\csname r@thm:symprelations469\endcsname{{47.5}{79}{}{proposition.47.5}{Dirichlet-GRH-v930-Volume-IX.pdf}}
\expandafter\def\csname r@thm:tower138\endcsname{{24.28}{23}{}{proposition.24.28}{Dirichlet-GRH-v930-Volume-V.pdf}}
\expandafter\def\csname r@thm:trace71\endcsname{{34.2}{65}{}{proposition.34.2}{Dirichlet-GRH-v930-Volume-VI.pdf}}
\expandafter\def\csname r@thm:weightchart362\endcsname{{36.21}{23}{}{proposition.36.21}{Dirichlet-GRH-v930-Volume-VII.pdf}}
\expandafter\def\csname r@thm:wild98\endcsname{{25.1}{24}{}{proposition.25.1}{Dirichlet-GRH-v930-Volume-V.pdf}}
\makeatother
\begin{document}
\setcounter{part}{10}
\fancyhead[L]{\small\color{muted}Dirichlet GRH \textperiodcentered\ Volume XI}
\fancyhead[R]{\small\color{muted}Strict cohomology at exceptional weights}
\thispagestyle{plain}\begin{center}
{\LARGE\bfseries The generalized Riemann hypothesis\\[2pt]for Dirichlet \(L\)-functions}\\[6pt]
{\small 27 November 2026 \quad\textperiodcentered\quad Edition 930}\end{center}
\part{Strict cohomology at exceptional weights}\label{part:XI}
\begin{abstract}\noindent We identify the twists at which two ordinary characters coincide. Bounds on conductors and the spacing of weights reduce the surviving cases to unit, regulator and class group calculations in the real cyclotomic tower and its quadratic extensions.\end{abstract}
\setcounter{section}{53}
\noindent The \hyperref[seriesguide]{guide in Volume~I} lists all twenty-seven volumes.
\begingroup\small\setlength{\parskip}{0pt}\tableofcontents\endgroup
\clearpage

\section{Strict orders at varying coefficient primes}\label{nav:chapter-varying-coefficients}\label{nav:volume-XI}
\subsection{The large-prime cohomology bound on the period fields}
\label{sec:periodfieldlarge425}

Continue with \(K=\QQ(\sqrt{-s})\), \(E_x\), the original
weight and the unramified coefficient ring \(O\) of
Section~\ref{sec:originalfixedhecke424}. Write \(D=J/O\) and
\(S(\xi)\) for classes in \(H^1(\QQ,D(\xi))\) unramified at
finite primes away from \(\ell\) and zero at \(\ell\). Let
\(B=\bigoplus_j B_j\) consist of the finite sums of homogeneous
elements in the completed marked solution ring. Here degree
is the weight for conjugation by
\(\operatorname{diag}(a,1,a^{-1})\). Rows and columns are
numbered \(0,1,2\).

\begin{proposition}[The marked presentation on the chosen quadratic field]\leavevmode
\label{prop:originalchart425}
One has \(B_0=T\). When \(S(\Psi^{-1})\ne0\), the augmented
graded presentation of \(B\) satisfies the numerical and
cotangent hypotheses of Proposition~\ref{prop:localmult399}.
Its negative weight-one cotangent quotient has length
\[
        \operatorname{length}_O S(\Psi^{-1})+
        \operatorname{length}_O S(\eta\Psi^{-1}),
\]
where \(\eta\) is the quadratic character of \(K\).
The fixed subscheme consists of the split section \(\operatorname{Spec}O\).
\end{proposition}
\begin{proof}
We check the dependence on the auxiliary quadratic character in
Sections~\ref{sec:original68}--\ref{sec:origin116}.
The prime \(\ell\) exceeds \(sN_*\), \(4t_*+5\) and the class
number of \(K\) throughout this range, for sufficiently large
\(V\). Both \(\alpha\) and \(\eta\) are unramified there,
\(\Psi(c)=1\), and \(\eta(c)=-1\). Choose
\(g_*\) in inertia at \(\ell\) mapping to a generator of
\(\mathbf F_\ell^\times\) under the residual cyclotomic
character. The three diagonal residual values are distinct.
Thus the ordinary word cuts out the unique ordered inertia
flag by Lemma~\ref{lem:flag69}. For \(r\ge2\) the dual
twisted strictly upper triangular module has inertia exponents
\(1-r\) and \(1-2r\), both nonzero modulo \(\ell-1\).
Theorem~\ref{thm:smooth69} gives the local chart. For \(r=1\)
the proof of Theorem~\ref{thm:first72} uses the unramified
quadratic algebra at \(\ell\). In the split case the two
adjacent characters are separate unramified twists of the
cyclotomic character. In the field case they form an induced
module from the unramified quadratic extension. In each case
the sum of the local degrees in Lemma~\ref{lem:kummer72} is
two, saturation removes the valuation coordinates, and the
outer obstruction vanishes by the residual exponent \(1-2\).
This gives the same local coordinate counts
\((1,2,3,4,2)\) in weights \(-2,-1,0,1,2\).

Away from \(\ell\) the inertia quotient has order prime to
\(\ell\), including at \(s\). In the relative cochain complex
of Lemma~\ref{lem:rel70} these places therefore contribute
zero Euler characteristic. The action at infinity is minus
reflected transpose as before. If \(a_j\) is the relative
tangent dimension in weight \(j\), the obstruction space has
dimension \(a_j+\mathbf1_{|j|\le1}\). Indeed the correction
in degree three is dual to the global invariants of the
dual twisted module. Inertia at \(\ell\) excludes all weights
except possibly \(j=r=1\). In that case the two characters
are \(\bar\alpha^{-1}\) and \(\bar\eta\bar\alpha^{-1}\).
The first is odd and the second is nontrivial on inertia
at \(s\), since \(s\nmid N_*\ell\). This proves the asserted
dimension by the Euler calculation in that lemma. The
small-extension argument of Theorem~\ref{thm:pres70} and
the eight marking entries now give the following numbers.
\[
 \begin{array}{c|rrrrr}
  j&-2&-1&0&1&2\\ \hline
  \text{variables}&a_{-2}+1&a_{-1}+2&a_0+3&a_1+4&a_2+2\\
  \text{relations, at most}&a_{-2}+1&a_{-1}+3&a_0+3&a_1+3&a_2+1 .
 \end{array}
\]
Homogeneous generators and the ambient formal chart are obtained
from the power series presentation by taking weight parts as
in Lemma~\ref{lem:algebra116}.

At a fixed lift all matrices on \(G_K\) are diagonal. The
quotients of their three coordinate characters by their
prescribed values reduce to one and are everywhere unramified.
At \(\ell\) this follows from the ordered flag and its distinct
residual characters, and at the other primes it follows by the
prime-to-\(\ell\) inertia condition. Over an Artinian local
test ring these quotients therefore factor through both the
class group of \(K\) and a finite \(\ell\)-group. They are
trivial since \(\ell\nmid|\mathrm{Cl}(K)|\). The proof of
Proposition~\ref{prop:fixed116} applied to all Artinian
truncations identifies the fixed quotient with \(O\). It also
shows that the equations in degree zero span the integral
degree-zero cotangent space.

For weight \(-1\) the two lower adjacent directions split
into the characters \(\Psi^{-1}\) and \(\eta\Psi^{-1}\).
The condition at \(c\) fixes the odd coboundary, and the
adjacent marking at \(g_*\) fixes the other coboundary:
its coefficient \((\Psi(g_*)+1)(\Psi(g_*)^{-1}-1)\) is
a unit. The lower long entry has character \(\eta\Psi^{-2}\);
its coboundary is fixed at \(c\). For either weight the
ordinary local word forces a lower cocycle to be a local
coboundary at \(\ell\): after removing its value at \(g_*\)
by a local coboundary, the ordered flag is the standard flag
and its lower entries vanish on inertia. The cocycle law and
the unit difference on \(g_*\) then force it to vanish on the
decomposition group. Conversely conjugating the split lift
realizes each such local coboundary on the selected component.
Away from \(\ell\) averaging over the finite inertia quotient
gives exactly the unramified condition. Applying this argument
over \(O/\ell^n\) for every \(n\), as in
Corollary~\ref{cor:dual68} and Theorem~\ref{thm:cot116},
identifies the \(D\)-duals of the two negative cotangents with
\(S(\Psi^{-1})\oplus S(\eta\Psi^{-1})\) and \(S(\eta\Psi^{-2})\).
The passage to \(D\) commutes with the strict conditions since
the invariants at \(\ell\) vanish and the invariants on each
other inertia group are zero or divisible. Restriction to
\(E_xK\) is injective on these groups by the prime-to-\(\ell\)
index. The finite characters become trivial there, and the
strict groups of twists \(r,2r\) are finite by
Theorem~\ref{thm:wild98}. Thus both negative cotangents are
torsion, supplying rational spanning and the stated length.

Finally \(Z=M(g_*)-M(cg_*^{-1}c)=\operatorname{diag}(z,0,-z)\)
with \(z\) a unit. Its square trace gives \(2z^2\), so the
specified square root \(z\) belongs to the closed trace ring
by Hensel's lemma. Its spectral projectors are therefore
polynomials in \(Z\) over \(T\). The adjunction relation
\(M(ch^{-1}c)=J_3 M(h)^{\mathsf T}J_3\) puts each
weight-zero entry monomial into products of directed cycles,
as proved in Theorem~\ref{thm:trace71}. Inserting the projectors
expresses each cycle as a trace. Closure gives \(B_0=T\).
Moreover the entries of finitely many group generators generate
\(B\) over \(T\). Indeed their homogeneous differences from
the split values topologically generate the completion. For
any one weight, every monomial in those differences has a
weight-zero factor times one of finitely many remainders by
Lemma~\ref{lem:algebra116}. Grouping each homogeneous series
by its remainders gives the claimed algebraic generation.
\end{proof}

\begin{lemma}[Multiplicity at a nonsplit trace]\leavevmode
\label{lem:markmultiplicity425}
For each nonsplit geometric characteristic-zero character
\(\sigma\) of \(T\), put \(k_0=\overline J\) and
\(B^\sigma=B\otimes_{T,\sigma}k_0\). Then
\[
  \dim_{k_0} B^\sigma_0=1,\qquad
  \dim_{k_0} B^\sigma_1=
      \begin{cases}1&\text{at an irreducible trace},\\
                    0&\text{at a line and an irreducible plane}.
      \end{cases}
\]
These are also the dimensions obtained by localization in
\(T\otimes_O k_0\).
\end{lemma}
\begin{proof}
Corollary~\ref{cor:originaletale424} makes the indicated local
factor of \(T\otimes_O k_0\) equal to \(k_0\). Taking weights
commutes with this scalar extension, so \(B^\sigma_0=k_0\).
Write \(E^\sigma=E_M\otimes_{T,\sigma}k_0\) for the finite
matrix algebra of Lemma~\ref{lem:torsioncrys375}. The universal
group matrices have homogeneous entries, and therefore give
a homomorphism \(E^\sigma\to M_3(B^\sigma)\) preserving
trace. Each of the three spectral projectors \(P_i\) of \(Z\)
is the corresponding diagonal matrix unit in this map.

At an irreducible trace \(E^\sigma\) maps onto the full
endomorphism algebra on its semisimple realization, by
Burnside's theorem. Each \(P_i\) has trace one and is a
rank-one projector there. In the eigenbasis of \(Z\), choose
two elements \(a,b\) mapping to opposite matrix units between
the first and middle lines. Their order can be chosen so that
\(\operatorname{tr}(P_0aP_1bP_0)=1\). In the universal matrices
over \(B^\sigma\) this same trace is \(a_{01}b_{10}\). These
entries have degrees \(1,-1\), respectively, since elements
of \(E^\sigma\) are scalar linear combinations of group matrices.
Hence \(a_{01}\) is a unit of degree one, and multiplication
by it identifies \(B^\sigma_0\) with \(B^\sigma_1\).

At a line and an irreducible plane
Lemma~\ref{lem:planeunits379}, applied first over a finite
splitting field and then over \(k_0\), gives
\(E^\sigma=k_0\times M_2(k_0)\), with the group acting trivially
on the first factor. Here the finite algebra calculation is
essential: the proof uses the vanishing of extensions in both
directions between the plane and the line. On the trivial
factor \(Z=0\), so the middle projector \(P_1=1-Z^2/z^2\)
is one there. Its total trace is one. On the plane it is an
idempotent of trace zero in characteristic zero, hence zero.
Thus \(P_1\) is central in \(E^\sigma\). Its image in
\(M_3(B^\sigma)\) is the middle diagonal matrix unit, so
every group matrix has all four adjacent entries zero in
\(B^\sigma\). The generation assertion of the proposition now
shows that \(B^\sigma\) is generated in even degrees. This
proves the assertion on the full odd piece, including its
possible infinitesimal elements.
\end{proof}

\begin{theorem}[The integral strict length for one character]\leavevmode
\label{thm:strictlarge425}
Whenever \(S(\Psi^{-1})\ne0\), with \(I,h\) as in
Theorem~\ref{thm:originalorder424}, one has
\[
 \operatorname{length}_O S(\Psi^{-1})+
 \operatorname{length}_O S(\eta\Psi^{-1})
       \ \le\ 2\operatorname{length}_O(I/hI).
\]
In particular \(\operatorname{length}_O S(\Psi^{-1})\) is
bounded by the right hand side for every character in the range.
\end{theorem}
\begin{proof}
Proposition~\ref{prop:originalchart425} supplies the presentation,
the rational formal chart, and the two negative spanning conditions
in Proposition~\ref{prop:localmult399}. The fixed section is
unique. Theorem~\ref{thm:originalfinite424} gives integral
finiteness of \(B_0\). The Frobenius difference \(h\) vanishes at
the split character. At every nonsplit point it is nonzero by the
archimedean estimate in Theorem~\ref{thm:originalorder424}.
The characters nonzero on \(B_{-2}B_2\) form a subset of the
nonsplit points, since every nonzero-degree element has split
value zero. At each of these characters the largest of the two
local dimensions in Proposition~\ref{prop:localmult399} is one
by the lemma. Applying that integral intersection inequality
and enlarging the sum to all nonsplit points gives the bound
by \(2\sum_{\sigma\ne\mathrm{split}}v_\ell(\sigma(h))\).
The determinant formula in Theorem~\ref{thm:originalorder424}
identifies this sum without its factor two with the length of
\(I/hI\).
\end{proof}

\begin{corollary}[All large coefficient primes on the original fields]\leavevmode
\label{cor:periodlarge425}
For some constant \(C\) and every sufficiently distant
polynomial-conductor interval of Section~\ref{sec:missing},
one has, for all \(V\le x\le V^{3/2}\),
\[
 \mathcal L_{>}(x)=
 \sum_{\ell>P(x)}(\log\ell)\sum_{1\le r<t(x)}
      \operatorname{length}_{\ZZ_\ell}
      \bigl(\mathcal S_{\ell,r}(E_x)^{c=(-1)^r}\bigr)
   \ \le\ 2\log D(V)\ \ll_{\psi,\varepsilon} V^C .
\]
For any such \(\ell\) prime to \(D(V)\) all these complementary
groups are zero. The large-prime summands \(R_{\ell,r}\) of
Section~\ref{sec:missing} satisfy the same total bound.
\end{corollary}
\begin{proof}
Since \(\ell>[E_x:\QQ]\), the character decomposition in
Theorem~\ref{thm:levelstrict345} is integral after unramified
coefficient extension. It gives the length at \(\ell,r\) as
the sum of \(b_\alpha\operatorname{length}_O S(\Psi^{-1})\)
over the complementary residue orbits. The conductor conditions
in that theorem change none of the lengths here because the
finite inertia has order prime to \(\ell\). Apply the theorem
just proved to these scalar groups and sum the result using
Theorem~\ref{thm:originalorder424}, which uses a single \(D(V)\)
for the interval. Its proof gives the same comparison at each
individual \(\ell\) with \(v_\ell(D(V))\), proving the vanishing
assertion. For \(\ell>P(x)\), \(R_{\ell,r}\) is a direct summand
of the displayed complementary group by definition of the field
average. This gives its bound.
\end{proof}

Since \(V^C=o(\exp(x^\varepsilon/2))\) uniformly for \(x\ge V\),
the contribution from \(\ell>P(x)\) fits the allowance in
equation~\eqref{eq:length12}. For the odd primes at most \(P(x)\),
Proposition~\ref{prop:res23} leaves the quotients by the subgroup
killed by \(\ell^{b_0}\). Section~\ref{sec:missing} also requires
the integral comparison of projected theta periods over these
same \(E_x\).


\subsection{Variation of the twist in the original weight space}
\label{sec:regularresidual426}

Fix \(E_*,N_*,K,s,u\) and \(D(V)\) as in
Theorem~\ref{thm:originalorder424}. For this subsection allow
any prime \(\ell\ge11\) such that
\[
          \ell\nmid suN_*[E_*:\QQ]\,|\mathrm{Cl}(K)|.
\]
Let \(n=\ell-1\). For each twist \(r\ge1\) let \(j\) be its
representative in \(\{0,\dots,n-1\}\). We consider \(r=1\),
and also all \(r\) satisfying
\[
               2\le j\le n-2,\qquad n/\gcd(n,j)>4.        \tag{*}
\]
Write \(\bar\omega=\bar\kappa_\ell|_{I_\ell}\).
The significance of \((*)\) is that \(\bar\omega^j\) has
order greater than four and \(j\not\equiv\pm1\bmod n\).
The period fields and the test prime \(u\) are the same for
the small and large coefficient primes.

\begin{lemma}[The auxiliary ordinary ring at a residual weight]\leavevmode
\label{lem:baseweight426}
Let \(\alpha\) be a character of \(E_x\) with
\(\alpha(c)=(-1)^r\). At \(r=1\) put \(j=1\). For each of
the indicated \(j,\alpha\), there is a layer \(L=KF_a^+\)
of Proposition~\ref{prop:auxfield361} for which every
geometrically irreducible prime of the ordinary determinant
ring with residual character \(\bar\alpha\bar\kappa_\ell^j\)
contains \(I_K\). The weight map for any of the indicated
positive twists congruent to \(j\) uses this same ring.
\end{lemma}
\begin{proof}
We verify the input to the propagation argument
of Theorem~\ref{thm:globalhecke423} with the exponent \(j\).
Only the residue of the exponent is used to specify the
ordinary determinant and its finite inertia and diamond
characters. The choice of a classical point to start the
argument also needs an integral exponent of the specified
parity. Here \(j\) itself has that parity since \(n\) is even.

The character \(\alpha\) has order prime to \(\ell\) and is
unramified above it. At \(s\) the inertia character of \(K\)
is nontrivial while \(\alpha\) is unramified. Thus
\(K\cap E_*=\QQ\), \(\eta\alpha\ne1\), and its conductor is
\(s\,\operatorname{cond}(\alpha)\). Take the auxiliary prime
\(w\equiv1\bmod sN_*\ell\) with \(\ell\) a non-\(\ell\)-th
power modulo \(w\) by Lemma~\ref{lem:inertaux364}. In its
Kummer argument \(\ell\nmid sN_*\) gives valuation \(n\)
for \(\ell\) in the cyclotomic field. The other auxiliary
prime may be chosen in the progression
\(-1\bmod 8sN_*\ell\). The resulting \(L\) is disjoint
from \(E_*(\mu_\ell)\) as in Proposition~\ref{prop:auxfield361}:
the only possible common ramification is at \(\ell\), with
coprime ramification indices there. All designated finite
places split in \(L/F_a^+\) and inertia on \(G_L\) has
residual cyclotomic image of order \(n\). In particular one
can choose the element of \(G_L\) acting as complex conjugation
on \(E_*(\mu_\ell)\), used in the relative duality arguments.

For the ordinary adjoint use weight \(j+1\) in
Lemma~\ref{lem:cusp361}. Here \(2\le j+1<\ell-1\), so the
Bernoulli denominators through degree \(j+1\) are prime to
\(\ell\). Its auxiliary level
\(s\,\operatorname{cond}(\alpha)w\) is prime to \(\ell\).
Therefore the Eisenstein congruence, free diamond action,
and lifting of the indicated cuspidal character in that lemma
apply. The order of \(\bar\omega^j\) is greater than four
also for \(j=1\). Thus the quadratic self-twist excluded in
Theorem~\ref{thm:adjointseed361} is impossible. The resulting
adjoint is absolutely irreducible on every open subgroup
and ordinary with the required ordered residual characters.
Its classical ordinary deformations and the full Hecke
component of Theorem~\ref{thm:adjcomponent372} use this
seed. The ordinary flags are separated by inertia, and
the exclusion of a cyclic permutation of three lines in
Theorem~\ref{thm:slthree363} uses
\((\bar\omega^j)^3\ne1\).

For the strict classes on the layer the two characters in
Proposition~\ref{prop:strictclasses365} have inertia exponents
\(-ij\) and their Tate duals have exponents \(1+ij\),
\(i=1,2\). These four exponents are nonzero modulo \(n\).
For \(1+2j\) this follows from parity, for \(1+j\) from
\(j\ne n-1\), and for the other two from the order condition.
The local contribution at \(w\) is still one, so the proof
supplies both required strict classes. In
Proposition~\ref{prop:classrank365} restriction to
\(L E_*(\mu_\ell)\) has prime-to-\(\ell\) index over \(F_a^+\)
or \(L\), since that index divides
\(2[E_*:\QQ]n\). The bound \(h\) on class group generator
numbers along the cyclotomic tower follows there from
Ferrero--Washington. It may depend on \(\ell,E_*,w\).
Theorem~\ref{thm:commonframe365} then chooses \(a\) after
\(h\). The needed large real degree, local degrees and
intersection dimensions thus hold without any upper bound
on the conductor in terms of \(\ell\).

For clarity, the two residual exclusions in the subsequent
propagation have different origins. A hypothetically induced
plane would acquire a quadratic self-twist, interchanging
the two outer residual characters. This forces
\((\bar\omega^j)^4=1\), by Lemma~\ref{lem:planecurve417}.
A cyclic permutation of three lines in the absolutely
irreducible case would give a cubic self-twist, forcing
\((\bar\omega^j)^3=1\), as in the irreducible-curve argument
of Section~\ref{sec:irrcurvesupport409}. Both are excluded.
The residual projector constructions require the three
eigenvalues to be distinct, and the matrix and sign
projections require \(2\) and \(3\) to be units. The
Taylor--Wiles primes for the separated curves are chosen
after those curves, by the dual-space and Frobenius
argument of Proposition~\ref{prop:jointpatch417}.
For the integral Hecke modules at these primes use the
old-type coefficient lattices of Lemma~\ref{lem:typelattice427}.
This includes the primes dividing the order of a local finite
matrix group at the tame level.
The ordinary level, rank and complete-intersection arguments
apply to that field with those primes. The dimension bounds
in Sections~\ref{sec:weightintersections419} and
\ref{sec:markedintersections420} use \(d=[F_a^+:\QQ]\),
the degrees at \(\ell\) and \(h\) just obtained. With
these inputs the marked-chart propagation and the connected
lower fiber give Theorem~\ref{thm:globalhecke423} on the
chosen layer.

Finally for \(r\equiv j\bmod n\), the ordered residual
characters agree on the whole group. On local units the
change of exponent \(r-j\) is trivial on the prime-to-\(\ell\)
part of the cyclotomic character, whose order divides \(n\).
On the principal units it defines an arithmetic specialization
of the full weight algebra. The common cyclotomic normalization
is independent of \(r\). This proves the assertion about the ring.
\end{proof}

\begin{proposition}[Specialization at the actual twist]\leavevmode
\label{prop:weightshift426}
At any \(r=1\) or \(r\) satisfying \((*)\), the fixed-weight
ring \(T\) of the original quadratic field at exponent \(r\)
is finite over \(O\), with geometrically reduced generic
fiber. Its nonsplit points have the two sorts of classical
labels of Proposition~\ref{prop:alloriginalpoints424} at
the actual weight \(r\).
\end{proposition}
\begin{proof}
Use the upper ordinary problem from the lemma, while prescribing
\(\Psi=\alpha\kappa_\ell^r\) in the original problem. The
inertia element \(g_*\) separates the three original characters
since the order of \(\bar\omega^r=\bar\omega^j\) exceeds four.
The restriction-of-traces proof of
Proposition~\ref{prop:orighecke374} therefore gives a surjection
onto \(T\) at the arithmetic weight of \(r\): the polynomial
projector identities and equality of matrix spans there divide
only by residual unit differences. In characteristic \(\ell\)
the reducible trace argument of
Theorem~\ref{thm:reducibleminors378} applies. In the plane
case its rational ordinary pseudorepresentation has odd
determinant lifted by \(\eta\Psi^{-2}\), de Rham at \(\ell\).
Its reduction is nontrivial on \(I_\ell\) because
\((\bar\omega^j)^2\ne1\). Its local constituent lifting
the trivial residual character is unramified after the
prescribed twist. These are the hypotheses for the finiteness
argument using Pan over \(\QQ\) in that theorem. The
class-number argument for its one-dimensional summand uses
\(\ell\nmid|\mathrm{Cl}(K)|\). At a prime with irreducible
trace the preceding lemma puts its restriction in Hecke
support. The special-fiber argument of
Theorem~\ref{thm:originalfinite424} now gives finiteness.

For \(r>1\) the residual Tate dual of the strictly upper
triangular module has exponents \(1-j,1-2j\) on inertia.
Both are nonzero. The ordinary local chart is consequently
formally smooth by the obstruction calculation of
Theorem~\ref{thm:smooth69}. At arithmetic weight \(r\ge2\)
every successive positive difference of the three cyclotomic
exponents is at least two. In constructing an extension of an
earlier flag step by the next graded quotient, the Hom module
from the quotient to the earlier step therefore has positive
exponents and has no quotient with exponent one. Its Tate
dual has zero invariants. The local Euler characteristic and
the Bloch--Kato dimension formula make all such extensions
of crystalline steps crystalline, as in
Theorem~\ref{thm:crystalline121}. Density of integral
evaluations in the smooth chart and the fixed Hodge-type
crystalline quotient give that result on the whole chart.
At \(r=1\) the unit part of the Kummer chart gives exactly
the first-twist calculation of Theorem~\ref{thm:first72}.
The finite matrix module and limit argument of
Proposition~\ref{prop:tracecrystal375} uses the fixed weight
interval \([-r,r]\) for this triple. Liu's crystalline
limit theorem permits any fixed finite interval for odd
\(\ell\).~\cite{XI-VIII-liu375,XI-VIII-VI-kisin121}

The ordinary control theorem on the upper layer applies to
the dominant algebraic weight of \(r\). The cuspidal descent
in Proposition~\ref{prop:alloriginalpoints424} uses the tame
inertia type \(\alpha,1,\alpha^{-1}\) and the prime-to-three
degree \([L:K]\). It gives the desired original level for
irreducible points. For a line and a plane the odd descent
and primitive form computation in Propositions~\ref{prop:odddescent376}
and~\ref{prop:planelabels379} use distinct weights and the
absence of a quadratic self-twist, guaranteed by the residual
order greater than four. Thus the modular weight is \(2r+1\)
and the level divides \(sN_*\). The Selmer vanishing for the
line--plane extensions uses Kato's off-center twist \(r\)
of that same form, since \(1\le r\le 2r\) and \(2r\ne2r+1\).
The enormous-image proof at irreducible points uses the
three distinct Hodge--Tate weights and excludes the cubic
residual self-twist. Corollary~\ref{cor:originaletale424}
therefore computes these generic local rings at the shifted
weight as well.~\cite{XI-VIII-VI-kato120,XI-VIII-VI-nt119}
\end{proof}

\begin{theorem}[Full lengths for the indicated residual twists]\leavevmode
\label{thm:regularlength426}
Sum over the primes specified at the beginning of this
subsection and over \(1\le r<t(x)\) with \(r=1\) or \((*)\).
Then throughout the same interval
\[
 \sum_{\ell}\log\ell\sum_{r}
      \operatorname{length}_{\ZZ_\ell}
        \bigl(\mathcal S_{\ell,r}(E_x)^{c=(-1)^r}\bigr)
                       \ \le\ 2\log D(V).
\]
The same estimate holds at one prime with the factors
\(\log\ell\) removed and the right side replaced by
\(2v_\ell(D(V))\).
\end{theorem}
\begin{proof}
For \(r\) satisfying \((*)\), the five dual twisted global
matrix weights have exponents \(1-ar\), \(a=0,\pm1,\pm2\),
on inertia. They are all nonzero modulo \(n\): the two
even multiples cannot equal the odd number one modulo even
\(n\), and \(j\ne\pm1\) accounts for the adjacent weights.
For \(r=1\) the possible global invariants in weight one
are again excluded by parity and ramification at \(s\).
The proof of Proposition~\ref{prop:originalchart425} now gives
the same numerical presentation. Strictness at \(\ell\)
uses the unit differences for \(\bar\omega^{-j}\) and
\(\bar\omega^{-2j}\). Prime-to-\(\ell\) restriction to
\(E_xK\) proves finiteness of the negative cotangents even
when \(r>\ell\). The fixed-point argument uses the stated
class-number condition.

The proposition gives the finite \'{e}tale generic trace
algebra at this weight. At a line and plane its specialized
matrix algebra is the product of the scalar and full plane
algebras by the vanishing of the extensions just proved
and the opposite-matrix-unit calculation of
Lemma~\ref{lem:planeunits379}. Thus the proof of
Lemma~\ref{lem:markmultiplicity425} gives multiplicity one
in the maximum over the two graded dimensions. The integral
intersection inequality of Proposition~\ref{prop:localmult399}
bounds the scalar strict length by twice the sum of the
nonsplit Frobenius-difference valuations at \(u\). The
argument applies to every contributing scalar character.

For fixed \(r,\ell\) the character decomposition is integral
since \(\ell\nmid[E_x:\QQ]\). The local corrections at the
conductor vanish for the same reason. For each \(\ell\),
summing over the contributing characters and these \(r\)
uses a sublist of the very factors defining \(D(V)\) in
Theorem~\ref{thm:originalorder424}, since those factors are
indexed by the actual twists \(r\).
More explicitly, the diagonal value is
\(1+u^r\alpha(u)+u^{-r}\alpha(u)^{-1}\) and the plane value
is \(1+u^{-r}\alpha(u)a_u(f)\). The rank-three label has
the corresponding archimedean weight and the same level
bound as in that theorem. The test prime is different from
\(\ell\), all denominator-clearing powers of \(u\) are units,
and any unused factor has nonnegative valuation. Distinct
labels for different points and the separate residue-character
embeddings give the single-prime inequality. Multiplication
by \(\log\ell\) and summation prove the assertion.
\end{proof}

\begin{corollary}[The residual twists at smaller coefficient primes]\leavevmode
\label{cor:remainingresidues426}
At each of the indicated \(\ell\le P(x)\) and positive twists
covered by the theorem, \(R_{\ell,r}=0\) on every sufficiently
distant interval. For \(r\ge2\), a nonzero \(R_{\ell,r}\) at
one of these primes can therefore occur only if
\[
     \ell-1\mid 2r,\quad \ell-1\mid 3r,\quad
     \ell-1\mid 4r,\quad \ell-1\mid r-1,\quad
     \text{or}\quad \ell-1\mid r+1.
\]
In particular \(\ell\le4r+1\) there, and at each such prime
there are at most eight possible residue classes of twists.
\end{corollary}
\begin{proof}
The length of any one complementary group in the theorem
is at most \(2v_\ell(D(V))\le 2\log D(V)/\log 2\).
This is \(O_{\psi,\varepsilon}(V^C)\), uniformly on the
interval. For \(x\ge V\), \(b_0=\lfloor\exp(x^\varepsilon/8)\rfloor\)
eventually exceeds this bound. A finite \(\ZZ_\ell\)-module
of length at most \(b_0\) is killed by \(\ell^{b_0}\), by
its elementary divisors. This holds also for its direct
summand \(U_{\ell,r}\), proving the vanishing.

For the other twists the residue is \(1,-1\), or its
order in \(\ZZ/n\ZZ\) divides one of \(2,3,4\).
These are precisely the indicated divisibility alternatives.
There are at most
\(\varphi(1)+\varphi(2)+\varphi(3)+\varphi(4)=6\)
elements of orders at most four in a cyclic group. Adding
the two unit residues gives eight. For \(r\ge2\) all the
integers on the right of the divisibilities are positive
and at most \(4r\). This gives the prime bound.
\end{proof}

The exceptions on the prime side have a simple description
on the period interval. The odd part of \([E_*:\QQ]\) divides
the fixed degree \([K_\psi:\QQ]\), since the other extensions
adjoined to it have two-power degree. Apart from the two
moving primes \(q,d_0\) in \(N_*\), the excluded primes
therefore come from a fixed finite set. Bounding their tails
and the remaining residue classes is a further arithmetic
requirement in Section~\ref{sec:missing}.


\subsection{Distinct residual inertia characters at a coefficient prime}
\label{sec:ordinaryobstructions427}

Keep the primes \(\ell\), fields and test prime \(u\) of
Section~\ref{sec:regularresidual426}. Put \(n=\ell-1\). We shall
require only \(n\nmid 2r\) on the twist. Thus a separating
inertia element \(g_*\) still exists even when the residual
ratio of two adjacent ordinary characters is cyclotomic.
For \(r\ge2\) we use the full ordinary inertia word at \(\ell\)
in the original marked matrix problem. At the actual twist
\(r=1\) we keep the component of Theorem~\ref{thm:first72}.

\begin{lemma}[Integral coefficients for the prescribed tame types]\leavevmode
\label{lem:typelattice427}
On each layer \(L/F_a^+\) of Proposition~\ref{prop:auxfield361}
constructed with the prime-to-\(\ell\) conductor data above,
the three ordinary modules needed for the changes of level at \(Q\)
can be taken with coefficients in stable lattices in the old
compact types. They are finite free over their respective
weight and diamond group algebras, with the stated diamond
coinvariants. Their characteristic-zero systems have the prescribed
old inertia and zero monodromy.
\end{lemma}
\begin{proof}
The numerical proof of Lemma~\ref{lem:oldidempotent405}
uses \(\ell>N_*^{12}\) before averaging over a finite matrix
group at an old place. A prime coprime to \(N_*\) can divide
that group's order. For instance \(13\mid 3^3-1\).
We give the integral construction at any prime in the present
range. For each old compact type \(\sigma_v\) used in that
lemma, enlarge the coefficient field \(J'\) to realize it.
Its compact group acts through a finite quotient. The sum
of the translates of any full lattice is an invariant full
lattice. Take the tensor product \(W\) of their dual lattices.
At an unramified place use the trivial lattice.

At the old places use the full maximal compact group and its
action on \(W\). An automorphic form is an equivariant function
on the finite adelic double cosets with these coefficients.
At a representative \(g\) its possible values are precisely
the invariants in \(W\) under its arithmetic stabilizer
\(\Gamma_g\), with the ordinary algebraic coefficient included
when present. These stabilizers are finite. Their rational
points are integral at every finite place for some lattice,
and the real group is compact, so discreteness gives finiteness.
They have order prime to \(\ell\). In fact an element of order
\(\ell\) would act faithfully on \(L^3\), with a primitive
\(\ell\)-th root as an eigenvalue. Its degree over \(L\) is
\(\ell-1>3\) by the disjointness of
Proposition~\ref{prop:auxfield361}. Averaging over each
\(\Gamma_g\) consequently gives free coefficient modules and
commutes with reduction and extension of coefficients.

Consider any of the normal diamond covers at the coefficient
places or at \(Q\) with finite \(\ell\)-group quotient \(\Delta\).
The stabilizer at the lower level maps trivially to \(\Delta\),
since its order is prime to \(\ell\). Thus over each lower
double coset the lifts are a free transitive \(\Delta\)-set.
Evaluation at one lift, followed by right translation to all
the lifts, identifies its contribution with the regular
module tensored with that lift's invariant lattice. This
also gives the trace identification of coinvariants with the
lower level. Dualizing gives the corresponding assertion for
the dual modules. The argument works simultaneously for
the independent normal covers. The finite diamonds of order
prime to \(\ell\) select their prescribed character by an
integral idempotent.

For precision, the remaining change in the ordinary filtration
is the inclusion between the levels with indices \((b,b)\)
and \((b,c)\), \(c\ge b\), at a coefficient place. The
double-coset identity in Geraghty's Lemma~2.5.2 sends a form
at the deeper level to the shallower level by a power of the
positive diagonal operator. The coset representatives act
only at the coefficient place, hence act trivially on \(W\).
Invertibility of that operator on the ordinary part gives the
same control identity with these coefficients. The ordinary
projector commutes with the diamonds and the trace maps.
Its images in the finite free diamond modules are therefore
projective, hence free over the local group rings. Compatible
coinvariants give a common rank at every depth. Lifts of a
basis at the residual level give a surjection from the free
module of this rank over the completed group algebra by
compact Nakayama. Modulo each open diamond kernel it is a
surjection of equal-rank free modules, so an isomorphism.
These kernels have zero intersection on the source. This proves
freeness over the completed algebra, also with the diamonds
at \(Q\) retained. Localization at the common residual Hecke
system is a direct summand and preserves these conclusions.
This is the double-coset proof of Geraghty's Lemma~2.2.6 and
Proposition~2.5.3 with the old coefficient lattice carried
through the maps.~\cite{XI-VIII-geraghty363}

Over \(J'\), the equivariance condition at \(v\) selects
\(\operatorname{Hom}_{\mathrm{GL}_3(\mathfrak o_v)}
(\sigma_v,\pi_v)\) in the local automorphic representation.
The compact-type theorem in Lemma~\ref{lem:oldidempotent405}
therefore gives exactly its old inertial condition. In
particular the ordinary vector of the initial adjoint occurs.
The change-of-level operators and the integral identities
at \(Q\) act by finite double cosets with identity component
at the old places. Their convolution relations, and the
classical tests of Proposition~\ref{prop:ordinaryclass405},
apply to these modules. The determinant identities first
checked at the dense classical weights then hold integrally,
since the operators act on free modules over the weight ring.
This supplies the integral modules used in the support argument
of Lemma~\ref{lem:baseweight426} even when the finite matrix
group at an old conductor prime has order divisible by \(\ell\).
\end{proof}

\begin{lemma}[A local presentation with the ordinary inertia word]\leavevmode
\label{lem:ordinaryrelations427}
Let \(r\ge2\), \(n\nmid 2r\), and let \(W_1\) be the residual
module of the two upper adjacent matrix entries over
\(G_{\QQ_\ell}\), with its polarized action. Put
\(e=\dim_k H^2(\QQ_\ell,W_1)\). The full framed local ordinary
ring admits a presentation with
\[
    (1,\,2,\,3,\,4+e,\,2)
\]
variables of weights \((-2,-1,0,1,2)\) and at most \(e\)
relations, all of weight one. In the full original marked
global presentation this adds \(e\) to the variable count
and to the relation allowance at weight one in
Proposition~\ref{prop:originalchart425}.
\end{lemma}
\begin{proof}
The three distinct residual eigenvalues of \(g_*\) give the
unique ordinary flag by Lemma~\ref{lem:flag69}. Use its lower
unitriangular coordinates and put it in standard position.
These are three free coordinates of weights \(-2,-1,-1\).
The diagonal characters on inertia are fixed. The diagonal
variation is determined by its unramified Frobenius value
in the signed diagonal group, hence is formally smooth in
three coordinates of weight zero. Work relative to this
diagonal ring. The remaining matrices are upper unitriangular
corrections. At the residual diagonal their Lie module is
\(W_1\oplus W_2\), with \(W_2\) the upper long entry.
On the Tate dual of \(W_2\), inertia has exponent \(1-2r\).
It is nonzero modulo the even integer \(n\), so
\(H^2(\QQ_\ell,W_2)=0\) by local duality.

With the diagonal held fixed the framed first-order lifts
on each weight form \(Z^1(\QQ_\ell,W_i)\). The local Euler
characteristic gives
\(\dim Z^1=2\dim W_i+\dim H^2\), after adding the coboundary
dimension \(\dim W_i-\dim H^0\) to \(\dim H^1\).
Thus the two further variable counts are \(4+e\) and two.
For a small coefficient extension with square-zero kernel \(J'\)
annihilated by the maximal ideal, fix the lifted diagonal
and lift the upper entries as continuous cochains. Their
multiplication defect is a two-cocycle in
\((W_1\oplus W_2)\otimes_k J'\). Changing the entry lifts
changes it by a coboundary, and subtracting a correcting
cochain gives a lift exactly when its class vanishes. Keeping
the diagonal fixed preserves the required inertia characters.

Choose a power series presentation over the diagonal ring
using the computed relative tangent basis, with kernel \(J_0\)
and power series maximal ideal \(\mathfrak m\). The usual
small-extension argument in Theorem~\ref{thm:pres70}
injects the dual of \(J_0/\mathfrak mJ_0\) into this obstruction
space. Indeed a functional killing its component of the
universal obstruction allows the corresponding pushed-out
extension to lift. The relative variables can then only change
by square-zero elements. Every relation has zero relative
linear term at the residual point, so these changes cannot
alter its value in that kernel. The functional must therefore
vanish on every relation. One may take sufficiently high
Artinian truncations as in that theorem. This construction
is equivariant under the diagonal torus. Taking the weight
parts of lifts of a basis, followed by Nakayama's lemma,
gives at most \(e\) generating relations, all of weight one.

In the relative global calculation of Lemma~\ref{lem:rel70}
the local lift at \(\ell\) is given. Its ring need not be
smooth for that calculation: the obstruction is the defect
of a global lift with that specified local restriction.
The possible dual invariants in degree three have inertia
exponents \(1-ar\) with \(a=0,\pm1,\pm2\). Only \(a=\pm1\)
could contribute, and then \(r\equiv\pm1\bmod n\) is odd.
The residual characters of the corresponding dual pair,
after cancellation of the cyclotomic factor, are
\(\bar\alpha^{\pm1}\) and \(\bar\eta\bar\alpha^{\pm1}\).
The first is odd and the second ramifies at \(s\), so neither
has global invariants. The relative counts and the eight
marks in Theorem~\ref{thm:pres70} therefore apply as before.
Adding the local presentation gives the asserted global counts.
In the notation of Corollary~\ref{cor:present116} this replaces
\(R'\) by \(R'+e\), which is allowed in
Proposition~\ref{prop:localmult399}.
\end{proof}

\begin{proposition}[Crystallinity of modules for the ordinary matrix algebra]\leavevmode
\label{prop:filteredalgebra427}
Let \(r\ge2\). Write \(T\) for the original trace ring with
the full ordinary word and \(E_M\) for the \(T\)-span of the
group matrices on \(G_K\). Let \(T\to D'\) be a continuous
specialization to a finite local algebra over a finite extension
of \(\QQ_\ell\) with integral residue point. Every finite
\(E_M\otimes_T D'\)-module, as a coefficient-field representation,
is crystalline at all places of \(K\) above \(\ell\).
\end{proposition}
\begin{proof}
As in Lemma~\ref{lem:torsioncrys375}, \(E_M\) is a finite
\(T\)-module by finiteness of the Cayley--Hamilton algebra.
Its action and the action of the Galois group on any such
module are continuous. Fix a place above \(\ell\). Inside
\(E_M\) use the polynomial projectors \(e_0,e_1,e_2\) of
the local inertia element \(g_*\) and the three fixed
eigenvalues. They lie in \(E_M\), even if \(D'\) has
nilpotents. For each local group element \(g\) the following
identities hold in the original matrix algebra:
\[
  e_iM(g)e_j=0\quad(i>j),\qquad
  e_iM(g)e_i=\xi_i(g)e_i,\quad
  \xi_i(g)=t(e_iM(g)).
\]
Indeed the eigenlines of \(g_*\), in their ordinary order,
split the ordinary filtration. Each initial sum of these
lines is stable under the local group by uniqueness of the
flag. On each successive quotient the group acts by the
corresponding diagonal character, which gives both identities
in this eigenbasis. At a conjugate place use the reflected
dual flag in its ordinary order. The characters \(\xi_i\)
take values in \(T\), and restrict on inertia to
\(\kappa_\ell^{r},1,\kappa_\ell^{-r}\), respectively.

On a finite module \(U\) over the specialized algebra, the
partial sums of \(e_iU\) are therefore stable. Each successive
quotient, after removing its indicated scalar Tate twist,
is an unramified finite-dimensional representation, including
any unramified Jordan blocks. Such a representation is
crystalline of weight zero. Suppose an initial part \(U'\)
of the filtration is crystalline, and \(Q'\) is the next
quotient. For \(V'=\operatorname{Hom}(Q',U')\), all the
cyclotomic exponents are at least \(r\). This module is
crystalline, and \(H^0(K_v,(V')^\vee(1))=0\), since a
trivial subrepresentation would have weight zero and its
weights are at most \(1-r<0\) in this exponent convention.
The Bloch--Kato dimension formula gives
\(\dim H^1_f(K_v,V')=\dim H^0(K_v,V')+[K_v:\QQ_\ell]\dim V'\);
all its exponents are positive. The local Euler characteristic
and local duality give the same dimension for \(H^1(K_v,V')\).
Every extension of \(Q'\) by \(U'\) is consequently crystalline.
Induction proves the result. In particular this computation
applies to the trace realization, its dual-number lifts, and
extensions between two distinct simple constituents of the
finite algebra.~\cite{XI-VIII-VI-bk117}
\end{proof}

\begin{lemma}[An ordinary flag excludes a transitive permutation of lines]\leavevmode
\label{lem:monomialflag427}
Consider the normalized plane or the normalized irreducible
three-dimensional representation in the curve arguments of
Lemmas~\ref{lem:planecurve417} and~\ref{lem:connectedthree418},
or at the classical points in the corresponding arguments.
Suppose the residual inertia character \(\bar\Psi\) has order
greater than two. If the local ordinary characters have pairwise
ratios of infinite order on a common inertia subgroup, the
global group cannot permute transitively a decomposition of
the representation into lines.
\end{lemma}
\begin{proof}
In a stable lattice the subgroup \(\ker\bar\Psi\) has pro-\(\ell\)
image, because its reduction is unipotent after a triangular
residual change of basis. A permutation action on two or three
lines is continuous: the stabilizers of the lines have closed
finite index. Its order is prime to \(\ell\), so this action
factors through the cyclic image of \(\bar\Psi\). In the
transitive case the permutation group has order \(b\), where
\(b\) is the dimension two or three. Giving its generator a
primitive \(b\)-th root of unity as value defines a character
\(\delta\). In a basis of the lines numbered in cyclic order,
conjugation by the diagonal with successive powers of that
root gives the self-twist by \(\delta\) or its inverse.

On the specified compact inertia subgroup choose an element
whose local diagonal ratios all have nontrivial sixth powers.
Each ratio to the sixth power has infinite image. Its kernel
has Haar measure zero since it has arbitrarily many disjoint
translates of equal measure. A finite union of these kernels
therefore omits an element. The sixth power of the chosen
matrix fixes every global line and has distinct eigenvalues.
The lines of the decomposition are exactly its eigenlines.
Every step of the ordinary flag, being stable under this
matrix, is spanned by a subset of them. A permutation
preserving every step of a complete coordinate flag fixes
each line individually. The local decomposition group thus
acts trivially on the lines, so \(\delta\) is trivial there.

The reduction \(\bar\delta\) still has order \(b\).
For the plane it must interchange the two global residual
characters \(\bar\Psi,\bar\Psi^{-1}\), forcing
\(\bar\delta=\bar\Psi^{2}\) or \(\bar\Psi^{-2}\).
For the triple, the orbit of the trivial character under
multiplication has three elements and is the residual
triple, so \(\bar\delta=\bar\Psi\) or \(\bar\Psi^{-1}\).
Each conclusion is nontrivial on the local inertia group
under the stated order hypothesis. This contradicts local
triviality of \(\delta\).
\end{proof}

\begin{theorem}[The length bound at every twist with distinct inertia]\leavevmode
\label{thm:adjacentweights427}
For every \(x\) in the same sufficiently distant interval,
\[
 \sum_{\substack{\ell\ge11\ {\rm prime}\\
           \ell\nmid suN_*[E_*:\QQ]|\mathrm{Cl}(K)|}}
   (\log\ell)
   \sum_{\substack{1\le r<t(x)\\ \ell-1\nmid 2r}}
         \operatorname{length}_{\ZZ_\ell}
           \bigl(\mathcal S_{\ell,r}(E_x)^{c=(-1)^r}\bigr)
       \ \le\ 2\log D(V).
\]
For each one of these primes the inner sum is at most
\(2v_\ell(D(V))\).
\end{theorem}
\begin{proof}
First obtain the upper-field Hecke support at the residue
\(j\in\{1,\dots,n-1\}\) of \(r\). The hypothesis
\(n\nmid 2j\) ensures that \(\bar\omega^j,1,\bar\omega^{-j}\)
are distinct at every coefficient place of the layer. In
Lemma~\ref{lem:baseweight426} there are three issues when the
stronger restrictions on \(j\) are removed.

First, for \(j=n-1\) the starting Eisenstein series in the seed
argument has weight \(k=j+1=\ell-1\). Put \(\chi=\eta\alpha\)
and let \(D_0\) be its primitive conductor. In the expansion of
\[
 B_{k,\chi}=D_0^{\,k-1}\sum_{a=1}^{D_0}
                            \chi(a)B_k(a/D_0)
\]
the only Bernoulli number through degree \(k\) whose denominator
can contain \(\ell\) is \(B_k\). Its coefficient
\(D_0^{\,k-1}\sum_a\chi(a)\) is zero because \(\chi\)
ramifies at \(s\). All remaining terms are integral at \(\ell\),
and \(2k\) is a unit. Thus the constant term is integral,
which supplies the required congruence at this weight in
Lemma~\ref{lem:cusp361}. The ordinary unit root and the
prescribed tame characters follow as in that lemma.

Second, the Tate dual of \(\bar\Psi^{-1}\eta\) can have trivial
inertia at \(\ell\) when \(j=n-1\). Its global character then
equals \(\bar\alpha\bar\eta\), nontrivial at \(s\). For
\(\bar\Psi^{-2}\eta\) the dual inertia exponent \(1+2j\)
is odd and nonzero modulo \(n\). The two original characters
are nontrivial at \(\ell\), by \(n\nmid2j\). These are the
global invariant vanishings and the zero local \(H^0\) at
\(\ell\) entering the strict-class Euler count in
Proposition~\ref{prop:strictclasses365}. Its two strict classes
on the auxiliary layer are therefore available.

Third, use Lemma~\ref{lem:monomialflag427} in place of the
order-three and order-four exclusions in the argument of
Lemma~\ref{lem:baseweight426}. For a separated curve in
residue characteristic the ratios on inertia in the cyclotomic
kernel reduce to one. Each nontrivial such ratio has infinite
image since the principal units of a Laurent series field in
characteristic \(\ell\) have no torsion. The initial curve
of the auxiliary ordinary family has the same property by its
chosen principal-unit direction. At classical specializations
with distinct ordinary exponents, on an open subgroup of
inertia the ratios are nonzero powers of the cyclotomic
character in the parallel-weight cases. More generally,
distinct labeled weights give infinite character ratios,
since finite characters have weight zero at every embedding.
For the plane its local flag is obtained by intersecting the
full stable flag with the plane and discarding repeated steps.
Thus in each use of these exclusions the lemma applies. A
scalar action on an open subgroup is excluded already by the
same infinite ratios. The original two-dimensional seed with
residual characters \(1,\bar\eta\bar\Psi\) also has no quadratic
self-twist, since a swap would force \((\bar\omega^j)^2=1\).
This verifies the upper-field propagation, including at order
three or four on inertia. The lower-root dimension estimates
and the local degree of the auxiliary layer are unchanged in
the argument of Lemma~\ref{lem:baseweight426}.

We specialize at the actual weight \(r\). The trace minors
and ordered diagonal characters in Proposition~\ref{prop:weightshift426}
require just the three distinct residual inertia eigenvalues.
At the reducible special-fiber points, Pan's ordinary finiteness
argument uses the nontriviality of \(\bar\Psi^{-2}\) on inertia,
which is precisely \(n\nmid2r\). Thus finiteness of the entire
original trace ring follows from the two-quotient argument
in Theorem~\ref{thm:originalfinite424}. For \(r\ge2\)
Proposition~\ref{prop:filteredalgebra427} gives crystallinity
for all the modules used in the generic-fiber arguments.
The old inertia is finite of prime-to-\(\ell\) order by the
original tame condition. Hence the descent and the level
calculation of Proposition~\ref{prop:alloriginalpoints424} give
exactly the same lists at \(r\). In the non-CM and enormous-image
arguments for the plane and for an irreducible original point,
the corresponding local ordinary exponents are distinct, so
Lemma~\ref{lem:monomialflag427} again excludes the permutation
case. The adjoint Selmer vanishings and the two extension
vanishings used in Lemma~\ref{lem:planeunits379} therefore
apply. In particular the generic trace algebra is \'{e}tale
and the graded multiplicities are those of
Lemma~\ref{lem:markmultiplicity425}. At \(r=1\) these conclusions
already follow from Proposition~\ref{prop:weightshift426}.

For \(r\ge2\) the local presentation of
Lemma~\ref{lem:ordinaryrelations427} has exactly the allowed
surplus of variables over relations in each weight. The
identification of the two negative cotangents in
Proposition~\ref{prop:originalchart425} remains valid for the
full ordinary word. Indeed a negative root class, after its
local coboundary adjustment at \(g_*\), has standard ordered
flag and zero lower entries on inertia. Conversely any local
coboundary is a conjugate of the ordinary split lift. The
unramified and marked conditions away from this place, the
fixed section and the trace description of \(B_0\) follow from
the same identities. In particular the two negative cotangents
are torsion. Apply Proposition~\ref{prop:localmult399} with
the Frobenius difference at \(u\). For \(r=1\) apply
Theorem~\ref{thm:strictlarge425} with the local chart already
checked in Proposition~\ref{prop:weightshift426} and
Theorem~\ref{thm:regularlength426}. The sum of the valuations
over nonsplit points in all the resulting bounds is at most
\(v_\ell(D(V))\) for each prime: the labels use the actual
\(\alpha,r\), their multiplicities are one, and \(u\) is a
unit there. The character descent and summation of
Theorem~\ref{thm:regularlength426} prove both assertions.
\end{proof}

\begin{corollary}[The two remaining inertia residues]\leavevmode
\label{cor:smallorder427}
At primes \(\ell\le P(x)\) satisfying the prime conditions
of the theorem, the remaining \(R_{\ell,r}\) in
Section~\ref{sec:missing} can be nonzero only for
\[
                 r\equiv 0\quad\text{or}\quad
                 r\equiv(\ell-1)/2\pmod{\ell-1}.
\]
At these twists the summand for the trivial character of
\(E_x\), when it has the complementary sign, is zero.
\end{corollary}
\begin{proof}
The high-power truncation kills each group bounded by the
theorem, since \(2v_\ell(D(V))<b_0\) uniformly for sufficiently
large \(V\), as in Corollary~\ref{cor:remainingresidues426}.
Precisely the two displayed residues satisfy \(\ell-1\mid2r\).
The trivial character has the specified sign only for even
\(r\). Its strict summand descends integrally to \(\QQ\)
since \(\ell\nmid[E_x:\QQ]\). At this even twist its dual is
the \(\ell\)-primary part of \(K_{2r}(\ZZ)\), which vanishes
under this divisibility by Proposition~\ref{prop:qbound102}.
\end{proof}


\subsection{The two moving ramified coefficient primes}
\label{sec:movingcoefficient428}

Fix the polynomial-conductor family of
Theorem~\ref{thm:conductor232}. Write \(\ell=q\) or \(\ell=d_0\),
\(n=\ell-1\), and let \(\tau\) be the quadratic Dirichlet
character of \(\QQ(\sqrt{-\ell})\). We keep the same split test
prime \(u\), the same quadratic field \(K=\QQ(\sqrt{-s})\), and
the integer \(D(V)\) of Theorem~\ref{thm:originalorder424}.
The finite character of an ordinary Galois representation can
now be ramified at the coefficient prime itself. On both outer
lines its finite inertia character then equals \(\tau\).

\begin{lemma}[An unramified character giving the same residual problem]\leavevmode
\label{lem:quadraticseed428}
Let \(\alpha\) be a character of \(E_x\) of sign \((-1)^r\),
with \(1\le r<t(x)\). There are \(b\in\{0,1\}\) and a character
\(\alpha_0\) of conductor prime to \(\ell\) such that
\(\alpha=\alpha_0\tau^b\). Put \(j=r+bn/2\). Then
\[
       \bar\alpha\,\bar\kappa_\ell^r
        =\bar\alpha_0\,\bar\kappa_\ell^j,\qquad
       \alpha_0(c)=(-1)^j .
\]
There is an auxiliary layer \(L=KF_a^+\) on which all the
geometrically irreducible primes of the corresponding ordinary
determinant ring contain the Hecke kernel. The finite ordinary
diamonds on this ring agree with those prescribed by
\(\alpha\kappa_\ell^r\).
\end{lemma}
\begin{proof}
The construction of Theorem~\ref{thm:conductor232} gives
\(q,d_0>H^{401W}\), \(V^2\le H^{K_0}\) and \(W>K_0+2\).
Hence \(\ell-1>8t_*+8\) for sufficiently large \(H\), since
\(t_*\le H^{401(3K_0/4)}\). The primes \(s,u\), the class
number of \(K\), and the odd prime divisors of \([E_*:\QQ]\)
are fixed. Thus \(\ell\) divides none of them. In \(E_*\) the
only factor ramified at \(\ell\) is \(\QQ(\sqrt{-\ell})\).
Indeed all other factors in the displayed field of
Theorem~\ref{thm:conductor232} have conductor prime to \(\ell\).
Their compositum is disjoint from this ramified quadratic field,
by looking at inertia. Take \(b\) according to the restriction
of \(\alpha\) to this order-two inertia group. Then \(\alpha_0\)
factors through the other factors, whose compositum we denote
by \(E^\circ\). Write \(N^\circ=N_*/\ell\). Since the level in
that theorem contains each moving prime exactly once, the
conductor of \(E^\circ\) divides \(N^\circ\).

Both moving primes are \(3\bmod 4\). The unique quadratic
subfield of \(\QQ(\mu_\ell)\) is therefore
\(\QQ(\sqrt{-\ell})\), and \(\bar\tau=\bar\kappa_\ell^{n/2}\)
on \(G_\QQ\). Its value at complex conjugation is \(-1\)
and \(n/2\) is odd. This proves the two identities. For
\(b=0\) the size inequality gives \(j=1\) or
\(2\le j\le n-2\) with \(n/\gcd(n,j)>4\).
For \(b=1\) the latter alternative always holds. In fact
\(hj\equiv hr\) or \(n/2+hr\bmod n\) for \(1\le h\le4\),
and \(0<hr<n/2\). Neither residue is zero.

Apply the construction in Lemma~\ref{lem:baseweight426} to
\(\alpha_0,j,E^\circ,N^\circ\). Its proof uses the conductor
divisibility, \(\ell\nmid sN^\circ[E^\circ:\QQ]\), the
specified parity, and the ramification of \(\eta\) at \(s\).
They hold because \(E^\circ\subset E_*\) and \(\ell\nmid N^\circ\).
In particular \(L\) is disjoint from \(E^\circ(\mu_\ell)\),
which already contains \(E_*\). The auxiliary congruences
include splitting at \(\ell\) in \(L/F_a^+\), and the
cyclotomic inertia image on \(G_L\) still has order \(n\).
The class-rank bound and the choice of the depth use the field
\(L E^\circ(\mu_\ell)\). The old types away from \(\ell\)
are the same for \(\alpha_0\) and \(\alpha\), because \(\tau\)
is unramified there. Use the integral modules of
Lemma~\ref{lem:typelattice427} for those types. The Hecke
support conclusion of Lemma~\ref{lem:baseweight426} follows.

On the prime-to-\(\ell\) part of the local units, finite
character values of order prime to \(\ell\) are determined
by their reductions. Thus the residual identity gives equality
of the prescribed finite diamonds of the two ordered triples,
including after restriction to \(L\). On principal units
\(\tau\) is trivial, and the actual exponent \(r\) specifies
the arithmetic specialization of the variable weight ring.
The residual equality holds on whole decomposition groups,
so it also gives the same order of the residual diagonal
characters there.
\end{proof}

\begin{proposition}[The local ordinary modules of quadratic type]\leavevmode
\label{prop:tamecrystal428}
Suppose \(b=1\). Use the full ordered ordinary word in the marked
problem, with inertia characters
\(\tau\kappa_\ell^r,1,\tau\kappa_\ell^{-r}\).
The framed local ring over \(O\) is smooth with coordinate
counts \(1,2,3,4,2\) in weights \(-2,-1,0,1,2\).
At each \(v\mid\ell\) of \(K\), every finite module of the
specialized matrix algebra of Proposition~\ref{prop:filteredalgebra427}
becomes crystalline over \(K'_v=K_v(\sqrt{-\ell})\).
Extensions of such modules have the Bloch--Kato finite local
condition whenever the extension itself is a module for that algebra.
\end{proposition}
\begin{proof}
The separating element exists by the lemma. On the two
adjacent upper entries the residual Tate-dual inertia exponent
is \(1-j\), and on the long upper entry it is \(1-2j\).
Both are nonzero modulo \(n\). Hence the obstruction group
for lifting the upper entries with a given unramified diagonal
variation is zero by local duality. The flag and diagonal
are smooth in the three lower and three diagonal coordinates,
respectively. The cocycle calculation of
Lemma~\ref{lem:ordinaryrelations427} gives four adjacent
and two long upper variables. It uses only the residual
actions and works also when the actual exponent is one.

The ordered projectors of the separating element give the
filtration on every finite matrix module exactly as in
Proposition~\ref{prop:filteredalgebra427}. On a graded piece
the group acts by the prescribed inertial character times an
unramified representation, which can have unipotent Frobenius.
This piece is de Rham over \(K_v\) and crystalline over
\(K'_v\). We prove the latter statement for successive
extensions by induction. Let \(A\) be an initial part already
crystalline over \(K'_v\) and \(B\) the next graded piece.
The module \(V'=\operatorname{Hom}(B,A)\) is de Rham over
\(K_v\) by descent of de Rham periods, and every cyclotomic
exponent on it is strictly positive.

We claim \(H^0(K_v,(V')^\vee(1))=0\). For exponents greater
than one this follows from the nonzero Hodge--Tate weights
on the corresponding graded part of the dual twist. An
exponent equal to one requires \(r=1\) and two adjacent
steps. The corresponding finite inertia character of the
dual twist is \(\tau\), even with the unramified coefficient
blocks. Inertia of \(K_v\), which is unramified over
\(\QQ_\ell\), contains an element acting there by \(-1\).
Thus this graded part also has zero invariants. An invariant
in a filtered space with zero invariants on all graded pieces
is zero, proving the claim.

The Bloch--Kato dimension formula for de Rham coefficients
gives \(\dim H_f^1=\dim H^0+[K_v:\QQ_\ell]\dim V'\),
since all exponents of \(V'\) are positive. The local Euler
characteristic and the claim give the same value for
\(\dim H^1\). Thus the extension class restricts to
\(H_f^1(K'_v,V')\) by functoriality. Its end terms are
crystalline there, so its middle term is crystalline there
as well. This proves the induction, including the first
twist. For an arbitrary short exact sequence of finite
matrix modules its restriction to \(K'_v\) now splits after
tensoring with \(B_{\rm cris}\). Indeed the crystalline
period spaces give an exact sequence of vector spaces over
the maximal unramified period field, which admits a linear
splitting. Extension to \(B_{\rm cris}\) gives an equivariant
splitting. The class in the \(B_{\rm cris}\)-cohomology of
the Hom module over \(K_v\) is therefore zero by restriction
and corestriction, whose composite multiplies by two.
This is the finite local condition over \(K_v\).~\cite{XI-VIII-VI-bk117}
\end{proof}

\begin{proposition}[The original trace ring at a moving prime]\leavevmode
\label{prop:movinglabels428}
For each \(\alpha,r\) above, the fixed-weight trace ring \(T\)
is finite over \(O\) with geometrically reduced generic fiber.
Its nonsplit geometric traces occur with the original level
and weight in the lists defining \(D(V)\).
\end{proposition}
\begin{proof}
For \(b=0\), the original problem can be imposed with
\(E^\circ,N^\circ\). In the constructions of
Propositions~\ref{prop:originalchart425} and~\ref{prop:weightshift426}
the inertia kernel away from the coefficient prime is the
kernel of the prescribed character together with the quadratic
sign. At the coefficient prime it is the ordinary condition
alone that is imposed. Thus replacing \(N_*\) by \(N^\circ\)
here gives precisely the required trace calculation. The
resulting levels divide the original ones. We treat \(b=1\)
in detail.

Restrict the original integral matrices to \(G_L\). The
projector calculation of Proposition~\ref{prop:orighecke374}
recovers each diagonal character from a trace. Lemma~\ref{lem:quadraticseed428}
puts these in the same ordered weight component. On inertia
away from \(\ell\), \(\alpha\) and \(\alpha_0\) agree, so
restriction factors through the prescribed ramification
quotient. Put \(G=G_K\), \(H=G_L\), and take a separating
inertia element in \(H\). Its spectral projectors \(e_i\)
belong to the restricted matrix span. For \(g\in G\),
normality puts \(M(g)e_iM(g)^{-1}\) there as well. Its
\((k,i)\)-block is \(e_kM(g)e_i\) multiplied by
\(\operatorname{tr}(e_iM(g)^{-1})\), a unit since its
residue is a diagonal character value. This proves equality
of the two matrix spans over \(T\) by summing the blocks.
Thus all trace-minor ideals on \(G\) are extended from \(H\)
by multilinearity. In particular the four-by-four minors
generate the original augmentation ideal by
Proposition~\ref{prop:origideal373}. That proposition uses
the separated ordered diagonal and the splitting of all
allowed character lifts by the class-number condition. In
this application each ordered character divided by its split
value is unramified at \(\ell\), because their prescribed
inertia characters agree there, and at the other places by
the prime-to-\(\ell\) inertia condition. Its residual value
is one and \(\ell\nmid|\mathrm{Cl}(K)|\), as required.
The closed restricted trace ring then equals \(T\) by
Theorem~\ref{thm:restrictedtrace374}, giving the surjection
of Proposition~\ref{prop:orighecke374} at the actual weight.

At an irreducible special-fiber trace the restriction is
therefore in Hecke support. For a reducible trace the
special-fiber descent in Theorem~\ref{thm:reducibleminors378}
uses the distinct residual eigenlines and the class number
of \(K\) prime to \(\ell\). On the plane the twist by
\(\bar\Psi^{-1}\) makes one ordinary inertia character trivial.
The determinant lift \(\eta\Psi^{-2}\) is odd and de Rham,
and its inertia reduction is nontrivial since \(n\nmid2r\).
Thus the ordinary pseudodeformation finiteness theorem used
there applies with the ramification set dividing \(sN^\circ\ell\).
These two cases give finiteness of \(T\) by the special-fiber
argument in Theorem~\ref{thm:originalfinite424}.~\cite{XI-VIII-VI-pan40}

At an irreducible characteristic-zero point the upper ordinary
specialization is classical. In Geraghty's control theorem the
weight consists of a dominant algebraic part and a finite
diagonal character. Here the first is the weight belonging
to \(r\), and on local units the second has components
\(\tau,1,\tau\), restricted by the local norm on the upper
field. Their product with the algebraic character has exactly
the prescribed tame diamonds. Lemma~2.6.4 of that theorem
thus places the specialization on the finite-level classical
module, with any nilpotent kernel of the specialization map
killed by evaluation in a field. The old type coefficients
commute with this control map by Lemma~\ref{lem:typelattice427}.
The restricted representation is irreducible, so the transfer
is cuspidal. The degree \([L:K]\) is prime to three. The
determinant-one cyclic descent, including its uniqueness, is
that of Proposition~\ref{prop:alloriginalpoints424}. The
determinant on \(K\) is indeed one since the product of the
three prescribed inertia characters is one at every finite
place and \(\ell\nmid|\mathrm{Cl}(K)|\).~\cite{XI-VIII-geraghty363,XI-VIII-I-ac123}

At \(v\mid\ell\) the local parameter of the descent has zero
monodromy by Proposition~\ref{prop:tamecrystal428}. On its
ordinary filtration the inertial types of the graded pieces
are \(\tau,1,\tau\). Exactness of the potentially semistable
period functor and semisimplicity of finite inertia in
characteristic zero give this same inertial type on the
Frobenius-semisimplified Weil--Deligne parameter. Frobenius
preserves each inertia eigenspace. After semisimplification
it can be diagonalized on each, giving three Weil characters
trivial on principal local units. Local--global compatibility
at the coefficient prime applies since the rank is odd and the
weight is regular. This places the local component in the
same principal level as at any finite place of exponent one
in the argument of Theorem~\ref{thm:auxdescent361}. At the
other finite places that argument applies verbatim. The
weight at infinity is determined by \(r\). Thus the rank-three
label is at level \(N_*\).~\cite{XI-VIII-lgp361,XI-VIII-V-shin}

At a line and a plane the line is trivial by uniqueness of
the middle residual character and the class-number argument
of Theorem~\ref{thm:alllabels375}. The symmetric plane
adjunction gives its odd extension \(W\) over \(\QQ\),
of determinant \(\eta\), as in Proposition~\ref{prop:odddescent376}.
It is de Rham with weights \(r,-r\), since this is true upon
finite local extension. Its modularity and the non-CM
assertion follow there by Pan's theorem and the order of
the residual inertia character, which exceeds four. The form
for \(W\kappa_\ell^r\alpha^{-1}\) has weight \(2r+1\) and
nebentypus \(\eta\alpha^{-2}\). Twisting its plane parameters
at \(\ell\) by \(\alpha^{-1}\) removes \(\tau\) from both
inertia characters. Monodromy is zero. Since \(K_v/\QQ_\ell\)
is unramified, the Weil inertia group for the rational
parameter is the same. The primitive level is consequently
the conductor of \(\eta\alpha^{-2}\), as in
Proposition~\ref{prop:planelabels379}. At \(u\) the trace
formula \(1+u^{-r}\alpha(u)a_u(f)\) uses the actual \(\alpha\).
~\cite{XI-VIII-VI-pan40,XI-VIII-VI-saito40}

Finally, all extensions of simple modules and dual-number
lifts arising from the marked algebra satisfy the finite
local condition by Proposition~\ref{prop:tamecrystal428}.
Away from \(\ell\) their prescribed inertia is finite of
order prime to \(\ell\). For a split triple the decreasing
extension groups vanish after restriction to \(K E_x\) by
the negative-twist argument of Theorem~\ref{thm:reduciblefiber376}.
For a plane the off-diagonal groups vanish by Kato's theorem
over the same quadratic \(K\) for the modular form attached
to \(W\kappa_\ell^r\), at the twist \(r\). Its restriction
after the dual and twist is the indicated plane. These are
the vanishings giving the product matrix algebra at the plane
in Lemma~\ref{lem:planeunits379}. At an irreducible trace
the ordinary exponents are distinct and its residual inertia
order exceeds four. The enormous-image argument of
Lemma~\ref{lem:cyclointegral377} and the polarized adjoint
Selmer vanishing of Newton and Thorne therefore apply. The
same adjoint vanishing for the non-CM modular plane gives
its trace rigidity. The tangent argument of
Corollary~\ref{cor:originaletale424} now makes every
geometric generic local ring a field.~\cite{XI-VIII-VI-kato120,XI-VIII-VI-nt119}
\end{proof}

\begin{theorem}[Full length at both moving coefficient primes]\leavevmode
\label{thm:movingprimes428}
For sufficiently distant intervals and every \(V\le x\le V^{3/2}\),
at each \(\ell\in\{q,d_0\}\) one has
\[
  \sum_{1\le r<t(x)}
       \operatorname{length}_{\ZZ_\ell}
          \bigl(\mathcal S_{\ell,r}(E_x)^{c=(-1)^r}\bigr)
            \ \le\ 2 v_\ell(D(V)).
\]
Consequently \(R_{\ell,r}=0\) at both moving primes for every
twist on the interval.
\end{theorem}
\begin{proof}
For \(b=1\) the local chart has the same five weight counts
as in Proposition~\ref{prop:originalchart425}. For the
relative global presentation, none of the dual inertia
exponents \(1,1\pm j,1\pm 2j\) is zero, since
\(j=n/2+r\) and \(0<4r<n/2\). The away Euler terms and the
polarization at infinity are unchanged. The fixed section
and the negative cotangent modules use the actual \(\Psi\)
in that proposition. Its proof uses the units
\(\bar\Psi(g_*)^{\pm1}-1\) and \(\bar\Psi(g_*)^{\pm2}-1\),
the diagonal projectors, and prime-to-\(\ell\) inertia away
from \(\ell\). Each holds here. At the coefficient place the
local lower-root coboundary is equivalent to the strict local
condition, by subtracting its value at the separating element
and using the ordinary flag. Restriction to \(K E_x\) still
has index prime to \(\ell\) and proves finiteness of the
negative cotangents. For \(b=0\) use the same presentation
at \(r\ge2\), and the first-twist chart of Theorem~\ref{thm:first72}
at \(r=1\), with the reduced tame modulus \(N^\circ\).

The proposition and the calculation at its plane points
give the multiplicities of Lemma~\ref{lem:markmultiplicity425}
in both cases. Apply Proposition~\ref{prop:localmult399}
and sum the resulting valuations at \(u\). The lists in
Theorem~\ref{thm:originalorder424} are indexed by the original
\((\alpha,r)\). Distinct trace characters for that pair give
distinct labels, since the embedded label determines all
good traces and these determine the trace character. The
factors clearing denominators are powers of \(u\ne\ell\).
Summing over the residue conjugates and the twists thus uses
at most \(v_\ell(D(V))\), with the factor two in the length
bound. The scalar character descent of
Theorem~\ref{thm:levelstrict345} is exact also at these
ramified primes since \(\ell\nmid[E_x:\QQ]\). In particular
local restriction and corestriction at \(\ell\) are injective
in the required direction because the local degree is a unit.
This proves the displayed inequality. The moving primes are
at most \(N_*\ll V^C\), hence eventually at most \(P(x)\).
Since \(2v_\ell(D(V))\ll V^C<b_0\), truncation kills all
their remaining summands.
\end{proof}

\begin{corollary}[Support of the remaining odd-prime tails]\leavevmode
\label{cor:unboundedprimes428}
There is a finite set of primes depending only on the fixed
data such that, outside it, a nonzero \(R_{\ell,r}\) with
\(\ell\le P(x)\) requires \(\ell-1\mid 2r\). In particular
\(\ell\le 2r+1\). Above \(P(x)\) the total complementary
length with weights \(\log\ell\) is \(O_{\psi,\varepsilon}(V^C)\)
as in Corollary~\ref{cor:periodlarge425}.
\end{corollary}
\begin{proof}
Include the primes below eleven and the prime factors of
\(su\,2F[K_\psi:\QQ]\,|\mathrm{Cl}(K)|\) in the finite set.
Outside it, the only prime divisors of \(N_*[E_*:\QQ]\)
are \(q,d_0\). The theorem treats both, and
Corollary~\ref{cor:smallorder427} treats the other primes
with \(\ell-1\nmid2r\). Divisibility by the positive integer
\(\ell-1\) gives the last inequality for any remaining
positive twist.
\end{proof}

\section{Colliding ordinary characters}\label{nav:chapter-coincident-characters}

\subsection{Frobenius at the two coincident inertia weights}
\label{sec:distinguishedfrobenius435}

Keep the interval and the fields of Theorem~\ref{thm:conductor232}.
Write \(K=\QQ(\sqrt{-s})\), write \(\eta_K\) for its quadratic
character, and take a prime \(\ell\ge 11\) prime to
\(2suN_*[E_*:\QQ]\,|\mathrm{Cl}(K)|\), with \(u\) as in
Theorem~\ref{thm:originalorder424}. Let \(\alpha\) be a character
of \(E_x\) with \(\alpha(c)=(-1)^r\), where
\(2\le r<t(x)\) and \(\ell-1\mid 2r\). Set
\(\Psi=\alpha\kappa_\ell^r\) and let \(O\) be an unramified
coefficient DVR containing the values of \(\alpha\), with residue
field \(k\). The two outer characters of the rank-three diagonal
\((\Psi,1,\Psi^{-1})\) coincide on residual inertia at \(\ell\).
For a substantial part of this family the ordered flags and the
reducible traces can still be determined at \(\ell\).

\begin{lemma}[A separating element on every auxiliary layer]\leavevmode
\label{lem:onefrobenius435}
Suppose \(\alpha(\ell)^4\ne1\). Fix a place above \(\ell\)
in an algebraic closure of \(\QQ\). There is an element \(h\)
in its decomposition group which belongs to \(G_K\) and to
every \(G_{KF_a^+}\) in Proposition~\ref{prop:auxfield361}
for the fixed auxiliary primes. It has \(\kappa_\ell(h)=1\),
and \(b=\bar\Psi(h)\), \(1\), \(b^{-1}\) are pairwise distinct.
For a fixed \(\ell\) one choice works for every indicated
\(\alpha,r,x\). Conjugate elements have the same three values.
\end{lemma}
\begin{proof}
The local cyclotomic extension of \(\QQ_\ell\) is totally
ramified. Multiply an arithmetic Frobenius lift by an inertia
element to obtain a lift \(\phi\) with \(\kappa_\ell(\phi)=1\).
Take \(h=\phi^{2\ell}\). The cyclotomic factor \(F_a\) of
\(F_a^+\) is fixed already by \(\phi\). On the other factors,
the unramified local degree divides \(2\ell\). Indeed the two
quadratic factors are unramified at \(\ell\) and the factor of
prime conductor \(w\) has degree \(\ell\). This also shows that
\(h\) fixes \(K\). We have \(b=\bar\alpha(\ell)^{2\ell}\).
Reduction preserves the order of \(\alpha(\ell)\), which is
prime to \(\ell\). Thus \(b^2\ne1\) is equivalent to
\(\alpha(\ell)^4\ne1\). All characters used here are global
abelian characters, which proves the last assertion.
\end{proof}

Use the marked matrix problem on \(G_K\) of
Section~\ref{sec:original68} with its symmetric adjunction and
the same inertia kernels away from \(\ell\). Use \(h\) for the
element in the adjacent-entry mark of Corollary~\ref{cor:dual68}.
At places above
\(\ell\) prescribe a flag stable under the full decomposition
group. Its ordered inertia characters are \(\Psi,1,\Psi^{-1}\)
and its residual graded characters on the full local group
are \(\bar\Psi,1,\bar\Psi^{-1}\). At a conjugate place use
the dual flag in the reversed order. Write \(M\) for the
matrices on \(G_K\). The adjunction in the marked coordinates is
\(M(cg^{-1}c)=J_3 M(g)^{\mathsf T}J_3\), where \(J_3\) is
the anti-diagonal permutation matrix. We describe the local
conditions before passing to traces.

\begin{proposition}[Integral flags and the strict lower entries]\leavevmode
\label{prop:frobeniusflag435}
Under the hypothesis of Lemma~\ref{lem:onefrobenius435},
the ordered local condition is given by closed matrix equations
over \(O\). On its solutions the flag is unique. For a
square-zero deformation supported in one lower matrix entry at
the split diagonal, these equations hold exactly when the local
extension class is zero. This assertion holds over every finite
length \(O\)-module of coefficients.
\end{proposition}
\begin{proof}
Work first over a complete Noetherian local \(O\)-algebra
with residue field \(k\). Index the three rows by \(0,1,2\).
The characteristic polynomial of \(M(h)\) has three distinct
residual roots. Hensel's lemma gives roots \(\lambda_i\) with
the prescribed reductions. Their differences are units, and
Cayley--Hamilton gives three mutually orthogonal projectors
\[
   e_i=\prod_{j\ne i}
       \frac{M(h)-\lambda_j I}{\lambda_i-\lambda_j},
       \qquad \sum_{i=0}^2 e_i=I .
\]
Each has free image of rank one, by reduction and Nakayama's
lemma. The roots lie already in the completed algebra of
characteristic-polynomial coefficients. An ordered stable
flag, if it exists, is split by these projectors. Its first
line and first plane must be the first eigensummand and the
sum of the first two, since the residual ranks prescribe
which summands occur. It follows that the equations are
\[
 \begin{aligned}
    e_i M(g)e_j&=0 &&(i>j,\ g\in G_{K_v}),\\
    e_i M(\tau)e_i&=\Psi(\tau)^{1-i}e_i
                        &&(\tau\in I_v,\ 0\le i\le2).
 \end{aligned}
\]
Conversely these equations give the indicated stable flag and
inertia action. The construction uses only units as
denominators, commutes with change of coefficients, and gives
a closed ideal in the universal complete local ring.

For a lower entry use a square-zero ideal \(N\) and write its
matrix as \((1+z(g)E_{ij})d_3(g)\), \(i>j\), with \(d_3\)
the split diagonal. Its character ratio \(\theta_{ij}\)
satisfies \(\theta_{ij}(h)-1\in O^\times\) by the lemma.
Conjugate by a unique matrix \(1+yE_{ij}\), \(y\in N\), to
make \(z(h)=0\). Conjugation adds a coboundary and gives the
standard projectors at \(h\). The first row of equations now
forces \(z(g)=0\) on the entire local group. A local
coboundary conversely gives a conjugate of the split flag.
This works with \(O\oplus N\) for every such \(N\); the
normalization divides by the displayed unit. For a class
defined over \(\QQ_\ell\), restriction to \(K_v\) also
detects zero, since corestriction composed with restriction
is multiplication by \([K_v:\QQ_\ell]\in\{1,2\}\).
\end{proof}

Let \(T^{\rm ord}\) denote the completed trace algebra of this
fixed-weight problem. More explicitly, take the image of the
universal determinant ring for \(G_K\), unramified outside the
given finite set, in the framed marked matrix ring with the
displayed local equations. Since \(2,3\in O^\times\), the
coefficients of characteristic polynomials are polynomials
in traces of powers. Denote the group trace by \(t\). For
\(m\ge1\) let \(I_m\) be generated by
\(\det(t(g_i h_j))_{1\le i,j\le m}\), \(g_i,h_j\in G_K\).

\begin{theorem}[Finiteness on the reducible locus at coincident inertia weights]\leavevmode
\label{thm:reduciblecollision435}
If \(\alpha(\ell)^4\ne1\), then \(T^{\rm ord}/I_m\) is a
finite \(O\)-module for \(6\le m\le9\). At any geometric
field point its semisimple representation is reducible
exactly when \(I_m\) vanishes there. For every continuous local
map from the quotient to \(k'[[z]]\) with \(k'/k\) finite,
all group traces belong to \(k'\).
\end{theorem}
\begin{proof}
The determinant ring is complete Noetherian, and its universal
Cayley--Hamilton algebra is finite over it. Here the
required finiteness of residual first cohomology follows by
restriction to a finite extension killing the residual
representation. For the trivial residual module one may also
adjoin \(\mu_\ell\); Kummer theory then gives a bound from the
finite \(\ell\)-torsion of the class group and the finitely
generated group of units outside the given finite set.
Inflation--restriction adds only finite groups. The determinant
and Cayley--Hamilton finiteness theorems therefore apply just
as in Lemma~\ref{lem:torsioncrys375}. The trace ring, as an
image of a compact complete ring with finite residue field,
is a complete Noetherian quotient. The span of the universal
group matrices over it is finite by the same Cayley--Hamilton
theorem.~\cite{XI-VIII-chenevier364}

Let \(T^{\rm ord}\to k'[[z]]\) be a continuous local map with
\(k'/k\) finite. Take the semisimple representation of its
determinant over an algebraic closure of \(k'((z))\).
It factors through the specialized matrix algebra. Indeed
each matrix relation gives zero trace after multiplication
by every group matrix. In a semisimple realization these
matrices span the product of full matrix algebras on its
distinct constituents, and the trace pairing on each factor
is the usual pairing times its multiplicity. All these
multiplicities are units since \(\ell>3\). The pairing is
thus nondegenerate. The realization is defined over a finite
extension of \(k'((z))\), since the algebra through which it
factors is finite dimensional. It is continuous there by
finiteness of the integral matrix algebra, and its compact
image preserves a lattice. Its semisimplified reduction on
such a lattice is \( \bar\Psi\oplus1\oplus\bar\Psi^{-1}\).
The projectors and both local matrix identities in the
proposition descend to this realization.

If the representation is a sum of lines, match each character
with its reduction at the element \(h\). At every place above
\(\ell\) the prescribed inertia character on that line in
equal characteristic is exactly its residual one. In fact
the three idempotents label the three local constituents, and
the distinct reductions at \(h\) determine the label, also
at a conjugate place. Away from \(\ell\) the prescribed
inertia quotient has order prime to \(\ell\). Dividing the
character by the constant lift of its residual character
therefore yields an everywhere unramified character into
principal units. Such a character factors through the class
group of the imaginary quadratic field \(K\); its image is
trivial since \(\ell\nmid|\mathrm{Cl}(K)|\). Thus every line
is constant. The same argument applied to the determinant
shows that the determinant of the entire realization is one:
its inertia product at \(\ell\) and its residual character are
both trivial.

In the other reducible case write the representation as a
line \(\lambda\) and an absolutely irreducible plane \(P\).
By conjugate duality the unique line has
\(\lambda^c=\lambda^{-1}\). Each possible reduction extends
to \(G_\QQ\); the two outer characters have square different
from one at \(h\). Hence \(\bar\lambda=1\). The ordinary
projectors assign this line the middle inertia character,
including when inertia alone has three equal reductions.
The unramified character argument just given shows
\(\lambda=1\), so that \(\det P=1\).

We recall the odd descent of this plane to make the use of
Pan's theorem precise. The adjunction descends to the
semisimple group span: its kernel there is the radical of
the trace pairing, which adjunction preserves. It fixes the
plane factor, the unique simple factor of dimension four.
On it the adjunction is adjoint to a symmetric or an
alternating form, because an involution of a matrix algebra
fixing scalars is transpose followed by an inner
automorphism. The alternating possibility is excluded by
the integral identity \(2XYX=t(XY)X\) for skew matrices in
this rank-three symmetric adjunction, proved in
Proposition~\ref{prop:odddescent376}. Skew elements in the
quotient lift by taking half the difference with the
adjoint, and their line component is zero. For alternating
adjunction the traceless matrices \(X=\operatorname{diag}(1,-1)\)
and \(Y=E_{12}+E_{21}\) are skew, whereas \(XYX=-Y\) and
\(\operatorname{tr}(XY)=0\).

Write \(S\) for the symmetric form on the plane and \(J_0\)
for its alternating determinant form. Then \(S^{-1}J_0\)
intertwines \(P\) with its conjugate by \(c\): the adjunction
sends \(P(g)\) to \(P(cg^{-1}c)\), and \(\det P(g)=1\).
This intertwiner is invertible with trace zero. After a
finite scalar extension rescale it to square one. Assigning
it to \(c\) extends \(P\) to \(W\) on \(G_\QQ\) with
\(\det W=\bar\eta_K\). In a stable lattice the residual
semisimplification is a sum of two characters. An irreducible
residual plane would restrict semisimply to the normal
subgroup \(G_K\), since the translates of a simple subspace
span it. Its two distinct restricted characters extend to
\(G_\QQ\), so their two isotypic lines would each be stable
under \(c\), a contradiction. Twisting \(W\) by \(\bar\eta_K\)
if necessary, its two residual characters are therefore
\(\bar\Psi,\bar\eta_K\bar\Psi^{-1}\).

On \(G_{K_v}\) the first local flag line lies in \(P\):
its eigenvalue at \(h\) has reduction \(b\ne1\).
This line is preserved by \(\phi\), because \(\phi\) commutes
with \(h\) and the two eigenvalues of \(W(h)\) are distinct.
Together \(G_{K_v}\) and \(\phi\) generate \(G_{\QQ_\ell}\);
the quotient is unramified of degree one or two. Thus
\(W\bar\Psi^{-1}\) has reducible local trace. On its first
line inertia acts trivially, since \(I_{K_v}=I_{\QQ_\ell}\).
Its residual character on that line is one: among the two
global residual characters after twisting, only one takes
the value one at \(h\).

Apply Pan's Theorem~5.1.2(1) to \(F=\QQ\) with determinant
\(\chi=\eta_K\Psi^{-2}\). It is odd, unramified outside the
stated finite set and de Rham at \(\ell\). Its reduction extends
to \(G_\QQ\) and is nontrivial on \(G_{\QQ_\ell}\), since
\(\bar\chi(h)=b^{-2}\ne1\). These are the hypotheses of
Section~5.1.1 of that paper. The ordinary pseudodeformation
ring is finite over \(O[[1+\ell\ZZ_\ell]]\); this last
algebra acts by the local character reducing to one.
That character is unramified here, so the action factors
through the augmentation. Applying this to the trace of
\(W\bar\Psi^{-1}\) over the coefficient valuation ring shows
that its image is a finite domain in characteristic \(\ell\),
hence a finite field. If necessary enlarge the constant
residue field in this application. Restoring the
finite-valued twist and the line shows that the original
traces are algebraic over \(k'\). Inside \(k'[[z]]\) such
elements are constant.~\cite{XI-VIII-VI-pan40}

It remains to pass from curves to the entire quotient.
At a geometric field point the semisimple group span has
dimension nine in the irreducible case and at most five
in the reducible case. By the nondegenerate trace pairing,
in the irreducible case one can successively choose \(m\)
independent group matrices and then \(m\) independent
pairings with other group matrices. This gives a nonzero
\(m\)-minor. In the reducible case every such minor vanishes.
If the special fiber of \(T^{\rm ord}/I_m\) had positive
dimension, it would have a one-dimensional complete local
domain quotient by a prime. The normalization of this curve
is a complete equal-characteristic DVR with finite residue
extension, hence is \(k'[[z]]\). Here one uses the standard
normalization theorem for complete local domains, also
obtained by finite Noether normalization over \(k[[z_0]]\).
The induced trace has reducible generic representation by
the minor test. Its traces are constant by the preceding
argument, and topologically generate the curve algebra over
\(k\), contradicting its dimension. The special fiber is
therefore Artinian and has finite \(k\)-dimension. In the
complete quotient a power of the maximal ideal then lies
in \(\ell\) times the ring. Its \(\ell\)-adic and maximal-ideal
topologies agree, so lifts of a basis of the special fiber
generate the quotient over \(O\), by successive approximation.
\end{proof}

\begin{corollary}[The characters for which this Frobenius test fails]\leavevmode
\label{cor:undistinguished435}
Write \(E_x=B C_J\), \(B=K_\psi(\sqrt{-q},\sqrt{-d_0})\),
and put \(o=\operatorname{ord}(\psi)\). Define
\[
  c_0=v_2(4\operatorname{lcm}(o,2)),\qquad
  t_\ell=v_2(\ell^2-1)-3 .
\]
Every \(\alpha\) with \(\alpha(\ell)^4=1\) factors through
\(B C_{\min(J,t_\ell+c_0)}\). There are at most
\(4o\,2^{\min(J,t_\ell+c_0)}\) such characters, without
imposing a parity. In the range \(r\le x^{\eps/128}\),
if \(2^{t_\ell+c_0}\le 2^{\lfloor(\eps/128)\log_2 x\rfloor}\),
every character contributing to the complementary field
summand \(U_{\ell,r}\) in Section~\ref{sec:missing} passes
the Frobenius test of the theorem.
\end{corollary}
\begin{proof}
Every character of \(B C_J\) can be expressed as
\(\beta\delta\), with \(\beta\) a character of \(B\) and
\(\delta\) a character of \(C_J\). Indeed its Galois group
is a subgroup of the product of the two finite abelian
groups, and characters with values in all roots of unity
extend from a subgroup. This last fact follows by adjoining
one generator at a time and taking a root of the prescribed
value of its first power in the old subgroup. The exponent
of \(\operatorname{Gal}(B/\QQ)\) divides \(\operatorname{lcm}(o,2)\).
Thus \(\alpha(\ell)^4=1\) implies that the order of
\(\delta(\ell)\) divides \(2^{c_0}\).

The class of \(5\) generates
\((\ZZ/2^{J+2}\ZZ)^\times/\{\pm1\}\) of order \(2^J\).
Indeed \(v_2(5^{2^a}-1)=a+2\) for \(a\ge0\), by successive
squaring. Choose the sign of \(\ell\) so that \(u_\ell=\pm\ell\)
is \(1\bmod 4\). Then \(v_2(u_\ell-1)=t_\ell+2\), since
the other factor of \(\ell^2-1\) has valuation one. The
subgroups generated by \(5^{2^a}\) are precisely the images
of \(1+2^{a+2}\ZZ_2\), by their orders. For a character
\(\delta\) of order \(2^j\), \(j\le J\), this gives
\(\operatorname{ord}(\delta(\ell))=2^{\max(0,j-t_\ell)}\).
Consequently \(j\le t_\ell+c_0\), proving the factorization.
The two factors of the resulting field have degrees at most
\(4o\) and \(2^{\min(J,t_\ell+c_0)}\), respectively, which
proves the count. Under the last inequality the character
factors through \(L_x\) of Section~\ref{sec:missing}. The
average of \(\operatorname{Gal}(E_x/L_x)\) is the identity
on that character; its complementary projector is zero.
\end{proof}

Theorem~\ref{thm:reduciblecollision435} controls the reducible
part of the fixed-weight trace algebra. To bound the strict
length at these weights one still needs control of the
irreducible components and an upper bound for the integral
intersection lengths at the split section.


\subsection{Ordinary cuspidal adjoints at the original weights}
\label{sec:twoprimeadjoints436}

We construct an ordinary cuspidal form for each of the original
characters and weights, including those for which all three
residual Frobenius eigenvalues at \(\ell\) agree. A second
prime in the nebentypus rules out quadratic self-twists. Both
auxiliary primes can be chosen with polynomial size on the
period interval.

Keep \(K=\QQ(\sqrt{-s})\), \(E_x\), \(E_*\), \(N_*\), \(t_*\)
and \(L_*\) from Sections~\ref{sec:symplecticstrict429}
and~\ref{sec:distinguishedfrobenius435}. In particular
\(s\equiv3\bmod4\) is prime and \(s\nmid 2N_*\).
Here we allow every \(5\le\ell\le L_*\) prime to \(2sN_*\),
including primes dividing \([E_*:\QQ]\). Fix a geometric
character \(\alpha\) of \(E_x\), of sign \((-1)^r\), and
\(2\le r<t(x)\). Set \(\chi=\eta_K\alpha\), of primitive
conductor \(D=s\,\operatorname{cond}(\alpha)\), and
\(\Psi=\alpha\kappa_\ell^r\). In the proof, coefficients
are enlarged as needed at the chosen place over \(\ell\).

\begin{lemma}[Integrality of the Eisenstein constant term at every weight]\leavevmode
\label{lem:bernoulliall436}
Let \(\lambda\) be any nontrivial Dirichlet character of
conductor \(d\) prime to \(\ell\). Then \(B_{k,\lambda}/k\)
is integral at every place over \(\ell\) for every integer
\(k\ge1\).
\end{lemma}
\begin{proof}
For the Bernoulli polynomial \(B_k(X)\), set
\(P_k(X)=(B_k(X)-B_k(0))/k\). The identity
\(B_k(X+1)-B_k(X)=kX^{k-1}\) gives
\(P_k(m)=\sum_{j=0}^{m-1}j^{k-1}\) for each \(m\ge1\);
for \(k=1\) every summand is one. Also \(P_k(0)=0\).
The polynomial is continuous on \(\ZZ_\ell\), the nonnegative
integers are dense there, and \(\ZZ_\ell\) is closed.
Thus \(P_k(\ZZ_\ell)\subseteq\ZZ_\ell\). Since
\(\sum_{a=1}^d\lambda(a)=0\), the generalized Bernoulli
formula becomes
\[
       B_{k,\lambda}/k
           =d^{\,k-1}\sum_{a=1}^d\lambda(a)P_k(a/d).
\]
Each argument \(a/d\) lies in \(\ZZ_\ell\); all character
values are algebraic integers. This proves the assertion
also when \(\ell\) divides \(k\) or the order of \(\lambda\).
\end{proof}

For each \(\ell\) choose primes \(w_1,w_2\) such that
\[
  w_1\equiv1\bmod 8sN_*\ell,\qquad
  w_2\equiv1\bmod 8N_*\ell,\qquad \eta_K(w_2)=-1.
\]
The last condition specifies a nonsquare residue modulo \(s\).
The Chinese remainder theorem gives a unit residue class
for \(w_2\), distinct from that for \(w_1\). Linnik's
least-prime theorem bounds each chosen prime by an absolute
constant times \((8sN_*\ell)^{C_1}\), for an absolute \(C_1\).
They are independent of \(x,r,\alpha\) on this interval
for the fixed \(\ell\). Choose a character \(\vartheta_i\) of
order \(\ell\) and conductor \(w_i\) for each \(i=1,2\).

\begin{proposition}[An ordinary newform with two auxiliary characters]\leavevmode
\label{prop:auxiliarycusp436}
There is a normalized \(\ell\)-ordinary cuspidal newform \(f\)
of weight \(r+1\), exact level \(D w_1w_2\), and nebentypus
\(\chi\vartheta_1\vartheta_2\). At a place over \(\ell\) its
residual Galois semisimplification is
\(1\oplus\bar\chi\,\bar\kappa_\ell^r\).
\end{proposition}
\begin{proof}
Put \(k=r+1\ge3\) and \(M=Dw_1w_2\). The Eisenstein series
\(E_k(1,\chi)\) has the required parity. Its constant term
is \(-B_{k,\chi}/(2k)\), integral by the lemma since \(\chi\)
ramifies at \(s\). Its other coefficients are
\(\sum_{j\mid n}\chi(j)j^{k-1}\). Good reduction at \(\ell\)
and the expansion principle give integrality at all cusps,
as in Lemma~\ref{lem:cusp361}. Explicitly a vertical pole
on a component, after clearing its order, would give a
nonzero reduction on that component with zero expansion at
a cusp above infinity. One may check this on a full
prime-to-\(\ell\) level, with the conjugate expansions at
lifts of infinity on its components.~\cite{XI-VIII-katz361,XI-VIII-billerey361}

Write \(E=E_k(1,\chi)\) and reduce \(E(z)-E(w_1z)\)
at level \(D w_1\). It is cuspidal. At a cusp with primitive
column \((b,c)^{\mathsf T}\), put \(h_1=\gcd(c,w_1)\).
For \(E(w_1z)\) the column on level \(D\) is
\((w_1b/h_1,c/h_1)^{\mathsf T}\) and the constant term
acquires a factor \((h_1/w_1)^k\) by the slash formula.
The two columns agree modulo \(D\) since \(w_1\equiv1\bmod D\).
Their completions to determinant-one matrices therefore differ
modulo \(D\) by a translation at the cusp, which preserves
the constant term. The remaining factor reduces to one.
Pull back this cusp form to level \(M\). Its first Fourier
coefficient is one; it remains an eigenvector for all good
Hecke operators, including \(T_\ell\) with eigenvalue one
modulo the maximal ideal.

Here the desired characteristic-zero character can have an
\(\ell\)-part already at \(D\). Write the diamond group
\((\ZZ/M\ZZ)^\times=P\times H\) with \(P\) its \(\ell\)-Sylow
subgroup. The \(P\)-action on \(X_1(M)\) is free on geometric
fibers. A stabilizer in the smooth elliptic locus comes
from an automorphism of order divisible only by two or three.
At a cusp its order divides \(2n\) for an \(n\)-gon with
\(n\mid M\). Neither allows an element of order \(\ell\).
Thus \(X_1(M)\) is a finite \'{e}tale \(P\)-torsor over its
quotient. On that quotient descend the weight \(k\) cusp
line with the character of \(P\) prescribed by
\(\chi\vartheta_1\vartheta_2\), using the inverse character
for the equivariance on the line. The corresponding twist
reduces to the trivial line, since all \(\ell\)-power roots
of unity reduce to one. On \(H\) the projector for the
prescribed character is integral, and agrees on reduction
with the character of the cusp form.

On the special fiber the dual of the first cohomology of
this descended cusp line is the space of sections of
\(\omega^{2-k}\) by Serre duality and Kodaira--Spencer;
the cusp divisor cancels the divisor in the canonical line.
It vanishes because \(\omega\) has positive degree and \(k>2\).
Cohomology and base change consequently lift the residual
section to the prescribed character. The finite algebra of
good Hecke operators acts faithfully on the torsion-free
integral cusp space. Its residual eigencharacter has a
minimal prime below its maximal ideal in characteristic
zero, since a uniformizer is a nonzero divisor. Simultaneous
triangularization in characteristic zero gives a cuspidal
eigensystem there with the required reduction. The character
\(\chi\vartheta_1\vartheta_2\) is primitive of conductor
\(M\), so the normalized newform has exact level \(M\).
Its eigenvalue at \(\ell\) is a unit. Its good Frobenius
polynomials, reduction and Chebotarev give the asserted
residual semisimplification.~\cite{XI-VIII-katz361,XI-VIII-rapoport361,XI-VIII-V-newforms,XI-VIII-V-deligne}
\end{proof}

\begin{theorem}[Cuspidal adjoints for all the indicated original characters]\leavevmode
\label{thm:twoprimeadjoints436}
The Galois representation \(W_f\) of every newform supplied
by the proposition has Zariski-dense projective image in
\(\mathrm{PGL}_2\). For each \(\ell\) and each depth \(a\)
there is an auxiliary CM field \(L\), the same for all
\(x,r,\alpha\) on the interval, such that
\(U=\operatorname{Ad}^0(W_f)|_{G_L}\) is absolutely irreducible
on every open subgroup and is cuspidal automorphic in the
normalization of Theorem~\ref{thm:adjointseed361}. At places
away from \(\ell\) its inertia characters are
\(\alpha,1,\alpha^{-1}\), with zero monodromy. At places
above \(\ell\) it is ordinary crystalline with ordered
inertial characters \(\kappa_\ell^r,1,\kappa_\ell^{-r}\) and
ordered residual local characters \(\bar\Psi,1,\bar\Psi^{-1}\).
\end{theorem}
\begin{proof}
Suppose \(W_f\) has a nontrivial quadratic self-twist
\(\delta\). Reduction preserves the order of \(\delta\).
Multiplication by \(\bar\delta\) must interchange the two
residual characters, which are distinct since \(\chi\)
ramifies quadratically at \(s\). It follows that
\(\bar\delta=\bar\chi\,\bar\kappa_\ell^r\). This character is
unramified at \(w_2\) and takes value \(-1\) on its arithmetic
Frobenius: \(w_2\equiv1\bmod N_*\ell\) and \(\eta_K(w_2)=-1\).
The same is true of \(\delta\). On the other hand,
\(\dim W_f^{I_{w_2}}=1\). Indeed the local conductor exponent
is exactly one by conductor compatibility for a newform;
the determinant ramifies, so the codimension of inertia
invariants is at least one, and it plus the nonnegative
Swan conductor equals one. An intertwiner for \(\delta\)
must preserve this inertia-fixed line, and Frobenius
preserves it too. On the line the intertwiner and Frobenius
act by nonzero scalars. The proposed intertwining relation
would make their product equal its negative. This proves
that the self-twist is impossible.~\cite{XI-VIII-I-compat123}

The characteristic-zero representation \(W_f\) is absolutely
irreducible and has distinct Hodge--Tate weights. Its
projective image is infinite, since a scalar representation
on an open subgroup has equal weights. A proper connected
positive-dimensional subgroup of \(\mathrm{PGL}_2\) is solvable
and triangularizable. If it has a nontrivial unipotent
radical, its unique fixed line would be preserved by the
full image. The remaining possibility is a torus. Its
normalizer then permutes its two lines, transitively by
irreducibility, yielding a quadratic self-twist by conjugation
with the diagonal matrix with entries \(1,-1\). This too
was excluded. The projective image is therefore Zariski
dense, as is that of every open subgroup by connectedness.
The adjoint representation is irreducible there because its
invariant subspaces are ideals of \(\mathfrak{sl}_2\).

Choose a prime \(v\equiv-1\bmod 8sN_*\ell\), distinct from
\(w_1,w_2\). Let \(F_i\) be the real field of degree \(\ell\)
of conductor \(w_i\), and \(F_a\) the real degree-\(\ell^a\)
field in \(\QQ(\mu_{\ell^{a+1}})\). Set
\[
        F=\QQ(\sqrt{sv})F_1F_2F_a,\qquad L=KF .
\]
These fields depend only on the interval, \(\ell,a\), and
the choices of the three auxiliary primes. They have the
splitting at \(sN_*\ell\) used in
Proposition~\ref{prop:auxfield361} since \(-v\) is a square
at these places. The field \(L\) kills \(\eta_K,\vartheta_1,
\vartheta_2\). At an old ramified prime of \(W_f\), the
conductor exponents of \(W_f\) and its determinant agree:
both are the exponent of the primitive nebentypus conductor.
Each term of the Artin conductor bounds the corresponding
term for the determinant. Equality forces a one-dimensional
space of inertia invariants. In the Weil--Deligne
representation the two finite inertia characters are
therefore the trivial and the ramified determinant
characters. The nilpotent monodromy commuting with these
distinct characters is zero. This applies also at both new
primes; at other good primes inertia is trivial. Restricting
the two characters to \(G_L\) and taking their ratios proves
the claimed inertia description.

At \(\ell\) the ordinary crystalline representation of the
newform has unramified quotient with unit Frobenius root
reducing to one; the subcharacter follows from its
determinant. On \(G_L\) their ratio has inertia
\(\kappa_\ell^r\) and reduction \(\bar\Psi\). On
\(\mathfrak{sl}_2\) the strictly upper triangular line
and upper triangular plane give the asserted ordinary
flag, whether or not its residual characters coincide.
The Gelbart--Jacquet adjoint transfer of \(f\) is cuspidal,
since a quadratic automorphic self-twist would give one on
\(W_f\) by the good Frobenius polynomials. The transfer has
central character one and is self-dual. Along a tower of
cyclic extensions to \(L\) the degrees are two and \(\ell\).
Loss of cuspidality at one such step would give a self-twist
of that prime order by cyclic base change; the equality of
central characters forces this order to divide three. Thus
cuspidality persists. The archimedean parameters are regular
for \(r>0\), and at good primes the Satake parameters are
the two Frobenius ratios and one. These give the same Tate
normalization as in Theorem~\ref{thm:adjointseed361}.
~\cite{XI-VIII-ordinary361,XI-VIII-VI-saito40,XI-VIII-I-gj123,XI-VIII-I-ac123}
\end{proof}

\begin{corollary}[Fixed-weight rigidity at the constructed points]\leavevmode
\label{cor:seedrigidity436}
At \(U\) on \(L/F\), the polarized characteristic-zero
ordinary deformation problem of Section~\ref{sec:weightchart362}
has zero tangent space with the inertial weights and the
old inertia type fixed. Its completed local ring with all
weights varying is the formal weight ring of dimension
\(3[F:\QQ]\) over the coefficient field.
\end{corollary}
\begin{proof}
The closure of the image of the commutator of \(G_L\) under
the projective action is \(\mathrm{PGL}_2\), since the
original image is dense and this algebraic group equals
its derived group. Thus the image on the cyclotomic kernel
is already dense. In the adjoint action its span is the
full matrix algebra by irreducibility. For any nonzero
endomorphism \(X\), the open set
\(\operatorname{tr}(gX)\ne0\) therefore meets the open
set of elements with three distinct eigenvalues. Some
eigenprojector pairs nontrivially with \(X\). This verifies
the enormous-image hypothesis in Newton and Thorne's
Theorem~A as in Proposition~\ref{prop:dualselmer362}.
The polarized regular algebraic cuspidal representation
just constructed supplies its automorphic hypothesis, giving
vanishing of the polarized adjoint Bloch--Kato Selmer group.
~\cite{XI-VIII-VI-nt119}

For clarity, at every chosen place above \(\ell\) the base
representation is crystalline and the positive root exponents
are \(r,r,2r\), all at least two. On the dual strictly
upper triangular module twisted by one all exponents are
negative. Thus its local \(H^2\) vanishes by Tate duality,
and the fixed-inertia tangent lies in the finite local
condition by Lemma~\ref{lem:localweight362}. Away from
\(\ell\) the imposed inertia image is finite and averaging
is available in characteristic zero. The local excess of
the fixed-weight tangent dimension over the invariant
dimension is \(3[F_w:\QQ_\ell]\) at \(w\mid\ell\), zero
at other finite places, and \(-3\) at each real place.
Their sum is zero. The two global invariant spaces in
Poitou--Tate vanish as in Proposition~\ref{prop:dualselmer362}:
the commutant on the cyclotomic kernel consists of scalars,
complex conjugation negates the scalar Lie algebra, and
the Tate twist has infinite cyclotomic image. Hence both
fixed-weight Selmer and dual Selmer groups vanish. The
small-extension argument of Theorem~\ref{thm:weightchart362}
then gives a unique global lift for each lift of the weights,
up to strict conjugacy. Its local ordinary argument uses
distinct inertia eigenvalues in characteristic zero, where
the infinite cyclotomic powers already separate them. The
formal characters of the principal local units have the
asserted \(3[F:\QQ]\) independent parameters.
\end{proof}

\begin{theorem}[A simultaneous congruence bound for these adjoints]\leavevmode
\label{thm:actualseedcongruences436}
On the original interval there is a rational prime \(p\)
split in \(K\) and a positive integer \(D_{\rm ad}(V)\), with
\(p+\log D_{\rm ad}(V)\ll_{\psi,\eps}V^C\), such that all
the constructed forms satisfy
\[
 0<
 v_\ell\!\left(\operatorname{tr}\operatorname{Ad}^0(W_f)
    (\operatorname{Frob}_p)
    -\bigl(1+\alpha(p)p^r+\alpha(p)^{-1}p^{-r}\bigr)\right)
                  \le v_\ell(D_{\rm ad}(V)).
\]
Here \(v_\ell(\ell)=1\) and the place is the one giving the
congruence in Proposition~\ref{prop:auxiliarycusp436}.
\end{theorem}
\begin{proof}
Put \(M_\ell=sN_*w_1w_2\). All these integers and the number
of eligible primes are \(O_{\psi,\eps}(V^{C_2})\) by the
least-prime bound and the polynomial bounds for \(N_*,L_*\).
Choose a power \(2^j\) strictly above \(16,L_*\) and every
\(M_\ell\), taking the least such power. The least prime
\(p\equiv1\bmod 2^j s\) is polynomially bounded by the
same theorem. It exceeds these quantities, splits in \(K\),
and is a good prime different from every \(\ell\) under
consideration.

For each \(\ell\) and \(2\le r\le t_*\) let \(\mathcal F_{\ell,r}\)
be all normalized newforms of weight \(r+1\) and level
dividing \(M_\ell\), with all their coefficient embeddings.
Take each distinct normalized eigenform once, so the list
is stable under rational Galois conjugation. It has size
\(O((r+1)M_\ell^3)\): the forms are independent in the
classical space at principal level \(M_\ell\) by multiplicity
one at good primes. The valence formula bounds the dimension
there by \(1+(r+1)[\mathrm{SL}_2(\ZZ):\Gamma(M_\ell)]/12\),
and the index is at most \(M_\ell^3\). Equivalently one
may use the degree of the Hodge line and the degree bound
on the number of sections of a line bundle on a curve.

For such a form let \(\epsilon_f\) be its nebentypus.
Its adjoint trace is
\(u_f=a_p(f)^2/(\epsilon_f(p)p^r)-1\), so \(p^r u_f\)
is an algebraic integer. Deligne's theorem gives absolute
value \(p^{r/2}\) to each of the two Frobenius roots at
every complex embedding. Their ratios have modulus one,
giving \(|u_f|\le3\). For every \(\zeta^{[E_*:\QQ]}=1\) put
\[
       b(f,\zeta)=p^r\bigl(u_f-1-\zeta p^r-\zeta^{-1}p^{-r}\bigr).
\]
It is an algebraic integer. Its magnitude is positive,
since \(p^r-1-p^{-r}>3\), and is at most \(7p^{2r}\).
The absolute value of the product of these integers over
all listed \(\ell,r,f,\zeta\) is a positive rational integer
by Galois invariance; use one for an empty product.
The degree \([E_*:\QQ]\), \(t_*\), \(p\), all levels and
the numbers of factors are polynomial in \(V\). The
logarithm of each magnitude is at most \(2r\log p+\log 7\),
which proves the claimed estimate for this product.

For the actual form the list contains \(b(f,\alpha(p))\).
All other integral factors have nonnegative valuation there;
\(p^r\) is a unit at \(\ell\). Finally \(\eta_K(p)=1\).
Taking the adjoint of the residual semisimplification in
the proposition shows that this particular trace difference
reduces to zero. This proves both bounds.
\end{proof}

These congruences concern actual points obtained from the
constructed forms. An upper bound for the full strict
cohomology also requires a comparison of its integral length
with congruences on all the components that carry it. The
comparison with the integral projected theta period remains
required on the same interval.


\subsection{Recovering the ordered local characters from Hecke traces}
\label{sec:orderedtraceconductor437}

At a prime where the three residual inertia characters agree,
the integral Hecke trace algebra need not separate the three
ordinary quotients. At the original positive weight their
inertia characters in characteristic zero do separate them.
We compute an explicit integer that controls the discrepancy
on the whole specialization, including its torsion.

\begin{lemma}[Three ordered roots over a trace domain]\leavevmode
\label{lem:ordereddomains437}
Let \(B\) be a domain, let \(s_1,s_2,s_3\in B\) have nonzero
pairwise differences, and let \(t_0,t_1,t_2\in B\).
Use the three equations \(\sum_i s_i^j u_i=t_j\) for
\(0\le j\le2\) to define \(u_i\in\operatorname{Frac}B\).
Suppose they are the roots with multiplicity of a monic
cubic over \(B\) with invertible constant coefficient.
Set \(\delta_i=\prod_{j\ne i}(s_i-s_j)\) and
\(C=B[u_1,u_2,u_3]\). Then all \(u_i\) are units of \(C\).
It is generated as a \(B\)-module by
\[
           u_1^i u_2^j,\qquad 0\le i\le2,\quad 0\le j\le1,
       \qquad \delta_1^2\delta_2 C\subset B .
\]
\end{lemma}
\begin{proof}
For \(\{i,j,k\}=\{1,2,3\}\), subtract \((s_j+s_k)t_1\)
from \(t_2+s_js_k t_0\). Only \(\delta_i u_i\) remains.
The determinant of the three equations is a Vandermonde
product and is nonzero in the fraction field. The monic cubic
reduces all powers of \(u_1\) to degrees at most two; division
by \(X-u_1\) gives a monic quadratic for \(u_2\) over
\(B[u_1]\). The remaining root is \(t_0-u_1-u_2\).
These facts give the six generators, without requiring their
independence. Each becomes integral in \(B\) after
multiplication by \(\delta_1^2\delta_2\), by the first
identity. The product of the three roots is a unit of \(B\),
so their inverses already belong to \(C\).
\end{proof}

For the newform of Proposition~\ref{prop:auxiliarycusp436}
use the same fields \(F,L\) of Theorem~\ref{thm:twoprimeadjoints436},
at a chosen depth \(a\). Let \(d=[F:\QQ]\), let \(g_\ell\)
be the number of places of \(F\) above \(\ell\), and choose one
place of \(L\) above each of them. The completions at a chosen
pair coincide. Write \(\Lambda\) for the full three-character
ordinary weight algebra over a sufficiently large coefficient
DVR \(O\), with the prime-to-\(\ell\) diamonds of the adjoint
fixed. The local unit groups have no \(\ell\)-power torsion:
their ramification indices at \(\ell\) are \(\ell^a\), while
\(\QQ_\ell(\mu_\ell)\) is totally ramified of degree \(\ell-1\).
Consequently \(\Lambda\) is a power series ring in \(3d\)
variables. Use the cyclotomic normalization of the traces in
which the constructed point is \(U=\operatorname{Ad}^0(W_f)|_{G_L}\).

\begin{proposition}[Interpolation of Frobenius on the ordered quotients]\leavevmode
\label{prop:heckeorder437}
There is an ordinary good Hecke domain \(B\) through this
point, finite over and containing \(\Lambda\), with continuous
group traces \(t(g)\in B\) on \(G_L\). On any such full
component through the point one can adjoin all the ordered
local characters at \(\ell\) within \(\operatorname{Frac}B\),
obtaining a finite algebra \(B^{\rm ord}\). It has at most
\(6^{g_\ell}\) module generators over \(B\). For each
\(w\mid\ell\) choose an element \(\sigma_w\) of inertia
whose cyclotomic value is \((1+\ell)^{\ell^a}\). Let
\(s_{wi}\in\Lambda\) be the three ordered weight character
values there, let \(\delta_{wi}=\prod_{j\ne i}(s_{wi}-s_{wj})\),
and put
\[
        H=\prod_{w\mid\ell}\delta_{w1}^2\delta_{w2}\in\Lambda.
        \qquad H B^{\rm ord}\subset B .
\]
\end{proposition}
\begin{proof}
We specify the Hecke algebra in use. The group is the
definite unitary group of rank three over \(F\), quasi-split
at the finite places. The extension \(L/F\) is unramified
there. At a prime dividing \(sN_*\ell\) this follows from
the local splitting supplied by \(-v\) being a square in
the construction of \(F\); at other primes it follows by
base change from \(K/\QQ\). Labesse's descent to this group,
in the form used in Proposition~\ref{prop:ordinaryfamily363},
applies to the regular algebraic conjugate self-dual
cuspidal representation of Theorem~\ref{thm:twoprimeadjoints436}.
At \(\ell\) it has an ordinary refinement, since the two
Frobenius slopes of the newform are \(0,r\) and the adjoint
has their ordered differences. Take a tame level carrying
this descent; one can take it sufficiently small at an
additional split place. The dual ordinary module is finite
free over \(\Lambda\), and the control theorem identifies
its regular algebraic specializations with the corresponding
classical ordinary spaces. This is the definite ordinary
control theorem used in Proposition~\ref{prop:ordinaryfamily363}.
Its construction and level control are independent of
distinctions between residual inertia characters.~\cite{XI-VIII-geraghty363}

Form the commutative subalgebra of its endomorphisms generated
over \(\Lambda\) by good split Hecke operators. It is finite
and torsion-free over \(\Lambda\). Localize at the residual
system of the adjoint and quotient by a minimal prime through
the adjoint in characteristic zero. The resulting domain
\(B\) contains \(\Lambda\), since its nonzero elements act
as nonzero divisors before passage to the minimal prime.
Dominant algebraic weights with sufficiently large successive
gaps and the specified finite diamonds are Zariski dense:
on small logarithmic neighborhoods one can vary every labeled
integral exponent through an arithmetic progression, and
the matrix of the local embeddings has full rank. Increasing
successive exponents to retain the gaps leaves a dense set
in each variable. Points above these weights outside any
proper closed subset are dense in \(B\). Indeed a nonzero
element of the finite domain generates an ideal meeting
\(\Lambda\) nontrivially, by its minimal polynomial over the
fraction field and clearing denominators. Every such point
occurs in the specialized ordinary space, by faithful support
and Nakayama's lemma.

At these classical points the good Hecke generators give
the normalized Galois traces. To obtain \(t(g)\) for arbitrary
\(g\), approximate the conjugacy class of \(g\) by good split
Frobenius classes using Chebotarev. For each finite list of
classical points one can approximate the traces there to
arbitrary coefficient precision by continuity of the
classical representations. The compact Hecke algebra has
closed image in the product of their valuation rings.
This gives the trace tuple in the algebra; the density just
proved makes it unique in \(B\). The compact image is
homeomorphic to the algebra with its maximal-ideal topology,
which also gives continuity. Newton's identities supply the
characteristic coefficients since \(3!\in O^\times\).
This is the trace construction of
Proposition~\ref{prop:ordinaryfamily363}.

The only ramified factor at \(\ell\) in the construction of
\(L\) is the degree-\(\ell^a\) cyclotomic layer. All other
factors are unramified locally. Total ramification of the
cyclotomic extension therefore allows the indicated
\(\sigma_w\) in inertia. At the adjoint its ordered values
are \((x_\ell,1,x_\ell^{-1})\) with
\(x_\ell=(1+\ell)^{\ell^a r}\). They are distinct in
characteristic zero; every Vandermonde factor is consequently
nonzero in \(B\).

Choose an arithmetic Frobenius lift \(\phi_w\). In the lemma
take \(t_j=t(\phi_w\sigma_w^j)\). On the dense regular classical
points where the weight values are distinct, the ordinary
Galois filtration has the three ordered inertia characters
of the weight, with the fixed common normalization. This is
local ordinary compatibility in
Proposition~\ref{prop:ordinaryclass405}. We can restrict the
tests to the open set on which a nine-by-nine trace minor
nonzero at the adjoint is nonzero; their transfer is then
cuspidal, by the rank of the semisimple group span. At each
test point the trace of \(\phi_w\sigma_w^j\) is the sum of
the three products on the flag quotients, regardless of
off-diagonal entries. The three solutions \(u_{wi}\) over
\(\operatorname{Frac}B\) therefore agree with the correctly
ordered values at \(\phi_w\). Their elementary symmetric
functions agree with the characteristic coefficients of
\(\phi_w\) on this dense open set, and hence in the field.
The constant coefficient is a unit of \(B\), by reduction
at the residual system. This proves every hypothesis of the
lemma without dividing by a residual eigenvalue difference.

Adjoin these three values at every \(w\). Products of the
local module generators generate the resulting algebra,
and multiplication by \(H\) carries each of them into \(B\).
It remains to see that these values give the full characters.
Local reciprocity identifies the abelian local group with
the product of its compact units and a procyclic factor
generated by the image of a Frobenius lift, after making the
corresponding choice of splitting. The values on inertia
are given by the weights. On the procyclic factor there is
a continuous character sending its generator to \(u_{wi}\):
the unit group of a finite complete local algebra with finite
residue field is profinite. The resulting characters agree
with the classical ones on the dense test set, proving the
assertion for every local group element.
\end{proof}

\begin{theorem}[Integral loss on the entire original weight specialization]\leavevmode
\label{thm:orderedconductor437}
Specialize \(\Lambda\to O\) to the original three inertia
weights \(\kappa_\ell^r,1,\kappa_\ell^{-r}\).
The kernel and cokernel of
\[
      B\otimes_\Lambda O\longrightarrow
                   B^{\rm ord}\otimes_\Lambda O
\]
are annihilated by
\(\ell^{\,6g_\ell(1+a+v_\ell(r))}\), with \(v_\ell(\ell)=1\).
The same bound holds after any further scalar extension
from \(O\), including every finite precision.
One has \(g_\ell\le2\ell^2\). In particular at \(a=0\) the
annihilating exponent is polynomial in \(V\), uniformly for
all the characters and twists in Section~\ref{sec:twoprimeadjoints436}.
\end{theorem}
\begin{proof}
Write \(i:B\to B^{\rm ord}\) for the inclusion.
Multiplication by \(H\) is a \(B\)-linear map
\(j:B^{\rm ord}\to B\) by the proposition. Both \(ij\)
and \(ji\) are multiplication by \(H\). The two identities
persist under every scalar extension, so its image annihilates
the kernel and cokernel even when the extension is not flat.

At the prescribed weights the three specialized differences
are, up to sign, \(x_\ell-1, 1-x_\ell^{-1}\) and
\(x_\ell-x_\ell^{-1}\). The factors \(x_\ell\) and \(x_\ell+1\)
are units. All differences have valuation
\(1+a+v_\ell(r)\). Indeed raising a principal unit to an
exponent prime to \(\ell\) preserves the valuation of its
difference from one. Raising it to the \(\ell\)-th power
increases that valuation by one for these rational principal
units, by the binomial theorem and \(\ell>2\). Each
\(\delta_{wi}\) has two such factors, so the specialized
\(H\) is the stated power of \(\ell\) times a unit. Finally
\([F:\QQ]\le2\ell^{a+2}\), and \(F\) has ramification index
\(\ell^a\) at \(\ell\). The local degrees sum to \([F:\QQ]\),
giving \(g_\ell\le2\ell^2\). At depth zero use
\(\ell\le L_*\), \(r\le t_*\), both polynomially bounded on
the interval, to obtain the assertion.
\end{proof}

\begin{corollary}[The ordinary component through the adjoint]\leavevmode
\label{cor:seedcomponent437}
Suppose additionally that \(\ell\ge11\) is prime to
\([E_*:\QQ]\). Use the stable compact type at the old places
of Lemma~\ref{lem:typelattice427}, with hyperspecial level
at the other places away from \(\ell\).
Let \(R\) be the polarized ordinary determinant ring over
\(L/F\) with three labeled full local characters, the old
inertia kernels of \(U\), and freely varying ordered
weights as in Section~\ref{sec:adjcomponent372}. At the
characteristic-zero point of \(U\), \(R\) has a unique
irreducible component. Its integral domain is \(B^{\rm ord}\)
for any minimal good Hecke component \(B\) through that
point with these types. Its local ring at the point is
regular of dimension \(3d\).
\end{corollary}
\begin{proof}
The old conductor places are split in \(L/F\). The adjoint
has inertia type \(\alpha\oplus1\oplus\alpha^{-1}\) there
and zero monodromy. The compact types and their stable
coefficient lattices consequently include it. For ordinary
control with these coefficients the arithmetic stabilizers
have order prime to \(\ell\): an element of order \(\ell\)
in \(\mathrm{GL}_3(L)\) would have a primitive \(\ell\)-th
root as eigenvalue, of degree \(\ell-1>3\) over \(L\).
Here \(L\cap\QQ(\mu_\ell)=\QQ\) since ramification at
\(\ell\) in \(L\) has \(\ell\)-power index. The double-coset
and lattice proof of Lemma~\ref{lem:typelattice427} applies,
and the old compact type gives precisely this inertia and
zero monodromy at the cuspidal classical test points.
At every other finite place away from \(\ell\) the descent
and the test representation are unramified.

The determinant and local characters on \(B^{\rm ord}\)
give a surjection \(R\to B^{\rm ord}\). To check its
equations, one may invert \(H\), where
\(B^{\rm ord}[1/H]=B[1/H]\). At the dense cuspidal classical
tests the polarized trace identities and the identities
for every prescribed tame inertia kernel hold, by local
compatibility. For the local ordinary word,
each \(M(g_i)-\xi_i(g_i)I\) carries the \(i\)-th flag step
to the preceding step. The product in step order therefore
vanishes, also after multiplication and taking the trace
against any group element. The associated local characteristic
polynomial identities hold as well. All these are equalities
in the domain \(B^{\rm ord}\) by density. Surjectivity follows
because \(R\) supplies the weights, all the good Hecke
coefficients, and precisely the \(u_{wi}\) adjoined above.

The augmentation at the adjoint extends to \(B^{\rm ord}\)
with the same characteristic-zero residue field \(J\), since
the ordered Frobenius values there are those of \(U\) and
are integral units in a sufficiently large \(O\). The domain
\(B^{\rm ord}\) is integral and finite over the normal domain
\(\Lambda\). The prime of this weight in
\(\Lambda[1/\ell]\) has height \(3d\), since its completion
has the translated weight variables as a regular set of
parameters. Going down and incomparability give the same
height for the adjoint prime of \(B^{\rm ord}\).

At the corresponding prime of \(R\) its tangent space over
\(J\) injects into the \(3d\)-dimensional space of weight
tangents. Here one uses the characteristic-zero argument of
Lemma~\ref{lem:detangent372}, with the following checks at
coincident residual weights. A determinant over the dual
numbers lifting the absolutely irreducible \(U\) is given
by a representation unique up to strict conjugacy, by
Chenevier's theorem. The nonsingular trace pairing makes
the tame kernel identities actual matrix identities and
does the same for the ordinary word. The eigenvalues of
\(\sigma_w\) are distinct in \(J\), and their differences
are units over the dual numbers over \(J\). In its eigenbasis,
put \(\sigma_w\) in every position of the word except the
\(i\)-th. The factors on the left are units on rows of
index at least \(i\), those on the right on columns of index
at most \(i\). This forces triangularity in the prescribed
order with diagonal \(\xi_i\), and recovers those labeled
characters from the matrices. Conjugate duality for this
absolutely irreducible determinant gives the polarized lift
over \(F\), by the symmetric pairing argument of the same
lemma. Continuity follows by recovering matrix coordinates
from the traces against a fixed basis of nine group matrices.
Corollary~\ref{cor:seedrigidity436} gives the asserted
injection. The hypotheses concern this one adjoint, for
which its automorphic realization was proved above.

The local ring of \(R\) now has a quotient of dimension
\(3d\) and embedding dimension at most \(3d\). It is regular,
hence a domain. The prime kernel of the local surjection is
zero, since a positive-height prime would lower the dimension.
Its contraction is the unique minimal prime of \(R\) contained
in the adjoint prime. The original surjection identifies
the quotient by that contraction with \(B^{\rm ord}\).
\end{proof}

This identifies the component passing through the chosen
automorphic point even when all three residual local characters
coincide. It remains to determine which further components
meet the split section carrying the strict classes and to
bound their total integral intersection there.


\subsection{The fiber of the split trace at repeated local characters}
\label{sec:splitwordfiber438}

At the weight and on the auxiliary field \(L/F\) of
Theorem~\ref{thm:twoprimeadjoints436}, consider the split determinant
with diagonal \(\Psi,1,\Psi^{-1}\), restricted to \(G_L\).
It has the prescribed polarization, tame inertia and ordered
inertia weights. We use the cyclotomic normalization of that
theorem. Restoring the common twist multiplies the factors
of the ordinary word below by units and carries each local
characteristic polynomial to its twist. We keep one place of \(L\) over each \(w\mid\ell\)
of \(F\). Write \(G_w,I_w\) for its local group and inertia.
The coefficient ring \(O\) is the ring of integers of a finite
extension of \(\QQ_\ell\) containing all the diagonal values,
with residue field \(k\) and uniformizer \(\varpi\).
Here \(\ell\ge5\) and \(\ell\nmid 2sN_*\).

The source defines an ordinary pair by a global determinant,
three local characters \(\xi_{w,i}\) with prescribed reductions
and torsion values, and the equations
\begin{equation}
  t\bigl(h(g_1-\xi_{w,1}(g_1))
           (g_2-\xi_{w,2}(g_2))(g_3-\xi_{w,3}(g_3))\bigr)=0 .
  \label{eq:ordinaryword438}
\end{equation}
Here \(g_i\in G_w\), \(h\in G_L\), group elements are read in
the group algebra, and \(t\) is extended linearly. See
\srcloc{Development/AutomorphicArithmetic/AutomorphicVarieties/WeilGroupCharacters/UniversalOrdinaryDeterminantRing}{94}
and lines 127--147 there. The source's residual data in this
definition have three distinct \emph{global} characters. Their
restrictions to \(G_w\) can coincide. For an actual stable
flag with ordered line quotients the characteristic polynomial
on \(G_w\) also satisfies
\begin{equation}
  \operatorname{char}_{D}(g;X)
         =\prod_{i=1}^3(X-\xi_{w,i}(g))\qquad (g\in G_w).
  \label{eq:factorlocal438}
\end{equation}
Indeed in an adapted basis the matrix is triangular, also over
a local ring with nilpotents. We compute the effect of this
additional identity. The quotients considered below keep the
given local weights at every place.

\begin{theorem}[The complete fiber of the ordinary word equations]\leavevmode
\label{thm:splitwordfiber438}
Fix the global determinant to be the constant residual split
determinant over \(k\), after any extension of its coefficient
algebra. At each \(w\) partition \(\{1,2,3\}\) by equality of
the three residual characters on the \emph{full} \(G_w\).
Denote the partition by \(\mathcal B_w\).
In the problem defined by \eqref{eq:ordinaryword438}, the
complete ring for the remaining local characters is
\[
 k[[z_{wi}:w\mid\ell,\ 1\le i\le3]]
 \Big/\Big(\prod_{i\in B}z_{wi}: w\mid\ell,\ B\in\mathcal B_w\Big).
\]
It is reduced and equidimensional of dimension
\(\sum_{w\mid\ell}(3-|\mathcal B_w|)\).
There are \(\prod_{w,B}|B|\) minimal primes.
\end{theorem}
\begin{proof}
Write \(\beta_i\) for the residual diagonal characters at one
place and choose an arithmetic Frobenius lift \(\phi\). A lift
of \(\beta_i\) to a complete Noetherian local \(k\)-algebra
with residue field \(k\) and fixed inertia equals
\(\beta_i(g)(1+z_i)^{n(g)}\). Here \(n(g)\) is the image in
the pro-\(\ell\) quotient \(\ZZ_\ell\) of the unramified
Galois group, normalized by \(n(\phi)=1\).
Indeed \(1+\mathfrak m\) in each Artinian quotient is a
finite \(\ell\)-group, and passage to the limit gives precisely
the continuous powers of \(1+z_i\). An unramified variation
kills the torsion in the local abelianization, since its
unramified pro-\(\ell\) quotient is torsion-free. These
parameters therefore respect the specified torsion values.

With the global determinant fixed, its trace of the product
in \eqref{eq:ordinaryword438} can be computed on the split
diagonal. On a line of local character \(\beta\) the product
acts by \(P_\beta=\prod_i(\beta(g_i)-\xi_i(g_i))\).
Taking the testing element \(h\) in the local group gives the
sum of \(|B|\beta(h)P_\beta\) over the distinct characters
\(\beta\), where \(B=\{i:\beta_i=\beta\}\).
Their joint image is a finite abelian group of order prime to
\(\ell\). Multiplication by \(\beta(h)^{-1}\) and averaging
over representatives of that image selects \(|B|P_\beta\).
The factor \(|B|\) is invertible in \(k\). Thus the equations
are equivalent to the vanishing of every \(P_\beta\).
Conversely this vanishing kills every diagonal entry and
therefore every test with global \(h\).

For \(i\in B\) the ideal of differences
\(\beta(g)-\xi_i(g)\) is \((z_i)\), by evaluation at \(\phi\).
For \(i\notin B\) it is the unit ideal, by evaluation at a
group element on which the residual characters differ.
The \(g_i\) vary independently, so the ideal of the products
is exactly \((\prod_{i\in B}z_i)\). The only remaining data
of the pair are the characters at the chosen places, hence
the independent sets of equations give the asserted ring.
The ideal just found is the intersection of the coordinate
ideals obtained by choosing one variable from each block.
Indeed a monomial avoiding at least one variable in every
block fails to belong to the coordinate ideal generated by
those avoided variables. Membership in these monomial ideals
is coefficientwise for power series. All the coordinate ideals
in the intersection are incomparable primes of the same
height, one for each indicated choice. This gives all the
claims.
\end{proof}

\begin{proposition}[A triple local coincidence on the common intervals]\leavevmode
\label{prop:dyadiccollision438}
On every sufficiently distant interval of
Theorem~\ref{thm:conductor232} there are a single allowed
\(\ell\), an even character \(\alpha\) of all \(E_x\) and
a single twist \(r\) in the required range such that the
residual global diagonal has three distinct characters on
\(G_L\), while at every \(w\mid\ell\) all three are trivial.
This holds at every auxiliary depth \(a\). The constant-trace
fiber of \eqref{eq:ordinaryword438} then has dimension
\(2g_\ell\), where \(g_\ell=\#\{w\mid\ell\text{ in }F\}\).
\end{proposition}
\begin{proof}
Fix four distinct primes congruent to one modulo \(32\)
and prime to the conductor of the fixed \(\psi\). Such primes
exist by Dirichlet's theorem. Only \(q,d_0,s\) can remove
further primes from this list under the condition
\(\ell\nmid 2sN_*\), by the conductor in
Theorem~\ref{thm:conductor232}. Select a remaining prime.
On sufficiently large intervals all four primes are below
\(L_*\), and all their predecessors \(\ell-1\) are below
\(t(x)\). The dyadic layer \(C_J\) then contains
\(\QQ(\zeta_{32})^+\), since \(J\) grows uniformly with \(x\).
Take a faithful character of this real cyclic field of
degree eight and set \(r=\ell-1\). Its value at \(\ell\)
is one, so its restriction to \(G_{\QQ_\ell}\) is trivial.
The residual cyclotomic factor to the power \(r\) is trivial
everywhere. These facts give the assertion at the local places.

The field \(L\) is the compositum of two quadratic fields
and abelian extensions of \(\ell\)-power degree. The
\(2\)-primary part of its Galois group has exponent at most
two. Its intersection with the cyclic field of degree eight
therefore has degree at most two. Reduction at the odd prime
\(\ell\) preserves the order of the eighth roots of unity.
It follows that \(\bar\Psi|_{G_L}\) has order at least four,
which makes the global triple distinct. The dimension is now
given by Theorem~\ref{thm:splitwordfiber438}. For example
one can vary a single unramified label freely in \(k[[z]]\)
and keep the two others and the global determinant constant.
\end{proof}

\begin{theorem}[The finite fiber imposed by local characteristic polynomials]\leavevmode
\label{thm:flagfiber438}
In the fiber of Theorem~\ref{thm:splitwordfiber438}, impose
\eqref{eq:factorlocal438} for every \(w\) and \(g\).
The resulting ring is the tensor product over \(k\), over
all blocks \(B\in\mathcal B_w\), of
\[
 k[z_{wi}:i\in B]/
 \bigl(e_1(z_B),\ldots,e_{|B|}(z_B)\bigr),
\]
where \(e_j\) is the elementary symmetric polynomial of
degree \(j\). It has length \(\prod_{w,B}|B|!\). The
dimension of its cotangent space is
\(\sum_w(3-|\mathcal B_w|)\).
\end{theorem}
\begin{proof}
At a place let \(\Delta\) be the joint residual character
image. The map \(G_w\to\Delta\times\ZZ_\ell\) given by the
residual characters and \(n\) is onto. Indeed take powers
of a lift of any \(\delta\in\Delta\) with exponents congruent
to one modulo \(|\Delta|\) and tending to zero in
\(\ZZ_\ell\). Its images tend to \((\delta,0)\) in the
compact image. Surjectivity onto \(\ZZ_\ell\) gives the
other factor after multiplication by such elements.

For every positive integer \(j\), take local elements with
images \((\delta,j)\) and compare their traces. Orthogonality
on \(\Delta\) gives for each block of size \(m\)
\[
       \sum_{i\in B}(1+z_i)^j=m .
\]
At \(j=1,\ldots,m\) successive expansion gives zero for the
power sums of \(z_i\) of degrees \(1,\ldots,m\).
Newton's identities and \(\ell>m\) give \(e_1=\cdots=e_m=0\).
Conversely these identities imply \(z_i^m=0\), since
\(\prod_{i\in B}(X-z_i)=X^m\) and each \(z_i\) is a root.
All positive power sums vanish by Newton's identities in
degrees below \(m\) and this nilpotence thereafter. Expansion
gives the displayed equality for all nonnegative integer
exponents and then for all \(n\in\ZZ_\ell\) by continuity.
Thus the local traces agree at every element, including its
second and third powers, which gives the full characteristic
polynomial by Newton's identities. The product equation
from the ordinary word is already \(e_m=0\) in each block.

The block ring represents ordered factorizations of \(X^m\)
into monic linear factors. Adjoin the first root of \(X^m\),
divide by its monic linear factor, and adjoin a root of the
remaining monic polynomial. Continuing gives free extensions
of ranks \(m,m-1,\ldots,1\). The resulting ring represents
the same factorizations, so its dimension is \(m!\).
All variables are nilpotent and the only linear equation
is \(\sum_{i\in B}z_i=0\). Tensoring the block algebras proves
the formula for both invariants. The product is already
complete because it is Artinian and local.
\end{proof}

\begin{corollary}[The integral contribution of ordered labels]\leavevmode
\label{cor:flagcotangent438}
At the original inertia weight over \(O\), let \(C\) be the
ordinary determinant algebra with \eqref{eq:factorlocal438}
imposed, or any closed quotient through the split section.
Only the determinant and the three ordered characters at
each chosen place are recorded. Let \(B\) be the closed
\(O\)-subalgebra of \(C\) generated by all global traces.
Write \(J_B,J_C\) for the augmentation ideals at the split
section, put \(b=1+a+v_\ell(r)\), and normalize \(v_\ell(\ell)=1\).
The map
\[
              J_B/J_B^2\longrightarrow J_C/J_C^2
\]
has cokernel of \(O\)-length at most
\[
  2\,v_\varpi(\ell)\, b\,
       \sum_{w\mid\ell}(3-|\mathcal B_w|).
\]
Its kernel is killed by \(\ell^{12g_\ell b}\).
These assertions apply to the image of the determinant and
its ordered characters in any family of actual local flags
through the split section, including families over rings with
nilpotents. At depth zero the sum of the displayed upper
bounds, divided by \(v_\varpi(\ell)\) and weighted by
\(\log\ell\), over the indicated primes, twists and characters
of any \(E_x\), is polynomially bounded in \(V\) throughout
the common interval.
\end{corollary}
\begin{proof}
Let \(u_{wi}=\xi_{w,i}(\phi_w)\) for an arithmetic Frobenius
lift. The algebra \(C\) is topologically generated over \(B\)
by these values. Indeed the determinant is determined by
traces when \(3!\) is a unit, and all inertia characters are
fixed. Every unramified character is determined continuously
by Frobenius. Equivalently one may use a presentation of the
determinant and character rings with these generators before
imposing the equations. The coefficient ring and the completed
deformation rings here are compact with finite residue field.
The finite spans of monomials over the closed ring \(B\)
used below are therefore closed.

Take \(\sigma_w\) as in
Proposition~\ref{prop:heckeorder437}. Its inertia weights
\(s_i\) are \(x,1,x^{-1}\) with \(x=(1+\ell)^{\ell^a r}\).
The trace coefficient of \eqref{eq:factorlocal438} at
\(\phi_w\sigma_w^j\) gives
\(t(\phi_w\sigma_w^j)=\sum_i u_{wi}s_i^j\).
For \(\{i,j,k\}=\{1,2,3\}\), multiplication by the
coefficients of \((X-s_j)(X-s_k)\) and summation give
\((s_i-s_j)(s_i-s_k)u_{wi}\in B\).
Each difference has valuation \(b\), by the binomial
calculation of Theorem~\ref{thm:orderedconductor437}.
Thus \(\ell^{2b}u_{wi}\in B\) without any condition on
the components of \(C\).

By \eqref{eq:factorlocal438} at \(\phi_w\), the three \(u_{wi}\)
are the roots of a split monic cubic over \(B\) with unit
constant coefficient. One can eliminate the third using
their sum, reduce powers of the first using the cubic
and those of the second using the remaining quadratic.
This holds over arbitrary commutative rings by monic
division. The resulting six monomials at each place have
total degree at most three. Inverses of the roots are
already in this algebra by the unit product. Taking products
over the places gives \(C\) as the finite span of these
monomials and gives \(\ell^{6g_\ell b}C\subset B\).
If \(y\in J_B\cap J_C^2\), multiplication by
\(\ell^{12g_\ell b}\) multiplies both factors in each product
from \(J_C^2\) by \(\ell^{6g_\ell b}\). Each factor then lies
in \(J_B\), proving the kernel assertion.

The cokernel is generated over \(O\) by the classes of
\(u_{wi}-u_{wi}^{0}\), where the superscript denotes
augmentation. The interpolation calculation kills each
class by \(\ell^{2b}\). Modulo \(\varpi\), this cokernel
is the cotangent space of \(C/(\varpi,J_B C)\) at its
residue point. In fact its presentation is
\(J_C/(J_C^2+J_B C+\varpi J_C)\), since modulo \(J_C^2\)
a coefficient in \(C\) acts on \(J_B\) by augmentation.
The ring in this fiber is a quotient of the ring in
Theorem~\ref{thm:flagfiber438}. It fixes every trace
to its residual split value, hence fixes the determinant
as well. The dimension of its cotangent space is at most
the number in that theorem. Nakayama's lemma and the
annihilator give the length bound. Finally \(g_\ell\le2\ell^2\);
at depth zero \(b\le1+\log_\ell t_*\). There are at most
\(n_*=[E_*:\QQ]\) characters on each field, at most \(t_*\)
twists and at most \(L_*\) primes, all polynomially bounded
in \(V\). This proves the uniform assertion.
\end{proof}

\begin{theorem}[Elementary divisors at the original split trace]\leavevmode
\label{thm:smithlabels438}
Keep the global determinant equal to the original split
determinant over \(O\), for any \(\alpha,r,\ell,a\) in this
subsection. Choose \(\phi_w\) of cyclotomic value one. Let
\(y_w\) be the \(\ell\)-power part of the root of unity
\(\alpha(\phi_w)\), set \(x=(1+\ell)^{\ell^a r}\), and put
\[
   h_w=\min\{v_\varpi(x-1),v_\varpi(y_w-1)\},
        \qquad v_\varpi(0)=+\infty .
\]
In the fixed-weight algebra satisfying
\eqref{eq:ordinaryword438} and \eqref{eq:factorlocal438},
the cotangent module of this fiber at the split section is
\[
     \bigoplus_{w\mid\ell}\ \bigoplus_{B\in\mathcal B_w}\
             \bigoplus_{j=1}^{|B|-1} O/\varpi^{j h_w}O .
\]
For any quotient through the split section this gives the
sharper bound \(\sum_{w,B}\binom{|B|}{2}h_w\) for the
cokernel length in Corollary~\ref{cor:flagcotangent438}.
For the actual family of
Proposition~\ref{prop:dyadiccollision438} the displayed
cotangent module is
\[
             \bigoplus_{w\mid\ell}
               \bigl(O/\ell^{a+1}O\ \oplus\
                     O/\ell^{2(a+1)}O\bigr).
\]
\end{theorem}
\begin{proof}
Work one place at a time. Write \(\lambda=\Psi|_{G_w}\),
\(\tau\) for the Teichm\"uller lift of its reduction, and
\(d=\lambda/\tau\). The local field is an unramified extension
of the completion of the degree-\(\ell^a\) cyclotomic layer.
The cyclotomic extension above it is totally ramified.
Consequently the prescribed \(\phi_w\) exists by multiplying
a Frobenius lift by inertia. On inertia the principal-unit
image of \(d\) is topologically generated by \(x\), and
\(d(\phi_w)=y_w\). Indeed \(\alpha\) is unramified and
the image of the principal cyclotomic character is
\((1+\ell)^{\ell^a\ZZ_\ell}\). In particular
\(d(g)\equiv1\bmod\varpi^{h_w}\) for every \(g\).

Use the exponents \(e_i=1,0,-1\). Write the \(i\)-th
local character as \(\lambda^{e_i}(1+z_i)^n\), with \(n\)
the unramified coordinate used above. Modulo quadratic terms
its variation is \(\lambda(g)^{e_i}n(g)z_i\).
The image of \((\tau,d,n)\) is the direct product of
\(\operatorname{im}\tau\) with the image of \((d,n)\).
In fact the first group has finite prime-to-\(\ell\) order
and the second is pro-\(\ell\). Powers congruent to one
modulo the first order and to zero modulo arbitrarily large
powers of \(\ell\) isolate the first group in the compact
image. Complementary multiplication isolates the second.
We can therefore average the linearized trace relations
over \(\operatorname{im}\tau\) using its character idempotents
over \(O\). This separates precisely the stated blocks.
For a block \(B\) the relations are the rows
\[
              n(g)\bigl(d(g)^{e_i}\bigr)_{i\in B}
\]
acting on its generators \(z_i\). Comparing traces at the
second and third powers of each local element gives all
characteristic coefficients, as \(3!\) is invertible.

There are three available values of \(d(g)\), all with
\(n(g)\) a unit, whose pairwise differences have valuation
exactly \(h_w\). When \(v_\varpi(x-1)\le v_\varpi(y_w-1)\)
use \(y_w,y_w x,y_w x^2\), obtained with \(n=1\) by varying
inertia. Otherwise use \(y_w,y_w^2,y_w^3\), obtained on
the first three powers of \(\phi_w\). Here \(1,2,3\) are
units, and the difference between two consecutive powers
has the same valuation as the difference two powers apart,
since the sum of a principal unit and one is a unit.

A singleton block has a unit row entry and contributes zero.
A block of size two occurs when \(\bar\lambda\) has order two,
and its exponents are \(1,-1\). Multiplying each row by the
unit \(d(g)\) gives the two columns \(d(g)^2,1\), with row
factor \(n(g)\). Every two-row minor is divisible by
\(\varpi^{h_w}\), and two of the displayed values attain
that valuation since their sum is a unit. Its nonunit
elementary divisor is therefore \(\varpi^{h_w}\).
For a block of size three the same unit scaling gives the
columns \(d(g)^2,d(g),1\). Every two-row minor has valuation
at least \(h_w\) and every three-row determinant at least
\(3h_w\), by the Vandermonde product and the row factors.
The displayed choices attain both minima. The entries
generate the unit ideal. Smith normal form now gives the
three valuations \(0,h_w,2h_w\). These minor computations
apply to the whole span of the rows, since it is a submodule
of a finite free module over a DVR and its minor ideals
are generated by the minors of the generating rows.

It remains to check the ordinary word at first order.
On the \(i\)-th split global line the \(i\)-th factor has
zero constant term, so only its own variation can contribute.
This variation is a multiple of \(z_i\). Each other index
in its block supplies a factor divisible by \(\varpi^{h_w}\),
because the Teichm\"uller characters agree and the remaining
values are congruent to one modulo that power. Thus the linear
term is a multiple of \(\varpi^{(|B|-1)h_w}z_i\), already
zero by the elementary divisors. The global polarization
and tame equations hold on the fixed determinant. Local
torsion values are fixed by inertia. The independent places
therefore give the direct sum asserted.

Quotienting the original deformation algebra through its
split section gives a quotient on the fibers with all traces
fixed to their split values. Their augmentation cotangent
modules are the respective cokernels in the corollary, by
the same ideal computation over \(O\). This gives the length
bound. In Proposition~\ref{prop:dyadiccollision438} each
place has one block, \(y_w=1\) and \(r=\ell-1\), giving
\(h_w=(a+1)v_\varpi(\ell)\) as claimed.
\end{proof}

These calculations concern a specific cost in passing from
global traces to local character labels. In particular the
positive-dimensional fiber of the word equations cannot
supply a finiteness argument at the split trace for these
characters. Recovering the total strict cohomology length
requires control of the global trace relations that detect
extensions. The signed period on the prescribed analytic row
in Section~\ref{sec:missing} remains to be compared with the
integral arithmetic period on each \(E_x\).


\subsection{Reducible traces when inertia is scalar on the descended plane}
\label{sec:scalarinertiaplanes439}

Fix one of the common intervals of Theorem~\ref{thm:conductor232}.
Use the same auxiliary quadratic field \(K=\QQ(\sqrt{-s})\),
characters \(\alpha\) of \(E_x\), sign condition, and twists
\(2\le r<t(x)\) as in Section~\ref{sec:twoprimeadjoints436}.
Take \(5\le\ell\le L_*\) prime to \(2sN_*\) and suppose
\(\ell-1\mid 2r\). Put \(\Psi=\alpha\kappa_\ell^r\). Bars
denote reduction in the residue field \(k\) of a coefficient
ring \(O\) containing the character values. We use Pan's two
ordinary finiteness theorems over \(\QQ\). Passing from the
plane over \(K\) to a representation over \(\QQ\) requires
checking its inertia before either theorem can be applied.

\begin{lemma}[A stable local line in the descended plane]\leavevmode
\label{lem:unipotentplane439}
Let \(V\) be a continuous two-dimensional representation of
\(G_{\QQ_\ell}\) over a finite extension of \(k((z))\).
Suppose every inertia element has both eigenvalues one.
After a finite extension of coefficients \(V\) has a
\(G_{\QQ_\ell}\)-stable line. The characters \(\xi_1,\xi_2\)
on the line and quotient are unramified. For any global
representation extending \(V\) its trace satisfies
\[
  \operatorname{tr}\bigl((V(\sigma)-\xi_1(\sigma)I)
                  (V(\tau)-\xi_2(\tau)I)M\bigr)=0
\]
for all local \(\sigma,\tau\) and any matrix \(M\) of the
global representation.
\end{lemma}
\begin{proof}
If inertia has a nonidentity element \(u\), choose a basis
in which \(u=I+E_{12}\). For every other inertia matrix \(A\),
both \(uA\) and \(A\) have trace two, which gives \(A_{21}=0\).
The diagonal entries of \(A\) are then both one. Consequently
the common inertia fixed space is the first line, unique
because \(u\ne I\). Normality of inertia makes this line
stable under the decomposition group. If inertia is trivial,
choose an eigenline of an arithmetic Frobenius lift after
splitting its characteristic polynomial. The unramified
quotient is procyclic and the stabilizer of a line is closed,
so this line also suffices. In both cases inertia acts
trivially on the two quotients. The second matrix difference
in the formula maps into the line, and the first kills the
line. This proves the identity, including after any
multiplication and trace.
\end{proof}

Consider lifts on \(G_K\) of the split residual diagonal
\((\bar\Psi,1,\bar\Psi^{-1})\), unramified outside a finite
set containing the primes over \(2sN_*\ell\). Use the
symmetric adjunction in a fixed basis from
Proposition~\ref{prop:odddescent376}. Thus the matrix of
\(cg^{-1}c\) is the adjoint of the matrix of \(g\) for the
fixed invertible symmetric antidiagonal form, with the
common cyclotomic twist removed. Here \(c\) is a complex
conjugation. Away from \(\ell\) impose finite inertia image
by fixing inertia kernels contained in the kernels of the
split type, of finite index in the local inertia groups.
For instance one may use the kernels of \(\alpha\) on
inertia there. At each \(v\mid\ell\) impose only
\[
  \operatorname{char}_{M(\sigma)}(X)
     =(X-\Psi(\sigma))(X-1)(X-\Psi(\sigma)^{-1})
          \qquad(\sigma\in I_v).
\]
These are closed equations in the framed matrix deformation
ring and hold on all actual flags of the original weights.
Let \(T\) be the closed \(O\)-algebra of traces of the
resulting universal matrices. Further matrix markings or
flag conditions give trace images which are quotients of
\(T\). Let \(I_m\subset T\) be generated by the size-\(m\)
minors of \((t(g_i h_j))\), where \(g_i,h_j\in G_K\).

\begin{theorem}[Finiteness of the entire reducible trace quotient]\leavevmode
\label{thm:scalarinertiaplanes439}
Suppose \(\bar\Psi|_{G_K}\) has order greater than two or
is trivial. For \(6\le m\le9\), \(T/I_m\) is a finite
\(O\)-module. Its geometric field points are exactly the
points of \(T\) with reducible semisimple representation.
In particular this holds for every globally distinct residual
triple on \(G_K\) in the stated range, regardless of the
values of \(\alpha(\ell)\) or \(\eta_K(\ell)\).
\end{theorem}
\begin{proof}
The framed ring is a quotient of the usual complete matrix
deformation ring unramified outside the indicated set.
Finiteness of residual first cohomology follows by Kummer
theory, as in Theorem~\ref{thm:reduciblecollision435}.
The complete determinant ring is therefore Noetherian and
its Cayley--Hamilton algebra is module finite. Its image in
the framed ring is \(T\), since \(3!\) is a unit and the
determinant coefficients are recovered from traces. The
image is closed by compactness. The span of group matrices
is finite as a \(T\)-module by the same Cayley--Hamilton
theorem.~\cite{XI-VIII-chenevier364}

At a geometric field point the semisimple group span is a
product of matrix algebras acting on the distinct simple
constituents. Its trace pairing is nondegenerate since it
is the ordinary matrix trace pairing on each factor times a
multiplicity between one and three. The span has dimension
nine in the absolutely irreducible case and at most five
otherwise. Group elements span it in both variables of the
pairing, which proves the minor assertion. A semisimple
realization factors through the specialized integral matrix
span. Indeed an identity of group matrices has trace zero
against all group matrices and therefore acts as zero on
the semisimple realization, by the nondegeneracy just proved.
The adjunction also passes to this semisimple span since its
trace radical is preserved.

It suffices to test a continuous map from \(T/I_m\) to
\(k'[[z]]\) with \(k'/k\) finite. Over a finite extension of
its fraction field take the semisimple realization just
described. Its continuity and the existence of a stable
lattice follow from the finite integral matrix span as in
Theorem~\ref{thm:reduciblecollision435}. Its residual
semisimplification is the prescribed diagonal. Character
values in residue fields are read as constants in the
equal-characteristic coefficient DVR, enlarged when needed.

Each one-dimensional constituent is the constant lift of its
reduction. In fact for any \(\sigma\in I_v\), \(v\mid\ell\),
its value is one of the three constant roots of the displayed
inertia polynomial in characteristic \(\ell\), hence is
determined by its reduction. At other places inertia has
finite image. A finite-order unit in an equal-characteristic
DVR with finite residue field is its constant residual lift:
the principal units are pro-\(\ell\) and have no nontrivial
\(\ell\)-power torsion, by the power identity in characteristic
\(\ell\). The character divided by its constant residual lift
is consequently everywhere unramified. The class group of
\(K\) is finite, so its image in these principal units is
trivial. The same reasoning for the determinant gives
determinant one, since its inertia at \(\ell\) and its
reduction are trivial. This treats a sum of three lines.

In the remaining case there is a unique line constituent
\(\lambda\) and an absolutely irreducible plane \(P\).
Polarization gives \(\lambda^c=\lambda^{-1}\).
Every possible residual character on this line extends to
\(G_\QQ\), hence is fixed by \(c\) on \(G_K\). The order
hypothesis forces \(\bar\lambda=1\): for order greater
than two neither outer character squares to one, and in
the trivial case all three are one. The preceding argument
gives \(\lambda=1\) and \(\det P=1\). In particular the
residual semisimplification of the plane on \(G_K\) consists
of \(\bar\Psi,\bar\Psi^{-1}\) with their multiplicities.

We give the odd descent explicitly. The adjunction fixes the
plane factor in the semisimple span. On a full matrix algebra
over a splitting field an involution fixing scalars is adjoint
to a symmetric or alternating form, by composing it with
transpose and using innerness of a matrix algebra automorphism.
Proposition~\ref{prop:odddescent376} gives the polynomial
identity \(2XYX=t(XY)X\) for skew elements of the original
rank-three symmetric matrix algebra. Skew elements in its
semisimple quotient lift by antisymmetrization, since two
is invertible. The identity rules out alternating adjunction
on the plane: in such a basis \(X=\operatorname{diag}(1,-1)\)
and \(Y=E_{12}+E_{21}\) are skew, but \(XYX=-Y\) and
\(\operatorname{tr}(XY)=0\). Both elements have zero line
component here.

Let \(S\) be the resulting symmetric form matrix on the plane
and \(J\) its standard alternating form, preserved by \(P\)
because the determinant is one. The adjunction identity
gives \(P(cgc)=S^{-1}P(g)^{-{\mathsf T}}S\).
The matrix \(S^{-1}J\) therefore intertwines \(P\) with its
conjugate by \(c\). It is invertible with trace zero, so
after a finite scalar extension it can be rescaled to square
one, with determinant \(-1\). Let \(c\) act by this matrix.
This extends \(P\) to a continuous representation \(W\)
of \(G_\QQ\) of determinant \(\eta_K\).

After taking a stable lattice, every irreducible constituent
of its residual semisimplification restricts semisimply to
\(G_K\): the translates of a simple submodule for the normal
subgroup span a semisimple module invariant under \(G_\QQ\).
When the two restricted characters are distinct they already
extend to \(G_\QQ\), so both isotypic lines in a purported
irreducible residual plane would be stable. When they are
equal, twisting off their global extension would give an
irreducible residual plane of a group of order two, also
impossible. Thus the residual characters on \(G_\QQ\) are
\(\bar\Psi,\bar\eta_K\bar\Psi^{-1}\), after twisting \(W\)
by \(\eta_K\) if necessary, using its determinant to fix
their product. In the repeated case the restrictions are
trivial and the determinant gives the pair \(1,\bar\eta_K\).

Now set \(V=W\bar\Psi^{-1}\), using the global constant
lift of \(\bar\Psi\) in characteristic \(\ell\).
Since \(K/\QQ\) is unramified at \(\ell\), its local inertia
is the same as that over \(\QQ_\ell\). The characteristic
polynomial on the plane there is the prescribed polynomial
with the line factor \(X-1\) removed. Its two roots are
\(\bar\Psi(\sigma),\bar\Psi(\sigma)^{-1}\).
They agree for every inertia element by \(\ell-1\mid2r\)
and \(\ell\nmid N_*\). This proves the unipotent inertia
hypothesis of Lemma~\ref{lem:unipotentplane439} for \(V\).
It gives two unramified local characters and a stable line
over a finite extension. The compact image preserves a
lattice, whose intersection with the line is saturated.
Thus both characters take values in the coefficient DVR.

Apply Pan's finiteness theorems to the rational field with
fixed determinant \(\chi=\eta_K\Psi^{-2}\). This character
is odd, unramified outside a finite set, and de Rham at
\(\ell\), being a finite character times a cyclotomic
power. The residual trace of \(V\) is \(1+\bar\chi\).
If \(\bar\chi|_{G_{\QQ_\ell}}\ne1\), the local trace is
the sum of two characters with the prescribed distinct
reductions. Pan's Theorem~5.1.2(1) applies to its ordinary
pseudodeformation. If \(\bar\chi|_{G_{\QQ_\ell}}=1\), mark
the actual stable line of the lemma. Its character reduces
to one, the quotient character is the determinant divided
by that character, and the matrix identity of the lemma
is exactly the ordered trace identity in Pan's Definition~5.3.2.
These verify the additional hypotheses for his
Theorem~6.1.2(1). In either case the applicable ring is
finite over \(O[[1+\ell\ZZ_\ell]]\), after finite coefficient
extension if needed. Its weight character is unramified on
our curve, since both local characters are unramified.
The map to the equal-characteristic coefficient DVR therefore
kills the weight augmentation ideal and \(\varpi\). The
image is a finite field. Undoing the constant twist, adjoining
the trivial line and restricting to \(G_K\) proves that all
original traces in \(k'[[z]]\) are constant.~\cite{XI-VIII-VI-pan40}

If the special fiber of \(T/I_m\) had positive dimension, a
one-dimensional complete local domain quotient and its
normalization would give just such a curve with nonconstant
trace. Indeed the normalization is a complete DVR finite
over \(k[[z_0]]\) by complete Noether normalization, and
the traces topologically generate the quotient. Thus the
special fiber is Artinian. A power of the maximal ideal
lies in \(\varpi(T/I_m)\); completeness and lifting a basis
of the special fiber show that \(T/I_m\) is a finite
\(O\)-module, including its torsion.
\end{proof}

\begin{corollary}[Integral finiteness with the characters of a flag]\leavevmode
\label{cor:flagfinitered439}
In the range of the theorem, adjoin to the trace image of
any of these fixed-weight symmetric matrix problems the
three ordered local characters of an actual flag at each
chosen place over \(\ell\). More generally one may adjoin
such labels with the local characteristic-polynomial
identities \eqref{eq:factorlocal438}. The quotient of the
resulting determinant and character algebra by \(I_m\) is
finite over \(O\) for \(6\le m\le9\). The conclusion includes
the characters and twists of Proposition~\ref{prop:dyadiccollision438}
on the same period intervals.
\end{corollary}
\begin{proof}
All inertia values are fixed and each local unramified
quotient is procyclic. Therefore only the ordered Frobenius
values must be adjoined. Their local cubic has coefficients
in the trace image and invertible constant coefficient.
Successive monic division gives at most six module generators
per place, as in Corollary~\ref{cor:flagcotangent438}, with
no assumption on the differences of its roots. This finite
span is compact and contains the continuous character
values, so it also contains their closed algebra. Reducing
by \(I_m\) leaves a finite module over a quotient of \(T/I_m\).
The order-eight character of the cited proposition retains
order at least four on \(G_L\), hence also on \(G_K\).
Its twist obeys the collision congruence. All hypotheses
therefore hold for that same character and field choice.
\end{proof}

Finiteness here controls the reducible part of the trace
problem at a fixed weight. An interval bound for the entire
strict group still requires a quantitative bound on the
integral extension classes, including their intersections with
the irreducible locus. The period comparison in
Section~\ref{sec:missing} uses those same \(E_x\).



\subsection{Quadratic residual characters on one real subfield of the interval}
\label{sec:realquadratictrace442}

Keep the symmetric rank-three trace algebra \(T\), the ideals
\(I_m\), the imaginary quadratic field \(K=\QQ(\sqrt{-s})\)
and the finite inertia conditions of
Section~\ref{sec:scalarinertiaplanes439}. Take the same interval
with largest period field \(E_*\) and conductor divisor \(N_*\).
We treat the case in which the two outer residual characters
agree globally and are nontrivial. When \(\ell-1\mid r\), they can
be removed over a subfield of \(E_*\) selected before \(x\) or
\(\ell\) is chosen. A trace identity brings the conclusion back
to \(G_K\).

\begin{lemma}[Integral traces over a quotient of exponent two]\leavevmode
\label{lem:quadratictrace442}
Let \(\rho:G\to\mathrm{GL}_3(A)\), where \(A\) is a commutative
ring in which \(6\) is invertible. Let \(\Gamma\) be normal with
\(G/\Gamma\) of exponent dividing two. For \(g\in G\), denote the
elementary characteristic coefficients of \(\rho(g^2)\) by
\(s_1,s_2,s_3\), in order. Then \(t=\operatorname{tr}\rho(g)\)
satisfies
\[
       \bigl((t^2-s_1)^2-4s_2\bigr)^2-64s_3t^2=0.
\]
All three coefficients \(s_i\) belong to the
\(\mathbb Z[1/6]\)-algebra generated by traces on \(\Gamma\).
In particular, on a curve over a finite
field, constancy of the restricted traces implies constancy of
every trace. If \(B\) is the closed restricted trace algebra and
is finite over a complete coefficient DVR \(O\) with finite residue
field, a Hausdorff trace algebra topologically generated over \(B\)
by \(d\) further traces is generated by at most \(8^d\) elements
as a \(B\)-module.
\end{lemma}
\begin{proof}
Write \(t,b,\delta\) for the three elementary coefficients of
\(\rho(g)\). The elementary symmetric polynomials of the three
squared roots give \(s_1=t^2-2b\), \(s_2=b^2-2\delta t\) and
\(s_3=\delta^2\). These are polynomial identities of matrices over
any commutative ring. Indeed they hold for the universal matrix
over its characteristic-zero rational function field by
triangularization over a splitting field, hence as identities
in the integral polynomial ring in the entries. Substitution
and squaring give the formula. Newton's identities express
\(s_1,s_2,s_3\) by traces of \(\rho(g^2),\rho(g^4),\rho(g^6)\)
with denominators dividing six. All three group elements belong
to \(\Gamma\).

In an equal-characteristic complete DVR with finite residue
field, an element algebraic over a finite constant field is a
constant. Indeed it belongs to a finite field and satisfies
\(X^{q}=X\) for a power \(q\) of the characteristic. Reduction is
injective on these roots since two roots with the same residue
have a difference of positive valuation equal to its \(q\)th
power. The unique lifts of the residue field inside the coefficient
DVR give those constants. Apply this to the monic equation for
\(t\). Finally powers of each additional trace reduce to powers
of degree at most seven by its monic equation. The span of the
resulting \(8^d\) monomials is compact as the image of a finite
power of \(B\). It is closed and contains the algebraic generators
and all their polynomial combinations. It contains their closure.
\end{proof}

Define \(H\subset E_*\) as the compositum of the real quadratic
subfields of \(E_*\), taking \(H=\QQ\) if the list is empty.
Put \(L=KH\) and write \(\eta\) for the quadratic character of
\(L/H\).

\begin{proposition}[Odd descent after the residual character is killed]\leavevmode
\label{prop:evenfieldplane442}
The field \(H\) is abelian and totally real, of degree at most
\(16\), and is unramified at every \(\ell\nmid N_*\).
Suppose \(5\le\ell\le L_*\), \(\ell\nmid 2sN_*\) and
\(\ell-1\mid r\). Let \(\Psi=\alpha\kappa_\ell^r\) with the sign
and field of Section~\ref{sec:twoprimeadjoints436}.
Suppose \(\bar\Psi|_{G_K}\) has order at most two. For every
equal-characteristic DVR curve of the reducible trace quotient
in Section~\ref{sec:scalarinertiaplanes439}, all traces on
\(G_L\) are constant.
\end{proposition}
\begin{proof}
The group of \(E_*/\QQ\) embeds into a product of four cyclic
groups by the description of \(E_*\) in Theorem~\ref{thm:realtwist404}.
A subgroup of such a product has at most four generators.
Projection onto one cyclic factor supplies a cyclic image. A lift
of a generator, together with generators for the kernel obtained
by induction, generates the subgroup. Its group of quadratic
characters therefore has order at most \(2^4\). The even quadratic
characters cut out \(H\), giving all the asserted field properties.
Since \(K\) is imaginary quadratic, \(H\cap K=\QQ\), and
\(\operatorname{Gal}(L/K)\) is elementary abelian of exponent at
most two.

The character \(\bar\Psi\) equals \(\bar\alpha\) on \(G_\QQ\) by
the congruence on \(r\). It is even by the sign condition.
Consequently its fixed field is totally real, is contained in
\(E_x\), and meets \(K\) in \(\QQ\). Restricting it to \(G_K\)
preserves its order. It follows that \(\bar\Psi|_{G_L}=1\).
In particular, the three residual characters are all trivial
on \(G_L\).

Take a semisimple realization on the curve, passing to a finite
extension of the coefficient fraction field as in
Theorem~\ref{thm:scalarinertiaplanes439}. The representation
over \(G_K\) is reducible and has a stable lattice. Its restriction
to \(G_L\) is semisimple over a splitting field. Indeed the
translates of any simple \(G_L\)-submodule in a simple
\(G_K\)-summand span that entire summand by normality and
irreducibility. A sum of simple submodules is semisimple.
Thus the restricted representation consists of three lines or
of a line and an absolutely irreducible plane.

Every one-dimensional constituent on \(G_L\) is trivial. To see
its reduction, filter a stable lattice by intersection with a
composition series on the generic fiber. The successive lattice
quotients are torsion free. Reduction of this filtration preserves
the composition factors with multiplicity; the residual
semisimplification restricted from \(G_K\) has only trivial
constituents. The line thus reduces to one, and it has finite inertia image away from
\(\ell\). At primes over \(\ell\) each inertia matrix has
characteristic polynomial \((X-1)^3\) in characteristic
\(\ell\), by the original inertia equations. Its character is
therefore one on this inertia. A finite-order principal unit in
the coefficient DVR equals one, since the units congruent to one
are pro-\(\ell\) and are torsion free in characteristic \(\ell\).
At the other primes this proves the same triviality. The remaining
character factors through the finite class group of \(L\) by
global class field theory and is again a principal unit, so is
trivial. The same argument gives trivial determinant on \(G_L\).
This settles the sum of three lines.

Write \(P_0\) for the plane in the remaining case. It has
determinant one. Choose a complex conjugation \(c_0\) in \(G_H\).
It acts nontrivially on \(K\) and \(G_H=G_L\rtimes\langle c_0\rangle\).
We may use the original conjugation of the marked symmetric
problem for \(c_0\), since every complex conjugation fixes the
totally real field \(H\). The adjunction preserves the simple
plane factor, the unique factor of matrix size two. On this
factor it is adjoint to a symmetric form. For clarity, the
polynomial identity of Proposition~\ref{prop:odddescent376}
on skew elements passes to the span of \(G_L\) and its semisimple
image by antisymmetrizing lifts. Alternating adjunction on the
plane would make both
\(X=\operatorname{diag}(1,-1)\) and \(Y=E_{12}+E_{21}\) skew.
With their line components zero the identity would say
\(2XYX=\operatorname{tr}(XY)X\), whereas \(XYX=-Y\) and
\(\operatorname{tr}(XY)=0\).

Let \(S\) be the matrix of this symmetric form and \(J_2\)
the alternating matrix with upper-right entry one. Polarization
gives \(P_0(c_0gc_0)=S^{-1}P_0(g)^{-{\mathsf T}}S\);
determinant one gives \(P_0(g)^{-{\mathsf T}}=J_2P_0(g)J_2^{-1}\).
The intertwiner \(S^{-1}J_2\) is invertible of trace zero.
After scaling over a finite extension it squares to the identity
and has determinant \(-1\). Using it as the value at \(c_0\)
extends the plane to a continuous representation \(W\) of \(G_H\)
with determinant \(\eta\). The sum of a \(G_L\)-stable lattice and
its translate by \(c_0\) is a \(G_H\)-stable lattice. On each
irreducible residual constituent the restriction to the normal
subgroup \(G_L\) is semisimple by the translate argument above.
The residual semisimplification on \(G_L\) is trivial, so these
constituents factor through the order-two quotient. The determinant
then gives residual trace \(1+\bar\eta\).

For every \(v\mid\ell\) of \(H\), the extension \(L/H\) is
unramified there because \(K/\QQ\) is unramified at \(\ell\).
Thus \(W\) has unipotent inertia at \(v\). If there is a
nonidentity unipotent, choose a basis in which it is \(I+E_{12}\).
The trace of its product with any inertia matrix equals the trace
of that matrix. Hence the lower-left entry of every inertia matrix
is zero. All inertia matrices fix the first line pointwise, and
their common fixed line is unique and stable under the decomposition
group. If inertia is trivial choose an eigenline of a Frobenius
lift after a finite extension. These arguments give a stable line
and two unramified local characters simultaneously at all
\(v\mid\ell\). The induced lattice on the line is saturated.

We apply the finiteness assertions of Pan with base field \(H\)
and fixed characteristic-zero determinant
\(\eta=\eta_K|_{G_H}\), using the finite set of primes containing
the ramification of the representation and all primes above
\(\ell\). The determinant extends from \(\QQ\), is finite order
and de Rham at \(\ell\), and is odd at every real place of \(H\).
The coefficient prime is unramified in this abelian totally real
field. All places \(v\mid\ell\) have the same local alternative
for \(\bar\eta\) because they are conjugate over \(\QQ\) and \(\eta\)
is the restriction of a rational character. If their local
character is nontrivial, the local trace is the sum of characters
with distinct reductions, one trivial, and Theorem~5.1.2(1)
of Pan applies. If they are all trivial, mark the local stable
lines just constructed. For local elements \(\sigma,\tau\), the
matrix \(W(\tau)-\xi_2(\tau)I\) maps into the line with character
\(\xi_1\), where \(\xi_2\) is the quotient character. Left
multiplication by \(W(\sigma)-\xi_1(\sigma)I\) kills this image.
Taking the trace after any further matrix multiplication gives
the ordered identity in Pan's Definition~5.3.2. His
Theorem~6.1.2(1) therefore applies in this case.
Both theorems give a ring finite over the product weight algebra,
the completed tensor product over \(O\) of the completed group
algebras of the pro-\(\ell\) unit groups at \(v\mid\ell\).
All weight characters on the curve are trivial on those units by
local reciprocity and unramifiedness. The map into the coefficient
DVR kills the augmentation ideal of the weight algebra and the
coefficient uniformizer. Its image is a finite domain over the
residue field, hence a field of constants. This proves constancy
of the traces of \(W\), and restriction to \(G_L\) and addition
of the trivial line prove the proposition.~\cite{XI-VIII-VI-pan40}
\end{proof}

\begin{theorem}[The reducible trace quotient at twists divisible by \(\ell-1\)]\leavevmode
\label{thm:allzeroweights442}
At every \(x\) of the same interval and every original character
of sign \((-1)^r\), let \(5\le\ell\le L_*\) be prime to \(2sN_*\).
If \(2\le r<t(x)\) and \(\ell-1\mid r\), the full reducible
quotient \(T/I_m\), \(6\le m\le9\), is finite over \(O\).
This includes the scheme structure at the split trace and the
case of repeated global residual characters. The corresponding
quotient with the local characters of genuine flags adjoined is
also finite over \(O\).
\end{theorem}
\begin{proof}
If \(\bar\Psi|_{G_K}\) has order greater than two,
Theorem~\ref{thm:scalarinertiaplanes439} applies directly.
In the remaining cases the proposition makes every restricted
trace constant on each equal-characteristic curve through the
split trace. Apply Lemma~\ref{lem:quadratictrace442} with
\(G=G_K\) and \(\Gamma=G_L\). Each original trace is constant.
The rank characterization of \(I_m\) in the proof of
Theorem~\ref{thm:scalarinertiaplanes439} uses only dimensions
of semisimple matrix spans, so those are exactly the curves
that must be considered here.

If the complete Noetherian local special fiber had positive
dimension, divide it by a prime of dimension one and normalize.
This gives a complete DVR finite over a power series ring in one
variable over a finite field. The map from the special fiber to
this DVR has one-dimensional image, whereas every trace in the
image is constant. The traces topologically generate and the
coefficient residue field is finite, a contradiction. The
special fiber is therefore Artinian. A power of the maximal
ideal of \(T/I_m\) is contained in \(\varpi(T/I_m)\); lifting
a basis modulo \(\varpi\) and iterating the expansion proves
module finiteness by completeness. This step retains nilpotent
elements and coefficient torsion. For a genuine flag each local
Frobenius character value is a root of the local characteristic
polynomial. Successive monic division bounds the rank of the
ordered-root algebra by six per place, with fixed inertia values,
as in Corollary~\ref{cor:flagfinitered439}. This proves the last
assertion.
\end{proof}

The field \(H\) works for the whole interval. The theorem proves
finiteness at these weights. The bound in Section~\ref{sec:missing}
requires a bound on the lengths of strict extension classes as the
field, twist and coefficient prime vary, including their
contributions on irreducible components.


\subsection{A ramified plane at an inert coefficient prime}
\label{sec:inertquadraticplane443}

We return to the second congruence for coincident weights in
Section~\ref{sec:scalarinertiaplanes439}. Suppose
\(\ell-1\mid 2r\) and \(\ell-1\nmid r\), with \(\ell\ge5\)
prime to \(2sN_*\). Write \(\omega=\bar\kappa_\ell\) and suppose
the even character \(\beta=\bar\alpha\omega^r\) has order two.
Its restriction to inertia at \(\ell\) is nontrivial. Indeed
\(r\equiv(\ell-1)/2\bmod(\ell-1)\), \(\alpha\) is unramified there
and the local cyclotomic character maps inertia onto
\(\mathbb F_\ell^\times\). We shall impose \(\eta_K(\ell)=-1\),
so that \(K\) has a unique prime over \(\ell\).

\begin{lemma}[The full decomposition group of the plane]\leavevmode
\label{lem:inertplane443}
Let \(P_0\) be the absolutely irreducible plane in a reducible
semisimple representation on an equal-characteristic curve of
the symmetric trace problem. Suppose that its line constituent
has residual character \(\beta|_{G_K}\). At the unique place
\(w\mid\ell\) of \(K\) one has, after a finite coefficient extension,
\[
                   P_0|_{G_{K_w}}=1\oplus\beta|_{G_{K_w}}
                     \quad\text{up to isomorphism}.
\]
\end{lemma}
\begin{proof}
The line is its constant residual lift. At \(\ell\), every inertia
value is a root of the prescribed characteristic polynomial and
all roots are constants in characteristic \(\ell\). At other places
the inertia image is finite. After division by the residual
character, the resulting principal-unit character factors through
the finite class group, as in Theorem~\ref{thm:scalarinertiaplanes439}.
The same argument gives determinant one for the entire
representation. The plane thus has determinant \(\beta|_{G_K}\)
and residual semisimplification \(1\oplus\beta|_{G_K}\).

Write \(c_0\) for the complex conjugation in the symmetric
adjunction. The adjunction on the plane is symmetric by the
skew-matrix argument in Proposition~\ref{prop:evenfieldplane442}.
With form matrix \(S\) and standard alternating matrix \(J_2\), put
\(B=S^{-1}J_2\). The two-dimensional alternating identity gives
\[
        P_0(c_0 g c_0)=\beta(g)\,B P_0(g)B^{-1}\qquad(g\in G_K).
\]
Consequently for any \(h\in G_\QQ\setminus G_K\), conjugation by
\(h\) on the group arguments gives a representation isomorphic to
\(\beta P_0\). Changing \(c_0\) to \(h=g_0c_0\) merely changes
the intertwiner to \(P_0(g_0)B\).

In a fixed rational decomposition group at \(\ell\), let \(D\)
be \(G_{K_w}\) and \(I\) its inertia. The extension is unramified
quadratic, so \(I\) is also the rational inertia group. Choose
\(h\) in this rational decomposition group outside \(D\).
It belongs to the nontrivial global coset, normalizes \(D,I\),
and preserves \(I_0=\ker(\beta|_I)\). Every element of \(I_0\)
acts unipotently on the plane. If there is a nonidentity such
element, the trace-two calculation in
Lemma~\ref{lem:unipotentplane439} gives a unique line fixed by
all of \(I_0\). It is stable under \(I\), with character either
\(1\) or \(\beta|_I\), since \(I/I_0\) has order two. Conjugating
the arguments by \(h\) preserves this character. Twisting by
\(\beta\) changes it to the other character. The two
representations of \(I\) would therefore be isomorphic with
different characters on their unique \(I_0\)-fixed lines,
a contradiction. Hence \(P_0(I_0)=1\).

The nontrivial coset of \(I/I_0\) has distinct eigenvalues
\(1,-1\) on the plane by the original inertia polynomial after
dividing out the line. The two eigenspaces are stable under \(D\),
which acts trivially on this order-two quotient. Thus
\(P_0|_D=\xi\oplus\beta\xi^{-1}\) with \(\xi\) unramified.
Conjugation by \(h\) fixes \(\xi\), as \(D/I\) lies in the abelian
unramified quotient of the rational decomposition group. It fixes
\(\beta\) as well. In \(\beta P_0\) the unique unramified
constituent is \(\xi^{-1}\). The isomorphism on \(D\) gives
\(\xi^2=1\). Its reduction is the trivial constituent of
\(1\oplus\beta|_D\), since the latter's other constituent
ramifies on inertia. To see this directly on a stable lattice,
the idempotents \((1\pm P_0(i))/2\) for \(i\in I\setminus I_0\)
split it integrally. A quadratic character reducing to one in
odd characteristic is trivial. This proves the formula on all of
\(D\).
\end{proof}

\begin{lemma}[A finite twist of the polarized plane]\leavevmode
\label{lem:quadtwist443}
There is a finite-order character \(\theta\) of \(G_K\) with
values of two-power order in a finite extension of the constant
field such that
\(\theta^{c_0}/\theta=\beta|_{G_K}\).
For the plane of Lemma~\ref{lem:inertplane443},
\(\theta P_0\) extends to an odd representation \(W\) of
\(G_\QQ\). Its residual representation is absolutely irreducible,
is induced from \(\theta\) over \(K\), and has projective fixed
field \(KH_\beta\), where \(H_\beta\) is the real quadratic
field cut out by \(\beta\).
\end{lemma}
\begin{proof}
Use the finite id\`ele class quotient
\(A=\mathbb A_{K,f}^\times/K^\times\). It is profinite.
Indeed the valuations map it onto the finite ideal class group
with open kernel \(\widehat{\mathcal O}_K^\times/\mathcal O_K^\times\).
The latter is profinite since an imaginary quadratic field has
finite units. The diagonal \(K^\times\) is closed in the finite
id\`eles, since in each open fiber of the valuation map its
intersection consists of the generators of a single fractional
ideal and hence is finite. This gives the Hausdorff quotient
used here. By class field theory, finite characters of \(A\)
are characters of \(G_K\). Restriction of \(\beta\) to \(G_K\)
corresponds to \(\beta\circ N_{K/\QQ}\) on the id\`eles.

The \(c_0\)-fixed subgroup of \(A\) is the image of the rational
finite id\`eles. To see this, if \(x\) represents a fixed class,
\(\overline x/x=b\in K^\times\) has norm one. Hilbert's
Theorem 90 gives \(b=\overline y/y\) for \(y\in K^\times\).
The id\`ele \(x/y\) is fixed in each conjugate system of local
fields, hence is rational. On these fixed classes
\(\beta\circ N_{K/\QQ}\) is trivial, since the norm of a rational
id\`ele is its square. On the closed subgroup
\(D_0=(c_0-1)A\) the rule
\[
        \overline x/x\longmapsto \beta(N_{K/\QQ}x)
\]
therefore defines a character. It is continuous since the
continuous surjection from compact \(A\) to Hausdorff \(D_0\)
is a quotient map. Take an open subgroup \(U\) of \(A\) whose
intersection with \(D_0\) is in its kernel. Its character on the
image of \(D_0\) in \(A/U\) extends to the finite abelian group
\(A/U\) with values of two-power order. For example, extend
one cyclic quotient generator at a time. If its order modulo
the previous subgroup is \(e\), prescribe a value whose \(e\)th
power is the value already assigned to its \(e\)th power.
Every such root can be chosen in the group of all two-power
roots of unity in \(\overline{\mathbb F}_\ell\), since raising to
any positive integer is surjective on that group. This gives
\(\theta\) and its asserted conjugation law.

The identity of Lemma~\ref{lem:inertplane443} now says that
\((\theta P_0)^{c_0}\) is conjugate to \(\theta P_0\) by \(B\).
The matrix \(B=S^{-1}J_2\) has trace zero and is invertible.
Rescale it over a finite extension to have square one and
determinant \(-1\), and let \(c_0\) act by this matrix. The
semidirect product \(G_\QQ=G_K\rtimes\langle c_0\rangle\)
gives \(W\), continuous and odd. Summing a \(G_K\)-stable
lattice and its translate by \(c_0\) supplies a stable lattice.

The two residual characters on \(G_K\) are \(\theta,\theta\beta\).
They are distinct and \(c_0\) exchanges them. A reducible residual
representation on \(G_\QQ\) would have two character factors,
each with \(c_0\)-invariant restriction to \(G_K\). This proves
absolute residual irreducibility. Its restriction to the normal
subgroup is semisimple by the translate argument of
Proposition~\ref{prop:evenfieldplane442}. A vector on the
\(\theta\)-line and its translate by \(c_0\) exhibit the induced
representation. On \(G_K\) its projective kernel is \(\ker\beta
\cap G_K\), and each element outside \(G_K\) interchanges the two
lines. This proves the field assertion.
\end{proof}

\begin{theorem}[Finiteness at the ramified collision for inert primes]\leavevmode
\label{thm:inertquadratic443}
At every \(x\) on the original interval suppose
\(5\le\ell\le L_*\), \(\ell\nmid 2sN_*\), \(\eta_K(\ell)=-1\),
\(2\le r<t(x)\), \(\ell-1\mid 2r\) and \(\ell-1\nmid r\).
Assume \(\bar\alpha\ne1\) and \(\beta=\bar\alpha\omega^r\)
has order two. Then \(T/I_m\) is a finite \(O\)-module for
\(6\le m\le9\), also after adjoining the local characters of
genuine flags. The auxiliary quadratic field \(K\) and the
period fields are the ones fixed on the interval.
\end{theorem}
\begin{proof}
It suffices to prove constancy on each equal-characteristic
curve, by the special-fiber argument of
Theorem~\ref{thm:allzeroweights442}. Each line constituent is
its constant residual lift by the first paragraph of
Lemma~\ref{lem:inertplane443}. If the representation is a
sum of lines this proves the claim.

Suppose first that there is an irreducible plane and the line
constituent is trivial. The plane has determinant one and the
symmetric adjunction descends it to \(G_\QQ\) with odd
determinant \(\eta_K\), by the explicit intertwiner in
Theorem~\ref{thm:scalarinertiaplanes439}. Twist the descended
plane by the global constant \(\beta\). Its determinant is still
\(\eta_K\) and its residual characters are \(1,\eta_K\).
Indeed its residual constituents restrict semisimply to \(G_K\)
where both factors are now trivial; each consequently factors
through the order-two quotient. Inertia at \(\ell\) is unipotent
after this twist, since the two roots left on the plane before
twisting were both \(\beta\). Lemma~\ref{lem:unipotentplane439}
gives unramified local characters, with distinct reductions
because \(\eta_K(\ell)=-1\). Pan's Theorem~5.1.2(1) over
\(\QQ\) gives finiteness over the ordinary weight algebra.
The weight acts by augmentation on this curve, so its trace
image is constant, as required.~\cite{XI-VIII-VI-pan40}

It remains to treat the line \(\beta|_{G_K}\). Use \(W\) of
Lemma~\ref{lem:quadtwist443}. Its residual image has two-power
order in the induced basis. We verify absolute irreducibility
after restriction to \(G_{\QQ(\zeta_\ell)}\). The field
\(KH_\beta\) is biquadratic. The fields \(K\) and its third
quadratic subfield both ramify at \(s\), since \(H_\beta\)
is unramified there. Moreover \(\bar\alpha=\beta\omega^{-r}\)
is a nontrivial quadratic character of conductor prime to
\(\ell s\). It must ramify at some finite prime dividing \(N_*\),
since \(\QQ\) has no nontrivial everywhere finite-unramified
quadratic extension. At that prime \(\omega^r\) is unramified,
so \(H_\beta\) ramifies there too. Thus none of the three
quadratic fields is contained in \(\QQ(\zeta_\ell)\). Their
compositum meets \(\QQ(\zeta_\ell)\) in \(\QQ\). The projective
image of the restricted residual group is consequently still the
full group of order four. It contains the nontrivial diagonal
ratio and an element exchanging that diagonal element's two
eigenlines. No line is stable under their lifts.

This restricted residual image is adequate. Indeed its order
is prime to \(\ell\), so averaging kills its positive-degree
cohomology, including with adjoint coefficients. Its invariant
trace-zero endomorphism space vanishes by absolute irreducibility
and \(\ell\ne2\). All group elements are semisimple. Their span
over an algebraic closure is the full matrix algebra by absolute
irreducibility, verifying the semisimple spanning condition
for adequacy.

Choose a finite set of primes away from \(\ell\) containing the
other ramified primes of \(W\), including any ramification of
\(\theta\). The inertia group at \(\ell\) equals the inertia
group in \(K_w\). By Lemma~\ref{lem:inertplane443} its matrices
in \(W\) have the two characters \(\theta,\theta\beta\) and
split over the fraction field. If an inertia element has trivial
residual matrix, both constant character values are one, so the
matrix itself is the identity. Hence the extension cut out by
\(W\) over the residual fixed field is unramified above \(\ell\).
This is the local condition in the unramified deformation ring
of Allen and Calegari, which allows the residual representation
itself to ramify at \(\ell\).

Apply their Theorem~1 over the totally real field \(\QQ\), with
the residual field enlarged to contain the matrix entries. The
representation is two-dimensional and odd, the prime is odd,
and adequacy over \(\QQ(\zeta_\ell)\) was just verified. Their
unramified deformation ring, allowing the specified finite set
away from \(\ell\), is finite over the Witt vectors of its
residue field. After choosing an integral basis, \(W\) on our
curve gives a specialization of that ring. Since the curve has
characteristic \(\ell\), the image is contained in a finite field
of constants. In particular its traces are constant. Restricting
to \(G_K\), undoing the constant twist and adjoining the constant
line recovers the traces originally under consideration.
This proves the assertion on curves. Completeness gives the
integral finiteness, and the ordered roots of the local
Frobenius polynomial give the assertion on flags as in
Theorem~\ref{thm:allzeroweights442}.~\cite{XI-VIII-finiteunr443}
\end{proof}

The conductor hypothesis is used to keep all three quadratic
subfields of the projective residual field outside
\(\QQ(\zeta_\ell)\). Estimating the strict cohomology lengths
on the entire interval still requires quantitative control on
the corresponding integral deformation rings and their
irreducible components.


\subsection{The conductor of the character used in descent}
\label{sec:rayconductor444}

The deformation argument in Theorem~\ref{thm:inertquadratic443}
uses a finite character \(\theta\) of \(G_K\). Its construction
there leaves its conductor unspecified. We determine which primes
are needed and bound the degree of the residual field. Throughout
this subsection \(K=\QQ(\sqrt{-s})\) is the field of
Proposition~\ref{prop:quadratic345}, \(c\) is its nontrivial
automorphism and \(\mathfrak s=(\sqrt{-s})\) is the prime above
\(s\). Write \(h_K\) for the ordinary class number.

\begin{proposition}[The obstruction at the discriminant prime]\leavevmode
\label{prop:rayobstruction444}
Let \(M\) be a positive integer prime to \(s\), and let \(\beta\)
be an even quadratic Dirichlet character of conductor dividing
\(M\). Write \(\beta_K\) for its restriction to \(G_K\), and put
\(\delta=0\) if \(\beta(s)=1\) and \(\delta=1\) if
\(\beta(s)=-1\). There exists a finite character \(\theta\) of
two-power order satisfying
\[
         \theta^c/\theta=\beta_K,\qquad
         \mathfrak f_\theta\mid
            \mathfrak f=(M)\mathfrak s^\delta .
\]
A solution of conductor dividing \((M)\) exists exactly when
\(\beta(s)=1\). The character values can be taken in a finite
extension of any field of odd positive characteristic.
\end{proposition}
\begin{proof}
Let \(A=\mathrm{Cl}_{\mathfrak f}(K)\) be the group of fractional
ideals prime to \(\mathfrak f\) modulo the ideals \((b)\) with
\(b\equiv1\bmod\mathfrak f\). The congruence is in the
localizations at the primes dividing \(\mathfrak f\).
By the norm compatibility of Artin reciprocity, \(\beta_K\)
acts on an ideal \(\mathfrak a\) by
\(\beta(N_{K/\QQ}\mathfrak a)\). The value on a positive rational
number prime to \(M\) is taken by reducing its numerator and
denominator modulo \(M\). This rule factors through \(A\), since
the norm of a generator congruent to one modulo \((M)\) is one
modulo \(M\). We use reciprocity and the ray class exact
sequence in their usual finite-modulus form.~\cite{XI-VIII-rayclass444}

We prove that \(\beta_K\) is trivial on the classes fixed by \(c\).
Let \(\mathfrak a\) represent such a class. Then
\(\overline{\mathfrak a}/\mathfrak a=(b)\) with
\(b\equiv1\bmod\mathfrak f\). Taking ideal norms gives
\(b\bar b=1\), since the field norm is positive. Choose
\(y\in K^\times\) with \(b=\bar y/y\) by Hilbert's Theorem~90.
At a split prime \(p\mid M\), the two valuations of \(y\) agree
since \(b\) is a unit. An inert \(p\mid M\) has one valuation,
and every prime over \(M\) has ramification index one. We can
therefore make \(y\) a unit at every prime over \(M\) by
multiplying by a rational number with the prescribed valuations
at the primes dividing \(M\). If \(\delta=1\), write locally at
\(\mathfrak s\)
\(y=(\sqrt{-s})^j u\) with \(u\) a unit. Conjugation acts
trivially on the residue field. Thus
\(\bar y/y\equiv(-1)^j\bmod\mathfrak s\). As \(s\) is odd and
\(b\equiv1\bmod\mathfrak s\), \(j\) is even.
Multiplication by \(s^{-j/2}\) makes \(y\) a unit at
\(\mathfrak s\) also. Rational rescalings preserve \(\bar y/y\).

The residue of \(y\) modulo \((M)\) is now a unit fixed by \(c\).
It lies in the image of \((\ZZ/M\ZZ)^\times\). Indeed,
\(\mathcal O_K=\ZZ+\ZZ(1+\sqrt{-s})/2\), and the difference
between the second basis element and its conjugate is
\(\sqrt{-s}\), a unit modulo \((M)\). Comparing the two
coefficients forces the second coefficient of a fixed residue
to vanish modulo \(M\). This applies to \(y\) by clearing
denominators prime to \(M\). Hence \(y\bar y\) is a square unit
modulo \(M\), and \(\beta(N(y))=1\).

The ideal \(\mathfrak b=\mathfrak a/(y)\) is fixed by conjugation.
Its split prime factors occur with equal conjugate exponents.
Each inert prime has square norm. Hence its norm is a rational
square times a power of \(s\). When \(\delta=1\), this last
exponent vanishes because \(\mathfrak b\) is prime to
\(\mathfrak s\). When \(\delta=0\), \(\beta(s)=1\).
It follows that \(\beta_K(\mathfrak a)=1\).

On the subgroup \(D=\{\bar a/a:a\in A\}\) the rule
\(\bar a/a\mapsto\beta_K(a)\) is therefore a well-defined
homomorphism. Extend it to \(A\) one cyclic quotient generator
at a time. For a generator of order \(d\) modulo the current
subgroup, its value must have \(d\)-th power equal to a prescribed
two-power root of unity. Such a root exists among the two-power
roots of unity in the algebraic closure of any field of odd
characteristic. The resulting character of the finite group
\(A\) has precisely the required quotient under \(c\).
Class field theory gives \(\mathfrak f_\theta\mid\mathfrak f\).

For necessity, work in the ray group of \((M)\) alone. The ideal
\(\mathfrak s\) is prime to \((M)\) and fixed by conjugation.
Every quotient \(\theta^c/\theta\) on this group takes value
one on its class. The value of \(\beta_K\) there is
\(\beta(N\mathfrak s)=\beta(s)\). This proves the equivalence.
\end{proof}

\begin{theorem}[Ramification and degree on the period interval]\leavevmode
\label{thm:rayconductor444}
Under the hypotheses of Theorem~\ref{thm:inertquadratic443},
put \(M=N_*\ell\). For the character just constructed, let
\(\bar W=\operatorname{Ind}_{G_K}^{G_\QQ}\theta\) and let
\(F_{\bar W}\) be the fixed field of its kernel. Its Artin conductor \(C(\bar W)\)
and field degree satisfy
\[
      C(\bar W)\mid s^{1+\delta}M^2,\qquad
      [F_{\bar W}:\QQ]\le 2 h_K(s-1)^\delta M^2 .
\]
The exponents of \(s\) and \(\ell\) in \(C(\bar W)\) are exactly
\(1+\delta\) and \(2\), respectively. On inertia over \(\ell\)
the character \(\theta\) has exact order
\(2^{1+v_2(\ell-1)}\).
\end{theorem}
\begin{proof}
The character \(\bar\alpha\) has conductor dividing \(N_*\) and
\(\omega^r\) is quadratic of conductor \(\ell\). Thus \(M\) is
a multiple of the conductor of \(\beta\) and
\(v_\ell(M)=1\). Use Proposition~\ref{prop:rayobstruction444}.
The ray class field \(K_{\mathfrak f}\) is stable under \(c\)
because its defining modulus is stable. Both \(\theta\) and
\(\theta^c\) factor through \(\operatorname{Gal}(K_{\mathfrak f}/K)\).
The restriction of \(\bar W\) to \(G_K\) is their direct sum,
so \(F_{\bar W}\subset K_{\mathfrak f}\). The ray class exact
sequence bounds its degree by
\[
     2|\mathrm{Cl}_{\mathfrak f}(K)|
       \le 2h_K|(\mathcal O_K/\mathfrak f)^\times|
       \le 2h_K M^2(s-1)^\delta .
\]
Here \(|\mathcal O_K/(M)|=M^2\), and the residue field of
\(\mathfrak s\) has \(s\) elements.

We compute the conductor using the finite complex lifts of
the two-power roots of unity. The kernel on every subgroup is
the same in the induced bases in both characteristics, and the
inertia characters may be computed with either realization.
At a split prime \(p\mid M\), the local representation is the
sum of the two characters at the primes over \(p\). Each
conductor exponent is at most \(v_p(M)\). At an inert
\(p\mid M\) the quadratic base change is unramified. Its
inertia and higher ramification groups agree with those over
\(\QQ_p\). On these groups \(\bar W\) again restricts to two
characters, both of conductor exponent at most \(v_p(M)\).
In either case their exponents add, giving at most \(2v_p(M)\).
Every prime outside \(sM\) is unramified.

At \(\ell\), inertia for \(\theta\) factors by local reciprocity
through \(\mathbb F_{\ell^2}^\times\). For a generator \(u\)
of this group, \(c\) sends \(u\) to \(u^\ell\), and its norm
generates \(\mathbb F_\ell^\times\). The restriction of
\(\beta\) to rational local units is the nontrivial quadratic
character there. Consequently
\[
              \theta(u)^{\ell-1}
               =(\theta^c/\theta)(u)=-1 .
\]
For a root of unity of two-power order this equation fixes its
order to be \(2^{1+v_2(\ell-1)}\). Both inertial summands are
nontrivial and tame, giving conductor exponent two.

At \(s\) the inertia image is tame since \(s\) is odd and the
induced image is a two-group. If \(\delta=0\), inertia over
\(K_{\mathfrak s}\) acts trivially. Rational inertia maps onto
the quadratic quotient and exchanges the induced lines. Its
nontrivial image is an involution with eigenvalues \(1,-1\),
so the fixed subspace has dimension one. If \(\delta=1\),
\(\theta\) must ramify at \(\mathfrak s\), by necessity in
Proposition~\ref{prop:rayobstruction444}. Its inertia character
factors through \(\mathbb F_s^\times\), on which conjugation
acts trivially. Some element of inertia over \(K_{\mathfrak s}\)
acts on both lines by the same nontrivial scalar. The fixed
subspace has dimension zero. In the tame case the conductor
exponent is the codimension of the fixed subspace, proving
the asserted exponent at \(s\).
\end{proof}

\begin{corollary}[One ray field for all points of the interval]\leavevmode
\label{cor:rayinterval444}
For each eligible \(\ell\le L_*\) both \(C(\bar W)\) and
\([F_{\bar W}:\QQ]\) are bounded by a fixed power of \(V\),
uniformly in \(x\in[V,V^{3/2}]\), \(r\) and \(\alpha\). There
is also one explicitly bounded auxiliary field containing these
residual fields for all the eligible primes together.
\end{corollary}
\begin{proof}
The prime \(s\) and its field are fixed by
Proposition~\ref{prop:quadratic345}. The estimates there and in
Lemma~\ref{lem:commondyadictwist429} bound \(N_*\) and \(L_*\)
by fixed powers of \(V\). Apply Theorem~\ref{thm:rayconductor444}.
For the simultaneous assertion take the set \(\mathcal L\) of all
primes \(5\le\ell\le L_*\) inert in \(K\) and prime to \(2sN_*\).
Put \(P=\prod_{\ell\in\mathcal L}\ell\), using \(P=1\) for
the empty set. Every modulus used above divides
\(\mathfrak F=(N_*P)\mathfrak s\). Functoriality of ray
class fields embeds all \(F_{\bar W}\) in \(K_{\mathfrak F}\).
Since \(P\le \max(2,L_*)^{L_*}\), the same ray class estimate gives
\[
  \log[K_{\mathfrak F}:\QQ]\le
    \log\!\bigl(2h_K(s-1)\bigr)+2\log N_*
                    +2L_*\log\max(2,L_*) .
\]
All these constructions use the \(N_*\) of the original interval
and the character \(\bar\alpha\) of the actual \(E_x\).
\end{proof}

The bounds control the residual representations entering the
deformation argument. The length of a strict cohomology group
also depends on how far its classes lift through powers of
\(\ell\). Bounding those lifts and comparing the integral
projected theta periods over \(E_x\) remain necessary for
Section~\ref{sec:missing}.


\subsection{Reducible traces for every pair of colliding weights}
\label{sec:allcollisiontraces445}

There are two ways that the residual characters on the plane
behave at a coefficient prime. At a split prime its decomposition
group preserves a line. At an inert prime the two characters on
inertia are exchanged after descent to \(\QQ\), and the entire
local trace is prescribed. We use these facts to treat the
remaining reducible quotients of Section~\ref{sec:scalarinertiaplanes439}.

\begin{lemma}[Local characters after quadratic descent]\leavevmode
\label{lem:quadraticdecomposition445}
Suppose \(5\le\ell\le L_*\), \(\ell\nmid 2sN_*\),
\(\ell-1\mid2r\), and \(\ell-1\nmid r\). Suppose the even
character \(\beta=\bar\alpha\omega^r\) is quadratic. On an
equal-characteristic curve of the reducible trace quotient, let
\(P\) be the absolutely irreducible plane accompanying the line
\(\beta_K\). Choose \(\theta\) by
Proposition~\ref{prop:rayobstruction444} with \(M=N_*\ell\).
There is an odd extension \(W\) of \(\theta P\) to \(G_\QQ\),
with absolutely irreducible residual representation
\(\operatorname{Ind}_{G_K}^{G_\QQ}\theta\).

If \(\ell\) splits in \(K\), \(W|_{G_{\QQ_\ell}}\) has a stable
line. The characters on the line and the quotient are trivial
on principal units under local reciprocity. If \(\ell\) is
inert, then, for the unique \(w\mid\ell\),
\[
      W|_{G_{\QQ_\ell}}\simeq
          \operatorname{Ind}_{G_{K_w}}^{G_{\QQ_\ell}}
                  \bigl(\theta|_{G_{K_w}}\bigr).
\]
In particular every trace on this rational decomposition group
is a constant in the coefficient DVR.
\end{lemma}
\begin{proof}
The character and determinant argument in
Theorem~\ref{thm:scalarinertiaplanes439} gives the line
\(\beta_K\) and \(\det P=\beta_K\). The residual characters of
\(P\) are \(1,\beta_K\). They are distinct because \(\beta_K\)
ramifies over \(\ell\). With the symmetric plane form \(S\)
and alternating matrix \(J_2\), put \(B=S^{-1}J_2\).
Polarization gives
\(P(cgc)=\beta_K(g)BP(g)B^{-1}\) for a complex conjugation
\(c\). Together with \(\theta^c/\theta=\beta_K\), this makes
\(B\) an intertwiner for \(\theta P\) and its conjugate.
The trace of \(B\) is zero and \(B\) is invertible. Scaling it
over a finite coefficient extension to square one makes it an
involution of determinant \(-1\). Giving \(c\) this action
defines \(W\). The sum of a stable lattice over \(K\) and its
translate by \(c\) is a stable lattice over \(\QQ\).
The two distinct residual characters \(\theta,\theta^c\) on
\(G_K\) are exchanged by \(c\). A residual stable line over
an algebraic closure would restrict to a character fixed by
\(c\), which is impossible. The restriction of this irreducible
residual representation to the index-two subgroup is semisimple
because the translates of a simple submodule span the representation.
Taking a nonzero vector on the \(\theta\)-line and its translate
by \(c\) gives the induced residual representation.

First suppose \(\ell\) splits. Identify \(D=G_{\QQ_\ell}\)
with the decomposition group at one of the two primes in \(K\).
Let \(I\) denote inertia and let \(I_0=\ker(\beta|_I)\).
The characteristic polynomial for \(P(i)\) has roots \(1,\beta(i)\)
for every \(i\in I\). Every element of \(I_0\) thus acts
unipotently. If some \(g\in I_0\) acts nontrivially, write
\(P(g)=1+E_{12}\) in a basis. For every \(h\in I_0\), the
equality \(\operatorname{tr}P(gh)=\operatorname{tr}P(h)=2\)
forces \(P(h)_{21}=0\). The common fixed space is exactly
the first line. Normality of \(I_0\) in \(D\) preserves that
line. If \(P(I_0)=1\), the inertia image has order two with
distinct eigenvalues. The decomposition group acts trivially
on the quotient \(I/I_0\), so both eigenlines are stable.
On any of the lines obtained this way the inertia character
and its quotient factor through \(I/I_0\). Twisting by the
finite two-power character \(\theta\) preserves triviality of
both characters on wild inertia. The principal units correspond
to the image of wild inertia under local reciprocity. Intersect
the stable line with the integral lattice. The intersection is
saturated, so both characters are integral.

If \(\ell\) is inert, Lemma~\ref{lem:inertplane443} gives
\(P|_{G_{K_w}}=1\oplus\beta\). Consequently \(W\) has the
two characters \(\theta,\theta\beta\) on \(G_{K_w}\), which
are already distinct on inertia. An element \(h\) outside
\(G_{K_w}\) in \(G_{\QQ_\ell}\) interchanges the two lines,
since its conjugation sends \(\theta\) to \(\theta\beta\).
In the basis \(v,W(h)v\) with \(v\) on the \(\theta\)-line,
\(G_{K_w}\) acts diagonally and \(W(h)^2v=\theta(h^2)v\).
These are precisely the matrices of the stated induction.
Every element outside \(G_{K_w}\) has trace zero, and the two
characters inside it are constant. This proves the assertion.
\end{proof}

\begin{theorem}[Finiteness of the reducible quotient at every collision]\leavevmode
\label{thm:allcollisiontraces445}
For every \(x\in[V,V^{3/2}]\) on the original intervals, take
any of the characters of \(E_x\) of sign \((-1)^r\). Suppose
\[
  5\le\ell\le L_*,\qquad \ell\nmid2sN_*,\qquad
  2\le r<t(x),\qquad \ell-1\mid2r .
\]
For each \(6\le m\le9\), the entire reducible trace quotient
\(T/I_m\) of Section~\ref{sec:scalarinertiaplanes439} is
finite as an \(O\)-module. This holds also after adjoining
the local characters of actual ordinary flags.
\end{theorem}
\begin{proof}
If \(\ell-1\mid r\), apply Theorem~\ref{thm:allzeroweights442}.
Let \(\ell-1\nmid r\). The character
\(\beta=\bar\alpha\omega^r\) is even and has order two on inertia
at \(\ell\). Its fixed field is totally real, hence disjoint
from the imaginary quadratic \(K\). Its order on \(G_K\)
equals its global order. If that order exceeds two, apply
Theorem~\ref{thm:scalarinertiaplanes439}. It remains to treat
a quadratic \(\beta\), allowing \(\bar\alpha=1\).

Take a semisimple realization over the fraction field of a
complete DVR finite over \(k[[z]]\), enlarging that field to
split the constituents as in that theorem.
Every line constituent is its constant residual lift. This
settles a sum of three lines. For a line and an absolutely
irreducible plane the line must be \(1\) or \(\beta_K\).
If it is \(1\), odd descent of the determinant-one plane gives
a rational representation with determinant \(\eta_K\).
Twist by \(\beta\). Its residual semisimplification restricts
trivially to \(G_K\), so has characters \(1,\eta_K\), by its
determinant. Its local inertia is unipotent, since the two
inertial roots before twisting both equal \(\beta\).
Lemma~\ref{lem:unipotentplane439} supplies a stable line and
unramified characters. With fixed odd determinant \(\eta_K\)
the ordinary pseudodeformation theorem of Pan applies over
\(\QQ\). Use his Theorem~5.1.2(1) if \(\eta_K\) is nontrivial
at \(\ell\), and Theorem~6.1.2(1) with the actual line if
\(\eta_K\) is locally trivial. The weight evaluates at
augmentation, so its finite algebra has constant image on
the curve. This gives constant traces on the plane after
undoing the twist.~\cite{XI-VIII-VI-pan40}

For the line \(\beta_K\), take \(W\) from
Lemma~\ref{lem:quadraticdecomposition445}. We use two
finiteness results of Zhang over \(\QQ\), stated in Appendix~A
of the cited preprint. We give the relevant local and global
hypotheses here. Choose the finite set \(\Sigma\) to include
\(\ell\) and every finite prime ramifying in \(W\). The residual
representation \(\bar W\) of \(G_{\QQ,\Sigma}\) is continuous,
odd and absolutely irreducible. Enlarge the finite residue
field to contain its entries. The determinant of \(W\) is
constant on \(G_K\), where it is \(\theta^2\beta_K\), and on
the second coset it is determined by \(\det W(c)=-1\).
Consequently it is the constant residual determinant on the
whole group. For the characteristic-zero fixed-determinant
deformation problem we may take the determinant of a regular
modular lift of \(\bar W\) as in Zhang, enlarging \(\Sigma\)
to include its ramification. Such a lift is supplied by
Serre's modularity theorem in the form used there. On an
equal-characteristic curve every lift of the fixed determinant
reduces to this same constant character. Thus \(W\) in an
integral basis is a point of the specified deformation problem.

At a split prime let \(\xi_1,\xi_2\) be the actual characters
on the integral line and the quotient from the lemma. Their
reductions are distinct on inertia. For \(\sigma,\tau\in
G_{\QQ_\ell}\) we have
\[
  \operatorname{tr}W=\xi_1+\xi_2\ \text{ on }G_{\QQ_\ell},
  \qquad (W(\sigma)-\xi_1(\sigma)I)
         (W(\tau)-\xi_2(\tau)I)=0 .
\]
The factor on the right maps into the stable line, which the
left factor kills. These are exactly the equations defining
the ordinary deformation ring with the specified residual
subcharacter. Zhang's Theorem~A.3.1 makes this ring finite
over \(O[[1+\ell\ZZ_\ell]]\), with weight given by \(\xi_1\)
on principal units. The lemma makes this weight trivial.
Modulo the coefficient uniformizer and the weight augmentation
ideal the ring is a finite-dimensional algebra over a finite
field. Its image in the coefficient DVR is a finite domain
and hence a field of constants.~\cite{XI-VIII-zhang445}

At an inert prime apply Zhang's Proposition~A.1.4. It says
that the global fixed-determinant deformation ring for
\(\bar W\), allowing the finite set \(\Sigma\), is finite over
the local fixed-determinant pseudodeformation ring at \(\ell\).
By the lemma the local trace equals its residual trace on
the full rational decomposition group, including Frobenius
elements outside \(G_{K_w}\). In dimension two with \(2\)
invertible, trace and fixed determinant determine the local
pseudorepresentation. Its map to the curve therefore factors
through the residue field. Finiteness again forces every
global trace to be constant. For \(\bar\alpha\ne1\) this
conclusion also follows from the separate argument in
Theorem~\ref{thm:inertquadratic443}.~\cite{XI-VIII-zhang445}

In both cases restrict \(W\) back to \(G_K\), divide out
the constant twist and adjoin the constant line character.
Every trace of the original representation is constant on
every equal-characteristic curve of \(T/I_m\). A
positive-dimensional special fiber would have a one-dimensional
complete domain quotient and a finite DVR normalization on
which the traces generate a positive-dimensional image.
Thus the special fiber is Artinian. Lifting a basis and
using completeness gives finiteness over \(O\), including
nilpotents and torsion. The flag assertion follows from
Corollary~\ref{cor:flagfinitered439} because the ordered
Frobenius values are successively integral by monic division
and their inertia values are prescribed.
\end{proof}

The theorem applies at all the stated coefficient primes on
the original interval. Its proof gives finiteness at each
weight. The sum in Proposition~\ref{prop:res23} requires a
quantitative bound on integral lengths, including contributions
where reducible and irreducible components meet.


\subsection{Conductor bounds at the two colliding residues}
\label{sec:collisionconductors457}

Let \(E=E_x\) be the period field at a point of the interval of
Section~\ref{sec:missing}, and let \(N_*\) be the conductor modulus
for that interval. Write \(\ell\ge5\) for a prime. Let \(\alpha\) be a
character of \(E/\QQ\) of order \(d\) and primitive conductor \(C\),
where \(\ell\nmid dC\). Take \(1\le r<t(x)\) with
\(\alpha(c)=(-1)^r\). We treat the two residues
\(r\equiv0,(\ell-1)/2\bmod{\ell-1}\) in
Corollary~\ref{cor:smallorder427}.
The scalar strict group \(S(\alpha,r)\) and its subgroup
\(S^{N_*}(\alpha,r)\) are defined in Section~\ref{sec:levelstrict345}.
The former uses divisible coefficients, zero localization at \(\ell\),
and unramified classes at every other finite prime.

Put \(O=\ZZ_\ell[\alpha]\), let \(J\) be its fraction field,
and write \(k=O/\ell O\). Since \(d\) is prime to \(\ell\), \(O\)
is unramified and reduction is faithful on \(\alpha(G_{\QQ})\).
Write \(\theta=\alpha^{-1}\kappa_\ell^{-r}\),
\(D=(J/O)(\theta)\), and \(\tau=\bar\alpha^{-1}\omega_\ell^{-r}\),
where \(\omega_\ell\) is the reduction of the cyclotomic character.
Let \(L\) be the fixed field of \(\ker\tau\) and let
\(\Delta=\operatorname{Gal}(L/\QQ)\). This is a cyclic group of
order \(m\) prime to \(\ell\). The parity condition makes \(L\)
totally real. When \(m>1\), put \(f=\operatorname{ord}_m(\ell)\).

\begin{lemma}[A class number divisor at a nonsplit coefficient prime]\leavevmode
\label{lem:collisionclass457}
Suppose that \(\tau\) is nontrivial on a decomposition group at
\(\ell\). If \(S(\alpha,r)\ne0\), then
\[
 \ell^f\ \bigm|\ 
 \left|\Cl(L)\big/\big\langle[\lambda]\mathrel{:}\lambda\mid\ell
                                      \big\rangle\right|.
\]
\end{lemma}
\begin{proof}
Every continuous cocycle with values in \(D\) has finite image by
compactness. A nonzero element of the strict group therefore has
a multiple of order \(\ell\). Both \(D^{G_{\QQ}}\) and
\(D^{G_{\QQ_\ell}}\) are zero since some local element acts by a
scalar whose difference from one is a unit. The exact sequence
for multiplication by \(\ell\) on \(D\) now lifts the element of
order \(\ell\) uniquely to a class with values in \(k(\tau)=D[\ell]\).
Its localization at \(\ell\) is zero with these finite coefficients
by the same exact sequence locally.

On \(G_L\) a cocycle of this finite class is a homomorphism \(u\).
It is nonzero, since inflation from \(\Delta\) has zero source by
averaging over \(m\). Away from \(\ell\), the divisible cocycle on
inertia is a coboundary. On inertia intersected with \(G_L\),
\(\kappa_\ell=1\) and \(\bar\alpha=1\), so also \(\alpha=1\).
There the coboundary is zero and hence so is \(u\). At every place
over \(\ell\) the finite cocycle itself is locally a coboundary, so
\(u\) is zero on the whole decomposition group over \(L\).
Global class field theory thus factors \(u\) through the class group
in the statement. At real places an odd-order quotient kills
complex conjugation, so the ordinary class group suffices.

For \(g\in G_{\QQ}\) and \(h\in G_L\), the cocycle formula gives
\(u(ghg^{-1})=\tau(g)u(h)\). Its nonzero image in the additive
group of \(k\) is consequently stable under the field generated
by the \(m\)-th roots in \(\tau(G_{\QQ})\) over \(\mathbf F_\ell\).
This field has degree \(f\), since a primitive \(m\)-th root lies
in \(\mathbf F_{\ell^a}\) exactly when \(m\mid\ell^a-1\).
The image has order a positive power of \(\ell^f\), proving the
divisibility.
\end{proof}

\begin{theorem}[A conductor test at a colliding residue]\leavevmode
\label{thm:smallcollision457}
At \(r\equiv(\ell-1)/2\bmod{\ell-1}\), set \(e=2\) and
\(m=\operatorname{lcm}(d,2)\). At \(r\equiv0\bmod{\ell-1}\),
suppose \(\alpha(\ell)\ne1\) and set \(e=1,\ m=d\).
In either case put
\[
 Y=\max\!\left\{1,\frac{m!}{m^m}\,
         C^{(m-1)/2}\ell^{\,m(1-1/e)/2}\right\}.
\]
If \(\ell^f>Y(1+\log Y)^{m-1}\), then \(S(\alpha,r)=0\).
At the zero residue \(\ell\nmid |\Cl(L)|\) also suffices, and
there \(L\) is the field cut out by \(\alpha\).
\end{theorem}
\begin{proof}
At the half residue the residual cyclotomic factor is the quadratic
character of conductor \(\ell\). The Chinese remainder theorem
allows the arguments of this character and of \(\alpha\) to vary
independently modulo \(\ell\) and \(C\). Thus their product has
image of order \(\operatorname{lcm}(d,2)\), and the inertia group
at \(\ell\) has order two. At the zero residue the cyclotomic
factor is trivial, and \(\alpha(\ell)\ne1\) gives nontrivial
decomposition group at \(\ell\). These observations verify the
local hypothesis of Lemma~\ref{lem:collisionclass457} in both cases.

The characters of \(L\) are the powers of the finite character
\(\tau\), after lifting its values to roots of unity. Their
prime-to-\(\ell\) conductors divide \(C\), and the trivial one
has conductor one. Exactly \(m(1-1/e)\) of the characters are
nontrivial on inertia at \(\ell\), all of conductor exponent one
there. The conductor--discriminant formula gives
\[
                  |\operatorname{disc} L|
                       \le C^{m-1}\ell^{m(1-1/e)}.
\]
Minkowski's ideal class bound for a totally real field supplies
an integral ideal of norm at most \(Y\) in each class. The number
of ideals of norm \(n\) is at most \(d_m(n)\), the number of
ordered positive factorizations of \(n\) into \(m\) factors.
Indeed at a rational prime the Euler factor for each prime above
it is coefficientwise at most \((1-X)^{-1}\), and there are at
most \(m\) such primes. After fixing the first \(m-1\) factors,
the number of choices for the last one is at most \(Y\) divided
by their product. Enlarging the first \(m-1\) sums to \([1,Y]\)
and bounding the harmonic sum by \(1+\log Y\) proves
\[
                     |\Cl(L)|\le Y(1+\log Y)^{m-1}.
\]
Each stated condition is now incompatible with the divisibility
from the lemma.
\end{proof}

\begin{corollary}[Quadratic and prime-order characters on the period intervals]\leavevmode
\label{cor:fixedcollision457}
At the half residue every quadratic \(\alpha\) of the specified
parity has \(S(\alpha,r)=0\) if
\[
                   \ell>4N_*(1+\log(2N_*))^2 .
\]
More generally fix \(m=\operatorname{lcm}(d,2)\) and suppose
\(\operatorname{ord}_m(\ell)>m/4\). There is an effective constant
\(K_m\) such that \(\ell>K_m(2N_*)^{4m}\) gives the same
conclusion for every such \(\alpha,r,x\). For an odd prime
order \(d\) this includes all \(\ell\equiv1\bmod4\) which
are primitive roots modulo \(d\) and satisfy this bound.
For cubic characters these primes are \(\ell\equiv5\bmod12\).
Whenever \(\ell\nmid[E_x:\QQ]\), the indicated summands have
length zero in the original period-field cohomology.
\end{corollary}
\begin{proof}
For \(d=2\) we have \(m=2,\ f=1\). A nontrivial quadratic
character has \(C\ge3\), so \(Y=\sqrt{C\ell}/2>1\).
With \(H=1+\log(2C)\), at the real value \(\ell_0=4CH^2\)
the ratio \(Y(1+\log Y)/\ell_0\) equals
\((1+\log C+\log H)/(4H)<1\). Here \(1+\log C<H\) and
\(\log H<H\). As a function of \(\ell\ge\ell_0\) the ratio
decreases, since it is a positive constant times
\((1+\log Y)/Y\). Apply the theorem and use \(C\le N_*\).

For general \(m\ge2\), let \(B_m\) be the maximum of
\((1+z)^{m-1}\exp(-z/(2m))\) for \(z\ge0\). This is an
effectively computable finite number. The ideal count is at most
\[
 B_mY^{\,1+1/(2m)}
          \ \le\ B_m N_*^{\,m/2}\ell^{\,m/4+1/8},
\]
because \(Y\le N_*^{(m-1)/2}\ell^{m/4}\) and
\((m-1)(1+1/(2m))<m\). The integer \(f>m/4\) satisfies
\(f\ge m/4+1/4\). Thus \(K_m=B_m^8\) suffices. If \(d\) is
an odd prime, \(m=2d\) and a primitive root \(\ell\bmod d\)
has \(f=d-1>d/2\). Such a character is even, and
\(\ell\equiv1\bmod4\) makes the half residue even. When \(d=3\)
the extra congruence is \(\ell\equiv2\bmod3\).

The bounds use the common modulus \(N_*\), so they hold
simultaneously on the prescribed interval for all positive
twists in the stated residue. Vanishing of \(S\) implies
vanishing of its subgroup \(S^{N_*}\). The character length
identity of Theorem~\ref{thm:levelstrict345} identifies the
stated zero terms when the field degree is prime to \(\ell\).
\end{proof}

The zero residue with \(\alpha(\ell)=1\) has a different local
condition. At the half residue the estimates above still leave
coefficient primes below the conductor bound and characters for
which the residue degree is small. These terms occur in the
weighted sum of Section~\ref{sec:missing}.

\section{Exceptional weights in the cyclotomic tower}\label{nav:chapter-dyadic-exceptional-weights}

\subsection{Uniform exponent bounds for characters of growing order}
\label{sec:collisiondepth458}

Keep \(\alpha,C,\ell,r,L,m,f\) from
Section~\ref{sec:collisionconductors457}. The bounds here use the
primitive conductor \(C\) of the individual character. In
particular, conductor factors belonging to a different character
of \(E_x\) do not enter the estimate.

\begin{lemma}[The class number divisor at higher levels]\leavevmode
\label{lem:collisiondepth458}
Suppose \(\tau\) is nontrivial on a decomposition group at
\(\ell\). For \(1\le a\le 1+v_\ell(r)\), an element of order
\(\ell^a\) in \(S(\alpha,r)\) implies
\[
    \ell^{af}\ \bigm|\ 
       \left|\Cl(L)/\langle[\lambda]\mathrel{:}\lambda\mid\ell
                                      \rangle\right|.
\]
\end{lemma}
\begin{proof}
Let \(\widetilde\tau\) be the lift of \(\tau\) to a character of
order \(m\) with values in \(O^\times\). For an \(\ell\)-adic unit,
the quotient by the root of unity lifting its residue is in
\(1+\ell\ZZ_\ell\) when the unit is in \(\ZZ_\ell^\times\).
By the binomial theorem an \(\ell\)-th power increases the
congruence depth in this subgroup by at least one. Applying this
to the cyclotomic character gives
\(\theta\bmod\ell^a=\widetilde\tau\bmod\ell^a\) through the
stated depth. In particular \(G_L\) acts trivially on
\(D[\ell^a]\).

The global and local invariants of \(D\) at \(\QQ\) and
\(\QQ_\ell\) vanish by the unit difference in the proof of
Lemma~\ref{lem:collisionclass457}. The finite lift of an
element of order \(\ell^a\) is unique and is locally zero at
\(\ell\). Restriction to \(G_L\) is injective on its global
finite cohomology group by averaging over \(\Delta\). Hence the
restricted homomorphism has exact additive order \(\ell^a\).
As in that lemma it kills all inertia away from \(\ell\) and
all decomposition groups over \(\ell\). Away from \(\ell\) this
uses that \(\theta=1\) on inertia intersected with \(G_L\)
even for the divisible coefficients.

Identify the target with \(O/\ell^a O\). A value of exact order
\(\ell^a\) is represented by a unit in \(O\). By equivariance,
its multiples by \(\ZZ_\ell[\widetilde\tau(\Delta)]\) lie in
the image. This ring is the ring of integers of the unramified
extension of degree \(f\). Indeed a primitive \(m\)-th root
modulo \(\ell\) generates the field of \(\ell^f\) elements,
and the roots of \(X^m-1\) lift uniquely. The ring modulo
\(\ell^a\) has cardinality \(\ell^{af}\), and multiplication
by the unit value embeds it in the image. This proves the
divisibility.
\end{proof}

\begin{theorem}[Conductor thresholds independent of the character order]\leavevmode
\label{thm:uniformdepth458}
Let \(1\le a\le1+v_\ell(r)\). Either of the following conditions
implies \(\ell^{a-1}S(\alpha,r)=0\).
\begin{enumerate}
\item At \(r\equiv(\ell-1)/2\bmod{\ell-1}\), one has
      \(af\ge m/2\) and \(\ell>4096C^4\).
\item At \(r\equiv0\bmod{\ell-1}\), one has
      \(\alpha(\ell)\ne1\), \(af\ge m/4\), and
      \(\ell>256 C^{8/3}\).
\end{enumerate}
For \(\ell\le P(x)\) prime to \([E_x:\QQ]\), each indicated
character component contributes zero to \(R_{\ell,r}\) on
sufficiently distant original intervals.
\end{theorem}
\begin{proof}
Let \(e=2\) at the half residue and \(e=1\) at the zero residue.
Use the discriminant bound of Theorem~\ref{thm:smallcollision457}
and Minkowski's theorem. Every class has an integral representative
of norm at most
\(Y_0=C^{(m-1)/2}\ell^{m(1-1/e)/2}\).
For every integral ideal with norm at most \(Y_0\), its
contribution to \(Y_0^{4/3}\zeta_L(4/3)\) is at least one.
The Euler factor comparison with \(m\) copies of the rational
zeta function gives \(\zeta_L(4/3)\le\zeta(4/3)^m\le4^m\),
using \(1+\int_1^\infty u^{-4/3}\,du=4\). Consequently
\[
          |\Cl(L)|^{1/m}
             \le 4 C^{2/3}\ell^{\,2(1-1/e)/3}.
\]
In the first case the right hand side is less than \(\ell^{1/2}\);
in the second it is less than \(\ell^{1/4}\). In each case
\(|\Cl(L)|<\ell^{af}\). The local hypotheses of
Lemma~\ref{lem:collisiondepth458} hold by
Theorem~\ref{thm:smallcollision457}. Any element of order at
least \(\ell^a\) in an \(\ell\)-primary torsion group has a
multiple of order precisely \(\ell^a\). The lemma therefore
proves the exponent bound, including vanishing when \(a=1\).

For the application, \(a-1\le v_\ell(r)\le\log_2 t(x)\).
This is smaller than \(b_0(x)=\lfloor\exp(x^\varepsilon/8)\rfloor\)
for all sufficiently large \(x\), since \(t(x)=\lfloor x^{401}\rfloor\).
The bound passes to \(S^{N_*}\). At primes prime to the field
degree, descent after unramified extension of coefficients
identifies these scalar groups with the character components
of the strict group over \(E_x\). This is the prime-to-degree
decomposition used in Theorem~\ref{thm:levelstrict345} and in
Lemma~\ref{lem:des28}. The quotient by the subgroup killed by
\(\ell^{b_0}\) therefore has zero component there.
\end{proof}

\begin{corollary}[Half and zero residues in the growing real dyadic tower]\leavevmode
\label{cor:dyadicdepth458}
In Section~\ref{sec:missing} let \(l\) be the least integer with
\(2^l\ge x^{2+8\varepsilon}\), and take a faithful character
\(\alpha\) of \(C_l=\QQ(\zeta_{2^{l+2}})^+\), inflated to \(E_x\).
For every sufficiently large \(x\) on the stated intervals and every
prime \(\ell\equiv5\bmod8\) in \([x^{16},2x^{16}]\) the following
hold for \(1\le r<t(x)\).
\begin{enumerate}
\item If \(\ell-1\mid r\), then \(S(\alpha,r)=0\).
\item If \(r\equiv(\ell-1)/2\bmod{\ell-1}\) and \(\ell\mid r\),
      then \(\ell S(\alpha,r)=0\).
\end{enumerate}
Both sets of twists occur on each interval and these components
give zero in the truncated cohomology \(R_{\ell,r}\).
\end{corollary}
\begin{proof}
The character has order \(m=2^l\) and conductor \(C=2^{l+2}\):
faithfulness prevents descent to the preceding dyadic level.
Thus \(C<8x^{2+8\varepsilon}\). Since \(0<\varepsilon<1/52\),
both conductor thresholds in Theorem~\ref{thm:uniformdepth458}
are below \(x^{16}\) for large \(x\). For \(\ell\equiv5\bmod8\)
one has \(v_2(\ell-1)=2\), and for every \(j\ge0\),
\(v_2(\ell^{2^j}-1)=j+2\). Indeed each successive squaring
introduces a factor \(\ell^{2^j}+1\) of valuation one.
Odd powers preserve the valuation. It follows that
\(\operatorname{ord}_{2^l}(\ell)=2^{l-2}=m/4\) for \(l\ge3\).
The class of \(\ell\) modulo \(2^{l+2}\) differs from both
\(1\) and \(-1\) already modulo eight. Its Frobenius in \(C_l\)
is nontrivial, so \(\alpha(\ell)\ne1\).

Take \(a=1\) at the zero residue and \(a=2\) at the half
residue with \(\ell\mid r\). The ramification and parity
conditions hold since \(C\) is a power of two, \(\alpha\) is
even and \((\ell-1)/2\) is even. The odd part of the degree of
\(E_x\) divides \([K_\psi:\QQ]\), independently of \(x\),
and \(2x^{16}<P(x)\) for large \(x\). This verifies the
remaining hypotheses.

The prime number theorem for the fixed progression \(5\bmod8\)
gives primes throughout the indicated dyadic range for large
\(x\). For each such prime the two types occur already at
\(r=\ell-1\) and \(r=\ell(\ell-1)/2\), respectively, which are
positive and less than \(t(x)\). Both exceed \(x^{\varepsilon/128}\)
and \(\ell\) exceeds the fixed primes in \(\Pi\). Hence the field
average \(J_{\ell,r}\) from Section~\ref{sec:missing} is zero at
these examples. Their elimination follows from the scalar
cohomology bounds themselves.
\end{proof}


\subsection{A single possible character orbit in a dyadic layer}
\label{sec:collisionorbits459}

Fix a character \(\alpha\) as in Section~\ref{sec:collisionconductors457}.
Choose one of the two residues there, with \(\alpha(\ell)\ne1\) in
the zero residue. For a twist \(r\) in that residue of the required
parity, write \(\tau=\bar\alpha^{-1}\omega_\ell^{-r}\) as before.
The field \(L\) cut out by \(\tau\) depends only on the residue,
\(\alpha\) and \(\ell\). Let \(H\) be the \(\ell\)-primary subgroup of
\(\Cl(L)/\langle[\lambda]\mathrel{:}\lambda\mid\ell\rangle\).
Its \(\Delta=\operatorname{Gal}(L/\QQ)\)-action will keep track of
the characters simultaneously.

\begin{lemma}[Faithful characters of the same residual field]\leavevmode
\label{lem:collisionorbits459}
As \(\beta\) ranges over the generators of the cyclic character
group \(\langle\alpha\rangle\), its residual twists in this residue
range bijectively over all faithful characters of \(\Delta\).
The correspondence commutes with raising character values to the
\(\ell\)-th power. Their coefficient rings are all isomorphic to
the ring of integers of the unramified extension of \(\QQ_\ell\)
of degree \(f=\operatorname{ord}_m(\ell)\), where \(m=|\Delta|\).
\end{lemma}
\begin{proof}
In the zero residue \(m=d=\operatorname{ord}(\alpha)\) and the
correspondence sends \(\alpha^u\) to \(\tau^u\) for every
\(u\) prime to \(d\). In the half residue the cyclotomic factor
has order two. If \(d\) is even, every \(u\) prime to \(d\) is
odd. If \(d\) is odd, each unit modulo \(d\) has exactly one
odd lift modulo \(2d\). Thus in both cases the units \(u\bmod m\),
\(m=\operatorname{lcm}(2,d)\), give all generators once and
preserve the quadratic factor when raising \(\tau\) to the
\(u\)-th power. Also \(\ell\) is odd, which gives equivariance.
The groups of \(d\)-th and \(m\)-th roots generate the same
coefficient ring since adjoining a sign does not enlarge it.
Its residue degree is \(f\) by lifting roots of unity of order
prime to \(\ell\). The hypotheses exclude \(m=1\).
\end{proof}

For each orbit \(\mathcal O\) of the faithful residual characters
under the \(\ell\)-th power map, choose a representative
\(\beta_{\mathcal O}\) on \(E_x\) by the lemma. Write \(O\) for
its coefficient ring, using an embedding for each choice.
The module length in the next formula is the length over this
ring. Twists need only have the same residue and parity.

\begin{theorem}[A simultaneous class group bound for different twists]\leavevmode
\label{thm:commonorbits459}
For each \(\mathcal O\) choose a positive twist \(r_{\mathcal O}\)
in the fixed residue and a depth
\(1\le a_{\mathcal O}\le 1+v_\ell(r_{\mathcal O})\). Then
\[
   f\sum_{\mathcal O}\operatorname{length}_O
       \bigl(S(\beta_{\mathcal O},r_{\mathcal O})
                            [\ell^{a_{\mathcal O}}]\bigr)
                         \ \le\ v_\ell(|H|).
\]
In fact the finite group in each summand is dual, at this depth,
to the \(\mathcal O\)-component of \(H\).
\end{theorem}
\begin{proof}
Fix an orbit and its \(\beta,r,a\). Let \(A=(O/\ell^a)(\widetilde
\tau_\beta)\), where \(\widetilde\tau_\beta\) lifts the residual
twist to roots of order prime to \(\ell\). The congruence in
Lemma~\ref{lem:collisiondepth458} identifies \(A\) with the
\(\ell^a\)-torsion of the divisible coefficient module. The
unit difference at \(\ell\) gives its unique finite global lift
for every strict class at this depth. Restriction to \(G_L\)
identifies \(H^1(\QQ,A)\) with
\(\operatorname{Hom}(G_L,A)^\Delta\), using continuous
homomorphisms. Indeed the first and second cohomology of
\(\Delta\) vanish by averaging, and \(G_L\) acts trivially.

Under this restriction the strict classes have precisely the
homomorphisms from \(H\) as their images. One inclusion is the
local calculation of Lemma~\ref{lem:collisiondepth458}.
Conversely pull a \(\Delta\)-equivariant map from \(H\) back
to \(G_L\) by reciprocity and take its unique preimage under
the restriction isomorphism. On inertia away from \(\ell\) its finite class
restricts to zero over \(L\). Restriction is injective on local
first cohomology in this direction because the inertia quotient
has order dividing \(m\). The same argument with decomposition
groups at \(\ell\) gives zero localization there. Thus the
divisible global class satisfies every required local condition.

To take the separate characters, choose a generator of \(\Delta\).
Then \(\ZZ_\ell[\Delta]=\ZZ_\ell[X]/(X^m-1)\) is the product of
unramified integer rings indexed by Frobenius orbits. In fact the
distinct irreducible factors modulo \(\ell\) lift by Hensel's
lemma and are pairwise coprime. For an orbit \(\mathcal O\) the
corresponding summand \(H_{\mathcal O}\) of \(H\) is a module
over its ring. For a faithful orbit that ring is \(O\) by
Lemma~\ref{lem:collisionorbits459}, with the identification
specified by the chosen character. The equivariant maps in the
preceding paragraph identify with
\(\operatorname{Hom}_O(H_{\mathcal O},O/\ell^a)\). On a cyclic
summand \(O/\ell^j\) their length is \(\min(a,j)\). Finally
\(\log_\ell |H|\) is the sum of all component lengths weighted
by their residue degrees. Summing over the distinct faithful
components proves the inequality.
\end{proof}

\begin{corollary}[One possible cyclic group at all half-residue twists]\leavevmode
\label{cor:dyadicorbit459}
Take the growing layer \(C_l\) and the coefficient primes
\(\ell\equiv5\bmod8\) of Corollary~\ref{cor:dyadicdepth458}.
For a fixed \(l,\ell\), there are two Frobenius orbits of faithful
characters of \(C_l\). At most one of them has nonzero strict
cohomology at any half-residue twist. If it exists, the same
orbit occurs at every such twist and its scalar group is cyclic
over \(O\). At \(\ell\mid r\) this group is \(O/\ell\).
For \(1\le r<t(x)\) the only possible contribution of
these faithful characters to the truncated length, when their
field average is zero, equals
\[
    f\max\{0,k_{\ell,l}(r)-b_0(x)\},
    \qquad f=2^{l-2}.
\]
Here \(k_{\ell,l}(r)\) is the length of the one surviving scalar
group, or zero in its absence. In particular the expression
vanishes when \(\ell\mid r\).
\end{corollary}
\begin{proof}
Here \(m=2^l,\ f=m/4\) and the number of faithful characters is
\(\varphi(m)=m/2\). The identity
\(\operatorname{ord}_{2^l}(\ell)=f\) gives the two orbits of
size \(f\). The estimate in the proof of
Theorem~\ref{thm:uniformdepth458} gives
\(|H|\le|\Cl(L)|<\ell^{m/2}=\ell^{2f}\).
Since any nonzero faithful component of \(H\) has order at
least \(\ell^f\), at most one of the two such components is
nonzero, and its length over \(O\) is then one. The finite
duality in Theorem~\ref{thm:commonorbits459} at \(a=1\)
identifies the first layers for every twist in the residue with
the duals of the same components. A nonzero \(\ell\)-primary
group has nonzero \(\ell\)-torsion, so this detects all
surviving orbits. Each scalar group is finite by
Theorem~\ref{thm:levelstrict345}. Over the discrete valuation
ring \(O\) the dimension of its first torsion layer counts its
cyclic summands, by the elementary divisor theorem. It is
therefore cyclic whenever nonzero. When \(\ell\mid r\), the
exponent bound of Corollary~\ref{cor:dyadicdepth458} makes
that cyclic group \(O/\ell\).

The two finite-level local conditions agree:
\(S^{N_*}(\beta,r)=S(\beta,r)\) by
Lemma~\ref{lem:locallevel345}, since \(\beta\) has order prime
to \(\ell\). The prime is also prime to the degree of \(E_x\).
An orbit of size \(f\) contributes \(f\) times the scalar
length by the integral character decomposition. A cyclic group
of length \(k\) modulo its subgroup killed by \(\ell^{b_0}\)
has length \(\max\{0,k-b_0\}\). This proves the formula when
the average is zero, as it is for the examples of
Corollary~\ref{cor:dyadicdepth458}.
\end{proof}

The first layer determines the orbit even at twists prime to
\(\ell\). At those twists the larger exponents of its cyclic
group still enter the comparison of Section~\ref{sec:missing}.


\subsection{Spacing of the weights with large strict cohomology}
\label{sec:spacingdyadic464}

Fix a faithful character of a real dyadic layer \(C_j\), and an odd
prime \(\ell\equiv5\bmod 8\). Assume \(j\ge3\) and
\(\ell>4096(2^{j+2})^4\). Consider positive weights in the residue
\(r\equiv(\ell-1)/2\bmod{\ell-1}\). The residual field \(L\) and
group \(\Delta=\operatorname{Gal}(L/\QQ)\) of
Lemma~\ref{lem:collisionorbits459} are common to all these weights
and to all faithful characters of the layer. Put \(f=2^{j-2}\).
The proof of Corollary~\ref{cor:dyadicorbit459} gives
\[
 H=\bigl(\Cl(L)/\langle[\lambda]\mid\lambda\mid\ell\rangle\bigr)
                [\ell^\infty],\qquad |H|<\ell^{2f}.
\]
Consequently at most one faithful Frobenius orbit has a nonzero
component in \(H\). That component, if present, has length one over
the unramified coefficient ring \(O\) of degree \(f\). We shall
use this length at a finite-order character to compare the groups
at different integer weights.

\begin{lemma}[One compact module for all weights in a fixed residue]\leavevmode
\label{lem:collisionfamily464}
Choose a faithful orbit, a character \(\alpha\) representing it as
in Lemma~\ref{lem:collisionorbits459}, and the lift
\(\widetilde\tau\) of its residual twist to roots of unity of
order prime to \(\ell\). Put \(J=\operatorname{Frac}(O)\) and
\(\Lambda=O[[T]]\). There is a compact \(\Lambda\)-module \(B\)
such that \(B/TB\) is isomorphic to this orbit's component of
\(H\). For \(t_r=(1+\ell)^{-r}-1\) there are isomorphisms
\[
  S(\alpha,r)\simeq
       \operatorname{Hom}_{O,\mathrm{cont}}
                 \bigl(B/(T-t_r)B,J/O\bigr).
\]
If the component of \(H\) is zero then \(B=0\). If it is nonzero
then \(B\) is cyclic over \(\Lambda\).
\end{lemma}
\begin{proof}
Let \(\QQ_\infty\) be the cyclotomic \(\ZZ_\ell\)-extension of
\(\QQ\), and put \(L_\infty=L\QQ_\infty\). Since \(|\Delta|\) is
prime to \(\ell\), the Galois group is
\(\Delta\times\Gamma\), where \(\Gamma\) is identified with
\(1+\ell\ZZ_\ell\) by the principal-unit part of the cyclotomic
character. Let \(\gamma\) correspond to \(1+\ell\). From the
maximal abelian pro-\(\ell\) quotient of \(G_{L_\infty}\), divide
out the closed subgroup generated by the images of all inertia
groups away from \(\ell\) and all decomposition groups at
\(\ell\). Denote the resulting compact group by \(X\). Conjugation
induces commuting actions of \(\Delta\) and \(\Gamma\) on \(X\).
The factor of \(\ZZ_\ell[\Delta]\) corresponding to
\(\widetilde\tau\) is \(O\), with \(\delta\) acting by
\(\widetilde\tau(\delta)\). This follows by factoring the prime-to-
\(\ell\) cyclotomic polynomials into unramified factors, as in
Theorem~\ref{thm:commonorbits459}. Let \(B\) be the image in
\(X\) of the idempotent of this factor and let \(T\) act by
\(\gamma-1\).

For \(t\in\ell O\), give \(J/O\) the action \(D_t\) on which
\(\Delta\) acts by \(\widetilde\tau\) and \(\gamma\) by \(1+t\).
The latter action extends continuously to \(\Gamma\) by the
binomial series. The coefficient action is trivial on
\(G_{L_\infty}\). Restriction identifies global classes for \(D_t\)
with continuous equivariant homomorphisms from \(G_{L_\infty}\)
to \(J/O\). Indeed in the inflation--restriction sequence the
outer groups in degrees one and two for \(\Delta\times\Gamma\)
vanish. Averaging over \(\Delta\) kills its positive cohomology,
and a nontrivial value of \(\widetilde\tau\) differs from one by
a unit, so \(D_t^\Delta=0\).

At a rational prime away from \(\ell\) the image of inertia in
\(\Delta\times\Gamma\) has order prime to \(\ell\), since
\(\QQ_\infty\) is unramified there. Restriction on its first
inertia cohomology therefore has zero kernel after passage to
\(L_\infty\). At \(\ell\), the image of a decomposition group
in \(\Delta\times\Gamma\) contains its projected prime-to-\(\ell\)
factor with trivial \(\Gamma\)-coordinate. To see this, take
powers of an element whose exponent tends to zero
\(\ell\)-adically and is one modulo \(|\Delta|\), and use
closedness. That finite factor acts nontrivially through
\(\widetilde\tau\), since the tame quadratic twist ramifies at
\(\ell\). Averaging over it gives zero first cohomology of the
decomposition image, just as globally. Thus a class is zero at
\(\ell\) exactly when its restricted homomorphism kills every
decomposition group above \(\ell\). These arguments identify
the strict classes with the equivariant homomorphisms from
\(X\). The real places impose no condition at an odd prime.
Equivariance takes the indicated idempotent component and imposes
\(T=t\). At \(t=t_r\) the coefficient action is exactly
\(\alpha^{-1}\kappa_\ell^{-r}\), giving the asserted formula.

At \(t=0\) the coefficient action already factors through
\(\Delta\). Restriction to \(G_L\), with the same averaging
argument globally, on inertia away from \(\ell\), and on the
decomposition image at \(\ell\), identifies the strict group
with the equivariant homomorphisms \(H\to J/O\). Here global
reciprocity identifies the unramified abelian group with the
class group, and imposing splitting at the primes over \(\ell\)
kills their ideal classes \cite{XI-VIII-rayclass444}. We have therefore
identified the continuous \(O\)-duals of \(B/TB\) and of the
indicated component of \(H\). Dualizing identifies the compact
modules. Explicitly, the trace pairing identifies \(J/O\) with
the \(\ZZ_\ell\)-dual of \(O\) with values in
\(\QQ_\ell/\ZZ_\ell\), since \(O\) is unramified. The claimed
biduality follows from ordinary finite abelian duality on the
finite quotients and passage to inverse limits.

In particular \(B/(\ell,T)B\) has dimension zero or one over
\(O/\ell O\). In dimension one choose a lift of a basis vector.
The submodule it generates is compact as the image of
\(\Lambda\), hence closed. In dimension zero take the zero
submodule. Its quotient \(Y\) satisfies \(Y=(\ell,T)Y\).
Every open additive subgroup contains an open subgroup stable
under \(O\) and \(\Gamma\). Indeed one can first intersect its
inverse images under multiplication by a finite \(\ZZ_\ell\)-basis
of \(O\) including one, obtaining an \(O\)-submodule, and then
intersect its finitely many \(\Gamma\)-translates. The stabilizer
is open by continuity of the action and compactness. On any
finite stable quotient \(Q\), the operator \(T\) is nilpotent
modulo \(\ell\), since the action of \(\Gamma\) factors through
a cyclic group of \(\ell\)-power order. Thus \(Q=(\ell,T)Q\)
forces \(Q/\ell Q=0\) and then \(Q=0\). The finite quotients
separate points, so \(Y=0\).
\end{proof}

\begin{theorem}[Lengths of the faithful dyadic strict groups]\leavevmode
\label{thm:collisionbranch464}
For an orbit whose component of \(H\) is nonzero, write
\(k(r)=\operatorname{length}_O S(\alpha,r)\). Either \(k(r)=1\)
at every indicated weight, or there exist \(z\in\ell O\) with
\(v_\ell(z)=1\) and \(b\in\{1,2,\ldots\}\cup\{\infty\}\) such
that at every such weight
\[
                    k(r)=\min\{b,v_\ell(t_r-z)\}.
\]
Here \(v_\ell(\ell)=1\) and \(v_\ell(0)=\infty\). For every
integer \(a\ge1\), the weights with \(k(r)>a\), if any, lie in
one residue class modulo \((\ell-1)\ell^a\).
\end{theorem}
\begin{proof}
Write \(B=\Lambda/I\) using Lemma~\ref{lem:collisionfamily464}.
Since \(B/TB\) has length one, the ideal formed by the constant
terms of elements of \(I\) is \(\ell O\). Choose
\(F(T)=\ell+cT+T^2G(T)\) in \(I\). Every \(P\in I\) satisfies
\(P(t)\in\ell O\) for \(t\in\ell O\). If \(c\in\ell O\), then
\(F(t)\) has valuation exactly one for every such \(t\).
Evaluation of \(I\) thus gives precisely \(\ell O\), so every
specialized dual has length one.

If \(c\) is a unit, \(F\) has a root \(z\in\ell O\). For
completeness, start with zero modulo \(\ell\). Given a root
modulo \(\ell^n\), the correction modulo \(\ell^{n+1}\) is
unique because \(F'\) is congruent to \(c\) modulo \(\ell\)
at that point. Completeness gives the root, and the equation
modulo \(\ell^2\) gives \(v_\ell(z)=1\). Division by \(T-z\)
in \(O[[T]]\) gives \(F(T)=(T-z)V(T)\), where the coefficient
of \(T^i\) in \(V\) is
\(\sum_{n\ge i+1}[T^n]F\,z^{n-i-1}\). The series converge
in \(O\) and the constant coefficient of \(V\) is congruent
to \(c\) modulo \(\ell\). Thus \(T-z\in I\). Evaluation at
\(z\) identifies \(\Lambda/(T-z)\) with \(O\) by the same
division formula. The image of \(I\) in this ring is then
\(\ell^b O\), with the zero ideal written as \(b=\infty\).
It is proper because \(B\ne0\). Imposing \(T=t_r\) now gives
\(O/(\ell^b,t_r-z)\). Its dual has the stated length. The
strict groups at the positive weights under consideration are
finite by Theorem~\ref{thm:levelstrict345}.

If two distinct weights \(r,s\) have \(k(r),k(s)>a\), the formula
implies \(v_\ell(t_r-t_s)\ge a+1\). On the other hand
\[
 v_\ell(t_r-t_s)=v_\ell((1+\ell)^{|r-s|}-1)=1+v_\ell(r-s).
\]
Indeed a power prime to \(\ell\) preserves the valuation of a
principal unit minus one, and raising it to the odd prime
\(\ell\) increases the valuation by exactly one, by the binomial
theorem. Hence \(\ell^a\mid r-s\). The weights already agree
modulo \(\ell-1\), which proves the assertion.
\end{proof}

\begin{corollary}[An interval bound with one possible exceptional weight]\leavevmode
\label{cor:onedeepweight464}
On an interval \([V,V^{3/2}]\) of Section~\ref{sec:missing},
fix \((\ell,j)\). Include every weight \(r\) in the half residue
for which \(1\le r<t(x)\) at some point of the interval where
the field contains the prescribed dyadic layer of index \(j\)
as in Corollary~\ref{cor:dyadicdepth458} or
Corollary~\ref{cor:shortcenter460}, and where
\(\ell\equiv5\bmod8\), \(x^{16}\le\ell\le2x^{16}\).
Call this set \(\mathcal R\). At sufficiently large \(V\),
there is at most one \(r\in\mathcal R\) at which the faithful
dyadic components can survive the truncation at \(b_0(x)\)
anywhere on the interval.

More precisely, suppose that \(\mathcal R\ne\varnothing\).
Write \(D=\max\mathcal R-\min\mathcal R\), take \(k(r)=0\)
in the absence of a surviving orbit, and choose \(r_0\) with
maximal \(k(r_0)\). Then
\[
 f\log\ell\sum_{\substack{r\in\mathcal R\\r\ne r_0}}
             \max\{0,k(r)-1\}
       \ \le\ \frac{fD\log\ell}{(\ell-1)^2}.
\]
\end{corollary}
\begin{proof}
In both field families \(2^j<2x^{2+8\varepsilon}\), so the
conductor hypothesis at the beginning of this subsection follows
from \(\ell\ge x^{16}\) for large \(V\). The prime is prime to
the field degree by Corollary~\ref{cor:dyadicdepth458}. For
fixed \(\ell,j\), the scalar group is independent of the choice
of \(x\) realizing the weight. The stronger level condition and
the scalar one agree by Lemma~\ref{lem:locallevel345}, and
Theorem~\ref{thm:levelstrict345} gives descent on these same
fields. The \(f\) conjugates in the orbit have the same length,
since coefficient conjugation fixes \(t_r\).

For each \(a\ge1\) a nonempty set of weights with \(k(r)>a\)
contains \(r_0\). Its elements are spaced by at least
\((\ell-1)\ell^a\), by Theorem~\ref{thm:collisionbranch464}.
On deleting \(r_0\) at most
\(\lfloor D/((\ell-1)\ell^a)\rfloor\) elements remain.
Sum over \(a\ge1\) and use
\(\sum_{a\ge1}\ell^{-a}=1/(\ell-1)\). This proves the displayed
inequality. Finally put \(t_*=\lfloor V^{(3/2)401}\rfloor\).
For large \(V\), every \(b_0(x)\) is at least
\(\lceil\log_2(t_*+1)\rceil\), since
\(b_0(x)=\lfloor\exp(x^\varepsilon/8)\rfloor\).
With this last ceiling as \(a\), the spacing exceeds
\(D<t_*\). Hence each \(r\ne r_0\) has \(k(r)\le b_0(x)\)
wherever it occurs.
\end{proof}

The estimate leaves the value of \(k(r_0)\) open. Its size enters
the truncated sum even when there is only one such weight. The
other character components and the upper bound for the projected
Gram determinant in Theorem~\ref{thm:primitivegram463} also enter
the comparison on these intervals.


\subsection{Full orders at weights outside the class group exceptions}
\label{sec:nonsplitcounts466}

Use the polynomial-conductor interval of Theorem~\ref{thm:conductor232},
its largest field \(E_*\), its degree \(n_*\), and its modulus \(N_*\).
Put \(t_*=\lfloor V^{1203/2}\rfloor\) and \(t_-=\lfloor V^{401}\rfloor\).
Consider primes \(\ell\ge5\) prime to \(n_*N_*\). For a fixed residue
\(s=0\) or \(s=(\ell-1)/2\bmod{\ell-1}\), take the characters
\(\alpha\) of \(E_*\) with \(\alpha(c)=(-1)^s\). We require that
\(\tau=\bar\alpha^{-1}\omega_\ell^{-s}\) be nontrivial on the
decomposition group at \(\ell\). This includes every character of
the indicated parity in the half residue. In the zero residue it
means \(\alpha(\ell)\ne1\).

Fix a cyclic subgroup of the character group of \(E_*\), and let
\(\alpha\) run through its generators that meet these conditions.
Their residual fields \(L\) are equal by Lemma~\ref{lem:collisionorbits459}.
For a Frobenius orbit \(\mathcal O\) of the resulting characters
of \(\Delta=\operatorname{Gal}(L/\QQ)\), write \(O\) for its
unramified coefficient ring, \(f\) for its residue degree, and
\[
 h_{\mathcal O}=\operatorname{length}_O
   \left(\bigl(\Cl(L)/\langle[\lambda]\mid\lambda\mid\ell\rangle\bigr)
                    [\ell^\infty]_{\mathcal O}\right).
\]
For a representative \(\alpha\) put
\(\lambda_{\mathcal O}(r)=\operatorname{length}_O S(\alpha,r)\)
when \(r>0\) is in the chosen residue. The value is independent
of the representative because coefficient conjugation preserves
the local conditions.

\begin{lemma}[The class group length bounds the number of large specializations]\leavevmode
\label{lem:classdet466}
Set \(h=h_{\mathcal O}\). If \(h=0\), all the scalar groups vanish.
For \(h>0\) the weights with \(\lambda_{\mathcal O}(r)>ha\)
occupy at most \(h\) residue classes modulo
\((\ell-1)\ell^a\), for every integer \(a\ge1\). In particular,
among \(1\le r<t_*\), at most \(h\) weights have
\[
 \lambda_{\mathcal O}(r)>
          h\bigl(1+\lfloor\log_\ell t_*\rfloor\bigr).
\]
\end{lemma}
\begin{proof}
The construction in Lemma~\ref{lem:collisionfamily464} applies to
this residual field. Indeed \(\Delta\) again has order prime to
\(\ell\), and the lift of \(\tau\) differs from one by a unit on
some member of the local decomposition image. These two facts
give the global and local restriction isomorphisms there.
Away from \(\ell\) the cyclotomic extension is unramified.
At the finite-order specialization \(T=0\), reciprocity over
\(L\) again identifies the dual with the indicated quotient of
the class group. Thus a compact module \(B\) over
\(\Lambda=O[[T]]\) has \(B/TB\) of length \(h\), and the
specialization \(B/(T-t_r)B\), where \(t_r=(1+\ell)^{-r}-1\),
is dual to the strict group. The compact generation argument in
that lemma works for any finite number of generators by lifting
a basis of \(B/(\ell,T)B\). In particular \(h=0\) implies \(B=0\).

For \(h>0\), choose a surjection \(\Lambda^g\to B\) in this way,
and write \(K\) for its kernel. In \(O^g\) the constant terms
of \(K\) form a submodule with quotient \(B/TB\), by right
exactness. This submodule has rank \(g\). Lift an \(O\)-basis
of it to \(g\) columns in \(K\). Their determinant \(F(T)\)
satisfies \(v_\ell(F(0))=h\) by the elementary divisor theorem.
At every \(t_r\) the specialized columns are relations for
\(B/(T-t_r)B\). Taking determinants over \(O\) gives
\(\lambda_{\mathcal O}(r)\le v_\ell(F(t_r))\), with the
right side read as infinity if the determinant vanishes.

Suppose \(h+1\) weights \(r_i\) with
\(\lambda_{\mathcal O}(r_i)>ha\) lie in distinct classes
modulo \((\ell-1)\ell^a\). Since their common residue modulo
\(\ell-1\) is fixed, they are distinct modulo \(\ell^a\).
Write \(z_i=t_{r_i}\). The binomial identity from
Theorem~\ref{thm:collisionbranch464} gives
\(v_\ell(z_i-z_j)=1+v_\ell(r_i-r_j)\le a\) for \(i\ne j\).
Each \(z_i\) is in \(\ell O\). Successively dividing \(F(T)\)
by the linear factors \(T-z_i\) gives
\[
 F(T)=Q(T)\prod_{i=1}^{h+1}(T-z_i)+R(T),\qquad
       Q\in O[[T]],\quad R\in O[T],\quad \deg R\le h.
\]
The divisions are integral by the convergent coefficient formula
in Theorem~\ref{thm:collisionbranch464}. Interpolate \(R\) at
the distinct \(z_i\) and evaluate at zero. The term corresponding
to \(i\) is
\[
             F(z_i)\prod_{j\ne i}\frac{-z_j}{z_i-z_j}.
\]
Its valuation exceeds \(ha+h-ha=h\). The term containing
\(Q(0)\) has valuation at least \(h+1\). This contradicts
\(v_\ell(F(0))=h\). Taking
\(a=1+\lfloor\log_\ell t_*\rfloor\) now gives the last claim,
since \(\ell^a>t_*\).
\end{proof}

\begin{theorem}[A uniform weight count for characters nontrivial at the coefficient prime]\leavevmode
\label{thm:nonsplitprefix466}
There is one set \(\mathcal N_V\subset[1,t_*]\cap\ZZ\) for the
whole interval such that, uniformly for \(t_-\le T\le t_*\),
\[
       |\mathcal N_V\cap[1,T]|
          \ll_{\psi,\varepsilon} \frac{T\log\log T}{\log T}.
\]
Outside this set, every orbit and prime specified above, at every
weight \(1\le r<t_*\) in its residue, satisfies
\(\lambda_{\mathcal O}(r)\le
h_{\mathcal O}(1+\lfloor\log_\ell t_*\rfloor)\).
\end{theorem}
\begin{proof}
Put in \(\mathcal N_V\) precisely those weights which exceed
their last bound for at least one such orbit and prime. Only
\(\ell\le2t_*+1\) can occur, because \(\ell-1\mid2r\).
For each \(\ell\) we first count the possible weights without
using this divisibility to estimate their number. Put
\(w(\ell)=\max\{v_2(\ell-1),v_2(\ell+1)\}\).

For a fixed cyclic subgroup of characters and one residue write
\(m=[L:\QQ]\), let \(C\mid N_*\) be the conductor of a
generator, and let \(H\) be the \(\ell\)-primary split class
quotient occurring in the definition of \(h_{\mathcal O}\).
All its faithful factors have residue degree
\(f=\operatorname{ord}_m(\ell)\). Since they are distinct
integral idempotent factors, \(f\sum_{\mathcal O}h_{\mathcal O}
\le\log|H|/\log\ell\). Every character of \(L\) has conductor
dividing \(C\ell\), so \(|\operatorname{disc}(L)|\le(C\ell)^m\)
by the conductor--discriminant formula. The field is totally real
by parity. Minkowski's bound gives every class an integral ideal
of norm at most \(\sqrt{|\operatorname{disc}(L)|}\). Counting
these ideals by the sum of their inverse squared norms gives
\[
  |H|\le|\Cl(L)|\le|\operatorname{disc}(L)|\zeta_L(2)
                         \le(2C\ell)^m.
\]
For the last inequality use the Euler factors and \(\zeta(2)<2\).

Let \(2^j\) be the two-part of the order of the character
\(\alpha\) before twisting. Its odd part divides
\(o=[K_\psi:\QQ]\), and the twist has order \(m\) divisible
by \(2^j\) and at most \(2o\,2^j\). Since \(2^j\mid\ell^f-1\),
one has \(f\ge2^{j-w(\ell)}\). In fact an odd power gives
\(v_2(\ell^f-1)=v_2(\ell-1)\), and for even \(f\) successive
squaring gives \(v_2(\ell^f-1)=w(\ell)+v_2(f)\).
Together with Lemma~\ref{lem:classdet466} the bounds just proved
allow at most
\[
 \sum_{\mathcal O}h_{\mathcal O}
       \ll_\psi 2^{w(\ell)}
                      \left(1+\frac{\log(2N_*)}{\log\ell}\right)
\]
exceptional weights for this cyclic subgroup and this residue.

At each fixed \(j\) there are \(O_\psi(1)\) cyclic subgroups to
count. Indeed the cyclic group of characters inflated from the
largest real dyadic layer has index at most \(4o\) in the full
character group. In each coset the elements killed by
\(o_{\rm odd}2^j\), where \(o_{\rm odd}\) is the odd part of
\(o\), number at most \(2^j\), since any two of them differ by
an element of the cyclic dyadic group killed by \(2^j\).
Each subgroup with two-part \(2^j\) has at least \(2^{j-1}\)
generators when \(j\ge1\), and at least one when \(j=0\).
Generators of distinct cyclic subgroups are disjoint.
There are \(O_{\psi,\varepsilon}(\log V)\) possible \(j\) since
\(n_*\ll_\psi V^{3+12\varepsilon}\), and only two residues.
As \(\log N_*\ll_{\psi,\varepsilon}\log V\), the bound for
all groups at this prime is \(O_{\psi,\varepsilon}
(2^{w(\ell)}(\log V)^2)\).

For \(t_-\le T\le t_*\) split the primes at
\(Y=T/(\log T)^4\). For \(a\ge1\) the integers with
\(w(\ell)\ge a\) lie in the two classes \(\pm1\bmod 2^a\).
Counting up to \(Y\) and summing \(2^a(Y/2^a+1)\) for
\(a\le\log_2(Y+1)\) shows that
\(\sum_{5\le\ell\le Y,\ \ell\ {\rm prime}}2^{w(\ell)}
\ll Y\log(2Y)\). The small primes therefore allow
\(O_{\psi,\varepsilon}(T/\log T)\) weights.

For larger primes a weight \(r\le T\) is a positive multiple of
\((\ell-1)/2\), and \(\ell\le2T+1\). Their contribution to
the count is at most
\[
       \sum_{Y<\ell\le2T+1,\ \ell\ {\rm prime}}
                       \frac{2T}{\ell-1}
                     \ll \frac{T\log\log T}{\log T}.
\]
Indeed there are \(O(\log\log T)\) dyadic intervals between
\(Y\) and \(2T+1\). On each, Chebyshev's upper bound for the
number of primes gives a reciprocal sum \(O(1/\log Y)\), and
\(\log Y\asymp\log T\). The two counts establish the claimed
bound for every prefix of the one fixed set.
\end{proof}

\begin{corollary}[A polynomial bound for the full orders at the other weights]\leavevmode
\label{cor:fullnonsplit466}
At each \(x\in[V,V^{3/2}]\), consider all triples consisting of
a prime \(\ell\), a weight \(1\le r<t(x)\) outside
\(\mathcal N_V\), and one of the specified character orbits
occurring on \(E_x\) in the residue of \(r\). Their full scalar
groups satisfy
\[
   \sum_{(\ell,r,\mathcal O)}
       f_{\mathcal O}\lambda_{\mathcal O}(r)\log\ell
            \ll_{\psi,\varepsilon}
                  n_*^2t_*^2(\log(2V))^3 .
\]
Each of these components contributes zero to \(R_{\ell,r}(E_x)\)
once \(V\) is sufficiently large.
\end{corollary}
\begin{proof}
Inflation of a character from \(E_x\) to \(E_*\) preserves its
scalar coefficient action over \(\QQ\), hence also its strict
group. The character decomposition is integral since
\(\ell\nmid n_*\). The comparison with the conductor-level group
is the one in Theorem~\ref{thm:levelstrict345} and
Lemma~\ref{lem:locallevel345}. From the bound for \(|H|\) in
the preceding proof and \(m\le2n_*\), each \(h_{\mathcal O}\)
is \(O_{\psi,\varepsilon}(n_*\log(2V))\). The theorem then
bounds each scalar length by
\(O_{\psi,\varepsilon}(n_*(\log(2V))^2)\).
For one prime and weight, the sum of the orbit degrees is at
most \(n_*\), the total number of characters. There are at most
\(2t_*+1\) primes and at most \(t_*\) weights, and
\(\log(2t_*+1)\ll\log(2V)\). Multiplication proves the
displayed inequality. Finally this polynomial bound on each
length is smaller than \(b_0(x)\) throughout a sufficiently
distant interval. All the indicated primes are then at most
\(2t_*+1<P(x)\), so truncation kills their components.
\end{proof}

The class group comparison in this argument uses the nontrivial
local residual action at \(\ell\). The groups on
\(\mathcal N_V\) and the zero-residue characters trivial on the
decomposition group at \(\ell\) are among the terms for which the
estimate in Corollary~\ref{cor:supportmass465} still needs an
upper bound.


\subsection{A finite class group determining the possible exceptional dyadic weight}
\label{sec:finiteexception467}

Return to the faithful characters of the real dyadic field \(C_j\)
and the primes and half-residue weights of
Section~\ref{sec:spacingdyadic464}. There are two faithful
Frobenius orbits. An orbit absent from the split class quotient
of \(L\) has zero strict group at every such weight by
Lemma~\ref{lem:collisionfamily464}. For the possible nonzero
orbit retain the notation \(O,J,\widetilde\tau,B,\gamma,t_r\)
from that lemma. In particular \(\gamma\) has cyclotomic
principal-unit value \(1+\ell\). Put \(u_r=1+t_r=(1+\ell)^{-r}\).

For an integer \(n\ge0\), let \(\QQ_n\) be the degree-\(\ell^n\)
subfield of the cyclotomic \(\ZZ_\ell\)-extension \(\QQ_\infty\).
Set \(L_n=L\QQ_n\) and define
\[
 H_n=\left(
     \bigl(\Cl(L_n)/\langle[w]\mid w\mid\ell\rangle\bigr)
           [\ell^\infty]\right)_{\mathcal O}.
\]
The subscript selects the same factor of \(\ZZ_\ell[\Delta]\)
as in Lemma~\ref{lem:collisionfamily464}. We use the action on
ideals of the element of \(\operatorname{Gal}(L_n/L)\) induced
by \(\gamma\). The fields \(L_n\) serve to calculate a scalar
strict group on the period field \(E_x\).

\begin{lemma}[Specialization through an ordinary ideal class group]\leavevmode
\label{lem:splitcyclolevel467}
There is an isomorphism respecting the action of \(\gamma\),
\[
        H_n\simeq B/\bigl((1+T)^{\ell^n}-1\bigr)B.
\]
With \(M^\vee=\operatorname{Hom}_{O,\mathrm{cont}}(M,J/O)\),
one therefore has for every indicated weight
\[
  \bigl(S(\alpha,r)[\ell^{n+1}]\bigr)^\vee
     \simeq H_n\big/
        \bigl(\ell^{n+1}H_n+(\gamma-u_r)H_n\bigr).
\]
\end{lemma}
\begin{proof}
Use the coefficient module \(D_0=(J/O)(\widetilde\tau)\) over
\(\QQ_n\), with trivial action in the cyclotomic
\(\ZZ_\ell\)-direction. Over \(L_\infty=L\QQ_\infty\) it is
trivial, and the quotient group over \(\QQ_n\) is
\(\Delta\times\langle\gamma^{\ell^n}\rangle\), where angular
brackets denote the closed subgroup. Restrict global classes as
in Lemma~\ref{lem:collisionfamily464}. Averaging over \(\Delta\)
gives zero first and second cohomology of this quotient because
\(D_0^\Delta=0\). At the prime above \(\ell\) the quadratic
inertia action remains nontrivial after the local extension
of degree \(\ell^n\), since a subgroup of odd index cannot lie
in the kernel of a quotient of order two. The finite part of
the local decomposition image therefore still has zero
invariants. Away from \(\ell\) the cyclotomic extension is
unramified, so the image of inertia has order prime to \(\ell\).
These are the local injectivity arguments in that lemma.
Consequently the strict classes of \(D_0\) over \(\QQ_n\) are
the continuous dual of the quotient of \(B\) imposing
\(\gamma^{\ell^n}=1\).

Instead restrict the same classes to \(L_n\). The quotient
over \(\QQ_n\) is now the finite group \(\Delta\), still of
order prime to \(\ell\). Restriction is again an isomorphism
onto the equivariant homomorphisms satisfying the same local
conditions. By global reciprocity these homomorphisms factor
through the ideal class group with all classes of primes above
\(\ell\) killed, and every homomorphism of its indicated component
occurs. Thus their group is \(H_n^\vee\). The two restrictions
commute with conjugation by the lift of \(\gamma\) which fixes
\(L\), and this lift acts trivially on \(D_0\). Dualizing gives
the first isomorphism with the asserted action.

The specialization formula of Lemma~\ref{lem:collisionfamily464}
identifies the dual of \(S(\alpha,r)[\ell^{n+1}]\) with
\(B/(\ell^{n+1},T-t_r)B\). Indeed the dual of the torsion
subgroup killed by a scalar is the quotient of the compact dual
by that scalar. In this quotient \((1+T)^{\ell^n}-1\) acts as
\(u_r^{\ell^n}-1\), of valuation \(n+1+v_\ell(r)\) by the
binomial theorem. Its action is therefore zero. We may first
take the quotient in the first isomorphism and then impose
\(\ell^{n+1}=0\) and \(\gamma=u_r\), giving the second formula.
\end{proof}

\begin{theorem}[A finite test for the full strict length to exceed a given depth]\leavevmode
\label{thm:scalarlevel467}
Let \(n\ge1\) and \(k(r)=\operatorname{length}_O S(\alpha,r)\)
for the possible surviving orbit. Then \(k(r)>n\) if and only if
\[
 \begin{gathered}
         H_n\simeq O/\ell^{n+1}O\quad\text{as an \(O\)-module},
         \qquad
         \gamma|_{H_n}=u_r\,\mathrm{id}_{H_n}.
 \end{gathered}
\]
If these two conditions fail for a weight, its full length is
the length of the finite quotient in
Lemma~\ref{lem:splitcyclolevel467}.
\end{theorem}
\begin{proof}
Suppose \(k(r)>n\ge1\). In the proof of
Theorem~\ref{thm:collisionbranch464} a relation with nonunit
linear coefficient would force \(k(r)=1\). The other
case of that proof identifies \(B\) with \(O/\ell^bO\), with
\(b=\infty\) denoting \(O\), and \(\gamma\) acts by \(1+z\)
for some \(z\in\ell O\) of valuation one. The formula for
\(k(r)\) gives \(b\ge n+1\) and \(1+z\equiv u_r\bmod
\ell^{n+1}O\). Raising \(1+z\) to the \(\ell\)-th power
increases the valuation of its difference from one by exactly
one, as \(\ell\) is odd and \(O\) is unramified. Induction
gives \(v_\ell((1+z)^{\ell^n}-1)=n+1\). The first isomorphism
of Lemma~\ref{lem:splitcyclolevel467} therefore gives precisely
the cyclic group and action asserted.

Conversely, under the displayed conditions both relations in
the second quotient of that lemma already hold on \(H_n\).
The dual of \(S(\alpha,r)[\ell^{n+1}]\) then has length
\(n+1\), forcing \(k(r)>n\). Failure of either condition
implies \(k(r)\le n\), in which case the entire scalar group
is killed by \(\ell^{n+1}\). The same quotient computes its
full length.
\end{proof}

\begin{corollary}[One class group at level twenty-five on the period interval]\leavevmode
\label{cor:leveltwentyfive467}
Fix \((\ell,j)\) and the set of weights \(\mathcal R\) in
Corollary~\ref{cor:onedeepweight464}. In particular each
eligible point has \(x^{16}\le\ell\le2x^{16}\). The group
\(H_{25}\) with its ideal action determines the full strict
orders for all but at most one weight of \(\mathcal R\).
All those determined by failure of the test have \(k(r)\le25\).
The remaining possible weight, if it exists, has \(k(r)\ge26\)
and is specified as follows.

First require \(H_{25}\simeq O/\ell^{26}O\), and write \(v\)
for the unit by which \(\gamma\) acts on it. Require also that
\(v\) belong to the image of \(1+\ell\ZZ_\ell\) in
\((O/\ell^{26}O)^\times\). Failure of either requirement leaves
no such weight. Otherwise determine \(a\bmod\ell^{25}\) by
\((1+\ell)^{-a}=v\bmod\ell^{26}O\). The possible weight is
the element of \(\mathcal R\), if present, satisfying
\[
          r\equiv a\bmod{\ell^{25}},\qquad
          r\equiv(\ell-1)/2\bmod{\ell-1}.
\]
The auxiliary field for this computation has degree
\(2^j\ell^{25}\) and conductor dividing \(2^{j+2}\ell^{26}\).
\end{corollary}
\begin{proof}
If the orbit is absent at level zero, all orders are already
one. Otherwise apply Theorem~\ref{thm:scalarlevel467} with
\(n=25\). An endomorphism of the indicated cyclic \(O\)-module
is given by a unique scalar modulo \(\ell^{26}\). The reduction
of \(1+\ell\ZZ_\ell\) has \(\ell^{25}\) elements and is
generated by \(1+\ell\), since
\(v_\ell((1+\ell)^{\ell^i}-1)=i+1\). The natural map of this
finite group into \((O/\ell^{26}O)^\times\) is injective by
unramifiedness. This gives the claimed exponent test. Its value
is independent of the generator chosen for the cyclic module.
Coefficient conjugation also preserves the test because every
\(u_r\) lies in \(\ZZ_\ell\).

The two congruences define one class modulo
\((\ell-1)\ell^{25}\). For each \(r\in\mathcal R\) there is
an \(x\) with \(x^{16}\le\ell\) and \(r<t(x)\), so all
such \(r\) lie between zero and \(\ell^{401/16}\), even if
they occur at different points of the interval. Since
\(401/16=25+1/16\) and \(\ell-1>\ell^{1/16}\) for
\(\ell\ge5\), the modulus exceeds this entire range. Thus
there is at most one matching integer. The theorem gives
its lower bound and the full orders at the other weights, by
the finite quotient already described.

Finally \(L\) has degree \(2^j\) prime to \(\ell\) and
conductor dividing \(2^{j+2}\ell\).
The field \(\QQ_{25}\) has degree \(\ell^{25}\) and is
contained in \(\QQ(\zeta_{\ell^{26}})\). Disjointness gives the
stated degree and the compositum gives the conductor bound.
These quantities are bounded by fixed powers of \(x\) on the
prescribed window. The descent in
Corollary~\ref{cor:shortcenter460} identifies the scalar
components on the same \(E_x\) used for the periods.
\end{proof}

If the congruences select a weight, the level-twenty-five test
gives a lower bound for its length. No bound for the excess above
that depth follows here. At any higher level the theorem gives
the same finite test, and finiteness of the strict group implies
eventual failure for each fixed weight. A uniform bound for that
level as \(\ell,j,r\) vary would give additional control of the
sum in Corollary~\ref{cor:supportmass465}.



\subsection{Full cohomology lengths at locally split weights}
\label{sec:locallysplit502}

Use the common fields \(E_x\subseteq E_*\) and the interval
\([V,V^{3/2}]\) of Theorem~\ref{thm:conductor232}. Write
\(n_*=[E_*:\QQ]\), let \(N_*\) be its cyclotomic modulus, and
put \(t_*=\lfloor V^{1203/2}\rfloor\), \(t_-=\lfloor V^{401}\rfloor\).
Fix a prime \(\ell\ge5\) with \(\ell\nmid n_*N_*\), and take
\(s=0\) or \(s=(\ell-1)/2\). For a character \(\alpha\) of
\(\operatorname{Gal}(E_*/\QQ)\) with \(\alpha(c)=(-1)^s\), let
\(\tau=\alpha^{-1}\omega^{-s}\), where \(\omega\) is the
Teichm\"uller lift of the cyclotomic character modulo \(\ell\).
Let \(L\) be its fixed-kernel field, \(\Delta=\operatorname{Gal}(L/\QQ)\),
\(O=\ZZ_\ell[\tau(\Delta)]\), and \(J=\operatorname{Frac}(O)\).
The field \(L\) is totally real, \(|\Delta|\) is prime to \(\ell\),
and \(O\) is unramified. Subscript \(\tau\) means the factor of
\(\ZZ_\ell[\Delta]\) determined by the coefficient Frobenius orbit
of \(\tau\). We allow the trivial character, in which case \(L=\QQ\).

For \(1\le r<t_*\) in the residue \(s\bmod{\ell-1}\), write
\(k_\tau(r)=\operatorname{length}_O S(\alpha,r)\). The group
\(S(\alpha,r)\) consists of the classes over \(\QQ\) with
coefficients \((J/O)(\alpha^{-1}\kappa_\ell^{-r})\) that are
unramified away from \(\ell\) and restrict to zero at \(\ell\).
Put \(t_r=(1+\ell)^{-r}-1\). Let \(\QQ_1\) be the first layer of
the cyclotomic \(\ZZ_\ell\)-extension of \(\QQ\), put \(L_1=L\QQ_1\),
and choose \(\gamma\) with principal cyclotomic value \(1+\ell\).
Its action on \(L_1\) fixes \(L\). Define
\[
 C_1=\left((\Cl(L_1)/\langle[w]\mid w\mid\ell\rangle)
                    [\ell^\infty]\right)_\tau,\qquad
 D_1=C_1/\Phi_\ell(\gamma)C_1,\qquad
 \Phi_\ell(Y)=1+Y+\cdots+Y^{\ell-1}.
\]
Fix a primitive \(\ell\)-th root \(\zeta\) and put
\(O'=O[\zeta]\), \(J'=\operatorname{Frac}(O')\). The action of
\(\zeta\) by \(\gamma\) makes \(D_1\) an \(O'\)-module. Set
\(h_\tau=\operatorname{length}_{O'}D_1\). Since \(O'\) and \(O\)
have the same residue field, \(h_\tau\le\operatorname{length}_O C_1\).

\begin{lemma}[A primitive cyclotomic specialization at a split prime]\leavevmode
\label{lem:primitiveclass502}
There is a finitely generated compact \(O[[T]]\)-module \(B\) with
\[
 S(\alpha,r)\simeq
   \operatorname{Hom}_{O,\mathrm{cont}}(B/(T-t_r)B,J/O),\qquad
 (B\otimes_O O')/(T-(\zeta-1))(B\otimes_O O')\simeq D_1 .
\]
These conclusions hold also when every prime of \(L\) above
\(\ell\) has local degree one. If \(\tau=1\), then \(B=0\).
\end{lemma}
\begin{proof}
Put \(L_\infty=L\QQ_\infty\), where \(\QQ_\infty/\QQ\) is
cyclotomic with group \(\Gamma=\langle\gamma\rangle\simeq\ZZ_\ell\).
Take the maximal abelian pro-\(\ell\) quotient of \(G_{L_\infty}\)
and kill the closed subgroup generated by inertia away from
\(\ell\) and all decomposition groups above \(\ell\).
Let \(B\) be its \(\tau\)-factor, with \(T=\gamma-1\).
Conjugation gives the action because \(\Delta\times\Gamma\) is
the quotient over \(\QQ\) and inner conjugations act trivially
on the abelian group. For the integral weights the coefficient
on this quotient is \(\tau\) on \(\Delta\) and \(1+t_r\) on
\(\gamma\). At the auxiliary specialization use \(J'/O'\)
with the same finite character and with \(\gamma\) acting by \(\zeta\).

Restriction from \(\QQ\) to \(L_\infty\) identifies the global
classes at each of these points with equivariant homomorphisms.
Indeed the first and second cohomology of the quotient vanish.
For \(\tau\ne1\), average over \(\Delta\), whose invariants
in the coefficient module vanish since some \(\tau(\delta)-1\)
is a unit. For \(\tau=1\), the second group vanishes by
\(\ell\)-cohomological dimension one of \(\Gamma\), and the
first is the cokernel of multiplication by \(t_r\) or by
\(\zeta-1\) on a divisible module. Both scalars are nonzero.

Away from \(\ell\) the inertia image in the quotient has
prime-to-\(\ell\) order, so inflation contributes no inertia
classes. At \(\ell\) the decomposition image is
\(\Delta_w\times\Gamma\). Its projections are onto the respective
local images, of coprime supernatural orders, and the cyclotomic
extension is totally ramified. For the product assertion one can
work in every finite layer, where a subgroup surjecting onto two
factors of coprime orders has the full product order. Its first
cohomology vanishes by averaging if \(\tau|_{\Delta_w}\ne1\),
and otherwise by the same surjective scalar multiplication on
the divisible module. Thus restriction detects zero at \(\ell\).
At the real place the coefficient prime is odd. Equivariance
on the resulting quotient gives the first isomorphism and gives
the continuous \(O'\)-dual of the auxiliary specialization of
\(B\). The scalar extension here is a finite direct sum as a
compact \(O\)-module.

For \(\tau\ne1\) compute that auxiliary group a second time,
by restriction to \(L_1\). Global averaging still kills degrees
one and two for \(\Delta\times C_\ell\). In the local quotient
\(\Delta_w\times C_\ell\) the only further case is the trivial
action of \(\Delta_w\). The first cohomology of \(C_\ell\)
is then \(\ker(1+\zeta+\cdots+\zeta^{\ell-1})/
(\zeta-1)(J'/O')=0\). Away from \(\ell\) inertia is again
of prime-to-\(\ell\) order. Class field theory over \(L_1\)
identifies the unramified abelian quotient split above \(\ell\)
with the stated quotient of its ideal class group
\cite{XI-VIII-rayclass444}. Equivariance selects \(C_1\) and imposes
\(\gamma=\zeta\) after extension of coefficients. The resulting
module is \(D_1\), since \(O'=O[Y]/(\Phi_\ell(Y))\) and the
additional relation identifies \(Y\) with the action of \(\gamma\).
Both calculations give its continuous dual. Dualizing recovers
the compact modules by finite abelian duality and inverse limits.
Over a finite extension of \(\ZZ_\ell\) the dualizing divisible
module can be taken as the fraction field modulo the integers,
since the inverse different is a principal fractional ideal.

For \(\tau=1\) the first cohomology of the global finite quotient
\(C_\ell\) is zero by the same calculation. Restriction therefore
injects the auxiliary strict group into the dual of the split
\(\ell\)-class group of \(\QQ_1\). This class group is zero.
Indeed its class field is Galois over \(\QQ\) with finite
\(\ell\)-group \(G\). The inertia at \(\ell\) has order \(\ell\)
and is cyclic. Its normal closure is \(G\), since the fixed field
of that closure would be everywhere unramified over \(\QQ\);
the odd degree excludes real ramification and Minkowski excludes
a nontrivial such field. A proper subgroup of a finite
\(\ell\)-group lies in a maximal subgroup, which is normal.
The inertia therefore equals \(G\). The split class field is
\(\QQ_1\) itself. Both auxiliary modules consequently vanish.

In all cases the finite auxiliary specialization surjects onto
\((B\otimes_O O')/(\zeta-1,T)(B\otimes_O O')\), which identifies
with \(B/(\ell,T)B\). Lift a basis of the latter to \(B\).
The submodule generated over \(O[[T]]\) by the lifts is closed
as the continuous image of a compact finite free module.
Its quotient \(Y\) satisfies \(Y=(\ell,T)Y\). In every finite
stable quotient \(T\) is nilpotent modulo \(\ell\), because
\(\gamma\) acts with \(\ell\)-power order. These quotients separate
points, by the open stable subgroup argument of
Lemma~\ref{lem:collisionfamily464}. Thus \(Y=0\), including the
case of an empty basis.
\end{proof}

\begin{theorem}[A bound for full lengths after removing at most one weight per root]\leavevmode
\label{thm:splitweights502}
Let \(I_s=\{1\le r<t_*:r\equiv s\bmod{\ell-1}\}\), \(N=|I_s|\),
and \(h=h_\tau\). There is a set \(D_\tau\subseteq I_s\) with
\(|D_\tau|\le h\) such that, if \(N>0\),
\[
 \begin{aligned}
 k_\tau(r)&\le h(1+\lfloor\log_\ell t_*\rfloor)
                          &&(r\in I_s\setminus D_\tau),\\
 \sum_{r\in I_s\setminus D_\tau} k_\tau(r)
          &\le h\left(N+\frac{N-1}{\ell-1}\right).
 \end{aligned}
\]
In particular all the strict groups vanish if \(D_1=0\).
\end{theorem}
\begin{proof}
If \(h=0\) the compact generation argument in the lemma gives
\(B=0\). The empty progression also requires no choice. Otherwise
put \(B'=B\otimes_O O'\), \(\pi=\zeta-1\), \(e=\ell-1\), and
normalize \(v(\ell)=1\), so \(v(\pi)=1/e\). Present \(B'\)
as a quotient of \(O'[[T]]^g\). At \(T=\pi\) its relations map
onto a full-rank sublattice of \((O')^g\) with cokernel of
length \(h\). Lift a basis of this sublattice to relation columns.
Their determinant \(F(T)\) satisfies \(v(F(\pi))=h/e\).
At \(t_r\) these columns are again relations, hence
\(k_\tau(r)\le v(F(t_r))\), allowing infinity when the determinant
vanishes. Indeed extension from \(O\) to \(O'\) multiplies
length by \(e\), and \(v\) divides the \(O'\)-length of a
determinant by that same factor.

By Weierstrass preparation over the complete discrete valuation
ring \(O'\), write \(F=\pi^\mu U\prod_i(T-\beta_i)\) over a
finite splitting field, where \(U\) is a unit power series,
the \(\beta_i\) have positive valuation, and roots are repeated
with multiplicity. The unit has valuation zero at all our
points. For \(v(\beta_i)<1\) one has
\[
 v(t_r-\beta_i)=v(\beta_i)\le e\,v(\pi-\beta_i).
\]
For \(v(\beta_i)<1/e\) the two distances have equal valuation.
For \(v(\beta_i)\ge1/e\) the right side is at least one.
Each root of valuation at least one contributes exactly \(1/e\)
at \(\pi\). There are therefore at most \(h\) such roots.
For each choose a weight \(r_i\in I_s\) maximizing
\(v(t_{r_i}-\beta_i)\), allowing an infinite maximum. Take
\(D_\tau\) to be the union of these choices. For \(r\ne r_i\)
the ultrametric inequality gives
\[
 v(t_r-\beta_i)\le v(t_r-t_{r_i})=1+v_\ell(r-r_i).
\]
The equality follows by the binomial theorem for the principal
unit \(1+\ell\). The constant \(\pi^\mu\), the roots of
valuation below one, and one unit of valuation per remaining
root cost at most \(e\,v(F(\pi))=h\). Thus outside \(D_\tau\)
the bound is \(h+\sum_i v_\ell(r-r_i)\), with the sum only
over the selected roots. This gives the first inequality.

Write the progression as \(r_{\min}+a(\ell-1)\)
for \(0\le a<N\), and write \(a_i\) for the index of \(r_i\).
Summing \(v_\ell(r-r_i)\) over the whole progression except
\(r_i\) gives \(v_\ell(a_i!)+v_\ell((N-1-a_i)!)\).
Counting multiples of powers of \(\ell\) bounds this by
\((N-1)/(\ell-1)\). Summing over at most \(h\) roots proves the
second inequality.
\end{proof}

\begin{corollary}[One sparse exceptional set for both inertia residues]\leavevmode
\label{cor:fullcollision502}
There is one set \(D_V\subseteq[1,t_*]\cap\ZZ\) such that
uniformly for \(t_-\le T\le t_*\),
\[
 |D_V\cap[1,T]|\ll_{\psi,\varepsilon} T/H(T),
\]
where \(H\) is defined in Section~\ref{sec:shiftedsupport465}.
At every \(x\in[V,V^{3/2}]\), take all the primes coprime
to \(n_*N_*\) and at least five, both residues above, and the
character orbits occurring on \(E_x\). With \(f_\tau=\operatorname{rank}_{\ZZ_\ell}O\)
denoting the unramified degree, their full lengths satisfy
\[
 \sum_{\substack{(\ell,r,\tau)\\1\le r<\lfloor x^{401}\rfloor\\
                              r\notin D_V}}
       f_\tau k_\tau(r)\log\ell
       \ll_{\psi,\varepsilon} n_*^2 t_*^3(\log(2V))^3 .
\]
Here a triple is included only when \(r\) is in the residue
defining \(\tau\). The bound includes the zero-residue characters
for which \(\alpha(\ell)=1\).
\end{corollary}
\begin{proof}
Inflating \(\alpha\) to \(E_*\) preserves the scalar group over
\(\QQ\), and the integral descent to \(E_x\) is that of
Theorem~\ref{thm:levelstrict345} and Lemma~\ref{lem:locallevel345}.
Take \(D_V\) to be the union of the sets in
Theorem~\ref{thm:splitweights502} for the characters of \(E_*\).
Only \(\ell\le2t_*+1\) occur since \(\ell-1\mid 2r\).
The field \(L_1\) has degree at most \(2n_*\ell\) and conductor
dividing \(N_*\ell^2\). By the conductor--discriminant formula
and the ideal-counting bound used in
Theorem~\ref{thm:nonsplitprefix466},
\[
 f_\tau h_\tau\log\ell\le \log|\Cl(L_1)|
                  \le 2n_*\ell\log(2N_*\ell^2).
\]
For \(\ell\le Y=V^{100}\) there are at most \(2n_*\) orbits
over the two residues. Hence the total number of their selected
weights is \(O_{\psi,\varepsilon}(n_*^2Y^2\log(2V))\).
The estimates \(n_*\ll_\psi V^{3+12\varepsilon}\) and
\(\log N_*\ll_{\psi,\varepsilon}\log V\), with
\(\varepsilon<1/52\), make this \(o(t_-/\log t_*)\).
Every weight selected for a larger prime belongs to the set in
Lemma~\ref{lem:shiftdivisors465}. Its count in every stated
prefix is \(O(T/H(T))\), which proves the first assertion.

For the mass, each \(k_\tau(r)\) outside the union is
\(O_{\psi,\varepsilon}(n_*t_*(\log(2V))^2)\) by the class number
bound and the first inequality of the theorem. At each prime
and weight the sum of the orbit degrees is at most \(n_*\).
There are at most \(2t_*+1\) primes and \(t_*\) weights.
Multiplying and using \(\log(2t_*+1)\ll\log(2V)\) gives
the asserted polynomial bound.
\end{proof}

For sufficiently large \(V\) every individual length bounded
here is smaller than \(b_0(x)\) on the entire interval.
At a weight in \(D_V\) the reference class group bounds the
number of possible exceptions. Section~\ref{sec:missing} still
requires control of the weighted sizes of those exceptional
groups and of the projected theta periods on these fields.


\subsection{Ideal classes at the primes testing the theta periods}
\label{sec:jointclasses503}

Fix one interval of Corollary~\ref{cor:onedeepweight464}. For
each pair \((\ell_i,j_i)\) from its dyadic family include the
exceptional weight \(r_i\) if its strict length \(k_i\) exceeds
25. There is at most one such weight and one such orbit per
pair by Corollary~\ref{cor:leveltwentyfive467}. Equivalently,
one may work with any subset of this finite list. Put
\(f_i=2^{j_i-2}\), let \(O_i\) be the corresponding unramified
integer ring, and let \(L_i\) be the totally real cyclic field
of the residual character \(\tau_i\) of order \(2^{j_i}\).
Choose an integer \(1\le n_i<k_i\) and set
\[
 K_i=L_i\QQ_{i,n_i},\qquad
 H_i=\left((\Cl(K_i)/\langle[v]\mid v\mid\ell_i\rangle)
                   [\ell_i^\infty]\right)_{\tau_i},
\]
where \(\QQ_{i,n_i}\) is the cyclotomic layer of degree
\(\ell_i^{n_i}\) over \(\QQ\). Theorem~\ref{thm:scalarlevel467}
gives \(H_i\simeq O_i/\ell_i^{n_i+1}O_i\). The group
\(\operatorname{Gal}(L_i/\QQ)\) acts through the lift of
\(\tau_i\), and the cyclotomic generator
acts on \(H_i\) by \((1+\ell_i)^{-r_i}\). Let \(M_i/K_i\)
be the abelian class field with group \(H_i\), using arithmetic
Frobenius as the Artin map on ideals. The defining subgroup is
stable under \(\operatorname{Gal}(K_i/\QQ)\), so \(M_i/\QQ\)
is Galois. All primes above \(\ell_i\) split in \(M_i/K_i\).

Let \(J/\QQ\) be any finite Galois field containing the
coefficients, torsion data, and normalizations of the chosen
finite list of elliptic theta tests on the interval, as in
Theorem~\ref{thm:commonsplit483}. Additional residual
splitting fields required for trace tests may be included in \(J\).
Write \(J^{\rm ab}\) for its largest subfield abelian over \(\QQ\).
Take a finite abelian field \(A/\QQ\) containing \(J^{\rm ab}\),
all the \(K_i\), and the period fields \(E_x\) on the interval.
Such an \(A\) exists because the \(K_i\) and \(E_x\) are abelian
and there are only finitely many distinct \(E_x\).
Put \(F=AJ\) and \(M=A\prod_i M_i\), where the product denotes
compositum.

\begin{lemma}[Independence over the abelian fields of the interval]\leavevmode
\label{lem:classdisjoint503}
Restriction gives \(\operatorname{Gal}(M/A)=\prod_i H_i\).
For uniquely determined integers \(0\le a_i\le n_i+1\), it gives
\[
 \operatorname{Gal}(MF/F)=\prod_i\ell_i^{a_i}H_i,\qquad
 \prod_i\ell_i^{f_i a_i}=[M\cap F:A]\ \mid\ [J:J^{\rm ab}].
\]
In particular enlarging the abelian field in this construction
costs no factor in the stated upper bound.
\end{lemma}
\begin{proof}
The conjugation action on \(H_i\) of a lift of a generator of
\(\operatorname{Gal}(L_i/\QQ)\) with trivial cyclotomic component
is multiplication by a primitive \(2^{j_i}\)-th root of unity.
Subtracting one gives a unit in \(O_i\), since the reduction of
a root of unity of order prime to \(\ell_i\) has the same order.
The commutators of this single lift with \(H_i\) thus fill \(H_i\).
The quotient by \(H_i\) is abelian. Consequently \(K_i\) is the
largest abelian subfield of \(M_i/\QQ\) and \(M_i\cap A=K_i\).

Let \(W=\prod_i H_i\) and \(D=\operatorname{Gal}(A/\QQ)\).
There is a well-defined action of \(D\) on \(W\) through its
quotients \(\operatorname{Gal}(K_i/\QQ)\). We describe every
\(D\)-stable subgroup of this finite abelian group. First
separate its primary parts. For a given prime \(\ell\) use the
subgroup \(D'\subset D\) of elements of order prime to \(\ell\).
On each factor with \(\ell_i=\ell\) it acts by \(\tau_i\).
Indeed the projection of \(D'\) to the cyclotomic layer is
trivial and its projection onto the prime-to-\(\ell\) group
\(\operatorname{Gal}(L_i/\QQ)\) is surjective. The residual
characters at this prime belong to distinct Frobenius orbits.
Their orders \(2^{j_i}\) differ, because the list contains at
most one entry for each \((\ell,j_i)\).

The ring \(\ZZ_\ell[D']\) has integral idempotents separating
these orbits. To see this, over an unramified coefficient
extension containing the character values use the character
idempotents with denominator \(|D'|\). Summing over one
Frobenius orbit gives coefficients in \(\ZZ_\ell\). On the
selected factor the image of the ring is \(O_i\), since the
values of \(\tau_i\) generate \(O_i\). A stable subgroup is
therefore the product of \(O_i\)-submodules of the separate
cyclic modules \(H_i\). Each such submodule is
\(\ell_i^{a_i}H_i\) with \(0\le a_i\le n_i+1\).

The restriction of \(\operatorname{Gal}(M/A)\) into \(W\) is
injective, and its image is \(D\)-stable by conjugation in
\(\operatorname{Gal}(M/\QQ)\). Each projection is surjective by
\(M_i\cap A=K_i\). The subgroup description gives the first
assertion. Restriction identifies \(\operatorname{Gal}(MF/F)\)
with \(\operatorname{Gal}(M/M\cap F)\), which is again stable
because \(F/\QQ\) is Galois. Its index in \(W\) is the stated
product. This index divides \([F:A]=[J:J^{\rm ab}]\), since
\(J\cap A=J^{\rm ab}\).
\end{proof}

\begin{theorem}[One prime representing the visible part of every exceptional class group]\leavevmode
\label{thm:classfrob503}
Fix \(0<\delta<1/4\). For all sufficiently large \(X\), a set
of rational primes \(q\in[X,2X]\) of positive limiting relative
density among those split in \(F\) has both of the following
properties. For every marked elliptic factor, at every place
over \(q\), its recurrence scalar has multiplicative order
greater than \(X^{1/4-\delta}\). For each \(i\), the
\(\tau_i\)-component of the class of a prime \(\mathfrak q_i\)
of \(K_i\) over \(q\) has additive order
\[
                  \ell_i^{\,n_i+1-a_i}\quad\text{in }H_i .
\]
This order is independent of the choice of prime above \(q\).
The \(O_i\)-multiples of the class fill
\(\ell_i^{a_i}H_i\). Every prime split in \(F\), away from
ramification in \(MF\), gives a class in this subgroup.
\end{theorem}
\begin{proof}
Choose \(g_i\) generating \(\ell_i^{a_i}H_i\) over \(O_i\);
use zero for the zero group. By the lemma the tuple defines one
element of \(\operatorname{Gal}(MF/F)\). Viewed in the finite
group \(\operatorname{Gal}(MF/\QQ)\), its conjugacy class gives
a positive Chebotarev density of rational primes. All split in
\(F\), and the dyadic-interval form of the theorem gives a
positive limiting relative density there
\cite[\S1A]{XI-VIII-I-tz50}. For each such prime choose a prime of
\(MF\) with Frobenius equal to the chosen tuple. The restrictions
to \(M_i/K_i\) are the arithmetic Frobenius elements of the
respective primes of \(K_i\), since their residue degrees over
\(\QQ\) equal one. Global reciprocity identifies them with the
desired projections of ideal classes. Conjugating the prime acts
on each \(H_i\) by a unit and preserves both the order and the
submodule generated. All primes split in \(F\) give elements in
the group of the lemma, which also proves the upper assertion.
Finally Theorem~\ref{thm:commonsplit483} discards at most
\(O(X^{1-2\delta})\) primes for the entire list of elliptic data.
For fixed fields this is \(o(X/\log X)\), so the remaining
relative density in the specified conjugacy class tends to one.
\end{proof}

\begin{corollary}[A single degree bound for the loss in the weighted lengths]\leavevmode
\label{cor:classmass503}
Take \(n_i=k_i-1\) and choose one prime \(q\) by the theorem.
Write \(\ell_i^{d_i(q)}\) for the order of the projected ideal
class at \(q\). For every subset \(S\) of the interval list and
arbitrary cutoffs \(b_i\ge0\), put \((y)_+=\max(y,0)\). Then
\[
 \begin{aligned}
 0\ \le\ 
 \sum_{i\in S} f_i\log\ell_i
     \big[(k_i-b_i)_+-(d_i(q)-b_i)_+\big]
       \ \le\ \log[J:J^{\rm ab}].
 \end{aligned}
\]
If \(J\) is abelian over a number field \(K\) abelian over
\(\QQ\), and \(K\subseteq A\cap J\), then \(d_i(q)=k_i\)
for every \(i\) with \(L_i\not\subseteq K\). In particular this
holds for all entries if \(K\) is imaginary quadratic.
\end{corollary}
\begin{proof}
The theorem gives \(k_i-d_i(q)=a_i\), and subtracting a fixed
cutoff and taking the positive part decreases this difference.
Sum the resulting bounds and take the logarithm of the index
in Lemma~\ref{lem:classdisjoint503}. In particular the same
prime works for the subset occurring on each \(E_x\) and for the
actual cutoff \(b_i=b_0(x)\).

In the last case \(F=AJ\) is abelian over \(K\), as a
compositum of two extensions abelian over \(K\). If
\(L_i\not\subseteq K\), choose an element of
\(\operatorname{Gal}(A/K)\) nontrivial on \(L_i\). Its action
on \(H_i/\ell_i H_i\) is multiplication by a nonidentity value
of the faithful residual character. Its action minus one is
therefore invertible on \(H_i\). On
\(\operatorname{Gal}(M\cap F/A)\) its action by conjugation is
trivial since \(M\cap F\) is abelian over \(K\). The quotient
\(H_i/\ell_i^{a_i}H_i\) in that group must be zero. Thus \(a_i=0\).
Every \(L_i\) here is totally real of degree \(2^{j_i}\ge8\),
giving the assertion for imaginary quadratic \(K\).
\end{proof}

The choice \(n_i=25\) in the theorem uses fixed layer indices
throughout the interval. The corollary uses the full depths and
allows its auxiliary prime to depend on them. To obtain the
bound of Section~\ref{sec:missing} from this realization one
still needs an upper bound for the orders \(d_i(q)\) and control
of the other components. The large orders of the elliptic
recurrence scalars at this prime provide the tests of
Corollary~\ref{cor:projectsymbol483}; their consequences for
the actual projected period determinant require its global
coefficient and rank estimates.

\section{Units, regulators and relative class numbers}\label{nav:chapter-dyadic-units-regulators}

\subsection{Global units modulo the coefficient prime and the opposite root}
\label{sec:oppositeunits472}

In the family of Section~\ref{sec:sympglobal469}, retain the
assumption \(S(\alpha,r)\ne0\) and write \(\tau=\bar\Psi\), the
character of the \emph{opposite} root. The cyclic field \(L\) cut
out by \(\tau\) is also the residual field for the negative root.
Its group \(\Delta\) has order \(2^j\); at \(\ell\) its inertia
subgroup \(I\) has order two, since \(\alpha\) is unramified and
\(\bar\kappa_\ell^r\) is the ramified quadratic character. The
character \(\tau\) restricts to the sign of \(I\). Let
\(H_{\rm rel}(\tau)\subset H^1(\QQ,k(\tau))\) consist of the
classes unramified at all finite primes different from \(\ell\),
with unrestricted localization at \(\ell\).

\begin{lemma}[The ray class group at the rational ideal \(\ell O_L\)]\leavevmode
\label{lem:rayconductor472}
Restriction to \(L\) gives a \(k\)-linear isomorphism
\[
 H_{\rm rel}(\tau)\simeq
   \operatorname{Hom}_{\Delta}\big(\Cl_{\ell O_L}(L),k(\tau)\big),
\]
where the maps on the right are additive group homomorphisms.
The dimension is at most one. Here the modulus has exponent
two at each prime over \(\ell\), independent of the weight.
\end{lemma}
\begin{proof}
Restriction to \(L\) identifies the global cohomology with the
\(\Delta\)-equivariant homomorphisms from \(G_L\): averaging over
\(\Delta\) kills the two outer terms in inflation--restriction.
At primes away from \(\ell\) the same argument over inertia shows
that the unramified conditions are exactly unramifiedness of the
restricted homomorphism. Real places have order at most two, so
they impose no distinction for these odd-primary homomorphisms.

Fix \(\lambda\mid\ell\). In \(L_\lambda\) the absolute
ramification index is two. Write \(U^a=1+\lambda^a O_{L_\lambda}\)
for \(a\ge1\). The \(\ell\)-adic logarithm gives
\((U^1)^\ell=U^3\). Indeed logarithm and exponential converge and
are mutually inverse on the maximal ideal, since the normalized
valuation of each element of that ideal is at least one, strictly
greater than \(2/(\ell-1)\). They identify taking the \(\ell\)-th
power with multiplication by \(\ell\), of valuation two.
Under local reciprocity the homomorphism therefore kills \(U^3\).
On \(U^2/U^3\simeq\lambda^2/\lambda^3\) inertia acts trivially.
To see this, the nonidentity tame inertia element acts by \(-1\)
on \(\lambda/\lambda^2\) via its faithful tame character and
fixes the residue field. Its square character on
\(\lambda^2/\lambda^3\) is one. On the target it acts as \(-1\),
so equivariance kills \(U^2/U^3\) as well. This proves the
conductor bound. Conversely a homomorphism of the stated ray
class group is unramified away from \(\ell\) and provides the
required global class by reciprocity \cite{XI-VIII-rayclass444}. Finally
localization at \(\ell\) on \(H_{\rm rel}(\tau)\) has strict
kernel zero by Lemma~\ref{lem:oppositeroot469}. Its target has
dimension one by the local Euler characteristic and the local
invariant vanishings in Theorem~\ref{thm:sympordinary468}.
\end{proof}

Let \(\kappa_\lambda=O_L/\lambda\), and let \(M_0\) be the least
common multiple of the integers \(|\kappa_\lambda^\times|\) for
\(\lambda\mid\ell\). It is prime to \(\ell\). Define
\[
 \begin{aligned}
 B_\ell&=\bigoplus_{\lambda\mid\ell}\lambda/\lambda^2,
 &e_\tau&=\frac{1}{|\Delta|}\sum_{\sigma\in\Delta}
                                  \tau(\sigma)^{-1}\sigma,\\
 b(u)&=\big(u^{M_0}-1\bmod\lambda^2\big)_{\lambda\mid\ell}
       \in B_\ell &&(u\in O_L^\times).
 \end{aligned}
\]
Here \(e_\tau\) acts after tensoring \(B_\ell\) from \(\mathbb F_\ell\)
to \(k\).

\begin{theorem}[A one-dimensional unit calculation for the ordinary root]\leavevmode
\label{thm:unitray472}
The space \(e_\tau(B_\ell\otimes_{\mathbb F_\ell}k)\) has
dimension one. The relaxed group \(H_{\rm rel}(\tau)\) vanishes
if and only if some global unit \(u\) has \(e_\tau b(u)\ne0\).
If every unit gives zero, \(H_{\rm rel}(\tau)\) has dimension one.
One may check any generating set of \(O_L^\times\). All of these
data are common to the half-residue weights for a fixed
\(\ell,j,\alpha\).
\end{theorem}
\begin{proof}
Consider first the ordinary class group of \(L\). Its faithful
\(\tau\)-component at \(\ell\) vanishes. Indeed imposing splitting
at the primes above \(\ell\) gives the quotient in
Lemma~\ref{lem:oppositeroot469}, with zero component there.
The subgroup generated by the omitted prime classes has trivial
\(I\)-action because \(\Delta\) is abelian and each such prime
is fixed by \(I\). It has zero \(\tau\)-component since \(\tau|_I\)
is nontrivial. One can take components integrally on the finite
\(\ell\)-primary groups by lifting \(\tau\) to its prime-to-\(\ell\)
roots of unity. A component with zero residual quotient is zero
by Nakayama's lemma.

The ray class exact sequence now shows that the only contribution
at \(\tau\) comes from \((O_L/\ell O_L)^\times\) modulo the global
units. Its \(\ell\)-primary subgroup is precisely \(B_\ell\):
the product of the residue unit groups has order prime to \(\ell\)
and the kernel is identified with the additive groups
\(\lambda/\lambda^2\) by subtracting one. Raising a unit to
\(M_0\) kills the prime-to-\(\ell\) part and acts as a unit scalar
on the \(\ell\)-part. The images of \(b(u)\) therefore generate
the contribution to be divided out. The map \(u\mapsto b(u)\)
is additive with respect to multiplication of units, as products
of two displacements lie in \(\lambda^2\). Taking components and
extending the residue field are exact here. In particular, no
extension from the ordinary class group changes this component.

To find its dimension before taking units, extend the coefficient
field temporarily to an algebraic closure of \(\mathbb F_\ell\).
The line \(\lambda/\lambda^2\) over \(\kappa_\lambda\) then splits
into one line per residue-field embedding. The group \(\Delta\)
acts transitively on the pairs consisting of a prime and such an
embedding; the stabilizer of one pair is its inertia group \(I\).
On its line the action is the nontrivial tame character of \(I\).
The whole module is therefore the representation induced from
this sign. Frobenius reciprocity gives multiplicity one for the
character \(\tau\), since it restricts to that sign. Dimensions
descend to \(k\), where \(\tau\) is defined. Thus quotienting by
the projected \(b(u)\) gives dimension zero or one as asserted,
and its equivariant dual computes the relaxed group by the lemma.
Finally, changing \(r\) in the prescribed residue preserves \(\tau\)
and \(L\), hence also the unit calculation.
\end{proof}

\begin{corollary}[Eliminating one parameter from the trace congruence]\leavevmode
\label{cor:eliminateupper472}
Write \(E\) for the positive-weight relation after diagonal
elimination in Corollary~\ref{cor:firsttorus469}. Its linear
coefficient on \(x\), reduced modulo \(\ell\), is nonzero if
and only if the unit in Theorem~\ref{thm:unitray472} exists.
In this case \(e=yE\) eliminates \(U\) over \(O\), and there are
series \(f'_1,f'_2\in O[[V,W]]\) giving
\[
       C=O[[V,W]]/(f'_1,f'_2),\qquad
       C/(V,W)C=O/\ell^{k(r)}O .
\]
In particular a characteristic-zero specialization of \(C\) with
\(V=W=0\) is then impossible.
\end{corollary}
\begin{proof}
Consider global square-zero deformations on the positive symplectic
root over \(k[\epsilon]/(\epsilon^2)\). They are the global
\(\tau\)-cocycles unramified away from \(\ell\), normalized to
have value zero at \(h\). On this upper root every such cocycle
satisfies the local ordinary condition: it fixes the coordinate
flag and induces the split diagonal characters. The normalization
selects exactly one representative of every relaxed class because
\(\tau(h)-1\) is a unit. Away from \(\ell\) the equivalence between
unramifiedness of a class and factorization through \(H_q\) is
the same inertia averaging used in Theorem~\ref{thm:symprelations469}.

In that presentation the only tangent coordinate of weight
\((0,1)\) is \(x\). Of the relations only \(E\) can have a linear
term in this coordinate. Consequently the relaxed group vanishes
exactly when that coefficient is a unit, giving the equivalence
by Theorem~\ref{thm:unitray472}. On setting \(v=w=0\) in the
first torus presentation, weight \(+1\) forces
\(E=xg(xy)\) for some \(g\in O[[U]]\). Here \(g(0)\) is the same
linear coefficient. Thus \(e(U,0,0)=Ug(U)\) with \(g(0)\) a
unit under the stated condition. The formal inverse function
theorem solves \(e=0\) uniquely as \(U=H(V,W)\). The solution
has \(H(0,0)=0\), as the unique solution at \(V=W=0\) in the
formal neighborhood of zero is \(U=0\). Substitution in the two
\(f_i\) proves the presentation; Theorem~\ref{thm:exactstricttrace470}
gives the indicated quotient. At a specialization with \(V=W=0\)
the coefficient ring therefore factors through \(O/\ell^{k(r)}O\).
\end{proof}



\subsection{Norms of units at the exceptional dyadic weights}
\label{sec:unitnorms505}

Fix a faithful character of \(C_j=\QQ(\zeta_{2^{j+2}})^+\) in
Section~\ref{sec:finiteexception467}, with \(j\ge3\),
\(\ell\equiv5\bmod8\), and the stated conductor bound. At the
half-residue weights write \(k(r)\) for its full strict length.
Let \(L\) be the residual field, \(\Delta=\operatorname{Gal}(L/\QQ)\),
and \(\tau\) its faithful character with the unramified coefficient
ring \(O\). Write \(L_a=L\QQ_a\) as in
Lemma~\ref{lem:splitcyclolevel467}. For any number field \(K\)
write \(U(K)\) for its group of global units, with its natural
Galois action. The subscript \(\tau\) on a finitely generated
abelian group with \(\Delta\)-action means to tensor with
\(\ZZ_\ell\) and take the idempotent of the orbit of \(\tau\).
Set \(C_a=\Cl(L_a)_\tau\).

\begin{lemma}[The ramified prime and the ambiguous class sequence]\leavevmode
\label{lem:ambiguous505}
There is exactly one prime above \(\ell\) at each level \(L_a\).
Its square is principal. For \(0\le a<n\) let
\(G=\operatorname{Gal}(L_n/L_a)\), let \(N\) be the norm, and
write \(j_{a,n}\) for extension of ideal classes. There is an
exact sequence of \(O\)-modules
\[
 0\longrightarrow H^1(G,U(L_n))_\tau
  \longrightarrow C_a \xrightarrow{\,j_{a,n}\,} C_n^G
  \longrightarrow \bigl(U(L_a)/N U(L_n)\bigr)_\tau
  \longrightarrow 0 .
\]
\end{lemma}
\begin{proof}
The quadratic inertia character at \(\ell\) in
\(\tau=\alpha^{-1}\omega^{-(\ell-1)/2}\) has order two, and
\(\tau^2=\alpha^{-2}\) cuts out \(C_{j-1}\). The image of
\(\ell\) generates \((\ZZ/2^{j+1}\ZZ)^\times/\{\pm1\}\).
Indeed \(\ell\equiv5\bmod8\) gives
\(v_2(\ell^{2^t}-1)=2+t\) for \(t\ge0\), by successive
squaring, and no power is \(-1\bmod4\). Thus the Frobenius
on \(C_{j-1}\) has order \(2^{j-1}\). The decomposition group
in \(L\) contains the inertia of order two and surjects onto
this whole unramified quotient. There is one prime over
\(\ell\), with ramification index two. Adjoining \(\QQ_a\)
adjoins a totally ramified local extension of coprime degree
\(\ell^a\). The local degree is the product by coprimality,
and its ramification index is a multiple of the coprime
indices \(2\) and \(\ell^a\). Its residue field contains that
of \(L\), of degree \(2^{j-1}\). The product of these lower
bounds is the full global degree, so the ramification index
is \(2\ell^a\) and there is still just one prime.
The prime in \(\QQ_a\) is generated
by the norm to \(\QQ_a\) of \(1-\zeta_{\ell^{a+1}}\). Indeed
this generates the unique prime in the full cyclotomic field
and the intermediate residue degree is one. Its ideal in
\(L_a\) is the square of the unique prime there. In particular
\(C_a\) equals the split class group factor \(H_a\) of
Lemma~\ref{lem:splitcyclolevel467}.

In \(L_n/L_a\) the unique prime over \(\ell\) is totally
ramified and all other finite primes are unramified. Each unit
of \(L_a\) is a local norm at those other primes, since units
in an unramified local extension map onto the base units.
The real places split. The product formula for local Artin
symbols makes the symbol of the unit at the last prime trivial
as well. By the local norm theorem and the Hasse norm theorem
for a cyclic extension, every base unit is the norm of a field
element. We need to distinguish which of them are norms of
units. The prime-degree ambiguous class formulas are given in
\cite{XI-VIII-ambiguous505}. The following ideal calculation works for
every cyclic degree \(\ell^{n-a}\) and retains the character action.

Write \(I_K,P_K\) for fractional and principal ideals.
The quotient \(I_{L_n}^G/I_{L_a}\), with lower ideals extended,
is a sum of cyclic groups of orders the ramification indices.
This follows by taking the product of the primes in each
orbit as the generator of its invariant ideal group. Only the
prime at \(\ell\) contributes, and \(\Delta\) fixes it. The
\(\tau\)-factor of this quotient vanishes. Also
\(H^1(G,I_{L_n})=0\), since each prime orbit gives a permutation
module with first cohomology \(\operatorname{Hom}(G_w,\ZZ)=0\).
Hilbert's theorem 90 applied to \(U(L_n)\to L_n^\times\to
P_{L_n}\) gives
\[
 P_{L_n}^G/P_{L_a}\simeq H^1(G,U(L_n)),\qquad
 H^1(G,P_{L_n})\simeq
       \bigl(U(L_a)\cap N L_n^\times\bigr)/N U(L_n).
\]
The second isomorphism uses \(H^2(G,D)=D^G/ND\) for cyclic
\(G\), first for units and then for field elements. The
intersection on the right is all of \(U(L_a)\) by the preceding
paragraph. Now take invariants in \(P_{L_n}\to I_{L_n}\to\Cl(L_n)\).
The kernel of the map to invariant classes is supplied by
invariant principal ideals, and its cokernel is \(H^1(G,P_{L_n})\).
After taking the \(\tau\)-factor the invariant ideals come from
\(L_a\), as proved above. Dividing by \(P_{L_a}\) gives the
displayed exact sequence. These operations are exact on the
required terms by flatness of \(\ZZ_\ell\) over \(\ZZ\) and
by the integral idempotent. The actions of \(\Delta\) and \(G\)
commute throughout.
\end{proof}

\begin{theorem}[Capitulation and norm indices between every pair of levels]\leavevmode
\label{thm:unitlength505}
Suppose some half-residue weight has \(k(r)>1\). Use the
scalar \(z\) and exponent \(b\) from the proof of
Theorem~\ref{thm:collisionbranch464}, so \(v_\ell(z)=1\),
\(1\le b\le\infty\) and \(k(r)=\min(b,v_\ell(t_r-z))\).
Put \(s_a=\min(b,a+1)\). For \(0\le a<n\) the kernel of
\(j_{a,n}\) and the \(\tau\)-component of \(U(L_a)/N U(L_n)\)
are cyclic \(O\)-modules, each of length
\[
                s_a-\max(0,s_n-(n-a)).
\]
At level zero this length is zero for \(n<b\) and one for
\(n\ge b\). Thus finite \(b\) is exactly the first level at
which the nonzero faithful base class capitulates. It is also
the first level at which the norms of units fail to generate
the faithful base unit component.
\end{theorem}
\begin{proof}
The scalar description gives \(B=O/\ell^bO\), where \(b=\infty\)
means \(B=O\). The action of \(\gamma\) is \(y=1+z\).
By the binomial theorem \(v_\ell(y^{\ell^a}-1)=a+1\).
Lemma~\ref{lem:splitcyclolevel467} therefore gives
\(C_a=O/\ell^{s_a}O\), with the same scalar action.
The invariants \(C_n^G\) are the kernel of \(y^{\ell^a}-1\)
on \(C_n\), hence have length \(s_a\).

The ideal norm \(C_n\to C_a\) is surjective. Indeed the
unramified class field over \(L_a\) for \(C_a\) has trivial
intersection with \(L_n\) over \(L_a\), since every nontrivial
intermediate field of \(L_n/L_a\) ramifies at its unique
prime above \(\ell\). Its base change to \(L_n\) is unramified.
Class field reciprocity, compatible with ideal norm under this
restriction, gives the surjection on class groups and then on
their \(\tau\)-factors. Extension of the norm of an ideal is
the product of all its \(G\)-conjugates, so \(j_{a,n}C_a\)
is the image on \(C_n\) of
\[
    \sum_{i=0}^{\ell^{n-a}-1}y^{i\ell^a}
              =\frac{y^{\ell^n}-1}{y^{\ell^a}-1}.
\]
This scalar has valuation \(n-a\). Its image thus has length
\(\max(0,s_n-(n-a))\). Both source and invariant target of
\(j_{a,n}\) are cyclic of length \(s_a\). Their kernel and
cokernel are cyclic of the claimed length, and the lemma
identifies the cokernel with the unit norm quotient. For
\(a=0\) the source has length one. Its entire nonzero class
line survives exactly when \(s_n=n+1\), or \(n<b\).
\end{proof}

\begin{corollary}[A weight-uniform test using units modulo the first power]\leavevmode
\label{cor:unitcap505}
For \(n\ge1\), form the norm on ordinary units and reduce
modulo \(\ell\) to obtain
\[
    e_\tau\bigl(U(L_n)\otimes\mathbb F_\ell\bigr)
         \longrightarrow e_\tau\bigl(U(L)\otimes\mathbb F_\ell\bigr),
\]
where \(e_\tau\) is the orbit idempotent. The target is a line
over \(O/\ell O\). If this map is zero, then \(k(r)\le n\)
at \emph{every} indicated weight on the common interval.
If a weight has \(k(r)>n\), the unit norm map onto this line
is surjective at each level \(1,\ldots,n\). In particular a
zero map at \(n=25\) rules out the exceptional weight for the
entire pair \((\ell,j)\) on the tested orbit.
\end{corollary}
\begin{proof}
The roots of unity in the totally real \(L\) have order
at most two. Dirichlet's unit theorem identifies the real span
of its free units, by equivariant logarithms of all absolute
values, with the hyperplane of sum zero indexed by the real
embeddings. The character of this representation is the
regular character of \(\Delta\) minus the trivial character.
Its trace character comes from the integral unit lattice and
is the same after scalar extension to characteristic zero at
\(\ell\). The nontrivial orbit therefore contributes a free
\(O\)-module of rank one to \(U(L)\otimes\ZZ_\ell\). Reduction
gives the stated line. Norm commutes with its idempotent.

If any \(k(r)>n\), the hypothesis of the theorem holds and
\(b>n\). Its conclusion with \(a=0\) gives a surjective norm
on completed faithful units at every level through \(n\).
Reduction proves the assertion and its contrapositive.
If some weight has length greater than one, the first failure
occurs exactly at finite \(b\). Indeed for \(n\ge b\)
the image has index \(|O/\ell O|\) and equals \(\ell\) times
the base lattice. Its reduction is then zero. Computing the
displayed map requires only the images of a set of fundamental
units of \(L_n\) in the base unit lattice modulo \(\ell\).
The assertion at level 25 follows for all admissible weights
by this same contrapositive.
\end{proof}

For every pair at which the displayed map vanishes at a depth
at most \(b_0(x)\), its faithful orbit contributes zero after
truncation at that \(x\). At pairs where these norm maps stay
surjective the scalar congruence for \(\gamma\) in
Theorem~\ref{thm:scalarlevel467} still enters. The argument
does not bound the number of successive surjective levels
uniformly over the fields of Section~\ref{sec:missing}.


\subsection{A base unit test for all zero-residue weights}
\label{sec:baseunit506}

In Section~\ref{sec:locallysplit502} take \(s=0\) and
\(\alpha\ne1\). Write \(\tau=\alpha^{-1}\),
\(L=\overline{\QQ}^{\,\ker\tau}\), \(\Delta=\operatorname{Gal}(L/\QQ)\),
and \(O=\ZZ_\ell[\tau(\Delta)]\). The field \(L\) is totally
real and for fixed \(\alpha\) it is independent of \(\ell\).
Take \(\ell\ge5\) prime to \(n_*N_*\). First suppose \(\ell\)
splits completely in \(L\). For \(n\ge1\) write
\(L_n=L\QQ_n\), \(G_n=\operatorname{Gal}(L_n/L)\), and
let \(h^0_\tau=\operatorname{length}_O\Cl(L)_\tau\) in this
subsection. Here the subscript on a group denotes its completed
orbit component as in Section~\ref{sec:unitnorms505}.

For a global unit \(u\) write \(u_v\in\ZZ_\ell^\times\) for
its image at \(v\mid\ell\). In the decomposition
\(\ZZ_\ell^\times=\mu_{\ell-1}\times(1+\ell\ZZ_\ell)\) denote
the second coordinate by \(\langle u_v\rangle\). Consider
\[
 \mathcal L(u)=\left(
     \frac{\log_\ell\langle u_v\rangle}{\log_\ell(1+\ell)}
                    \right)_{v\mid\ell}
     \ \in\ \bigoplus_{v\mid\ell}\ZZ_\ell .
\]
The logarithm here is its convergent principal-unit power
series. The map is equivariant for the permutation of places.
The target is a copy of the regular module \(\ZZ_\ell[\Delta]\),
and its \(\tau\)-component is a free rank-one \(O\)-module.
Define \(c_\tau\) by requiring the image on the completed
\(\tau\)-component of the units to be \(\ell^{c_\tau}\) times
this target component. We allow \(c_\tau=\infty\) for zero image.
Changing the identification of the target line multiplies it
by a unit and preserves this definition.

\begin{theorem}[Invariant ideal classes at a completely split coefficient prime]\leavevmode
\label{thm:splitamb506}
At every level \(n\ge1\) the ordinary class groups satisfy
\[
 \operatorname{length}_O\bigl((\Cl(L_n)_\tau)^{G_n}\bigr)
                         =h^0_\tau+\min(n,c_\tau).
\]
If \(h^0_\tau=c_\tau=0\), then \(\Cl(L_n)_\tau=0\) at every
level and \(S(\alpha,r)=0\) at every positive weight divisible
by \(\ell-1\).
\end{theorem}
\begin{proof}
The primes of \(L\) above \(\ell\) form a free transitive
\(\Delta\)-set. Each is totally ramified of degree \(\ell^n\)
in \(L_n\), and all other finite primes are unramified. As in
Lemma~\ref{lem:ambiguous505}, extending ideals multiplies the
generator of each invariant prime orbit by its ramification
index. Thus \(I_{L_n}^{G_n}/I_L\) is a copy of
\((\ZZ/\ell^n)[\Delta]\) and its \(\tau\)-factor has length \(n\).
Write \(U'=U(L_n)_\tau\) and \(U=U(L)_\tau\), and denote
the norm on units by \(N\). The equivariant logarithmic
embedding in Dirichlet's unit theorem, now at all real places,
gives the regular character of \(\Delta\times G_n\) minus its
trivial character on the rational unit lattice. Because \(\tau\)
is nontrivial, \(U'\otimes_O\operatorname{Frac}(O)\) is the
regular representation of \(G_n\) over \(\operatorname{Frac}(O)\).
Also \(U\) has rank one over \(O\). There is no \(\ell\)-torsion
in these units since the fields are totally real.

We record the resulting equality of lengths
\[
     \operatorname{length}_O H^1(G_n,U')
                 =\operatorname{length}_O(U/NU').
\]
Indeed \(U'\) is commensurable equivariantly with \(O[G_n]\)
inside the common rational representation. For the regular
lattice the two terms vanish by the cyclic cohomology formulas.
For a finite module \(Q\), with generator \(g\) of \(G_n\),
the quotients \(\ker N/(g-1)Q\) and \(Q^{G_n}/NQ\) have equal
orders. This follows by counting the kernels and images of
\(N\) and \(g-1\) separately. Their length difference is
additive in exact sequences by the cyclic Tate cohomology
sequence. Apply this to the intersection of the two lattices
and their finite quotients. Finally \((U')^{G_n}=U\) by
flatness and the orbit idempotent.

The group \((I_{L_n}^{G_n}/P_L)_\tau\) has length \(h^0_\tau+n\).
Passing to the subgroup of invariant classes represented by
invariant ideals removes \((P_{L_n}^{G_n}/P_L)_\tau\), of
length \(\operatorname{length}_O H^1(G_n,U')\) by Hilbert 90.
The cokernel of these ideal classes in all invariant classes is
\(\big((U(L)\cap NL_n^\times)/NU(L_n)\big)_\tau\), by the
ideal cohomology argument of Lemma~\ref{lem:ambiguous505}.
Here we retain the indicated intersection because several
primes ramify. Define \(U^{\rm norm}=(U(L)\cap NL_n^\times)_\tau\).
Together with the preceding equality of lengths this gives
\[
 \operatorname{length}_O\bigl((\Cl(L_n)_\tau)^{G_n}\bigr)
             = h^0_\tau+n-\operatorname{length}_O(U/U^{\rm norm}).
\]

The Hasse norm theorem for cyclic extensions tests membership of
a unit in \(NL_n^\times\) at every place. Away from \(\ell\)
all units are local norms, including at the split real places.
At \(v\mid\ell\) the local field is \(\QQ_\ell\) and the extension
is the degree-\(\ell^n\) cyclotomic layer. Local reciprocity on
its unit group is projection of the principal-unit exponent
modulo \(\ell^n\), up to sign. This follows from the action of
a local unit on the \(\ell\)-power roots of unity by its
cyclotomic exponent, with the prime-to-\(\ell\) factor removed.
Hence the norm condition is \(\mathcal L(u)=0\bmod\ell^n\)
at all coordinates. Completion and the integral idempotent
preserve this finite-image calculation. The length of the image
of the completed \(\tau\)-map modulo \(\ell^n\) is
\(\max(0,n-c_\tau)\), interpreted as zero for infinite \(c_\tau\).
Substitution proves the formula.

If both invariants vanish, \((\Cl(L_n)_\tau)^{G_n}=0\).
A nonzero finite \(\ell\)-group with an action of a cyclic
\(\ell\)-group has nonzero fixed elements. On its nonzero
subgroup killed by \(\ell\), \(g-1\) is nilpotent and thus has
nonzero kernel. Hence the whole component is zero at every
level. In particular its level-one split quotient \(D_1\)
in Lemma~\ref{lem:primitiveclass502} is zero. Theorem~\ref{thm:splitweights502} gives vanishing at every indicated
positive weight.
\end{proof}

\begin{corollary}[One class number product and a congruence on base units]\leavevmode
\label{cor:zeromass506}
There is a positive integer \(D_{\rm cl}(V)\) depending only on \(E_*\)
with
\[
                   \log D_{\rm cl}(V)\ll_\psi n_*\log(2N_*).
\]
For \(\ell\ge5\) coprime to \(n_*N_*D_{\rm cl}(V)\), a nonzero full
component at a positive zero-residue weight on any \(E_x\) of
the interval forces \(\ell\) to split completely in its
nontrivial residual field \(L\). It also forces
\[
   e_\tau\left(\frac{u_v^{\,\ell-1}-1}{\ell}\bmod\ell
                                   \right)_{v\mid\ell}=0
                         \qquad\text{for every }u\in U(L).
\]
Across \emph{all} primes \(\ell\ge5\) coprime to \(n_*N_*\), the
number of weights \(1\le r<t_*\) exceptional in the sense of
Lemma~\ref{lem:classdet466} for any zero-residue character
nontrivial at its coefficient prime is
\(O_\psi(n_*\log(2N_*))\) on the entire interval.
\end{corollary}
\begin{proof}
For each cyclic subgroup \(C\) generated by an even character of
\(\operatorname{Gal}(E_*/\QQ)\), take the field \(L_C\) cut out
by a generator and put \(D_{\rm cl}(V)=\prod_C|\Cl(L_C)|\). Its conductor
divides \(N_*\), and \(m_C=[L_C:\QQ]=|C|\). The discriminant
and ideal count in Theorem~\ref{thm:nonsplitprefix466} give
\(\log|\Cl(L_C)|\le m_C\log(2N_*)\). The odd part of the order
of any character divides \(o=[K_\psi:\QQ]\), because the other
fields adjoined have degrees powers of two. Thus
\(m_C/\varphi(m_C)\le\prod_{p\mid2o}p/(p-1)\).
The generators of distinct cyclic subgroups are disjoint and
there are at most \(n_*\) characters in total. This proves the
claimed bound on \(\log D_{\rm cl}(V)\).

For \(\ell\nmid D_{\rm cl}(V)\) every base \(\ell\)-class component is
zero. If the faithful residual character is nontrivial on the
local decomposition group, Lemma~\ref{lem:classdet466} gives
vanishing at every weight. If it is trivial on that group,
the prime splits completely in \(L\). A nonzero scalar group
then requires \(c_\tau\ge1\) by the theorem. This condition
says exactly that every projected local logarithm is zero
modulo \(\ell\). The binomial expansion gives
\[
 \frac{u_v^{\ell-1}-1}{\ell}
     \equiv (\ell-1)\,
        \frac{\log_\ell\langle u_v\rangle}{\log_\ell(1+\ell)}
                                                     \pmod\ell .
\]
The factor \(\ell-1\) is invertible. The globally trivial
character has zero scalar group by Lemma~\ref{lem:primitiveclass502}.
Integral character descent in Theorem~\ref{thm:levelstrict345}
identifies the components on every \(E_x\) with their scalar
groups for the inflated characters.

For the last assertion, the bound of Lemma~\ref{lem:classdet466}
leaves at most the length of one split class group component
per pair of prime and orbit. At a fixed field and prime,
the sum of these lengths multiplied by the residue degrees is
at most \(v_\ell(|\Cl(L_C)|)\), since the components have
orthogonal integral idempotents. Multiplying by \(\log\ell\)
and summing over all primes and all \(C\) gives at most
\(\log D_{\rm cl}(V)\). Since all residue degrees are at least one and
\(\log\ell\ge1\), this also bounds the number of exceptional
weights in their union.
\end{proof}

\begin{proposition}[The real quadratic layer \(\QQ(\sqrt2)\)]\leavevmode
\label{prop:pellunit506}
For the quadratic character of \(\QQ(\sqrt2)\), a nonzero
zero-residue component at a good \(\ell\ge5\) requires both
\[
 \ell\equiv 1\ \text{or }-1\bmod 8,\qquad
       (1+\sqrt2)^{\ell-1}\equiv1
                           \pmod{\ell^2\ZZ[\sqrt2]} .
\]
In particular its scalar group at \(\ell=7\) vanishes for
every positive multiple of six. This gives the corresponding
vanishing on the period fields whenever \(7\nmid n_*N_*\).
\end{proposition}
\begin{proof}
The field has discriminant eight, and Minkowski's bound
\(\sqrt8/2<2\) makes its class group trivial. A prime \(\ell\)
splits there exactly in the two displayed residue classes,
by the quadratic supplement to reciprocity. At the inert
primes under consideration Lemma~\ref{lem:classdet466} applies directly. At a split
prime test the actual unit \(u=1+\sqrt2\). Its conjugate is
\(-u^{-1}\). Thus its two Fermat quotients modulo \(\ell\) are
opposites, and the nontrivial quadratic idempotent fixes their
pair. Vanishing of that pair is equivalent to the displayed
congruence because the two distinct primes over \(\ell\) are
coprime. The base class group is already trivial, so the proof of
Corollary~\ref{cor:zeromass506} applies independently of the
class numbers of other subfields. At seven one has
\((1+\sqrt2)^6=99+70\sqrt2\), whose \(\sqrt2\)-coefficient is
nonzero modulo 49. The scalar calculation uses only the field
\(\QQ(\sqrt2)\), whose degree and conductor are prime to seven.
This proves the asserted vanishing. The
quadratic field occurs in every sufficiently large real dyadic
layer of the interval family.
\end{proof}

The Fermat quotient condition is necessary for a surviving
split orbit with trivial base class component. The proof supplies
no upper bound for its strict length when that condition holds.
Such lengths, the other inertia residue, and the projected
period comparison still enter Section~\ref{sec:missing}.


\subsection{Units supported at the coefficient prime and the real quadratic layer}
\label{sec:splitunits507}

Keep the nontrivial zero-residue character \(\tau\), its field
\(L\), and a completely split prime \(\ell\) as in
Section~\ref{sec:baseunit506}. Put \(S=\{v:v\mid\ell\}\),
\(U_S=\mathcal O_L[1/\ell]^\times\) and
\(\Cl_S(L)=\Cl(L)/\langle[v]\mid v\in S\rangle\). For
\(u\in U_S\) define
\[
 \mathcal L_S(u)_v=
 \frac{\log_\ell\langle\ell^{-v_\ell(u_v)}u_v\rangle}
      {\log_\ell(1+\ell)}\qquad(v\in S).
\]
Here \(L_v=\QQ_\ell\) and the angular brackets discard the
Teichm\"uller component. Write \(\ell^{c_{\tau,S}}O\) for the
image of the completed \(S\)-units under this map on the
\(\tau\)-component of the local exponent lattice, identified
with \(O\) as in Section~\ref{sec:baseunit506}.

\begin{theorem}[The constant specialization computed in the residual field]\leavevmode
\label{thm:sunitbase507}
For the compact module \(B\) of Lemma~\ref{lem:primitiveclass502}
there is an exact sequence
\[
 0\longrightarrow O/\ell^{c_{\tau,S}}O
   \longrightarrow B/TB
   \longrightarrow \Cl_S(L)_\tau\longrightarrow0.
\]
In particular \(c_{\tau,S}\le c_\tau\), where \(c_\tau\) is
the ordinary unit index of Theorem~\ref{thm:splitamb506}.
These indices are finite. Write
\(b_S=c_{\tau,S}+\operatorname{length}_O\Cl_S(L)_\tau\).
If \(b_S=0\), all the scalar groups at positive zero-residue
weights vanish. In general the weights at which the full
length exceeds \(a b_S\) occupy at most \(b_S\) residue classes
modulo \((\ell-1)\ell^a\) for every \(a\ge1\).
\end{theorem}
\begin{proof}
Use the constant coefficient in the cyclotomic direction,
\((J/O)(\tau)\) with \(J=\operatorname{Frac}(O)\). Restriction
from \(\QQ\) to \(L_\infty=L\QQ_\infty\) identifies its global
first cohomology with the equivariant homomorphisms on
\(G_{L_\infty}\). The outer groups in degrees one and two vanish
by averaging over the prime-to-\(\ell\) group \(\Delta\) with
nontrivial character, exactly as in
Lemma~\ref{lem:primitiveclass502}. Requiring the resulting
homomorphism to factor through the unramified, locally split
quotient defining \(B\) gives the dual of \(B/TB\), since
\(\gamma\) acts trivially on the coefficient in this calculation.

Instead restrict first to \(L\). The same averaging identifies
the global classes with equivariant homomorphisms on \(G_L\).
At \(v\nmid\ell\) their inertia condition is unchanged on
passing to \(L_\infty\), which is unramified at \(v\).
At \(v\mid\ell\) their restrictions to \(G_{L_{\infty,w}}\)
must be zero. Equivalently the local homomorphism on \(G_{L_v}\)
factors through the local cyclotomic \(\ZZ_\ell\)-extension.
Thus \(B/TB\) is the \(\tau\)-component of the maximal abelian
pro-\(\ell\) quotient of \(G_L\) with these two local conditions.
The equivalence follows by dualizing the same group of
equivariant homomorphisms in the two descriptions.

By reciprocity, this quotient is computed by finite ideles.
Keep the ideal valuations away from \(S\). At \(v\in S\) keep
only the local principal-unit exponent. Indeed reciprocity for
the cyclotomic extension of \(\QQ_\ell\) kills \(\ell\) and the
Teichm\"uller units and uses precisely this exponent up to a
common sign. The ideal map outside \(S\) gives a surjection
onto \(\Cl_S(L)\) before the pro-\(\ell\) quotient. Its kernel
is the local exponent lattice \(\bigoplus_{v\in S}\ZZ_\ell\)
modulo \(\overline{\mathcal L_S(U_S)}\). To verify this, an
outside ideal trivial in \(\Cl_S(L)\) can be removed by a
principal idele. Two choices differ by an \(S\)-unit, and a
principal idele with zero outside valuations comes from
precisely an \(S\)-unit. The group of outside ideals is discrete,
so taking closure changes only the local image in each such
fiber. The kernel is compact pro-\(\ell\) and the remaining
ideal class group is finite. Consequently passing to the
pro-\(\ell\) quotient retains the kernel and takes the
\(\ell\)-primary part of the class group. The character
idempotent gives the claimed exact sequence. This also
identifies the nontrivial character component of the logarithmic
class group described in Section~\ref{nav:III-74}.

Ordinary units belong to \(U_S\), which proves the inequality
for the indices. Leopoldt's theorem for abelian number fields
gives full rational rank to the diagonal \(\ell\)-adic logarithm
on the ordinary unit lattice \cite{XI-VIII-brumer330}. The nontrivial
orbit has rank one over \(O\) in this lattice by Dirichlet's
unit theorem. Its image in the local \(O\)-line is therefore
nonzero, so both indices are finite. The exact sequence gives
\(\operatorname{length}_O(B/TB)=b_S\). The proof of
Lemma~\ref{lem:classdet466} from this point uses only a finite
presentation over \(O[[T]]\), this constant length, and the
specialization formula for \(S(\alpha,r)\). Finite generation
and that formula are available here by
Lemma~\ref{lem:primitiveclass502}, and \(O[[T]]\) is Noetherian.
The determinant interpolation in Lemma~\ref{lem:classdet466}
therefore gives the stated residue-class bound, also when the
character is locally trivial. For \(b_S=0\), compact Nakayama
gives \(B=0\).
\end{proof}

\begin{proposition}[A prime where ordinary units alone leave a nonzero invariant class]\leavevmode
\label{prop:thirtyone507}
For \(L=\QQ(\sqrt2)\), its nontrivial character and \(\ell=31\),
the ordinary class group at the first cyclotomic layer has
invariant \(\tau\)-component of length one. The compact module
\(B\) is zero and the full strict group vanishes at every positive
weight divisible by \(30\).
\end{proposition}
\begin{proof}
Write \(\eta=1+\sqrt2\). The units are generated by \(-1\)
and \(\eta\). Indeed for a positive real unit \(1<u<\eta\)
the conjugate is \(u'=\pm u^{-1}\). In both cases
\(0<(u-u')/(2\sqrt2)<1\), whereas this is the integral
coefficient of \(\sqrt2\). Thus \(\eta\) is the least unit
greater than one and its powers account for all positive units.
The class number is one as in Proposition~\ref{prop:pellunit506}.
At the prime over 31 with \(\sqrt2\equiv116\bmod961\) one has
\(116^2\equiv2\bmod961\) and \(117^{30}\equiv1\bmod961\).
The latter follows by squaring successively, with residues
\(235,448,816,844\) at exponents \(2,4,8,16\).
The conjugate of \(\eta\) is \(-\eta^{-1}\), so the same
congruence holds at the other prime. Thus \(c_\tau\ge1\), and
Theorem~\ref{thm:splitamb506} gives invariant length one at
level one.

The element \(\pi=7+3\sqrt2\) has norm 31, so it belongs to
\(U_S\). At the specified prime its value is \(355\bmod961\),
which is a unit. The residues of its powers \(2,4,8,16\)
are \(134,658,514,882\), giving \(355^{30}\equiv900
=1+29\cdot31\bmod961\). Its local principal-unit exponent
is therefore nonzero modulo 31. At the other prime the
local value divided by 31 is the inverse of this unit since
\(\pi\bar\pi=31\). The two exponents are negatives, so their
pair already belongs to the nontrivial quadratic component.
Its reduction generates the target line, giving \(c_{\tau,S}=0\).
The split class group is also zero. The preceding theorem
now proves the vanishing at every such weight.
\end{proof}

\begin{corollary}[All coefficient primes through \(10^7\) in the real quadratic layer]\leavevmode
\label{cor:quadraticrange507}
For the quadratic character of \(\QQ(\sqrt2)\), every prime
\(5\le\ell\le10^7\) and every positive multiple \(r\) of
\(\ell-1\) satisfy \(S(\alpha,r)=0\). The conclusion transfers
to its component on each period field of the interval for which
the degree and modulus are prime to \(\ell\).
\end{corollary}
\begin{proof}
For inert primes use the trivial base class group and
Lemma~\ref{lem:classdet466}. For the split primes first apply
Proposition~\ref{prop:pellunit506}. Of the \(332213\) primes
in this range congruent to \(1\) or \(7\) modulo eight, only
\(31\) and \(1546463\) satisfy \(\eta^{\ell-1}=1\bmod\ell^2\)
in \(\ZZ[\sqrt2]\). The following exact finite calculation is
included for reproducibility. The pair \((a,b)\) denotes
\(a+b\sqrt2\). Each loop over the binary digits squares the
previous power and multiplies it by \(\eta\) when the digit is
one. The sieve marks the proper multiples starting with the
prime square, which exhausts the composites by their least
prime factors.
\begin{verbatim}
from math import isqrt
bound = 10**7
prime = bytearray([1]) * (bound+1)
prime[:2] = bytes(2)
for d in range(2, isqrt(bound)+1):
    if prime[d]:
        prime[d*d::d] = bytes(1+(bound-d*d)//d)
survivors = []
count = 0
for p in range(5, bound+1):
    if not prime[p] or p % 8 not in (1, 7):
        continue
    count += 1
    m = p*p
    a, b = 1, 0
    for bit in bin(p-1)[2:]:
        a, b = (a*a+2*b*b) % m, (2*a*b) % m
        if bit == '1':
            a, b = (a+2*b) % m, (a+b) % m
    if (a,b) == (1,0):
        survivors.append(p)
assert (count, survivors) == (332213, [31, 1546463])
\end{verbatim}

For the two survivors choose \(\pi=A+B\sqrt2\) of norm \(\ell\).
At the prime where \(\pi\) is a unit its residue modulo \(\ell\)
is \(2A\). From \(\pi(2A-\pi)=\ell\) we obtain the unit value
\(y=2A-\ell(2A)^{-1}\bmod\ell^2\), with the inverse taken
modulo \(\ell\). The last column in the table is
\((y^{\ell-1}-1)/\ell\bmod\ell\).
\[
 \begin{array}{r|r r|r|r}
 \ell&A&B& y\bmod\ell^2& (y^{\ell-1}-1)/\ell\bmod\ell\\ \hline
 31&7&3&355&29\\
 1546463&1601&713&534027699288&823555
 \end{array}
\]
The norm equalities follow from \(A^2-2B^2=\ell\). For an
integer check of the final column, compute
\(\big(\operatorname{pow}(y,\ell-1,\ell^2)-1\big)/\ell\)
with modular exponentiation. Both residues are nonzero.
At the conjugate prime the unit after division by \(\ell\)
is \(y^{-1}\) modulo \(\ell^2\). The argument of
Proposition~\ref{prop:thirtyone507} gives \(c_{\tau,S}=0\)
in each case. This proves vanishing for every indicated weight
in the whole range of primes. Descent to \(E_x\) uses
Theorem~\ref{thm:levelstrict345} and Lemma~\ref{lem:locallevel345}.
\end{proof}

The same \(S\)-unit test is defined at every larger split prime.
An interval proof of the residual bound in
Section~\ref{sec:missing} also requires uniform control at
the unbounded primes, the other characters and the other
inertia residue.


\subsection{Two unit polynomials for the zero-residue strict groups}
\label{sec:unitpolys523}

We compute the number of generators of the entire complementary
strict group at weights divisible by \(p-1\), outside a fixed
integer of controlled size for the interval. For a cyclic character
the calculation also determines the full order whenever the
constant specialization has length less than \(1+v_p(r)\).
Keep \(E_x\subseteq E_*\), \(n_*=[E_*:\QQ]\) and the cyclotomic
modulus \(N_*\) of Theorem~\ref{thm:conductor232}.

Let \(L\ne\QQ\) be a real cyclic subfield of \(E_*\),
\(\Delta=\langle g\rangle=\Gal(L/\QQ)\), and \(d=|\Delta|\).
For each nontrivial subfield \(K\subseteq L\), including \(L\), let
\(C_K\) be its conductor. Form the unit
\[
  w_K=N_{\QQ(\zeta_{C_K})/K}(1-\zeta_{C_K}),\qquad
  \epsilon_K=g(w_K)/w_K .
\]
Here \(g\) restricts to a generator on every subfield. Let \(U_L^0\)
be generated by \(-1\) and all conjugates of these \(\epsilon_K\).

\begin{lemma}[Units of bounded index on the interval]
\leavevmode\label{lem:explicitunits523}
Write \(I_L=[\mathcal O_L^\times:U_L^0]\) and \(h_L=|\Cl(L)|\).
There is a positive integer
\[
   J_U=2n_*N_*\prod_{\substack{\QQ\ne L\subseteq E_*^+\\
                          L/\QQ\ {\rm cyclic}}} I_L h_L,
 \qquad \log J_U\ll n_*^2\log(2n_*N_*).
\]
\end{lemma}
\begin{proof}
An automorphism above \(g\) takes \(\zeta_C\) to \(\zeta_C^a\) with
\(a\) prime to \(C\). The ratio \((1-\zeta_C^a)/(1-\zeta_C)\) is an
algebraic integer by the geometric sum, and its inverse is integral
by the same argument using an inverse of \(a\) modulo \(C\). Its norm
to \(K\) gives the asserted unit. At each embedding the logarithm of
the absolute value of this unit has modulus at most
\(2N_*\log(2N_*)\). Indeed each conjugate of \(1-\zeta_C\) has
absolute value between \(2\sin(\pi/C)\) and two, and there are at
most \(N_*\) factors in each of the two norms.

For a nontrivial complex character of \(\Delta\), let \(K\) be the
field cut out by the character. Its primitive Dirichlet character
\(\xi\) has conductor \(C=C_K\) and is even. The \(\xi\)-coefficient
of the logarithm vector of \(w_K\) is nonzero. To see this, unfold
the norm to get \(\sum_{a\bmod C}\bar\xi(a)\log|1-\zeta_C^a|\),
summing over units. Replace \(\zeta_C^a\) temporarily by
\(z\zeta_C^a\) for \(0<z<1\) and expand the logarithm. The sums with
frequencies \(j\) and \(-j\) agree by evenness. The sum at frequency
\(j\) is \(\xi(j)G\), where \(G=\sum_a\bar\xi(a)\zeta_C^a\).
For invertible \(j\) change variables. Otherwise primitivity gives
a unit congruent to one modulo \(C/\gcd(C,j)\) on which \(\xi\) is
nontrivial, so the sum is zero. Additive character orthogonality,
summed over the frequencies, gives \(|G|^2=C\). Letting \(z\) increase
to one now gives \(-G L(1,\xi)\), by convergence of the Dirichlet
series at one and Abel's limit theorem. The value is nonzero by
Dirichlet's nonvanishing theorem. Passing to \(g(w_K)/w_K\)
multiplies by a nonzero difference between a character value and one.
Repeating coordinates from \(K\) in the embeddings of \(L\) preserves
nonvanishing in this component. The span of the logarithm vectors
of the conjugates is invariant under \(\Delta\). The complex regular
representation of a cyclic group contains each character once, and
the vectors consequently span the full hyperplane of coordinate sum
zero. In particular \(I_L\) is finite by the unit theorem.

Delete one real coordinate from that hyperplane. For a nontorsion
unit some \(|\log|\sigma u||\) is at least \((\log 2)/2\). Otherwise
every conjugate of the nonzero algebraic integer \(u^2-1\) would
have modulus strictly less than one. The deleted coordinate is the
negative sum of the others. Thus in the remaining \(d-1\) coordinates
the logarithm lattice has minimum sup norm at least
\(c/d\), where \(c=(\log2)/2\). An open cube of side \(c/d\)
injects up to its boundary into its real torus, giving covolume at
least \((c/d)^{d-1}\). Select \(d-1\) independent displayed unit
vectors. Hadamard's bound on their determinant is at most
\((2\sqrt d\,N_*\log(2N_*))^{d-1}\). Dividing the two bounds gives
\(\log I_L\ll d\log(2dN_*)\).

The conductor--discriminant formula gives \(|\operatorname{disc} L|
\le N_*^d\). Minkowski's bound supplies an integral ideal of norm at
most \(N_*^{d/2}\) in each ideal class. The number of ideals of norm
at most \(Y\) is bounded by \(Y^2\zeta_L(2)\le Y^2\zeta(2)^d\), by
the Euler factors. Hence \(\log h_L\ll d\log(2N_*)\). There are at
most \(n_*\) cyclic subfields, since each corresponds to a cyclic
subgroup of the character group with a generator. Summing proves
the lemma.
\end{proof}

Suppose \(p\ge5\), \(p\nmid J_U\), and \(p\) splits completely in
\(L\). Choose a place \(v_0\mid p\) and write \(v_j=g^jv_0\). Choose
\(\beta\) with \((\beta)=v_0^{h_L}\). For a unit \(z\) away from the
places above \(p\), put \(z_j^\circ=p^{-v_j(z)}z_{v_j}\in\ZZ_p^\times\),
where \(v_j(p)=1\) and \(L_{v_j}=\QQ_p\). Define
\[
 \begin{split}
 \mathcal L_z(X)&=\sum_{j=0}^{d-1}
    \frac{\log_p((z_j^\circ)^{p-1})}{(p-1)\log_p(1+p)}X^j,
       \\[-2pt]
 f_z(X)&=\sum_{j=0}^{d-1}
       \left(\frac{(z_j^\circ)^{p-1}-1}{p}\bmod p\right)X^j .
 \end{split}
\]
The first polynomial has coefficients in \(\ZZ_p\), the second in
\(\mathbb F_p\). Their reductions differ by the unit factor \(p-1\),
as follows from \(\log_p(1+pa)\equiv pa\bmod p^2\).

\begin{theorem}[The first digits and the full lengths at divisible weights]
\leavevmode\label{thm:unitroots523}
Let \(\alpha\) be any faithful character of \(\Delta\),
\(\tau=\alpha^{-1}\), \(O=\ZZ_p[\tau(g)]\), and \(r>0\) a multiple of
\(p-1\). Normalize \(v_p(p)=1\) in this unramified coefficient ring.
Put \(s=1+v_p(r)\) and
\[
 b_\tau=\min\bigl(v_p(\mathcal L_{\epsilon_L}(\tau(g))),
                    v_p(\mathcal L_{\beta}(\tau(g)))\bigr).
\]
A zero value is assigned infinite valuation. Then \(b_\tau\) is finite.
The finite strict group \(S(\alpha,r)\) is cyclic over \(O\), and
for every \(1\le m\le s\) its subgroup killed by \(p^m\) has length
\(\min(m,b_\tau)\). If \(b_\tau<s\), its full length is \(b_\tau\).
It is nonzero precisely when the reduction of \(\tau(g)\) is a
common root of \(f_{\epsilon_L}\) and \(f_\beta\).
\end{theorem}
\begin{proof}
Identify the exponent lattice at the places above \(p\) with
\(\ZZ_p[X]/(X^d-1)\), by sending the place \(v_j\) to \(X^j\).
The action of \(g\) is multiplication by \(X\). Evaluation at
\(\tau(g)\) is the projection to its orbit factor \(O\), since
\(p\nmid d\). Under this identification the local principal-unit
logarithm map of Theorem~\ref{thm:sunitbase507} sends \(z\) to
\(\mathcal L_z\). The completed units are generated by the conjugates
of the \(\epsilon_K\) since \(p\nmid I_L\). If \(K\ne L\), the
logarithm vector is fixed by every element of \(\Gal(L/K)\), and
the faithful character component is zero. Conjugates of
\(\epsilon_L\) differ there by unit scalar multiplication.

The valuations at the \(v_j\) of the conjugates of \(\beta\) are the
vectors \(h_L e_j\). Since \(h_L\) is a \(p\)-adic unit, adjoining
these elements to the ordinary units generates the completed
\(S\)-units, where \(S=\{v_j\}\). Indeed their valuations contain
\(h_L\ZZ^d\), and the kernel of the valuation map consists of
ordinary units. Thus the image in \(O\) is generated by the two
values displayed in the theorem. The \(p\)-part of \(\Cl_S(L)\)
vanishes. Theorem~\ref{thm:sunitbase507} identifies \(B/TB\) with
\(O/p^{b_\tau}O\), with finite index, for the compact module \(B\)
of Lemma~\ref{lem:primitiveclass502}.

Write \(A_r=B/(T-t_r)B\), where \(t_r=(1+p)^{-r}-1\). This is the
dual of the strict group by that lemma. The odd-prime binomial
valuation gives \(v_p(t_r)=s\). Consequently
\[
  A_r/p^m A_r=B/(p^m,T)B\simeq O/p^{\min(m,b_\tau)}O
                      \qquad(1\le m\le s).
\]
The equality of quotients follows from equality of the two ideals
after reduction modulo \(p^m\). In particular the finite \(O\)-module
\(A_r\) has at most one generator, by reduction modulo \(p\). Its
dual is cyclic with the same length. If \(b_\tau<s\), a cyclic
module with a longer elementary divisor would have larger length
even after reduction modulo \(p^s\). Taking \(m=1\) and reducing the
two polynomial values proves the common-root assertion.
\end{proof}

\begin{corollary}[The number of generators on every common period field]
\leavevmode\label{cor:intervalroots523}
At every \(p\ge5\) outside \(J_U\), throughout the interval,
the dimension over \(\mathbb F_p\) of the subgroup killed by \(p\)
in the complementary sign of \(\mathcal S_{p,r}(E_x)\) equals
\[
 \sum_{\substack{\QQ\ne L\subseteq E_x^+,\ L/\QQ\ {\rm cyclic}\\
                            p\ {\rm splits\ completely\ in}\ L}}
   \deg\gcd\bigl(\overline\Phi_{[L:\QQ]}(X),
                      f_{\epsilon_L}(X), f_{\beta}(X)\bigr)
       \qquad (r>0,\quad p-1\mid r).
\]
The polynomials in each summand use its own \(L,g,v_0,\beta\).
In particular the same integer counts the generators at all such
weights. If it is zero the entire complementary group vanishes.
\end{corollary}
\begin{proof}
The weights are even, so the complementary characters factor through
\(E_x^+\). The trivial one has zero strict group by
Lemma~\ref{lem:primitiveclass502}. Each other character is faithful
on the cyclic quotient cut out by its kernel. When \(p\) does not
split completely in that field, the residual character is nontrivial
on its decomposition group. Lemma~\ref{lem:collisionclass457} would
force \(p\mid h_L\) if the group survived. At a split prime the
primitive \(d\)-th roots in characteristic \(p\) are distinct,
since \(p\nmid d\). The theorem gives one generator over the residue
field for each surviving Frobenius orbit. Its degree over
\(\mathbb F_p\) is the orbit size, exactly the degree of the
corresponding irreducible factor of the displayed gcd.

Finally the finite group of \(E_x/\QQ\) has order prime to \(p\).
Its integral orbit idempotents decompose the period-field group into
these scalar groups by the coefficient descent in the proof of
Theorem~\ref{thm:levelstrict345}. At the primes away from \(p\) the
additional conductor localizations cost zero by
Lemma~\ref{lem:locallevel345}, since every inertia character here
has order prime to \(p\). Summing the residue dimensions proves the
formula for the original group.
\end{proof}

\begin{proposition}[The cubic field of conductor seven at thirteen]
\leavevmode\label{prop:cubicunit523}
For each faithful character of \(L=\QQ(\zeta_7)^+\) the scalar
strict group at \(p=13\) vanishes at every positive multiple of
\(12\). This gives the same zero components on each period field
containing \(L\) with \(13\nmid n_*N_*\).
\end{proposition}
\begin{proof}
Put \(t=\zeta_7+\zeta_7^{-1}\), with minimal polynomial
\(t^3+t^2-2t-1\). Its discriminant is \(49\), so Minkowski's bound
is at most \(6\sqrt{49}/27<2\) and \(h_L=1\). Thirteen is \(-1\)
modulo seven and splits in the real subfield. With \(g(t)=t^2-2\),
one has \(w_L=2-t\) and \(\epsilon_L=2+t\). At the three local
embeddings in the order \(v_0,gv_0,g^2v_0\), the values of \(t\)
modulo \(169\) can be taken as \(34,163,140\). Indeed \(34\) solves
the cubic modulo \(169\) with simple reduction and \(g^{-j}(t)\)
gives the other two entries in this order. Raising the corresponding
values of \(2+t\) to the twelfth power modulo \(169\), subtracting
one and dividing by thirteen gives \((9,6,11)\bmod13\). The
polynomial \(9+6X+11X^2\) takes values \(9,5\) at \(3,9\), the two
primitive cube roots modulo thirteen. Thus this single ordinary unit
already generates each local exponent line over its coefficient
ring. Together with the trivial class group the exact sequence of
Theorem~\ref{thm:sunitbase507} gives \(B/TB=0\) for both characters,
so the specialized groups vanish. This calculation uses no index
assumption on the other fields of the interval.
\end{proof}

When the common root occurs and \(b_\tau\ge1+v_p(r)\), the formula
computes the indicated digits and leaves the possible excess length
at that weight. The primes dividing \(J_U\), the other weight
residues, and the comparison of signed projected periods still enter
Section~\ref{sec:missing}.


\subsection{Vanishing at zero-residue weights in the first three dyadic layers}
\label{sec:dyadicunits524}

Let \(C_j=\QQ(\zeta_{2^{j+2}})^+\). These are the actual real dyadic
subfields of the period construction. We examine the whole strict
group \(\mathcal S_{p,r}(C_j)\) at weights \(p-1\mid r\). The rational
and quadratic characters were treated in Lemma~\ref{lem:primitiveclass502}
and Corollary~\ref{cor:quadraticrange507}. We compute the remaining
faithful contributions for the layers of degrees four and eight.

\begin{lemma}[A single unit in the faithful part of a dyadic layer]
\leavevmode\label{lem:dyadicpolynomial524}
Let \(d=2^j\ge4\), \(q=4d\), \(n=2d\), \(L=C_j\), and suppose that
\(p\ge5\) is prime with \(p\equiv\pm1\pmod q\). For the calculation
use \(Z=\zeta_q\), \(t=Z+Z^{-1}\), and the following polynomial of
degree less than \(n\), with integer coefficients taken modulo \(p^2\),
\[
     A(Z)=(1+Z-Z^{n-1})^{2(p-1)}-1
                  \pmod{Z^n+1,\ p^2}.
\]
Each coefficient is divisible by \(p\). Write \(a_h\in\mathbb F_p\)
for its quotient by \(p\) modulo \(p\). For \(0\le i<d/2\) take the
representative \(e_i\) of \(3^i\bmod q\) in \([0,q)\), and put
\[
       H_{p,d}(X)=\sum_{i=0}^{d/2-1}
             (-1)^{\lfloor e_i/n\rfloor}a_{e_i\bmod n} X^i .
\]
If \(p\nmid h_L\) and \(\gcd(H_{p,d},X^{d/2}+1)=1\) in
\(\mathbb F_p[X]\), every faithful character of \(L/\QQ\) has zero
scalar strict group at every positive multiple of \(p-1\).
\end{lemma}
\begin{proof}
The group of the real field is cyclic of order \(d\), generated by
\(g:Z\mapsto Z^3\) modulo complex conjugation. Indeed successive
squaring gives \(v_2(3^{2^i}-1)=i+2\) for \(i\ge1\), and no power
of \(3\) is \(-1\bmod8\). Its element \(g^{d/2}\) takes \(t\) to
\(-t\), since \(3^{d/2}\equiv1+2d\bmod4d\). Here
\(w_L=2-t\) and \(\epsilon_L=(2-t^3+3t)/(2-t)=(1+t)^2\).
The ring \(\ZZ[Z]/(Z^n+1)\) realizes the integral cyclotomic
coordinates. The prime splits in \(L\), so the unit
\(\epsilon_L^{p-1}-1\) is in \(p\mathcal O_L\). It is therefore in
\(p\ZZ[\zeta_q]\). Since \(Z^{-1}=-Z^{n-1}\), reduction gives exactly
the polynomial in the statement and permits the coefficient division.

In the split ring \(\mathcal O_L/p\simeq\prod_{v\mid p}\mathbb F_p\),
the quotient by \(p\) just computed is the table of Fermat quotients
of \(\epsilon_L\). The inclusion into \(\ZZ[\zeta_q]/p\) is injective:
an element of \(L\) integral after division by \(p\) in the larger
field is already integral in \(L\). Write
\(b_i=Z^{e_i}+Z^{-e_i}\) for \(0\le i<d/2\). These lie in the
\((-1)\)-eigenspace of \(g^{d/2}\). Their pairs of exponents modulo
\(n\), up to sign, are disjoint because \(3\) generates the real
cyclotomic quotient and its order after factoring by \(g^{d/2}\)
is \(d/2\). Their coefficients on those distinct exponent pairs are
units. They form a basis of that eigenspace over \(\mathbb F_p\):
in the split ring the involution exchanges the places in pairs, so
the eigenspace has dimension \(d/2\). It is the direct sum of the
faithful character components, whose eigenvalues for \(g\) are
exactly the roots of \(X^{d/2}+1\).

Project to this eigenspace by \((1-g^{d/2})/2\). This retains the
odd exponents of the polynomial, since the involution acts there by
\(Z\mapsto-Z\). The coefficient of \(b_i\) is the coefficient at
\(Z^{e_i\bmod n}\) with the sign shown in the statement. Moreover
\(g b_i=b_{i+1}\), with \(g b_{d/2-1}=-b_0\). In this basis the
projection of the unit table is \(H_{p,d}\) in the module
\(\mathbb F_p[X]/(X^{d/2}+1)\). Its conjugates span that module if
the two polynomials are relatively prime, by B\'ezout's identity.
Consequently each faithful local logarithm line is generated modulo
\(p\) by ordinary units, hence integrally by Nakayama's lemma. With
\(p\nmid h_L\), Theorem~\ref{thm:sunitbase507} gives zero constant
specialization at each such character. Lemma~\ref{lem:primitiveclass502}
and compact Nakayama give the asserted vanishing.
\end{proof}

\begin{theorem}[Full strict vanishing in degrees four and eight]
\leavevmode\label{thm:octicrange524}
At every prime \(5\le p\le10^7\) and every positive multiple \(r\)
of \(p-1\),
\[
           \mathcal S_{p,r}(C_j)=0\qquad(0\le j\le3).
\]
For \(p\not\equiv\pm1\pmod{64}\) in this same range, the equality
holds at every \(j\ge0\).
\end{theorem}
\begin{proof}
Fukuda and Komatsu prove that no prime below \(1.2\cdot10^8\)
divides the class number of any \(C_j\)
\cite[Corollary 1.1]{XI-VIII-weberFK524}. At a nonsplit prime the residual
faithful character of \(C_j\) is nontrivial on the decomposition
group, so Lemma~\ref{lem:collisionclass457} gives zero. At the split
primes for \(d=4,8\), the following finite calculation gives the
number of primes tested and the number having a gcd of positive
degree in Lemma~\ref{lem:dyadicpolynomial524}, respectively:
\[
\begin{array}{c|rr}
 d&\text{split primes}&\text{positive gcd degree}\\ \hline
 4&166087&0\\
 8&82993&0
\end{array}
\]
Here are exact arithmetic instructions to reproduce the table. All
coefficient lists start with the constant term. For each indicated
prime reduce after each multiplication. In the quotient by
\(Z^n+1\), form an ordinary convolution \(c\) with an extra zero
at index \(2n-1\), and return \((c_i-c_{i+n})\bmod p^2\). Thus all
intermediate coefficients may be computed as integers. The following
pseudocode uses polynomial remainders over \(\mathbb F_p\) in the
last step.
\begin{verbatim}
prime[0..10000000] = true
for a from 2 through floor(sqrt(10000000)):
    if prime[a]: mark a*a, a*a+a, ... <= 10000000 false
for d in [4, 8]:
    count = 0; exceptional = 0; n = 2*d
    for p from 5 through 10000000:
        if not prime[p] or p mod (4*d) not in [1, 4*d-1]: skip p
        count = count+1
        u = 1+Z-Z^(n-1); b = 1; e = 2*(p-1)
        while e > 0:
            if e is odd: b = product(b,u)
            u = product(u,u); e = floor(e/2)
        b = b-1
        check that p divides each coefficient of b
        h = 1; H = 0
        for i from 0 through d/2-1:
            H = H + (-1)^floor(h/n)*(b[h mod n]/p)*X^i
            h = (3*h) mod (4*d)
        F = X^(d/2)+1; G = H reduced modulo p
        while G is nonzero: (F,G) = (G,remainder(F,G))
        if degree(F) > 0: exceptional = exceptional+1
    record (count, exceptional)
\end{verbatim}
The sieve removes exactly the composites, each of which has a prime
factor at most its square root. Repeated squaring preserves the
product of the accumulator and the remaining power of the base.
The loop over \(i\) reads the faithful basis in the order of the
lemma. The Euclidean algorithm terminates with a scalar multiple of
the gcd. This justifies the finite computation independently of any
choice of roots in a residue field.

Each character of \(C_3\) is either trivial or faithful on exactly
one of \(C_1,C_2,C_3\). The rational component vanishes by
Lemma~\ref{lem:primitiveclass502}. The quadratic component vanishes
by Corollary~\ref{cor:quadraticrange507}. The lemma and the table
account for the other two layers. All these fields are totally real
and of degree prime to \(p\). Character descent, or the
prime-to-\(p\) idempotents in Theorem~\ref{thm:levelstrict345}, gives
vanishing of the entire strict group. For the final assertion
Theorem~\ref{thm:tower138} identifies the group at any layer with
one of the four groups just treated.
\end{proof}

\begin{corollary}[The corresponding summand on the period fields]
\leavevmode\label{cor:realperiod524}
At \(x\) let \(C_j\subseteq E_x\) be any layer covered by
Theorem~\ref{thm:octicrange524}, and assume \(p\nmid[E_x:C_j]\).
For each of its weights the field average onto \(C_j\) has zero
image on \(\mathcal S_{p,r}(E_x)\).
\end{corollary}
\begin{proof}
The integral average has image the invariants under
\(\Gal(E_x/C_j)\). Restriction identifies these with the lower
strict group: global cohomology descends by averaging, and restriction
on the local conditions is injective by corestriction at the primes
over \(p\) and on inertia elsewhere. The lower group is zero by the
theorem.
\end{proof}

This removes the indicated real dyadic character components at all
their matching weights on the common intervals. Characters involving
the other factors of \(E_x\) and the other twist residues still
occur in the local sum of Section~\ref{sec:missing}.


\subsection{The two possible characters in the full real dyadic tower}
\label{sec:dyadictower525}

For an odd prime \(p\) put \(q=2^{\max(v_2(p-1),v_2(p+1))}\) and
\(d=q/4\). Thus
\(L=\QQ(\zeta_q)^+\) has degree \(d\). It is the highest real dyadic
layer in which \(p\) splits completely: the splitting condition at
conductor \(2^h\) is \(p\equiv\pm1\bmod 2^h\). For the primes
through \(10^7\) the class numbers at all levels are prime to \(p\)
by Fukuda and Komatsu \cite[Corollary 1.1]{XI-VIII-weberFK524}. Consequently
at a positive multiple of \(p-1\) every faithful character from a
level above \(L\) has zero strict group by Lemma~\ref{lem:collisionclass457}
for \(p\ge5\). We evaluate the logarithms on all levels up to \(L\)
simultaneously.

\begin{lemma}[A single transform for all the split layers]
\leavevmode\label{lem:splittower525}
Suppose \(d\ge4\). Choose a primitive \(q\)-th root \(z\) in the
unramified extension of \(\ZZ_p\) of residue degree one if
\(p\equiv1\bmod q\), two otherwise. Write \(g(\zeta_q)=\zeta_q^3\)
on the real field. Set \(u_i=2-z^{3^i}-z^{-3^i}\) for
\(0\le i<d\). These are in \(\ZZ_p^\times\). Put
\[
 c_i=\frac{u_i^{p-1}-1}{p}\bmod p,\qquad
 A_k=\sum_{i=0}^{d-1} c_i\omega^{ki},\qquad
          \omega=(z\bmod p)^4,\quad 0\le k<d .
\]
Let \(\tau\) be the lifted scalar character with reduction
\(\bar\tau(g)=\omega^{-k}\) and order \(e=d/\gcd(k,d)\ge4\).
On the unique subfield of degree \(e\) its faithful ordinary unit
\(\epsilon\) of Section~\ref{sec:unitpolys523} generates the local
logarithm component if \(A_k\ne0\).
\end{lemma}
\begin{proof}
The Frobenius of the unramified coefficient extension sends the
chosen Teichm\"uller root to its \(p\)-th power, which is itself or
its inverse. The traces \(z^{3^i}+z^{-3^i}\) are therefore rational
\(p\)-adic integers. Modulo \(p\), each \(u_i\) is nonzero since a
primitive root of order \(q\) raised to an odd exponent is different
from one. These values are the conjugates of
\(w_L=(1-\zeta_q)(1-\zeta_q^{-1})\) at a fixed place. The indexing
is by \(g^i\) acting on the element, hence by \(g^{-i}\) acting on
the place. This gives precisely the inverse in \(\bar\tau(g)\).

For a smaller nontrivial layer \(K\) the norm of \(w_L\) to \(K\)
is \(w_K\) with compatible choices of primitive roots. At a single
step this follows in the full cyclotomic field from
\((1-Z)(1+Z)=1-Z^2\); then take the norm to the real subfield.
The logarithm turns a norm into the sum over the fibers. The same
holds for Fermat quotients modulo \(p\) since
\((1+pa)(1+pb)\equiv1+p(a+b)\bmod p^2\). Thus the weighted sum for
an inflated character is its logarithm at the field of that
character, up to the common unit normalizing the logarithm. The
unit \(g(w_K)/w_K\) supplies the further factor
\(\omega^{-k}-1\), which is a unit for the indicated orders. The
value \(A_k\ne0\) consequently gives a generator modulo \(p\), hence
a generator of the rank-one local exponent lattice.
\end{proof}

\begin{theorem}[Full dyadic strict groups outside two coefficient primes]
\leavevmode\label{thm:dyadicall525}
Let \(3\le p\le10^7\) be prime, and let \(r>0\) be divisible by
\(p-1\). At every \(j\ge0\) one has \(\mathcal S_{p,r}(C_j)=0\)
unless \(p=18433\) or \(p=114689\). At the former prime only a
single character of degree-\(256\) quotient can contribute, and at
the latter only a single character of degree-\(1024\) quotient can
contribute. In the table, \(k\) follows the root convention specified
in the proof, \(e\) is the character order, and the last column is
\(\bar\tau(g)\) in \(\mathbb F_p\).
\[
 \begin{array}{r|rrrr}
 p& d& k& e&\bar\tau(g)\\ \hline
 18433  &512 &426 &256 &13924\\
 114689 &4096&1420&1024&34058
 \end{array}
\]
For each such prime the possible group is cyclic over \(\ZZ_p\) and
is obtained by restriction from the layer of degree \(e\) once
\(2^j\ge e\). If \(p\mid r\) its order is at most \(p\).
\end{theorem}
\begin{proof}
At three the vanishing follows from Theorem~\ref{thm:tower138}.
For \(p\ge5\) the rational and quadratic components are zero by
Theorem~\ref{thm:octicrange524}. It remains to test the characters
in the lemma. We give the complete finite calculation. All powers,
including those in a quadratic ring, may be taken by binary
exponentiation. For \(p\equiv1\bmod4\), take the smallest positive
quadratic nonresidue \(h\) and put
\[
       z_0=h^{(p-1)/q}\bmod p,\qquad z=z_0^p\bmod p^2 .
\]
For \(p\equiv3\bmod4\), work in \((\ZZ/p^m)[i]\), \(i^2=-1\), with
\(m=1,2\) as specified. Take the least positive integer \(h\) for
which \(h^2+1\) is a quadratic nonresidue. Such an \(h\) exists:
the equation \(v^2-h^2=1\) has \(p-1\) solutions by factoring the
left side, so \(\sum_h (\frac{h^2+1}{p})=-1\). The term at zero is
one. Take \(z_0=(h+i)^{(p^2-1)/q}\bmod p\), lift its two
coefficients to \([0,p)\), and take \(z=z_0^p\bmod p^2\) in that
ring. The norm of \(h+i\) is a nonsquare, hence \(h+i\) is a
nonsquare in the cyclic residue multiplicative group. Raising a
nonsquare to the group order divided by \(q\) gives exact order
\(q\) because \(q\) is a power of two. This also proves the order
assertion in degree one. The subsequent \(p\)-th power kills the
possible principal-unit error of the lifted root modulo \(p^2\)
and preserves exact order. In degree two it selects the conjugate
residue root, which defines the allowed choice of embedding.

For each prime in the sieve of Theorem~\ref{thm:octicrange524} the
following operations suffice. The symbol \(F(c,w)\) denotes the
list of the values of the polynomial with coefficient list \(c\) on
the consecutive powers of \(w\), starting with exponent zero.
\begin{verbatim}
if d < 4: skip this prime
s = z + z^(q-1) modulo p^2
for j from 0 through d-1:
    c[j] = (least_residue((2-s)^(p-1),p^2)-1)/p
    s = (s^3-3*s) modulo p^2
w = (z modulo p)^4
values = F(c,w)
for k from 1 through d-1, omitting d/2:
    if values[k] == 0: record (p,d,k,d/gcd(d,k))
\end{verbatim}
The initial \(s\) has zero imaginary part, by the order and Frobenius
properties proved above. The recurrence is the identity
\(x^3+x^{-3}=(x+x^{-1})^3-3(x+x^{-1})\). To implement \(F\), return
the input for length one. Otherwise set \(b=F(c_{\rm even},w^2)\)
and \(a=F(c_{\rm odd},w^2)\), with each subsequence in its original
order. For input length \(D\) the entries at \(i,i+D/2\) are
\(b_i+w^i a_i,b_i-w^i a_i\) for \(0\le i<D/2\). This follows by
separating the even and odd exponents and using \(w^{D/2}=-1\).
These instructions thus evaluate the sums of the lemma using only
integer arithmetic and pairs of residues.

There are \(166087\) primes with \(d\ge4\) in this range. In total
\(5211578\) character values of order at least four are tested. The
zero values give exactly the two rows displayed in the theorem. For
additional counts, the numbers tested at orders \(2^j\), for
\(j=2,\ldots,18\) successively, are
\[
 \begin{gathered}
 332174,331972,331752,331840,330432,331904,333056,331008,333312,\\
 331776,327680,311296,335872,327680,262144,196608,131072 .
 \end{gathered}
\]
Here the counts can also be checked directly by adding
\(\varphi(2^j)\) for each prime with \(2^j\mid d\).
For every nonzero test value, the class-number theorem and
Theorem~\ref{thm:sunitbase507} now give zero constant specialization
and hence vanishing at every indicated weight. The same two pairs
of prime and layer degree occur in Gras's computations of the
ordinary \(p\)-adic regulator \cite{XI-VIII-grasdyadic525}. The strict local
condition here also allows units supported at \(p\).

To bound the two remaining components we retain one more digit. Both
primes are \(1\bmod q\), so this calculation takes place in the
integers. For each displayed row lift \(z_0\) by \(z=z_0^{p^2}\bmod
p^3\), set \(w=z^{4k}\bmod p^2\), and run the same trace recurrence
modulo \(p^3\). For each \(u_i\) let
\(b_i=((u_i^{p-1}\bmod p^3)-1)/p\), taking the least residue first.
The resulting values are
\[
\begin{array}{r|rrrr}
 p&h& z_0 & w\bmod p &
 p^{-1}\sum_i w^i(b_i-p b_i^2/2)\bmod p\\ \hline
 18433&5&17660&12072&921\\
 114689&3&2187&2401&99260
\end{array}
\]
To reproduce the last column, accumulate the sum modulo \(p^2\),
using the inverse of two there, and then divide its least residue
by \(p\). The logarithm series gives
\(p^{-1}\log(1+p b_i)\equiv b_i-p b_i^2/2\bmod p^2\) since \(p\ge5\).
Changing to the normalization in Theorem~\ref{thm:unitroots523}
multiplies the entire sum by a unit. The norm to the smaller layer
and the factor from \(g(w_L)/w_L\) in the lemma preserve this
valuation. Thus the indicated ordinary unit has valuation exactly
one in each of the two logarithm lines. The additional \(p\)-units
can only enlarge its image.

The class group term is zero by Fukuda and Komatsu. Hence the
constant specialization \(B/TB\) of Theorem~\ref{thm:sunitbase507}
is either zero or \(\ZZ/p\) in the two possible characters. Their
coefficient rings are \(\ZZ_p\) because \(e\mid p-1\). Reduction of
\(B/(T-t_r)B\) modulo \(p\) gives at most one generator. For \(p\mid r\)
we may reduce modulo \(p^2\) and replace \(T-t_r\) by \(T\), as in
Theorem~\ref{thm:unitroots523}; the resulting quotient has order at
most \(p\), forcing the full cyclic group to have order at most \(p\).
The character decomposition over the totally real fields is integral
since the degrees are powers of two. Restriction on these towers of
degree prime to \(p\) identifies the individual characters with their
inflations, including the strict local conditions, as in
Corollary~\ref{cor:realperiod524}. This proves the theorem.
\end{proof}

\begin{corollary}[A single arithmetic progression for longer groups]
\leavevmode\label{cor:dyadicdepth525}
At each of the two primes either all indicated lengths are at most
one, or there exists \(z_p\in\ZZ_p^\times\) such that for every
positive \(r\) divisible by \(p-1\) and every \(j\),
\[
       \log_p|\mathcal S_{p,r}(C_j)|\le 1+v_p(r-z_p).
\]
Here \(v_p(0)=+\infty\). In particular lengths greater than \(a\)
are confined to at most one residue class modulo \((p-1)p^a\)
for each integer \(a\ge1\).
\end{corollary}
\begin{proof}
It suffices to consider the one possible field and character.
If \(B/TB=0\), compact Nakayama gives \(B=0\). Otherwise a single
generator gives \(B=\ZZ_p[[T]]/I\), since its reduction modulo
\((p,T)\) has dimension one. Constant terms of elements of \(I\)
form exactly \(p\ZZ_p\) by the constant specialization. Choose
\(F=p+T c(T)\in I\). For \(t\in p\ZZ_p\), each specialization has
length at most \(v_p(F(t))\). If \(c(0)\) is divisible by \(p\), this
valuation is one for all \(t\). Otherwise \(F\) has a unique root
\(u\in p\ZZ_p\): beginning at zero modulo \(p\), each lift of a
root modulo \(p^a\) changes the next digit with unit coefficient
\(c(0)\bmod p\). Completeness gives the root. For \(t\in p\ZZ_p\)
the integral divided difference \((F(t)-F(u))/(t-u)\) is a unit,
by the same reduction, whenever \(t\ne u\). Also \(u=-p/c(u)\)
has valuation one. Thus \(z_p=-\log(1+u)/\log(1+p)\) is a unit,
and \(v_p(t_r-u)=1+v_p(r-z_p)\) by the odd-prime logarithm. This
proves the inequality. Length greater than \(a\) forces
\(r\equiv z_p\bmod p^a\); combining with the weight residue proves
the final assertion.
\end{proof}

\begin{corollary}[The real dyadic image in the residual sum on the common interval]
\leavevmode\label{cor:actualdyadic525}
Include \(18433,114689\) in the fixed finite set \(\Pi\) of
Section~\ref{sec:missing} when selecting the common intervals. Let
\(C_j\subset E_x\) be a real dyadic layer. For every odd prime
\(p\le10^7\) and every weight in the residual sum divisible by
\(p-1\), the image of \(\mathcal S_{p,r}(C_j)\) in \(R_{p,r}(E_x)\)
is zero on every sufficiently distant such interval. The fields and
terminal primes can be those of Corollary~\ref{cor:sameprimes517}.
\end{corollary}
\begin{proof}
Only the two displayed primes require consideration. Every incoming
class is restricted from \(C_8\) or \(C_{10}\), whenever its supplying
layer is present, by the theorem. For \(r\le x^\eta\) both fixed
layers lie in the real dyadic part of \(L_x\) for large \(V\). For
\(r>x^\eta\) both lie in \(B(x)\), and the prime belongs to \(\Pi\).
In either case the field average \(J_{p,r}\) fixes the restricted
class, since it comes from the fixed field of that average. The
factor \(1-J_{p,r}\) in the definition of \(R_{p,r}\) therefore kills
it already before truncation. This argument concerns the image of
restriction even if a larger extension has degree divisible by \(p\).
The common field and prime selection allows an arbitrary fixed
finite set \(\Pi\) containing the required odd divisors, so the same
choices apply to the signed period calculation.
\end{proof}

The other character components of \(E_x\), other weights and the
larger coefficient primes must still be included in the full local
estimate. The two possible nonzero strict groups in the real dyadic
tower were removed here by the actual field averages.


\subsection{Relative regulators at the half-residue weights}
\label{sec:relativereg529}

Let \(C_j=\mathbb Q(\zeta_{2^{j+2}})^+\). We treat the residual
character at the second of the two weights in
Corollary~\ref{cor:smallorder427}. Suppose \(j\ge1\), put
\(2n=2^j\), and let \(\alpha\) be a faithful character of
\(C_j/\mathbb Q\). Its primitive conductor is \(c=8n\).
Let \(p\equiv1\pmod4\) be prime and write \(\eta_p\) for the
quadratic character of conductor \(p\). For
\(r\equiv(p-1)/2\pmod{p-1}\) the residual character of
Section~\ref{sec:collisionconductors457} is the reduction of
\(\lambda=\alpha^{-1}\eta_p\). Write \(L\) for the field fixed
by the kernel of \(\lambda\) and \(F=C_{j-1}\), with \(C_0=\mathbb Q\).
The two characters have coprime conductors, so \(L\) is real
cyclic of degree \(2n\), \(F\) is cut out by \(\lambda^2\), and
\(L/F\) is quadratic. We compare their ordinary class numbers
\(h_L,h_F\).

\begin{lemma}[The relative unit lattice at a ramified prime]
\leavevmode\label{lem:relregulator529}
For these fields, with the usual regulators and \(R_{\mathbb Q}=1\),
\[
       \frac{R_L}{R_F}
          \ \ge\ \frac{[n(\tfrac12\log p-\log 2)]^n}{n!}.
\]
Also extension of ideals injects \(\operatorname{Cl}(F)\) into
\(\operatorname{Cl}(L)\).
\end{lemma}
\begin{proof}
Pair the real embeddings of \(L\) over those of \(F\), choosing one
member \(\tau_i\) of each pair. Let \(s\) generate
\(\operatorname{Gal}(L/F)\). The kernel of the norm on units modulo
signs consists of units \(u\) with \(u\,s(u)=\pm1\), modulo
\(\{1,-1\}\). Their log vectors
\(a_i=\log|\tau_i(u)|\) form a lattice \(\Gamma\) of rank \(n\)
in \(\mathbb R^n\). Indeed the full logs are \((a,-a)\), and the
image of the norm on units contains the squares of units of \(F\).
Dirichlet's unit theorem gives the rank, discreteness and full span.

For a nonzero vector of this lattice one has \(u\notin F\), since
a unit of \(F\) with square \(\pm1\) is \(\pm1\). The relative
discriminant divides the discriminant \((u-s(u))^2\) of the
order with basis \(1,u\), as seen by passage to local integral
bases. At \(p\) the field \(F\) is unramified and the inertia
in \(L\) has order two, from the factor \(\eta_p\). Thus every
prime of \(F\) above \(p\) ramifies tamely in the quadratic
extension and occurs in the relative discriminant. Their product
has norm \(p^n\). Consequently
\[
 p^{n/2}\le\prod_{i=1}^n |\tau_i(u-s(u))|
                  \le 2^n\exp\!\left(\sum_i |a_i|\right).
\]
Every nonzero vector of \(\Gamma\) has \(\ell^1\)-norm at least
\(b=n(\frac12\log p-\log 2)>0\). The cross-polytope of this
radius has volume \(2^n b^n/n!\). Minkowski's convex body
theorem, first for smaller radii and then in the limit, gives
\(\det\Gamma\ge b^n/n!\).

Complete a basis of the relative units modulo signs by lifts of
a basis of the image of the norm. In the regulator matrix omit
the last conjugate row. Add the first \(n-1\) rows to the
corresponding conjugate rows. The bottom rows now give the logs
of norms, zero on the first \(n\) columns. The two diagonal
determinants in absolute value are \(\det\Gamma\) and \(R_F\)
times the index of the norm image in the units of \(F\) modulo
signs. This proves the regulator assertion, including \(n=1\).

The class number of \(F\) is odd. For completeness, if it were even,
the Hilbert class field for the ordinary ideal classes would give a
nontrivial maximal elementary abelian unramified two-extension
\(D/F\) split at the real places. It is Galois over \(\mathbb Q\)
by invariance of these conditions; its Galois group is a two-group.
The only possible finite ramification over \(\mathbb Q\) is at two,
which is totally ramified in \(F\). Inertia in \(D/\mathbb Q\)
there has order \([F:\mathbb Q]\), hence is proper. A maximal
subgroup containing it is normal of index two. Its fixed field
would be quadratic and unramified at every finite prime, contrary
to the quadratic discriminant formula. Finally the norm after
extension of ideal classes is squaring. Its kernel on the odd
class group of \(F\) is trivial.
\end{proof}

\begin{lemma}[The joint mean of the faithful Dirichlet values]
\leavevmode\label{lem:jointLvalues529}
Uniformly for \(p\to\infty\) and
\[
                   n\le\frac{(\log p)^2}{\log\log(3p)},
\]
one has
\[
 \left(\frac1n
    \sum_{\substack{a\bmod 2n\\a\ {\rm odd}}}
    |L(1,\alpha^a\eta_p)|^2\right)^{1/2}
    \le (1+o(1))\frac{\log p}{4\sqrt n}.
\]
\end{lemma}
\begin{proof}
All the indicated characters are primitive of conductor \(D=cp\).
Put \(y=\lceil\sqrt D\,\log^2(2D)\rceil\). For any primitive
nonprincipal character of conductor \(D\), its sum over an interval
is \(O(\sqrt D\log(2D))\), with an absolute constant. Indeed its
additive Fourier coefficients vanish at nonunits and have modulus
\(\sqrt D\) at units by the primitive Gauss sum identity. Expand the
indicator of a residue interval using its geometric sums, bounded
at frequency \(1\le h<D\) by \(D/(2\min(h,D-h))\), and divide by
\(D\). The zero frequency for the character vanishes, and full
periods may first be omitted. Partial summation therefore permits
truncating the Dirichlet series at \(y\) with error
\(O(1/\log(2D))\) per character.

Let \(G_0=(\mathbb Z/c\mathbb Z)^\times/\{1,-1\}\), of order
\(2n\). Form \(s_g=\sum_{m\le y,\ [m]=g}\eta_p(m)/m\) on
this group. For the representative \(1\le a<c/2\) of a class
\(g=\{a,c-a\}\), bounding the two harmonic sums separately gives
\[
 |s_g|\le \frac{2\log y}{c}+e_g,\qquad
 e_g=\frac1a+\frac1{c-a}+\frac2c,\qquad
                         \sum_g e_g^2\ll 1 .
\]
For each progression, terms after its first are bounded by the
harmonic sum up to \(y\) divided by \(c\). The squared error estimate
follows by summing reciprocal squares and the constant terms.
The truncated character values are precisely the Fourier sums of
the same \(s_g\) at the odd powers of \(\alpha\).
Orthogonality on \(G_0\), enlarged to all characters before applying
Parseval, gives their root mean square over the \(n\) odd powers at
most
\[
    \sqrt 2\,\|s\|_2\le
       \sqrt 2\left(\sqrt{2n}\frac{2\log y}{c}+O(1)\right)
                  =\frac{\log y}{2\sqrt n}+O(1).
\]
By the bound on \(n\),
\(\log y=\frac12\log p+O(\log\log(3p))\) and
\(\sqrt n/\log p\to0\) uniformly. The truncation error and the
absolute \(O(1)\) term are absorbed, which gives the assertion.
\end{proof}

\begin{theorem}[Vanishing in a growing range of dyadic layers]
\leavevmode\label{thm:halfdyadic529}
There is an absolute \(p_0\) such that for \(p\ge p_0\),
\(p\equiv1\pmod4\), and \(n\) in the range of
Lemma~\ref{lem:jointLvalues529},
\[
 \frac{h_L}{h_F p^{n/2}}
        \le \frac{n!}{n^n}\left(\frac1{\sqrt2}+o(1)\right)^n
        <1 .
\]
The error in the parenthesis is uniform. If
\(\operatorname{ord}_{2n}(p)\ge n/2\), then
\(S(\alpha,r)=0\) for every positive
\(r\equiv(p-1)/2\pmod{p-1}\), with the strict local conditions
of Section~\ref{sec:collisionconductors457}. In particular this
proves vanishing at every such weight on all the indicated layers
when \(p\equiv5\pmod8\), and on the quadratic and quartic layers
for both prime classes \(1,5\pmod8\).
\end{theorem}
\begin{proof}
The characters of \(L\) outside those of \(F\) are exactly the \(n\)
odd powers of \(\lambda\), whose conductors are all \(cp\).
The conductor--discriminant formula gives
\(|\Delta_L|/|\Delta_F|=(cp)^n\).
Divide the analytic class number formulas for these two totally real
fields, both with two roots of unity \cite[Theorem~19.12]{XI-VIII-sutherland529}. Factor their Dedekind zeta
functions into primitive Dirichlet functions. The common factors
cancel in the quotient of the residues at one, giving
\[
 \frac{h_L}{h_F}=\frac{(cp)^{n/2}}{2^n}\,
            \frac{R_F}{R_L}
        \prod_{\substack{a\bmod 2n\\a\ {\rm odd}}}
                 |L(1,\alpha^a\eta_p)|.
\]
The product of the absolute values equals the positive residue
ratio since nonreal characters are paired with their conjugates.
By the arithmetic--geometric mean the product is at most the
\(n\)-th power of the root mean square in the preceding lemma.
Insert its bound and the regulator bound. After division by
\(p^{n/2}\) the result is
\[
 \frac{n!}{n^n}
 \left(
    \frac{\sqrt{8n}}{2}\,
    \frac{(1+o(1))\log p}{4\sqrt n\,(\frac12\log p-\log2)}
 \right)^n.
\]
The parenthesis tends uniformly to \(1/\sqrt2\). Since
\(n!\le n^n\) the strict inequality follows.

Suppose the strict group were nonzero. Its first torsion layer
gives a nonzero equivariant homomorphism from \(\operatorname{Cl}(L)\)
to the additive residual character module, by the proof of
Lemma~\ref{lem:collisionclass457}. Its local hypothesis holds because
inertia at \(p\) acts by \(\eta_p\). The image has size at least
\(p^{\operatorname{ord}_{2n}(p)}\), since it is stable under a
primitive \(2n\)-th root after reduction. Classes extended from
\(F\) are fixed by \(s\). The character of \(s\) on the target is
\(-1\), and \(p\) is odd, so the homomorphism kills those classes.
By Lemma~\ref{lem:relregulator529} the available quotient has
order \(h_L/h_F\). This contradicts the class number bound whenever
the stated residue degree condition holds.

For \(p\equiv5\pmod8\) one has
\(v_2(p^{2^a}-1)=a+2\) for \(a\ge0\), by successive squaring.
Hence the order modulo \(2^j\) is \(2^{\max(j-2,0)}\), at least
\(n/2\). For \(j=1,2\) every \(p\equiv1\pmod4\) has order one
modulo \(2^j\), which also suffices.
\end{proof}

\begin{corollary}[The half-residue summands on the period fields]
\leavevmode\label{cor:periodhalf529}
On the fields and intervals of Section~\ref{sec:missing}, the
faithful components from the layers and primes in
Theorem~\ref{thm:halfdyadic529}, with \(C_j\subset E_x\) and
\(p\nmid [E_x:\mathbb Q]\), contribute zero to \(R_{p,r}\) for
every indicated \(1\le r<t(x)\). One may include the bounded
coefficient primes in the fixed set \(\Pi\) so that the inflated
contributions from the same degree range at these smaller primes
are removed by the field averages as well.
\end{corollary}
\begin{proof}
The integral character comparison of Theorem~\ref{thm:levelstrict345}
identifies the component over \(E_x\) with the scalar group subject
also to its conductor conditions. This is a subgroup of the strict
group just shown to vanish. Increase the absolute cutoff by the
fixed prime divisors of \([K_\psi:\mathbb Q]\) if necessary to
cover all possible odd divisors of the field degrees. For the
finitely many smaller primes the displayed degree range involves
only finitely many \(C_j\). All these fields lie in both \(L_x\)
and \(B(x)\) for sufficiently large \(x\). Include those primes
in \(\Pi\). The applicable field average then fixes every class
inflated from these subfields and \(1-J_{p,r}\) kills it.
This argument only uses invariance for the small primes, even when
the whole character decomposition over \(E_x\) fails to be integral.
The choice of \(\Pi\) is fixed before the intervals and is
compatible with the period primes in
Corollary~\ref{cor:whittakerterminal527}.
\end{proof}

\begin{corollary}[Every weight on a common growing subfield]
\leavevmode\label{cor:allweightlayer529}
For a sufficiently distant period interval \([V,V^{3/2}]\), let
\(D_V=C_j\) for the largest \(j\ge1\) with
\[
       2^{j-1}\le \frac{(\log V)^2}{2\log\log(3V)} .
\]
Its degree is of order \((\log V)^2/\log\log(3V)\).
For every \(x\) in the interval, \(V\le p\le P(x)\) with
\(p\equiv5\pmod8\), and every \(1\le r<t(x)\), the field
average onto \(D_V\) acts as zero on the actual residual quotient
\(R_{p,r}\) over \(E_x\). For \(p>P(x)\) in the same progression
the sum of the log orders coming from \(D_V\), over all such primes
and weights, is \(O_{\psi,\varepsilon}(V^C)\) for a fixed \(C\).
\end{corollary}
\begin{proof}
The field \(D_V\) lies in every \(E_x\) on the interval because
the available real dyadic degrees grow at least as a fixed positive
power of \(V\), by Theorem~\ref{thm:conductor232}. All the indicated
primes are prime to \([E_x:\mathbb Q]\) for large \(V\). The
field average commutes with the sign and field averages defining
\(R_{p,r}\). Its image on the full strict group is restriction
from \(D_V\) by descent through a prime-to-\(p\) extension.
We bound this latter group before forming the quotient.

At a multiple of \(p-1\), Theorem~\ref{thm:tower138} gives zero at
every real dyadic level when \(p\equiv5\pmod8\). At the half residue
the twist is even. The rational component vanishes by
Proposition~\ref{prop:qbound102}. The other components are faithful
on some dyadic subfield of \(D_V\). Each satisfies the size condition
of Theorem~\ref{thm:halfdyadic529}, because
\((\log z)^2/\log\log(3z)\) is increasing for sufficiently large
\(z\) and \(p\ge V\). They all vanish by that theorem and integral
character descent.
At every remaining weight, \(p-1\nmid2r\).
Lemma~\ref{lem:fullinteger517} bounds each full log order by
\(\log J(V)\ll V^C\), uniformly. This bounds its exponent in
powers of \(p\) by \(V^{O(1)}< b_0(x)\), so its image is killed
in the formation of \(R_{p,r}\) for \(p\le P(x)\).
Finally \(P(x)>2\lfloor V^{1203/2}\rfloor+1\) for large \(V\).
The product divisibility in the same lemma bounds the total log
order at the larger primes by \(\log J(V)\). This proves both claims.
\end{proof}

\begin{corollary}[The full real dyadic tower at primes three modulo eight]
\leavevmode\label{cor:fullreal529}
Write \(C_{J(x)}\) for the real dyadic field used to construct
\(E_x\), and \(e_D\) for the field average onto \(D\).
For \(V\le p\le P(x)\) with \(p\equiv3\pmod8\),
\(e_{C_{J(x)}}R_{p,r}=0\) at every \(1\le r<t(x)\).
More generally, for \(p\equiv3\pmod4\) and
\(p\not\equiv-1\pmod{64}\) in this range,
\[
       e_{C_{J(x)}}R_{p,r}=e_{C_{\min(3,J(x))}}R_{p,r}.
\]
The total log order from these real dyadic subfields at \(p>P(x)\)
is polynomially bounded in \(V\).
\end{corollary}
\begin{proof}
At separated weights use Lemma~\ref{lem:fullinteger517} and the
truncation argument of Corollary~\ref{cor:allweightlayer529}.
At the half residue, \((p-1)/2\) is odd. Complex conjugation
fixes the real subfield, hence acts as the identity on the image
of its field average. The complementary sign projector \((1-c)/2\)
kills that image. At the zero residue
Theorem~\ref{thm:tower138} gives zero for \(p\equiv3\pmod8\).
For the other indicated primes it gives descent of the entire
dyadic strict group to a layer of degree at most eight.
Consequently the field averages onto \(C_{J(x)}\) and
\(C_{\min(3,J(x))}\) have the same image on this group.
They are nested commuting idempotents, so this equality passes
to the sign, field and truncation quotients. The large-prime
bound follows on \(E_*\) from Lemma~\ref{lem:fullinteger517}
and descends to every \(C_{J(x)}\) by prime-to-\(p\) restriction.
\end{proof}

The estimate here includes twists prime to \(p\) as well as multiples
of \(p\). Its range grows with the coefficient prime. The highest
dyadic layers of the period fields, all the other residue characters
outside the earlier bounds, and the full marked period comparison
still enter Section~\ref{sec:missing}.


\subsection{Relative class numbers in complex dyadic layers}
\label{sec:complexdyadic542}

Write \(C_j=\mathbb Q(\zeta_{2^{j+2}})^+\). Let \(j\ge2\),
\(2n=2^j\), and let \(\beta\) be a faithful character of
\(C_j/\mathbb Q\). Its primitive conductor is \(c=8n\).
Let \(\xi\) be a quadratic character or the trivial character,
and let \(a\) be the odd part of its primitive conductor.
Take an odd prime \(p\nmid a\). Write \(\eta_p\) for the
quadratic character of conductor \(p\), and put
\[
             \lambda=\beta^{-1}\xi\eta_p,\qquad D=acp .
\]
Every odd power of \(\lambda\) is primitive of conductor \(D\).
Indeed an odd power of \(\beta\) still has conductor \(c\).
The dyadic conductor of a quadratic character divides eight,
and \(c\ge16\), so twisting leaves this conductor exponent
unchanged. The characters at the odd factors of \(a\) and
at \(p\) are nontrivial quadratic characters on separate factors.

\begin{lemma}[A geometric mean at one for all the odd powers]
\leavevmode\label{lem:geometricmean542}
With an absolute implied constant,
\[
  \left(\prod_{\substack{b\bmod 2n\\b\ {\rm odd}}}
         |L(1,\lambda^b)|\right)^{1/n}
                    \ll \bigl(\log(3D)\bigr)^{1/n}.
\]
\end{lemma}
\begin{proof}
First consider any primitive nonprincipal Dirichlet character
\(\chi\) of conductor \(D\). For real \(\sigma\ge1\), the
absolute value of its completed function \(\Lambda(\sigma,\chi)\)
is nondecreasing. Here is the zero argument. The entire
function of order at most one
\[
   \Lambda(1/2+w,\chi)\Lambda(1/2+w,\bar\chi)
\]
is even by the functional equations, including their conjugate
root numbers. Its zeros belong to \(|\operatorname{Re}w|\le1/2\)
by the Euler product and the functional equations. Pair each
nonzero zero \(z\) with \(-z\), counting multiplicities. By
Hadamard factorization the paired factors are \(1-w^2/z^2\)
and the central factor is \(w^{2h}\) for some \(h\ge0\).
The remaining factor is a constant because the entire function
is even and its exponential prefactor has degree at most one.
The paired product converges on compact sets, since the sum
of \(|z|^{-2}\) converges away from zero. At real \(w>1/2\)
each of \(1/(w-z)\) and \(1/(w+z)\) has positive real part.
Their sum is the logarithmic derivative of the paired factor.
The central factor contributes \(2h/w\). This proves the
monotonicity since the product of the two completed functions
is the positive number \(|\Lambda(1/2+w,\chi)|^2\) there.
Continuity gives the endpoint. Taking
\(\sigma=1+1/\log(3D)\), the conductor power and ratio of
gamma factors between \(1\) and \(\sigma\) are bounded by an
absolute constant for either parity. Hence
\[
                   |L(1,\chi)|\ll |L(\sigma,\chi)|.
\]

At this \(\sigma>1\) take logarithms of the absolutely
convergent Euler products. Prime powers of exponent at least
two contribute \(O(1)\) to the average of their real logarithms,
uniformly in all data. At a prime \(q\nmid D\) the average of
\(\lambda(q)^b\) over odd \(b\) is zero unless
\(\beta(q)^2=1\), by the finite geometric series. In that case
its real part is at most one. Enlarging the prime set to
\(q\nmid c\) in this nonnegative upper bound and using
orthogonality gives
\[
 \begin{split}
 \frac1n\sum_{b\ {\rm odd}}\log |L(\sigma,\lambda^b)|
   &\le \sum_{\substack{q\nmid c\\\beta(q)^2=1}}q^{-\sigma}+O(1)\\
   &=\frac1n\sum_{j'=0}^{n-1}
               \log |L_c(\sigma,\beta^{2j'})|+O(1).
 \end{split}
\]
The sum over \(q\) is over primes. The subscript \(c\) means
that every character is read modulo \(c\), so all its Euler
factors at two are omitted. The term for \(j'=0\) is bounded
by \(\log\zeta(\sigma)\le\log(1+1/(\sigma-1))\).

For the other \(j'\), a partial sum of the character has
absolute value at most \(c\) by periodicity and zero mean.
Partial summation bounds the tail after \(c\) of its Dirichlet
series by an absolute constant for \(1<\sigma\le2\).
The mean square of the first \(c\) terms, averaged over all
\(n\) characters, is bounded absolutely. In fact
orthogonality groups \(1\le m\le c\), \(m\) odd, by their
value under \(\beta^2\). Every fiber has
\(\varphi(c)/n=4\) members. Cauchy--Schwarz on each fiber
bounds the mean square by \(4\sum_{m\ge1}m^{-2}\).
It follows that the mean of \(\log^+|L_c(\sigma,\beta^{2j'})|\)
over the nonprincipal terms, divided by \(n\), is \(O(1)\).
Exponentiate the resulting bound and use the transfer from
\(\sigma\) to \(1\). This proves the assertion.
\end{proof}

\begin{theorem}[The half residue after a fixed quadratic twist]
\leavevmode\label{thm:complexhalf542}
Suppose \(\xi(-1)=\eta_p(-1)\), let \(L\) be the field cut out
by \(\lambda\), and put \(F=C_{j-1}\).
There is an absolute constant \(C\) for which
\[
 \left(\frac{h_L}{h_F p^{n/2}}\right)^{1/n}
 \le \frac{(n!)^{1/n}}{n}\,
      \frac{C\sqrt{a n}\,(\log(3acp))^{1/n}}
                         {\log p-2\log2}
             \qquad(p>4).
\]
For each fixed \(\xi\) there is \(b_\xi>0\) such that for all
sufficiently large \(p\) of this parity and every
\(2\le n\le b_\xi(\log p)^2\), the right hand side is less
than one. If in addition
\(\operatorname{ord}_{2n}(p)\ge n/2\), the strict group
\(S(\beta\xi,r)\) vanishes at every positive
\(r\equiv(p-1)/2\bmod(p-1)\).
\end{theorem}
\begin{proof}
The character \(\lambda\) is even of order \(2n\).
Its square cuts out \(F\), so \(L/F\) is quadratic and both
fields are totally real. At every prime over \(p\) this
extension is ramified of index two, whereas \(p\) is
unramified in \(F\). The relative regulator estimate of
Lemma~\ref{lem:relregulator529} applies also to this extension.
To verify the dependence on the quadratic twist, for a
nontrivial relative unit modulo signs let \(s\) be the
involution and take one embedding above each embedding of \(F\).
The ideal of the relative discriminant divides the ideal of
\((u-su)^2\), since \(u\notin F\) generates \(L/F\).
The product of the norms of primes of \(F\) above \(p\) is
\(p^n\). Consequently
\[
        p^{n/2}\le\prod_\tau|\widetilde\tau(u-su)|
              \le 2^n\exp\left(\sum_\tau
                                  |\log|\widetilde\tau u||\right).
\]
Here the relative unit is in the kernel of the norm modulo
signs, so that \(|\widetilde\tau su|=|\widetilde\tau u|^{-1}\).
The lattice of these \(n\) log coordinates has full rank
because norms of units contain the squares of units of \(F\).
The cross-polytope of radius \(n(\tfrac12\log p-\log2)\)
has no nonzero lattice point in its interior. Minkowski's
theorem gives covolume at least the \(n\)-th power of this
radius divided by \(n!\). The regulator determinant, with
one row omitted, gives \(R_L/R_F\) equal to this covolume
times the index of the norm image modulo signs, as computed in
Lemma~\ref{lem:relregulator529}. This proves the required
regulator inequality without any condition on \(a\).

Exactly the \(n\) odd powers of \(\lambda\) occur in the
quotient of the zeta functions for \(L\) and \(F\). Each
has conductor \(acp\). The conductor--discriminant formula
and the analytic class number formulas therefore give
\[
     \frac{h_L}{h_F}
       =\frac{(acp)^{n/2}}{2^n}\frac{R_F}{R_L}
                   \prod_{b\ {\rm odd}}|L(1,\lambda^b)|.
\]
Insert the two estimates and \(c=8n\), absorbing the fixed
numerical factor into \(C\). For fixed \(a\), take first
\(2\le n\le\log\log(3p)\). Then
\(\sqrt n(\log(3acp))^{1/n}/\log p\)
is \(O_a(\sqrt{\log\log(3p)/\log p})\).
For \(\log\log(3p)<n\le b_\xi(\log p)^2\), we may take
\(b_\xi\le1\), and the factor \((\log(3acp))^{1/n}\)
is at most four eventually. Choosing \(b_\xi\) small enough
in terms of \(a,C\) gives a strict upper bound less than one.

If \(S(\beta\xi,r)\ne0\), the argument of
Lemma~\ref{lem:collisionclass457} supplies an equivariant
homomorphism out of the ordinary class group of \(L\) with
image of order at least \(p^{\operatorname{ord}_{2n}(p)}\).
Inertia at \(p\) acts nontrivially on the residual character,
as required there. The involution of \(L/F\) acts as minus
one on this image, so the homomorphism kills classes extended
from \(F\). The class number of the real dyadic field \(F\)
is odd by Lemma~\ref{lem:relregulator529}, and extension of
its ideal classes is injective because norm after extension
is squaring. Thus the available quotient has order \(h_L/h_F\).
In the indicated range it has order less than \(p^{n/2}\),
contradicting the lower bound on the image.
\end{proof}

\begin{corollary}[Every weight on a growing complex subfield]
\leavevmode\label{cor:complexallweights542}
Let \(K_\xi\) be any fixed imaginary quadratic field contained
in \(K_\psi C_\infty\), with associated character \(\xi\).
There exists \(b>0\) such that, on each sufficiently distant
original interval, one can take the largest \(j\ge1\) with
\(2^{j-1}\le b(\log V)^2\) and put \(D_V=K_\xi C_j\).
Then \(D_V\) is contained in every actual \(E_x\) of the
interval, and its degree is of order \((\log V)^2\). At
every odd prime \(V\le p\le P(x)\) with \(p\equiv3,5\bmod8\)
and every \(1\le r<t(x)\),
\[
                 R_{p,r}(E_x)^{\operatorname{Gal}(E_x/D_V)}=0 .
\]
At \(p>P(x)\) the total logarithmic order of the complementary
groups on this subfield, summed over those primes and all these
weights, is \(O_{\psi,\varepsilon}(V^{A})\) for some \(A\).
If \(K_\xi\) is omitted, the same assertions hold for the
real subfields \(C_j\) on every original family.
\end{corollary}
\begin{proof}
Take \(b\) at most the two constants in the theorem for
\(\xi\) and the trivial character. The fixed quadratic field
is contained in a fixed base layer. The available real dyadic
degree in \(E_x\) grows at least as a fixed positive power
of \(V\), by Theorem~\ref{thm:conductor232}. This gives the
inclusions, and successive real dyadic degrees double. Since
the quadratic field is imaginary it has intersection
\(\mathbb Q\) with the real tower.

For \(p\ge V\) the odd prime \(p\) is prime to the degrees
of the actual fields for large \(V\). At a separated weight
\(p-1\nmid2r\), Lemma~\ref{lem:fullinteger517} bounds the
strict length on \(D_V\) by \(\log J(V)/\log p<b_0(x)\).
At \(p-1\mid r\) the complementary sign is plus. Its component
descends to \(C_j\) by the order-two projection and vanishes by
Theorem~\ref{thm:tower138} in both progressions.

At the half residue the sign is odd for \(p\equiv3\bmod8\)
and even for \(p\equiv5\bmod8\). The corresponding characters
are respectively \(\xi\beta'\) and \(\beta'\), with \(\beta'\)
a character of \(C_j\). If \(\beta'\) is faithful on a layer
of order \(2n'\ge4\), the theorem applies with
\(n'\le b(\log p)^2\). Indeed \(v_2(p^2-1)=3\) in both
progressions, and each successive squaring from \(p^2\)
increases the valuation of its difference from one by one.
The order modulo \(2n'\) is at least \(n'/2\), including
modulus four by inspection. If \(\beta'\) has order at most
two, the finitely many resulting nontrivial quadratic characters
vanish here by Proposition~\ref{prop:quarticconductor539}
at all sufficiently large \(p\). The trivial character at
an even half-residue weight is covered by
Proposition~\ref{prop:qbound102}. These scalar vanishings give
vanishing of the corresponding sign summand on \(D_V\) by
prime-to-\(p\) character descent. For the real field alone
the minus component is already zero by conjugation.

The field average onto \(D_V\) commutes with the sign, the
other field averages and the high-power truncation in
Proposition~\ref{prop:res23}. On the strict groups its image
is restriction from \(D_V\), by prime-to-\(p\) descent.
Thus in every case the indicated residual invariant group is
zero. For \(p>P(x)\), all weights in question are separated
for sufficiently large \(V\), since \(P(x)>2t(x)+1\).
The local divisibilities of Lemma~\ref{lem:fullinteger517}
on \(D_V\subset E_*\) bound the total logarithmic order by
\(\log J(V)\ll V^A\).
\end{proof}

This argument allows arbitrary fixed odd conductor in the
quadratic twist. On the original fields the two moving
ramified quadratic characters and primes in the other
residue classes still occur in the complementary estimate.


\subsection{Fixed characters at the limiting residue degree}
\label{sec:fixedcharacters543}

For a Dirichlet character \(\gamma\), let \(K_\gamma\) be the
fixed field of its kernel in the abelian closure of \(\QQ\).
We always take the primitive \(L\)-function of a character,
except when the modulus of an imprimitive Euler product is
indicated explicitly. We use the strict group \(S(\alpha,r)\)
of Section~\ref{sec:oddquartic539}. Its local condition at
the coefficient prime \(p\) is zero.

\begin{theorem}[A class number bound at the limiting residue degree]
\leavevmode\label{thm:fixedboundary543}
Let \(\alpha\) be a fixed primitive character of conductor \(M\)
and order \(d\). Write \(m=\operatorname{lcm}(d,2)=2n\) and
suppose \(n\ge2\). For an odd prime \(p\nmid Md\) let
\(\eta_p\) be the quadratic character of conductor \(p\).
If \(\alpha(-1)=\eta_p(-1)\), the field
\(L=K_{\alpha^{-1}\eta_p}\) is totally real and
\[
                    h_L\ll_{\alpha}
                    \frac{p^{n/2}}{(\log p)^{n-1}} .
\]
Consequently, at every sufficiently large prime of this parity
with \(\operatorname{ord}_m(p)\ge m/4\), one has
\(S(\alpha,r)=0\) for all positive
\(r\equiv(p-1)/2\bmod(p-1)\).
\end{theorem}
\begin{proof}
Put \(\lambda=\alpha^{-1}\eta_p\). The conductor factors are
coprime, so the two characters attain their values independently.
It follows that \(\lambda\) has order \(2n\) and
\(\lambda^2=\alpha^{-2}\) has order \(n\).
Write \(F=K_{\alpha^2}\), a fixed field of degree \(n\).
The parity assumption makes its quadratic extension \(L\) totally
real. The inertia at \(p\) has order two in \(L\) and trivial
image in \(F\).

We first prove
\begin{equation}
          \prod_{\substack{b\bmod 2n\\b\ {\rm odd}}}
                 |L(1,\lambda^b)|\ll_{\alpha}\log(3Mp).
          \label{eq:fixedLproduct543}
\end{equation}
Every odd power has conductor divisible by \(p\) and dividing
\(Mp\). Thus it is nonprincipal. Take
\(\sigma=1+1/\log(3Mp)\). The completed-function inequality
in the proof of Lemma~\ref{lem:geometricmean542} gives
\(\lvert L(1,\lambda^b)\rvert\ll
\lvert L(\sigma,\lambda^b)\rvert\) with an absolute
constant. Indeed the primitive conductor \(D_b\le Mp\)
contributes a factor \(D_b^{(\sigma-1)/2}\le\sqrt e\),
and the two possible gamma ratios are bounded on this interval.

At a prime \(q\nmid Mp\), the mean of \(\lambda(q)^b\) over
these \(n\) exponents is zero unless \(\alpha(q)^2=1\).
This follows by summing a geometric progression with ratio
\(\lambda(q)^2\). Its mean has real part at most one when
the ratio is one. Terms of exponent at least two in the
logarithms of the Euler products have sum bounded absolutely
by \(\sum_q\sum_{k\ge2}q^{-k}/k\). The Euler factors belonging
to the fixed primes \(q\mid M\) can increase the mean logarithm
by at most \(\sum_{q\mid M}-\log(1-1/q)\). All factors at
\(p\) are omitted. We can include \(q=p\) when taking an
upper bound for the resulting nonnegative sum of prime terms.
For \(q\nmid M\), orthogonality expresses the indicator of
\(\alpha(q)^2=1\) as the average of the \(n\) powers of
\(\alpha(q)^2\). Restoring prime powers of exponent at least
two in this new average costs another absolute constant. Hence
\[
  \frac1n\sum_{b\ {\rm odd}}\log |L(\sigma,\lambda^b)|
   \le \frac1n\sum_{i=0}^{n-1}
            \log |L_M(\sigma,\alpha^{2i})|+O_\alpha(1),
\]
where \(L_M\) omits the factors at \(q\mid M\).
Only \(i=0\) is principal. Its absolute value is at most
\(\zeta(\sigma)\le1+1/(\sigma-1)\).
Each nonprincipal term is bounded in absolute value before
taking its logarithm, uniformly in \(1<\sigma\le2\).
In fact the partial sums of a nonprincipal character modulo
\(M\) have absolute value at most \(M\), and partial summation
bounds its Dirichlet series by \(M\) there.
Multiply the inequality by \(n\), which is fixed, and
exponentiate. The transfer to \(s=1\) proves
\eqref{eq:fixedLproduct543}.

Exactly the odd powers of \(\lambda\) occur in
\(\zeta_L/\zeta_F\). The conductor--discriminant formula gives
\(\lvert\operatorname{disc}L\rvert/
\lvert\operatorname{disc}F\rvert=\prod_{b\ {\rm odd}}D_b
\le(Mp)^n\).
In the analytic class number formula the two totally real
fields both have two roots of unity. Taking the quotient of
their residues therefore gives
\[
   h_L=\frac{h_F}{2^n}
        \left(\prod_{b\ {\rm odd}}D_b\right)^{1/2}
        \frac{R_F}{R_L}
        \prod_{b\ {\rm odd}}|L(1,\lambda^b)| .
\]
The regulator argument of Lemma~\ref{lem:relregulator529}
applies because \(p\) is unramified in \(F\) and every prime
of \(F\) above \(p\) ramifies in \(L/F\). More explicitly,
the discriminant of \(1,u\) for a nontrivial relative unit
modulo signs has norm at least \(p^n\). The logarithms at one
embedding in every conjugate pair thus have total absolute
value at least \(n(\frac12\log p-\log2)\). The cross-polytope
and determinant calculation of that lemma gives
\[
     \frac{R_L}{R_F}\ge
     \frac{\bigl[n(\tfrac12\log p-\log2)\bigr]^n}{n!}
     \qquad (p>4).
\]
The field \(F\) and its class number are fixed. Inserting the
last inequality and \eqref{eq:fixedLproduct543} proves the
claimed estimate for all sufficiently large \(p\). A single
larger constant includes the remaining primes.

At a half-residue weight the residual scalar character is
the reduction of \(\lambda\). Its order stays \(m\) because
\(p\nmid m\). Its restriction to inertia at \(p\) is nontrivial.
Lemma~\ref{lem:collisionclass457} therefore makes
\(p^{\operatorname{ord}_m(p)}\) a divisor of \(h_L\) if
there is a nonzero strict class. Since \(n\ge2\), the class
number estimate gives \(h_L<p^{n/2}\) for all sufficiently
large \(p\). This contradicts the divisibility even when
\(\operatorname{ord}_m(p)=m/4\).
\end{proof}

Let \(D_8\) be the compositum of the character fields of order
dividing eight in \(K_\psi C_\infty\). To see that it is finite,
observe that the Galois group of this tower has an open procyclic pro-two
subgroup. Its subgroup of eighth powers has finite index
and every indicated character kills it. In particular \(D_8\)
has degree a power of two and is contained in \(K_\psi C_{j_0}\)
for some fixed \(j_0\).

\begin{corollary}[Odd weights on the subfield of exponent eight]
\leavevmode\label{cor:octicspectrum543}
On the intervals of Theorem~\ref{thm:conductor232}, with
\(t=\lfloor x^{401}\rfloor\), there is a fixed \(A\) such that
\[
     \sum_{p>2}\ \sum_{\substack{1\le r<t\\r\ {\rm odd}}}
       \log\big|\mathcal S_{p,r}(D_8)^{c=-1}\big|
       \ll_{\psi,\varepsilon} V^A+\exp(x^{2\varepsilon/3}).
\]
For each such \(r\) and every odd \(p\le P(x)\),
\(R_{p,r}(E_x)^{\operatorname{Gal}(E_x/D_8)}=0\)
throughout every sufficiently distant interval.
\end{corollary}
\begin{proof}
Each odd character of \(D_8\) has order two, four, or eight.
Choose \(p_0\) larger than the finitely many conductor primes,
the prime divisors of \([K_\psi:\QQ]\), the exceptional primes
of Lemma~\ref{lem:fullinteger517}, and the thresholds in
Proposition~\ref{prop:quarticconductor539} and the theorem
above for these characters. At an odd half-residue weight
one has \(p\equiv3\bmod4\). In both resulting residue classes
modulo eight the order of \(p\) modulo eight is two.
The theorem therefore kills every character of order eight
at this weight for \(p>p_0\). Orders two and four are covered
by Proposition~\ref{prop:quarticconductor539}. Over the
unramified coefficient ring containing the eighth roots of
unity the character projectors are integral since \(p\) is
odd. Prime-to-\(p\) restriction, globally and at each local
condition, identifies the sign summand with these scalar
summands as in Theorem~\ref{thm:oddquarticfield539}. Faithful
flatness gives zero before extending coefficients. The zero
residue has no odd weights. The product of the orders at
the remaining weights and \(p>p_0\) divides \(J(V)\) by
Lemma~\ref{lem:fullinteger517}, and \(\log J(V)\ll V^A\).

For \(2<p\le p_0\) restrict from \(D_8\) to the actual
\(B_J\). Corestriction kills the kernel after multiplication
by \([B_J:D_8]\), whose \(p\)-valuation is at most
\(v_p([K_\psi:\QQ])\). At each of these finitely many primes
the \(p\)-rank on \(D_8\) is bounded independently of \(r\).
Indeed the Kummer calculation in Theorem~\ref{thm:oddquarticfield539}
uses only the \(p\)-units and the \(p\)-torsion in the class
group of the fixed \(D_8(\mu_p)\) with \(p\) inverted.
Thus the kernel costs \(O_\psi(1)\) in logarithmic order per
weight, and Theorem~\ref{thm:log147} bounds the sum of the
logarithmic orders of the images by \(\exp(x^{2\varepsilon/3})\).
Absorb the \(O_\psi(t)\) losses by increasing \(A\).

Write \(H=\operatorname{Gal}(E_x/D_8)\). If \(p\nmid |H|\),
averaging identifies strict invariants with restriction
from \(D_8\) and gives exact invariants on the subsequent
quotients. For large \(V\) the displayed bound divided by
\(\log p\) is less than
\(b_0(x)=\lfloor\exp(x^\varepsilon/8)\rfloor\).
The quotient by the subgroup killed by \(p^{b_0(x)}\)
therefore removes the invariant summand. If \(p\mid |H|\),
then \(p\) lies in the fixed prime set \(\Pi\) of
Proposition~\ref{prop:res23}. Both \(B(x)\) and \(L_x\)
there contain the fixed \(D_8\) for large \(V\). Hence the
subgroup defining its average \(J_{p,r}\) is contained in
\(H\). The average acts as the identity on every \(H\)-invariant
and as zero on the residual quotient, proving the assertion.
\end{proof}

\begin{theorem}[One growing subfield down to the smallest coefficient primes]
\leavevmode\label{thm:allprimeslayers543}
Fix a quadratic field \(K_\xi\) as in
Corollary~\ref{cor:complexallweights542}, or take \(K_\xi=\QQ\).
Set
\[
        Y_V=\tfrac12c_{\psi,\varepsilon}(\log\log V)^{1/5}
\]
with the constant of Theorem~\ref{thm:log147}.
There is a choice of \(D'_V=K_\xi C_j\), common to all
\(x\in[V,V^{3/2}]\), with
\([D'_V:\QQ]\asymp_{\psi,\varepsilon,\xi}
(\log\log\log V)^2\), for which
\[
   R_{p,r}(E_x)^{\operatorname{Gal}(E_x/D'_V)}=0
           \qquad(1\le r<\lfloor x^{401}\rfloor)
\]
at every odd \(p\le Y_V\) and at every odd \(p\le P(x)\)
congruent to three or five modulo eight. For \(p>P(x)\) the
logarithmic order on \(D'_V\), summed over all such primes and
weights, is \(O_{\psi,\varepsilon}(V^A)\).
\end{theorem}
\begin{proof}
Choose \(b\) as in Corollary~\ref{cor:complexallweights542}
and take \(j\) largest with \(2^{j-1}\le b(\log Y_V)^2\).
Successive degrees in the real tower double and
\(\log Y_V=\tfrac15\log\log\log V+O_{\psi,\varepsilon}(1)\).
This proves the degree estimate, since adjoining the fixed
quadratic field changes the degree by at most two.
For all sufficiently large \(V\) we have
\(j=O(\log\log Y_V)\) and any fixed initial layer containing
\(K_\xi\) is available. Thus \(D'_V\) lies in \(B_J\) and
in each of the two fields \(B(x),L_x\) used for the
averages in Proposition~\ref{prop:res23}. In particular this
choice works throughout the original interval.

First take \(p\le Y_V\). If \(p\mid[E_x:\QQ]\), its odd
part comes from \([K_\psi:\QQ]\), so \(p\in\Pi\). The average
argument at the end of the preceding corollary applies
because both its candidate fields contain \(D'_V\).
If \(p\nmid[E_x:\QQ]\), restriction injects the complementary
strict group on \(D'_V\) into the corresponding one on
\(B_J\). Theorem~\ref{thm:log147} bounds their logarithmic
orders even after summing over all weights and primes up to
\(Y_V\), since \(x\ge V\). The resulting bound divided by
\(\log p\) is less than \(b_0(x)\) for large \(V\).
Prime-to-\(p\) invariants are exact, so the truncation again
removes the entire indicated summand.

Now take \(Y_V<p\le P(x)\) in either specified progression.
Eventually \(p\) exceeds every fixed prime excluded in
Lemma~\ref{lem:fullinteger517} and is prime to \([E_x:\QQ]\).
Here \(2^{j-1}\le b(\log p)^2\), the precise size condition
used at the half residue in Theorem~\ref{thm:complexhalf542}.
For each character in the required sign summand the dyadic
order is a power of two. At orders at least four,
\(v_2(p^2-1)=3\) gives the residue-degree bound in that
theorem. The characters of order at most two are covered by
Propositions~\ref{prop:quarticconductor539}
and~\ref{prop:qbound102}, as in the proof of
Corollary~\ref{cor:complexallweights542}. At a multiple of
\(p-1\) the sign is plus and descends to the real \(C_j\),
where Theorem~\ref{thm:tower138} applies. At all separated
weights Lemma~\ref{lem:fullinteger517} bounds the length by
\(\log J(V)/\log p<b_0(x)\). These arguments require only
the lower bound \(p>Y_V\) just verified, so they also cover
the primes below \(V\). Exactness of averaging gives the
claimed vanishing. Finally \(p>P(x)>2\lfloor
V^{1203/2}\rfloor+1\) on sufficiently distant intervals.
All weights under consideration are separated and the same
integer \(J(V)\) bounds the summed logarithmic order.
\end{proof}

The subfields in the last two statements belong to the fixed
tower \(K_\psi C_\infty\). Section~\ref{sec:missing} also
requires estimates on components involving the two quadratic
characters whose conductors vary with the interval.




\subsection{Half-residue weights over a fixed abelian field of two-power degree}
\label{sec:mixedhalf633}

The field cut out by the square of a character need not be a pure
cyclotomic layer. Working over that field gives a uniform estimate as
the order and the conductor both vary. For a finite Dirichlet character
\(\nu\), let \(K_\nu\) be its character field and let
\(\operatorname{cond}(\nu)\) be its primitive conductor, equal to one
for the trivial character. All \(L\)-values in this subsection have
their primitive Euler factors.

\begin{lemma}[An upper bound over the field of the squared character]
\leavevmode\label{lem:squarefield633}
Let \(\lambda\) be a Dirichlet character of order \(2n\), \(n\ge2\).
Put \(B=\operatorname{cond}(\lambda^2)\),
\(D=\max_{b\ {\rm odd}}\operatorname{cond}(\lambda^b)\), where
exponents are read modulo \(2n\), and set
\[
 H(B,n)=\min\left\{1+\log(2B),\
        \frac{B}{\varphi(B)}\left(1+\sqrt{\frac{\varphi(B)}{n}}\right)
                  \right\}.
\]
There is an absolute constant \(A\) such that
\[
 \left(\prod_{b\ {\rm odd}} |L(1,\lambda^b)|\right)^{1/n}
 \le A\,(\log(3D))^{1/n} H(B,n)^{\,1-1/n}.
\]
\end{lemma}
\begin{proof}
Write \(L=K_\lambda\), \(F=K_{\lambda^2}\), and
\(\sigma=1+1/\log(3D)\). For each odd power, the comparison of
completed functions proved in Lemma~\ref{lem:geometricmean542} gives
\(\lvert L(1,\lambda^b)\rvert\le A_1\lvert L(\sigma,\lambda^b)\rvert\).
Indeed the conductor factor from \(1\) to \(\sigma\) is at most
\(\sqrt e\), and the gamma ratios for both parities are bounded on
\([1,2]\).

The product of the absolute values at \(\sigma\) is
\(\zeta_L(\sigma)/\zeta_F(\sigma)\), by abelian factorization and
conjugate pairing. At a prime ideal of \(F\) of norm \(Q\), the
Euler factor of this quotient is \((1-Q^{-\sigma})^{-1}\),
\((1+Q^{-\sigma})^{-1}\), or \(1\), according as the prime splits,
is inert, or ramifies in \(L/F\). Each is bounded above by the
Euler factor of \(\zeta_F(\sigma)\). Thus the entire quotient is
at most \(\zeta_F(\sigma)\), including at primes dividing \(B\).

Factor this last zeta function using the \(n\) powers of
\(\nu=\lambda^2\). The principal factor is
\(\zeta(\sigma)\le1+1/(\sigma-1)\). For a nonprincipal primitive
character of conductor \(b_0\mid B\), periodicity bounds its
partial sums by \(b_0\). Summing the initial \(b_0\) terms and using
partial summation for the tail gives an upper bound
\(A_2(1+\log(2B))\) for the absolute value of each such factor.

For a second bound, first read all powers of \(\nu\) modulo \(B\).
For each nonprincipal power the tail after \(B\) is bounded by an
absolute constant, again by partial summation and periodicity.
Let \(T_i\) denote the sum of the first \(B\) terms for \(\nu^i\)
at \(\sigma\), including \(i=0\). Orthogonality gives
\[
 \frac1n\sum_{i=0}^{n-1}|T_i|^2
 =\sum_{z\in\operatorname{im}\nu}
    \left|\sum_{\substack{1\le a\le B,\ (a,B)=1\\ \nu(a)=z}}
                      a^{-\sigma}\right|^2
 \le \frac{\varphi(B)}{n}\sum_{a\ge1} a^{-2}.
\]
Every fiber has \(\varphi(B)/n\) members, which gives the inequality
by Cauchy--Schwarz. Restoring omitted Euler factors increases each
absolute value by a factor at most
\(\prod_{\ell\mid B}(1-1/\ell)^{-1}=B/\varphi(B)\).
The arithmetic mean over the \(n-1\) nonprincipal factors is therefore
at most \(A_3H(B,n)\), using \(n/(n-1)\le2\) and taking the better
of the two bounds. Their geometric mean satisfies the same estimate.
Insert the principal factor and take the \(n\)-th root.
\end{proof}

\begin{theorem}[A uniform necessary inequality for a half-residue class]
\leavevmode\label{thm:relhalf633}
Let \(\alpha\) have order \(d\) and conductor \(M\), let
\(2n=\operatorname{lcm}(2,d)\) with \(n\ge2\), and let \(p\ge5\)
be a prime not dividing \(Md\). Suppose that
\(r\equiv(p-1)/2\bmod(p-1)\), \(r>0\), and
\(\alpha(-1)=(-1)^r\). Put \(f=\operatorname{ord}_{2n}(p)\),
\(B=\operatorname{cond}(\alpha^2)\), and
\[
 M_{\rm o}=\left(\prod_{\substack{b\bmod 2n\\ b\ {\rm odd}}}
                     \operatorname{cond}(\alpha^b)\right)^{1/n}.
\]
If the full strict scalar group \(S(\alpha,r)\) of
Section~\ref{sec:levelstrict345} is nonzero, then, for an absolute \(A\),
\[
 p^{f/n}\ \le\
 A\,\frac{\sqrt{p M_{\rm o}}}{\log(p/4)}
       (\log(3Mp))^{1/n} H(B,n)^{\,1-1/n}.
\]
\end{theorem}
\begin{proof}
Let \(\eta_p\) be the quadratic character of conductor \(p\),
\(\lambda=\alpha^{-1}\eta_p\), \(L=K_\lambda\),
and \(F=K_{\lambda^2}\). The conductor factors of \(\lambda\)
are coprime, so that its order is \(2n\) and its odd powers have
conductors \(p\operatorname{cond}(\alpha^b)\). The parity condition
makes \(L\) totally real. Every prime over \(p\) ramifies in the
quadratic extension \(L/F\) because \(F\) is unramified at \(p\)
and inertia acts there on \(\lambda\) by \(\eta_p\).

Suppose a strict class survives. Its residual character is
\(\bar\lambda\), and inertia at \(p\) acts nontrivially.
The proof of Lemma~\ref{lem:collisionclass457} supplies an
equivariant homomorphism \(u\) from \(\operatorname{Cl}(L)\) to
an additive group of characteristic \(p\) whose image has at
least \(p^f\) elements. The involution \(s\) of \(L/F\) acts on
that image by minus one. For every ideal class \(a\), the class
\(a/s(a)\) has trivial norm and its image is \(2u(a)\). Thus
the kernel of the norm already maps onto the image, since two
is invertible there.

In this extension the class group norm is surjective. Indeed the
ordinary Hilbert class field \(H_F\) has \(H_F\cap L=F\) by the
ramification at \(p\). The compositum \(H_FL/L\) is unramified
at finite places and split at the real places. Reciprocity on
\(L\) maps its class group onto \(\operatorname{Gal}(H_FL/L)\);
restriction to \(H_F\) identifies this map with the ideal norm
followed by reciprocity on \(F\). This proves
\[
                 p^f\le|\ker N_{L/F}|=h_L/h_F .
\]

The regulator estimate used in Theorem~\ref{thm:complexhalf542}
depends only on total reality, quadratic relative degree, and full
ramification over the unramified rational prime \(p\). To recall its
content, for a nontrivial relative unit \(v\) modulo signs the order
with basis \(1,v\) has discriminant \((v-s(v))^2\). The norm of the
relative discriminant is at least \(p^n\). For one embedding
\(\tau\) of \(L\) above each embedding of \(F\), this gives
\[
 p^{n/2}\le\prod_\tau|\tau(v-s(v))|
       \le 2^n\exp\left(\sum_\tau|\log|\tau(v)||\right).
\]
Here \(v s(v)=\pm1\). The volume of the resulting empty open
cross-polytope and the norm determinant of
Lemma~\ref{lem:relregulator529} give
\(R_L/R_F\ge[n\log(p/4)/2]^n/n!\).
The conductor--discriminant and class number formulas, with two
roots of unity in both fields, now give
\[
 \frac{h_L}{h_F}
  =\frac{(p M_{\rm o})^{n/2}}{2^n}
        \frac{R_F}{R_L}\prod_{b\ {\rm odd}}|L(1,\lambda^b)|.
\]
Apply Lemma~\ref{lem:squarefield633}, use \(D\le Mp\), and take
the \(n\)-th root. The factor \((n!)^{1/n}/n\) is at most one.
\end{proof}

\begin{corollary}[Half-residue vanishing on the mixed subfield throughout the interval]
\leavevmode\label{cor:mixedhalf633}
Let \(K\) be the maximal subfield of \(K_\psi\) of two-power degree,
including \(K=\mathbb Q\). There is a constant \(b_\psi>0\) with
the following property on every sufficiently distant interval of
Theorem~\ref{thm:conductor232}. Take \(j\) largest with
\(2^j\le b_\psi(\log V)^2\), and put \(T_V=K C_j\).
This one field is contained in every \(E_x\) on the interval.
For every \(p\ge V\) congruent to three or five modulo eight and
every half-residue weight \(1\le r<t(x)\), one has
\[
       \mathcal S_{p,r}(T_V)^{c=(-1)^r}=0,\qquad
       R_{p,r}(E_x)^{\operatorname{Gal}(E_x/T_V)}=0 .
\]
\end{corollary}
\begin{proof}
The initial field \(K\) is fixed and cyclic; write \(a=[K:\mathbb Q]\).
The cyclotomic degree available in every \(E_x\) grows as a positive
power of \(V\) by Theorem~\ref{thm:conductor232}, so all contain
\(T_V\). A character of \(T_V\) is a product \(\alpha=\gamma\beta\)
with \(\gamma\) from \(K\) and \(\beta\) from \(C_j\). Indeed the
Galois group embeds in the product for the two factors, and a
character of a subgroup of a finite abelian group extends to the
whole group. One may extend a generator at a time by taking a root
of its prescribed power.

If \(\alpha\) has order \(2n\ge4\), then
\(\operatorname{ord}(\beta)\le\max(2n,a)\), since both factors
have two-power order. A nontrivial character of exact order \(2^i\)
in the real dyadic tower first occurs at \(C_i\) and has conductor
\(2^{i+2}\). Consequently \(M\le A_\psi n\), also \(B\le M\),
and every prime dividing \(B\) divides the fixed integer
\(2\operatorname{cond}(\psi)\). Hence \(H(B,n)\le A'_\psi\).
The possible \(2n\) are at most \(2^j\) once \(2^j\ge a\).
For \(p\equiv3,5\bmod8\), \(v_2(p^2-1)=3\) and every subsequent
squaring increases the valuation by one. This gives
\(\operatorname{ord}_{2n}(p)\ge n/2\), including modulus four.
By the theorem a nonzero group would require
\[
       1\le A''_\psi
            \frac{\sqrt n\,(\log(3A_\psi n p))^{1/n}}{\log(p/4)} .
\]
Put \(y=\log\log(3p)\). For \(2\le n\le y\) the right side is
\(O_\psi(\sqrt{y/\log p})\), which tends to zero uniformly.
If \(y<n\le b_\psi(\log p)^2\) and \(b_\psi\le1\), the factor
with exponent \(1/n\) is at most four for large \(p\). Choosing
\(b_\psi\) sufficiently small makes the right side strictly less
than one in this range as well.

The characters of order at most two in \(K C_\infty\) have bounded
conductor by the same argument for \(\gamma\beta\). For them
\(L=K_{\alpha^{-1}\eta_p}\) is real quadratic of discriminant
\(O_\psi(p)\) when the signs match. Minkowski gives a representative
of norm at most \(Y=O_\psi(\sqrt p)\) for every ideal class, with
\(Y\ge1\). The quadratic Euler factors bound the number of ideals
of each norm by the divisor function, so \(h_L\le Y(1+\log Y)<p\)
eventually. This contradicts Lemma~\ref{lem:collisionclass457},
including when \(\alpha\) is trivial.

All primes in question are prime to \([E_x:\mathbb Q]\) and to the
conductors of the characters of \(T_V\) for large \(V\). The
character decomposition with the strict local conditions from
Theorem~\ref{thm:levelstrict345} therefore proves the first equality.
Restriction identifies this group with the invariants of the matching
sign group on \(E_x\), since averaging is exact for the relative
group of order prime to \(p\), also locally. Field averages and the
truncation in Proposition~\ref{prop:res23} commute with this
projection. This proves the second equality at every indicated \(x\).
\end{proof}

For instance, a quartic character modulo \(17\) is even and its square
has conductor \(17\). The corollary applies to all its twists by the
indicated growing dyadic characters. The two varying auxiliary
quadratic conductors of \(E_x\) can enter \(M_{\rm o}\) in the
theorem. Their contribution to the full sum in
Section~\ref{sec:missing} still requires an estimate.


\subsection{Descent at the zero residue over a fixed real cyclic field}
\label{sec:zerodescent634}

At weights divisible by \(p-1\), the places over \(p\) give local
coordinates for a strict class killed by \(p\). To use these coordinates
one must control the ordinary class group of the field on which they
are taken. We give a uniform bound as the cyclotomic layer grows.

\begin{lemma}[A lower bound for the relative regulator in totally real degree two]
\leavevmode\label{lem:realunits634}
Let \(L/F\) be a quadratic extension of totally real fields with
\([F:\mathbb Q]=n\). Put \(c_0=\tfrac12\log(3\sqrt3/2)>0\).
Then \(R_L/R_F\ge(c_0 n)^n/n!\).
\end{lemma}
\begin{proof}
For every real number \(z\),
\[
       |z(z^2-1)|\le\frac{2}{3\sqrt3}\max(1,|z|)^4 .
\]
On \(0\le t=|z|\le1\), the maximum of \(t(1-t^2)\) is
\(2/(3\sqrt3)\). For \(t\ge1\), dividing by \(t^4\) gives the
same function at \(1/t\).
Let \(s\) be the involution of \(L/F\), and choose a nontrivial
unit \(u\) in the kernel of the norm on units modulo signs.
Then \(u\ne\pm1\), so \(u(u^2-1)\) is a nonzero algebraic integer.
Its absolute norm is at least one. In each pair of real embeddings
over \(F\), the two absolute values of \(u\) are reciprocal since
\(u\,s(u)=\pm1\). Taking the product of the scalar inequality over
all \(2n\) embeddings gives
\[
          1\le \left(\frac{2}{3\sqrt3}\right)^{2n}
                 \exp\left(4\sum_\tau|\log|\tau(u)||\right),
\]
where one member \(\tau\) of each pair occurs in the sum.
Thus every nonzero log vector of the relative units has
\(\ell^1\)-norm at least \(c_0 n\). These vectors form a full lattice
in \(\mathbb R^n\), since norm images contain the squares of the
units of \(F\). The cross-polytope estimate and the determinant
computation of Lemma~\ref{lem:relregulator529} give the result.
\end{proof}

\begin{theorem}[Large primes in the class groups of mixed real layers]
\leavevmode\label{thm:mixedclasses634}
Fix a totally real cyclic number field \(K\), allowing \(K=\mathbb Q\).
Write \(a=[K:\mathbb Q]\) and let \(a_{\rm odd}\) be its odd part.
Put \(U_j=K C_j\), with \(C_j=\mathbb Q(\zeta_{2^{j+2}})^+\),
and \(C_0=\mathbb Q\). There is \(A_K\), independent of \(j\), such
that for every \(j\ge0\) and every prime \(p\equiv3,5\bmod8\),
\[
               p>A_K(1+2^j)^{a_{\rm odd}}
           \quad\Longrightarrow\quad
                    \operatorname{Cl}(U_j)\{p\}=0 .
\]
The bound is linear in \(2^j\) when \(a\) is a power of two.
\end{theorem}
\begin{proof}
Enlarge \(A_K\) to exclude all primes dividing \(2a\) or
\(\operatorname{cond}(K)\). Here the conductor of \(K\) is the
least cyclotomic modulus containing it. Suppose that
\(\operatorname{Cl}(U_j)[p]\ne0\). Its Galois action is semisimple
over \(\overline{\mathbb F}_p\), because \(p\) is prime to the
abelian Galois group. Choose an occurring character, lift its values
to roots of unity of order prime to \(p\), and call it \(\alpha\).
Put \(m=\operatorname{ord}(\alpha)\), \(L=K_\alpha\). The character
also occurs in the class group of \(L\) after reduction to its
\(p\)-torsion and extension of scalars. In fact, on the \(p\)-part
of the class groups, extension followed by norm is multiplication
by \([U_j:L]\), a unit. Extension of the norm is the sum over
\(\operatorname{Gal}(U_j/L)\), which is that same unit times the
averaging idempotent. This identifies the class group of \(L\) at
\(p\) with the invariant subgroup and is compatible with taking
\(p\)-torsion.

Write \(\alpha=\gamma\beta\) with \(\gamma\) from \(K\) and
\(\beta\) from \(C_j\), as in Corollary~\ref{cor:mixedhalf633}.
The order of \(\beta\) divides the two-part of
\(\operatorname{lcm}(m,a)\). If \(m\) is odd or \(m=2\), this
puts \(L\) in one fixed finite extension of \(\mathbb Q\). There
are only finitely many such subfields and their class numbers have
bounded prime factors. Enlarge \(A_K\) to cover them.

It remains to consider \(m=2n\), \(n\ge2\). Set \(F=K_{\alpha^2}\)
and denote the involution of \(L/F\) by \(s\). The Frobenius orbit
of the faithful residual character has size
\(f=\operatorname{ord}_m(p)\): it is the least exponent for which
raising a primitive \(m\)-th root to \(p^f\) fixes it. All its
conjugates occur, as the class group action is defined over
\(\mathbb F_p\). They contribute dimension at least \(f\) to the
minus eigenspace of \(s\). The norm kills this space, because it
has the same value on a class and its \(s\)-conjugate, and two is
a unit on the \(p\)-torsion.

The image of the ideal norm from \(L\) has index at most two in
\(\operatorname{Cl}(F)\). To see this, form the compositum of \(L\)
with the ordinary Hilbert class field \(H_F\). Class field reciprocity
on \(L\) maps onto \(\operatorname{Gal}(H_FL/L)\), with restriction
to \(H_F\) given by the ideal norm. Its image in
\(\operatorname{Gal}(H_F/F)\) has index \([H_F\cap L:F]\le2\).
Thus \(p^f\le|\ker N_{L/F}|\le2 h_L/h_F\).

Put \(M=\operatorname{cond}(\alpha)\), \(B=\operatorname{cond}(\alpha^2)\).
The order bound on \(\beta\) and the conductor \(4\operatorname{ord}(\beta)\)
of a nontrivial real dyadic tower character give \(M\le C_K m\).
The primes dividing \(B\) belong to the fixed set dividing
\(2\operatorname{cond}(K)\), and \(B\le M\). It follows that
\(H(B,n)\) of Lemma~\ref{lem:squarefield633} is bounded in terms of
\(K\). The same is true of \((\log(3M))^{1/n}\). All the odd powers
of \(\alpha\) have conductor at most \(M\). As \(L\) and \(F\) are
totally real, dividing their class number formulas and applying that
lemma and Lemma~\ref{lem:realunits634} yields
\[
   (h_L/h_F)^{1/n}
     \le \frac{\sqrt M}{2}
              \left(\frac{R_F}{R_L}
                     \prod_{b\ {\rm odd}\bmod 2n}|L(1,\alpha^b)|
              \right)^{1/n}
     \le C'_K\sqrt m .
\]
Here the first inequality uses the conductor--discriminant formula
for the odd powers, and \((n!)^{1/n}/n\le1\) in the second.

Write \(m=2^v w\) with \(w\) odd. Then \(w\mid a_{\rm odd}\).
For the indicated primes the order modulo \(2^v\) is at least
\(2^{v-2}\), by \(v_2(p^2-1)=3\) and successive squaring, with
\(v=1,2\) covered directly. Hence \(f\ge m/(4w)\).
Combining the lower and upper bounds gives
\(p^{1/(2w)}\le p^{f/n}\le C''_K\sqrt m\).
Since \(m\mid\operatorname{lcm}(a,2^j)\), this implies
\(p\le A_K(1+2^j)^{a_{\rm odd}}\) after increasing \(A_K\).
This contradicts the assumed lower bound for \(p\).
\end{proof}

\begin{corollary}[Descent of the full strict group to a bounded cyclic field]
\leavevmode\label{cor:frobenius634}
Under the size and congruence conditions in
Theorem~\ref{thm:mixedclasses634}, let \(\sigma_p\) be Frobenius
in \(U_j/\mathbb Q\), and put \(P_{p,j}=U_j^{\langle\sigma_p\rangle}\).
For every positive \(r\) divisible by \(p-1\), restriction gives
\[
          \mathcal S_{p,r}(P_{p,j})
                      \ \simeq\ \mathcal S_{p,r}(U_j).
\]
The field \(P_{p,j}\) is cyclic of degree at most \(a\) and is
contained in \(K C_{v_2(a)}\). For \(j\ge v_2(a)\) it equals
\(P_{p,v_2(a)}\).
\end{corollary}
\begin{proof}
The prime \(p\) is unramified in \(U_j\). By
Lemma~\ref{lem:slope138}, triviality of its ordinary \(p\)-class
group embeds \(\mathcal S_{p,r}(U_j)[p]\), equivariantly, in the
permutation module on the primes over \(p\) modulo the constant
line. The decomposition group \(\langle\sigma_p\rangle\) fixes
every such prime since the extension is abelian. Its average
therefore fixes the entire \(p\)-torsion. Any nonzero element
in the complementary summand would have a multiple of order \(p\)
there. Thus it fixes the full strict group. Restriction from the
fixed field identifies the lower strict group with these invariants:
the relative group has order prime to \(p\), so restriction and
averaging give descent on global cohomology. Local restriction is
injective at \(p\), and at every other prime an inertia homomorphism
which vanishes upstairs factors through a group of order prime to
\(p\). This proves the isomorphism.

The projection of \(\sigma_p\) generates
\(\operatorname{Gal}(C_j/\mathbb Q)\). In fact choose \(u=\pm p\)
with \(u\equiv1\bmod4\). Then \(v_2(u-1)=2\), and
\(v_2(u^{2^i}-1)=2+i\) for \(i\ge0\), by repeated factorization.
This gives order \(2^j\) in the quotient modulo signs at conductor
\(2^{j+2}\). It follows that the quotient by \(\langle\sigma_p\rangle\)
is generated by the image of \(\operatorname{Gal}(U_j/C_j)\), a
subgroup of the cyclic group \(\operatorname{Gal}(K/\mathbb Q)\).
This proves cyclicity and the degree bound. For any character
\(\gamma\beta\) trivial on \(\sigma_p\), one has
\(\beta(p)=\gamma(p)^{-1}\). Since \(p\) generates the dyadic
quotient, the order of \(\beta\) equals that of \(\beta(p)\),
and divides \(2^{v_2(a)}\). Every such character therefore factors
through \(K C_{v_2(a)}\). If \(j\ge v_2(a)\), Frobenius there is the
restriction of the same \(\sigma_p\); its characters trivial on this
element are exactly the ones just identified. This proves the last
two assertions.
\end{proof}

\begin{corollary}[The remaining summand on the original intervals]
\leavevmode\label{cor:intervalzero634}
For the actual fields of Theorem~\ref{thm:conductor232}, take
\(K=K_\psi^+\), the maximal totally real subfield of \(K_\psi\).
For \(j\le J(x)\) satisfying the hypotheses of
Theorem~\ref{thm:mixedclasses634} and \(p\nmid[E_x:\mathbb Q]\),
write \(e_D\) for averaging onto a subfield \(D\) on \(R_{p,r}(E_x)\).
At every \(1\le r<t(x)\) divisible by \(p-1\),
\[
                  e_{K C_j} R_{p,r}(E_x)
                     =e_{P_{p,j}}R_{p,r}(E_x).
\]
In particular the assertion applies to the entire real subfield
of \(K_\psi C_{J(x)}\) whenever the size bound holds at \(j=J(x)\).

For \(T_V\) of Corollary~\ref{cor:mixedhalf633} one also has, for
every \(x\) on the interval and every \(V\le p\le P(x)\) congruent to
three or five modulo eight,
\[
      e_{T_V}R_{p,r}(E_x)=
      \begin{cases}
        e_{P'}R_{p,r}(E_x),&p-1\mid r,\\
        0,&p-1\nmid r ,
      \end{cases}\qquad 1\le r<t(x).
\]
Here \(P'\) is the decomposition field of \(p\) in \(T_V^+\);
it belongs to a finite list depending only on \(\psi\).
\end{corollary}
\begin{proof}
Since \(p-1\mid r\), the complementary sign is plus. Apply
Corollary~\ref{cor:frobenius634} to \(K C_j\) and restrict into
\(E_x\). The field averages, the sign projector and the quotient by
the subgroup of bounded exponent in Proposition~\ref{prop:res23}
all commute. Exactness of averaging gives the first equality after
these operations. The fixed field of complex conjugation on the
compositum with the real \(C_j\) is \(K_\psi^+ C_j\): adjoining
\(K_\psi\) to this field has degree at most two, and conjugation acts
nontrivially if \(K_\psi\) is imaginary.

For the second assertion the same argument identifies \(T_V^+\)
as the compositum of \(C_j\) with the real part of the fixed
two-power subfield of \(K_\psi\). Its degree condition at \(p\)
holds uniformly since \(2^j\ll_\psi(\log V)^2\) and \(p\ge V\).
At the zero residue apply the first argument with this fixed real
field as \(K\); the preceding corollary supplies the finite list.
At the half residue apply Corollary~\ref{cor:mixedhalf633}.
For all other \(r\), Lemma~\ref{lem:fullinteger517} bounds the
logarithmic order on \(T_V\) by \(\log J(V)\ll_{\psi,\varepsilon}V^C\).
Dividing by \(\log p\) gives a length smaller than \(b_0(x)\) on
every sufficiently distant interval. The truncation for \(p\le P(x)\)
then kills the summand. This accounts for every positive weight
in the stated range.
\end{proof}

The decomposition fields in the formula are totally split at \(p\).
The local unit conditions on these fixed fields can still contribute
at multiples of \(p-1\); the descent identifies the group carrying
that contribution.


\subsection{Descent at varying Frobenius orders in the real dyadic tower}
\label{sec:varyingdecomposition642}

At zero-residue weights the part of the strict group detected
above a decomposition field depends on the class group of the
real layer. We keep track of the order of Frobenius in that
layer for each odd coefficient prime. This lets us use a single
lower field across almost all the large primes of a common
interval, with an explicit bound on the exceptional prime and
weight pairs.

Let \(K/\QQ\) be totally real cyclic of degree \(a\), and write
\(a_{\rm o}\) for its odd part. Let
\(C_j=\QQ(\zeta_{2^{j+2}})^+\), \(U_j=KC_j\). For an odd
prime \(p\) put
\[
              h(p)=v_2(p^2-1)-3
                  =\max\{v_2(p-1),v_2(p+1)\}-2 .
\]
This integer is nonnegative. We use the usual Frobenius acting
by \(p\)-th power on unramified roots of unity. The real layer
and its Frobenius are also specified in the archive at
\srcloc{Formalism/FrobeniusOrbits/Statement}{156}.

\begin{theorem}[A class number bound with the dyadic Frobenius order]
\leavevmode\label{thm:variableclasses642}
There is a constant \(A_K\ge2\) such that for every odd prime
\(p\) and \(j\ge0\),
\[
        p>\bigl[A_K(1+2^j)\bigr]^{\,2^{h(p)}a_{\rm o}}
              \quad\Longrightarrow\quad
                 \operatorname{Cl}(U_j)\{p\}=0 .
\]
\end{theorem}
\begin{proof}
Increase \(A_K\) to exceed the prime divisors of
\(2a\operatorname{cond}(K)\). If the class group has nonzero
\(p\)-torsion, choose an occurring character over
\(\overline{\mathbb F}_p\) and write \(\alpha\) for its lift to
roots of unity of order prime to \(p\). Put \(m=\operatorname{ord}
(\alpha)\). Its component descends to \(L=K_\alpha\) by extension
and norm of ideals, since \(p\nmid[U_j:L]\). Write
\(\alpha=\gamma\beta\) with \(\gamma\) from \(K\) and \(\beta\)
from \(C_j\). Such a factorization follows by extending the
character of the subgroup
\(\operatorname{Gal}(U_j/\QQ)\) to the product of the two cyclic
Galois groups. If \(m\) is odd or \(m=2\), the order of \(\beta\)
divides the two-part of \(\operatorname{lcm}(m,a)\), which is
bounded in terms of \(K\). Here the characters are viewed on
the absolute Galois group of \(\QQ\), where the same order
bound follows from \(\beta=\alpha\gamma^{-1}\). These \(L\)
belong to one fixed finite extension. Enlarge \(A_K\) to cover
the finitely many prime factors of their class numbers.

Suppose \(m=2n\ge4\). Put \(F=K_{\alpha^2}\) and
\(f_p=\operatorname{ord}_m(p)\). The distinct \(p\)-power
conjugates of the faithful character in the class group of \(L\)
give at least \(f_p\) dimensions over \(\mathbb F_p\) in the
minus part for \(L/F\). This part is killed by the ideal norm,
whose image in \(\operatorname{Cl}(F)\) has index at most two.
Hence \(p^{f_p}\le 2h_L/h_F\), as in the proof of
Theorem~\ref{thm:mixedclasses634}. Its relative class number
estimate
\((h_L/h_F)^{1/n}\le C_K\sqrt m\) applies to this \(L/F\)
for every such \(p\). Indeed its proof uses the conductor of
\(\alpha\), the relative regulator bound of
Lemma~\ref{lem:realunits634} and the product of the complex
\(L\)-values in Lemma~\ref{lem:squarefield633}. None of those
quantities involves the congruence class of the coefficient
prime. We have proved \(p^{f_p/n}\le C'_K\sqrt m\).

Write \(m=2^v w\) with \(w\) odd, so \(w\mid a_{\rm o}\).
Choose \(u\in\{p,-p\}\) with \(u\equiv1\bmod4\). Then
\(v_2(u-1)=h(p)+2\). Successive squaring gives
\(v_2(u^{2^i}-1)=h(p)+2+i\), because each factor
\(u^{2^i}+1\) has valuation one. Raising to a further odd
power leaves this valuation unchanged. In the quotient modulo
signs at conductor \(2^v\) the order of \(p\) is consequently
\(2^{\max(0,v-2-h(p))}\) for \(v\ge2\). A power of \(u\)
is \(1\bmod4\), so cannot be \(-1\bmod 2^v\) there. At
\(v=1\) use order one. In all cases
\[
           f_p\ge\frac{m}{4\,2^{h(p)}w},
       \qquad p^{1/(2\,2^{h(p)}w)}\le C'_K\sqrt m .
\]
Now \(m\mid\operatorname{lcm}(a,2^j)\) and
\(m\le a(1+2^j)\). Raising to the indicated power and
increasing \(A_K\) proves the theorem.
\end{proof}

\begin{theorem}[The full strict group on the decomposition field]
\leavevmode\label{thm:variabledescent642}
Assume the prime bound of Theorem~\ref{thm:variableclasses642}.
For \(P_{p,j}=U_j^{\langle\operatorname{Frob}_p\rangle}\) and
every \(r>0\) divisible by \(p-1\), restriction gives
\[
      \mathcal S_{p,r}(P_{p,j})\simeq\mathcal S_{p,r}(U_j).
\]
Moreover,
\[
   [P_{p,j}:\QQ]\le a\,2^{\min(j,h(p))},\qquad
   P_{p,j}\subseteq K C_{\min(j,h(p)+v_2(a))}.
\]
For every \(j\ge h(p)+v_2(a)\), the decomposition field itself
equals \(P_{p,h(p)+v_2(a)}\), also independently of the prime
bound on the class number.
\end{theorem}
\begin{proof}
At the unramified \(p\) the class group bound and
Lemma~\ref{lem:slope138} embed the \(p\)-torsion of the strict
group in the permutation module on primes over \(p\) modulo
constants. Its decomposition group fixes all these primes by
commutativity. Its order is prime to \(p\), and averaging
fixes the full strict group since a nonzero complementary
\(p\)-primary summand would have nonzero \(p\)-torsion.
Restriction from the decomposition field identifies its strict
group with these invariants. Indeed averaging gives global
cohomological descent, restriction at \(p\) is injective, and
away from \(p\) an inertia homomorphism becoming zero upstairs
would factor through a group of order prime to \(p\).

The order of Frobenius on \(C_j\) is
\(2^{\max(0,j-h(p))}\) by the calculation at conductor
\(2^{j+2}\) above. This bounds its order on \(U_j\) from below.
Since \([U_j:\QQ]\le a2^j\), the degree bound follows.
For each character \(\alpha=\gamma\beta\) trivial on this
Frobenius, the order of \(\beta(p)=\gamma(p)^{-1}\) divides
\(2^{v_2(a)}\). The subgroup generated by \(p\) in the
cyclic dyadic group has index \(2^{\min(j,h(p))}\). Hence
the order of \(\beta\) divides \(2^{h(p)+v_2(a)}\).
Every such \(\alpha\) factors through the lower field displayed
in the theorem. Conversely every character there that is
trivial on its Frobenius inflates to one trivial on the upper
Frobenius. These equal character groups prove stabilization.
\end{proof}

\begin{corollary}[One lower field and the exceptional prime and weight pairs]
\leavevmode\label{cor:exceptionalweights642}
On a common interval of Theorem~\ref{thm:conductor232} take
\(K=K_\psi^+\), or any fixed real cyclic subfield of \(K_\psi\).
Fix \(A>0\) and choose \(j\) with \(2^j\le(\log V)^A\). Put
\[
 H=\left\lfloor\log_2
        \frac{\log V}{4a_{\rm o}(A+1)\log\log V}\right\rfloor,
 \qquad M=2^{H+3},\qquad W=K C_{\min(j,H+v_2(a))}.
\]
For all sufficiently large \(V\), for every \(x\in[V,V^{3/2}]\),
\(p\ge V\), and \(1\le r<t(x)\) with \(p-1\mid r\), one has
\[
     (e_{U_j}-e_W)R_{p,r}(E_x)=0
                    \qquad\text{if }p\not\equiv\pm1\bmod M .
\]
Here \(e_D\) is the average onto the field \(D\) in the original
local-length test. Uniformly for \(V\le R\le V^{602}\), the
number of pairs \((p,r)\) with \(p\ge V,\ 1\le r<R,\ p-1\mid r\)
in the two omitted progressions is
\[
                  O_{K,A}\!\left(R\,\frac{\log\log V}{\log V}\right).
\]
\end{corollary}
\begin{proof}
Eventually \(H\ge0\). Every \(U_j\) under consideration is in
every \(E_x\), since the available dyadic degree grows as a
positive power of \(V\). For these primes the field averages
exist integrally. The only odd primes dividing \([E_x:\QQ]\)
divide the fixed degree of \(K_\psi\). The congruence condition
in the display is exactly \(h(p)\le H\). In this case
\[
   2^{h(p)}a_{\rm o}\log[A_K(1+2^j)]<\log V
\]
for large \(V\), by the choice of \(H\). The two theorems
give \(\mathcal S_{p,r}(U_j)\) by restriction from
\(P_{p,j}\subseteq W\). Thus its image in
\(\mathcal S_{p,r}(E_x)\) coincides with the image from \(W\).
This equality survives the sign and field projections and
the quotient by bounded-exponent torsion defining \(R_{p,r}\),
because the averages commute and give direct summands.

For the count, write \(N(y)\) for the number of primes at most
\(y\) in the two residue classes. For \(y\ge V\) the two
classes are reduced, \(M\asymp_{K,A}\log V/\log\log V\), and
\(\log(y/M)\ge \tfrac12\log y\) eventually. The
Brun--Titchmarsh inequality gives \(N(y)\ll y/(M\log y)\),
since \(\varphi(M)=M/2\). We use the bound
\(2y/(\varphi(M)\log(y/M))\) per class, with \(y>M\).
\href{https://doi.org/10.1112/S0025579300004708}{H.~L. Montgomery
and R.~C. Vaughan, \emph{The large sieve}, Theorem~2 (1973).}
Each \(p\le R+1\) contributes at most \(R/(p-1)\) weights.
By partial summation,
\[
  \sum_{\substack{V\le p\le R+1\\p\equiv\pm1\bmod M}}\frac1{p-1}
   \ll \frac{N(R+1)}{R+1}
                +\int_V^{R+1}\frac{N(y)}{y^2}\,dy
   \ll \frac1M .
\]
In the last inequality \(\log(R+1)/\log V\) is bounded by
an absolute constant. Multiplying by \(R\) proves the claim.
\end{proof}

\begin{corollary}[The entire real dyadic layer at large coefficient primes]
\leavevmode\label{cor:realhighprime642}
On the same intervals take the actual index \(J=J(x)\).
Whenever \(x^{300}\le p\), \(p\not\equiv\pm1\bmod1024\), and
\(p-1\mid r\) with \(1\le r<t(x)\), one has
\[
       \mathcal S_{p,r}(C_J)
             \simeq \mathcal S_{p,r}(C_{\min(J,h(p))}),
                       \qquad h(p)\le7 .
\]
The corresponding field averages of \(R_{p,r}(E_x)\) agree.
\end{corollary}
\begin{proof}
Use \(K=\QQ\) in the two theorems. The congruence condition
is \(h(p)\le7\), hence \(2^{h(p)}\le128\). The actual
degree bound gives \(2^J\le[K_\psi C_J:\QQ]<2x^{2+8\varepsilon}\).
For \(0<\varepsilon<1/52\) we have
\(128(2+8\varepsilon)<300\), so the prime bound holds uniformly
at sufficiently large \(x\). In the cyclic group of \(C_J\)
the index of the Frobenius subgroup is exactly
\(2^{\min(J,h(p))}\). Its fixed field is the unique subfield
of that degree. Apply restriction and the same commuting
averages as in the preceding corollary.
\end{proof}

These results concern full strict groups at the specified weights.
The count of omitted pairs bounds their number. Their individual
lengths, the surviving classes over the decomposition fields,
and the components involving the two varying quadratic inputs
still enter the weighted sum in Section~\ref{sec:missing}.


\subsection{The relative class group at the half-residue weights}
\label{sec:relativealldigits643}

At a half-residue weight all the odd powers of a chosen Galois
character can be tested in one relative class group. Keeping
the different coefficient Frobenius orbits together determines
their combined torsion at each available precision. In particular
the bound for a relative class number can be used once for all
these orbits.

Let \(\alpha\) be a Dirichlet character of order \(m_0\) and
conductor \(M_0\). Fix \(p\ge5\) prime to \(m_0M_0\), and \(r>0\)
with \(r\equiv(p-1)/2\bmod(p-1)\) and \(\alpha(-1)=(-1)^r\).
Put \(2n=\operatorname{lcm}(2,m_0)\), let \(\eta_p\) be the
quadratic character of conductor \(p\), and set
\[
       \lambda=\alpha^{-1}\eta_p,\qquad
       L=K_\lambda,\quad F=K_{\lambda^2},\quad
       \Delta=\operatorname{Gal}(L/\QQ).
\]
Here \(K_\chi\) is the fixed field of \(\ker\chi\). The conductors
of the two factors of \(\lambda\) are coprime, hence its order
is \(2n\). Its parity is even, so \(L\) is totally real. Write
\(\iota\) for the involution of \(L/F\) and
\[
              C^-=\tfrac12(1-\iota)\operatorname{Cl}(L)\{p\}.
\]
Index odd powers by odd residue classes \(b\bmod 2n\). For one
representative in each orbit of \(b\mapsto pb\), put
\(O_b=\ZZ_p[\lambda^b(\Delta)]\), \(J_b=\operatorname{Frac}(O_b)\),
and \(S_b=S(\alpha^b,r)\), with coefficients in \(J_b/O_b\).
The same coefficient ring is generated by \(\alpha^b\) since
the quadratic factor has values in \(\ZZ_p\). If the remaining
members of an orbit are used, they have the corresponding
conjugate coefficient modules.

\begin{theorem}[All odd powers at one precision]
\leavevmode\label{thm:halfclassdigits643}
For \(1\le z\le1+v_p(r)\), there are isomorphisms of finite
modules
\[
 S_b[p^z]\simeq
  \operatorname{Hom}_{\ZZ_p[\Delta]}
       \bigl(C^-,(J_b/O_b)[p^z](\lambda^b)\bigr).
\]
In particular
\[
          \prod_{b\ {\rm one\ per\ orbit}} |S_b[p^z]|
                    = |C^-[p^z]|,\qquad |C^-|=(h_L/h_F)_p .
\]
The last subscript denotes the \(p\)-part of a positive integer.
If \(N_z\) counts the odd labels \(b\) whose full strict group
has an element of order \(p^z\), counting every character label,
then
\[
        N_z\le
        \log_p\!\frac{|C^-[p^z]|}{|C^-[p^{z-1}]|}
               \le \frac{\log_p(h_L/h_F)}{z}.
\]
\end{theorem}
\begin{proof}
On divisible coefficients the scalar action is
\(\theta_b=\alpha^{-b}\kappa_p^{-r}\). On \(p^z\)-torsion it
equals \(\lambda^b\): the Teichm\"uller factor of the power of
\(\kappa_p\) gives \(\eta_p\), and raising a member of
\(1+p\ZZ_p\) to \(r\) gives one modulo \(p^{1+v_p(r)}\).
This last assertion follows by factoring out the prime-to-\(p\)
part of \(r\) and then applying the binomial theorem successively.
At \(p\) inertia acts nontrivially on the residual character
because \(\eta_p\) is ramified and \(\alpha\) is unramified.
An inertia element therefore acts by a scalar whose difference
from one is a unit. Both the local and the global invariants
of the divisible coefficients vanish. The sequence for
multiplication by \(p^z\) consequently identifies the subgroup
of \(p^z\)-torsion in global \(H^1\) with \(H^1\) on the finite
coefficients, and the same assertion holds locally at \(p\).

The finite coefficients are trivial on \(G_L\), and
\(p\nmid|\Delta|\). Restriction identifies their global \(H^1\)
with the equivariant continuous homomorphisms on \(G_L\), by
averaging in inflation--restriction. A strict class gives a
homomorphism killing the decomposition groups over \(p\).
At a finite prime away from \(p\), inertia intersected with
\(G_L\) acts trivially even on the divisible coefficients:
both cyclotomic characters are unramified there and
\(\alpha^{-1}=\lambda\eta_p\) on this inertia. A finite
cocycle whose divisible class is unramified consequently
vanishes on that intersection. Conversely, an equivariant
homomorphism unramified everywhere and zero at the
decomposition groups over \(p\) gives a strict class. At each
finite place in this converse, the class restricted to inertia
away from \(p\), or to the full decomposition group at \(p\),
restricts to zero over \(L\). The kernel of that restriction
is zero by averaging over a subgroup of \(\Delta\).
Real inertia has order two. Class field theory now gives all
the homomorphisms from the ordinary class group killing the
classes of the primes over \(p\).

The element \(\iota\) belongs to inertia at \(p\), and
\(\Delta\) is abelian. Thus it fixes every such prime, whereas
\(\lambda^b(\iota)=-1\). Every homomorphism with these target
coefficients already kills the prime classes and the entire
plus part. This proves the first isomorphism, in both directions.
On the \(p\)-primary group the plus part identifies with the
class group of \(F\) at \(p\) by extension and norm: their two
composites are multiplication by two and \(1+\iota\),
respectively. The ideal norm on the full class groups is
surjective here. Indeed \(L/F\) is quadratic ramified above
\(p\), so its intersection with the ordinary Hilbert class
field of \(F\) is \(F\); reciprocity on their compositum gives
surjectivity as in Theorem~\ref{thm:relhalf633}. These facts
give the claimed order of \(C^-\).

Since \(p\nmid2n\), the roots of \(X^{2n}-1\) are simple
modulo \(p\). Factoring over \(\mathbb F_p\) and lifting the
coprime factors gives a product of unramified integer rings
for \(\ZZ_p[\Delta]\), with one ring for each coefficient
Frobenius orbit. The minus factor uses precisely the odd
labels. In the factor for \(b\) the target in the first
display is \(O_b/p^zO_b\) under the corresponding action.
On a cyclic summand \(O_b/p^uO_b\) of the class module the
homomorphisms to this target number
\(|O_b/pO_b|^{\min(u,z)}\). This is the order of the
\(p^z\)-torsion of the same summand. Multiplying gives the
product identity.

The orbit size equals \([O_b:\ZZ_p]\), the residue degree.
Passing from \(z-1\) to \(z\) adds exactly that many to the
base-\(p\) logarithm for each cyclic summand of length at
least \(z\). An orbit carrying an element of order \(p^z\)
therefore contributes at least its number of labels to the
increase. The full logarithmic order of each such class
summand is at least \(z\) times its contribution to the
increase. This proves both inequalities.
\end{proof}

\begin{corollary}[A bound uniform in the Frobenius order]
\leavevmode\label{cor:allhalfprimes643}
Fix a cyclic abelian field \(K\) and use characters of \(K C_j\).
There is \(c_K>0\) such that whenever \(V\) is sufficiently
large, \(2^j\le c_K(\log V)^2\), and \(p\ge V\), each family
above with \(n\ge2\) satisfies
\[
              N_z<\frac{n}{2z}
                        \qquad(1\le z\le1+v_p(r)).
\]
For \(n=1\) the full strict group is zero. In particular, at
least half the odd labels in every such family have zero full
strict group at every fixed half-residue weight of the required
parity. If \([K:\QQ]\) is a power of two, the bound by one
half also holds for the proportion of surviving labels among
all characters of \(K C_j\) of the required parity, whenever
that set is nonempty.
\end{corollary}
\begin{proof}
Let \(a=[K:\QQ]\). As in
Theorem~\ref{thm:variableclasses642}, write \(\alpha=\gamma\beta\)
with \(\gamma\) from \(K\) and \(\beta\) dyadic. The order of
\(\beta\) divides the two-part of \(\operatorname{lcm}(a,2n)\).
The conductor of a nontrivial real dyadic character of order
\(2^i\) is \(2^{i+2}\). Thus \(M_0\le A_K n\);
the primes dividing \(\operatorname{cond}(\alpha^2)\) belong
to a fixed finite set, and \(n\le a2^j\). For \(n\ge2\) the
quantity \(H(\operatorname{cond}(\alpha^2),n)\) of
Lemma~\ref{lem:squarefield633} is bounded in terms of \(K\).
The regulator and relative class number calculation in the
proof of Theorem~\ref{thm:relhalf633}, before the use of any
one residual component, therefore gives
\[
   \frac{(h_L/h_F)^{1/n}}{\sqrt p}
        \le C_K\,\frac{\sqrt n\,(\log(3A_K n p))^{1/n}}
                              {\log(p/4)} .
\]
Put \(y=\log\log(3p)\). For \(2\le n\le y\) the right side
is \(O_K(\sqrt{y/\log p})\). For \(y<n\le a c_K(\log p)^2\),
the logarithmic factor with exponent \(1/n\) is bounded by
four for large \(p\). Choosing \(c_K\) small enough makes
the displayed right side strictly less than one throughout
both ranges. Thus \(h_L/h_F<p^{n/2}\), and the theorem applies.
For \(n=1\), the conductor of \(\alpha\) is bounded in terms
of \(K\). The quadratic field \(L\) has discriminant \(O_K(p)\)
and class number \(O_K(\sqrt p\log(2p))<p\), by the ideal
count in the proof of Corollary~\ref{cor:mixedhalf633}.
The theorem at \(z=1\) gives zero \(p\)-torsion, hence zero
strict group. For the last assertion, partition the character
group into the generators of its cyclic subgroups. In a
two-group the generators of a nontrivial cyclic subgroup
are exactly the odd powers of any generator, and their parities
are equal. Apply the result to each eligible subgroup, counting
the trivial character separately.
\end{proof}

\begin{corollary}[Full exponent control on the mixed subfield]
\leavevmode\label{cor:mixedhalfexponent643}
In the original fields and intervals take \(K=K_\psi\), including
its odd-degree part. Let \(a_{\rm o}\) be the odd part of
\([K:\QQ]\), \(2^j\le c_K(\log V)^2\), \(U=KC_j\), and
\(h(p)=v_2(p^2-1)-3\). At \(p\ge V\) put \(z_p=2^{h(p)}a_{\rm o}\).
For every \(x\) on the interval and \(1\le r<t(x)\) satisfying
\[
       r\equiv (p-1)/2\bmod(p-1),\qquad
                         1+v_p(r)\ge z_p,
\]
the full complementary-sign strict group on \(U\) is killed
by \(p^{z_p-1}\). In the notation of Section~\ref{sec:missing},
\[
                         e_U R_{p,r}(E_x)=0 .
\]
\end{corollary}
\begin{proof}
Every eligible \(\alpha\) has \(\lambda\) of order \(2n=2^v w\)
with \(w\mid a_{\rm o}\). The coefficient orbit of \(b=1\)
has size \(\operatorname{ord}_{2n}(p)\). The dyadic order
calculation in Theorem~\ref{thm:variableclasses642} gives
\[
       \operatorname{ord}_{2n}(p)\ge
            \frac{2n}{4\,2^{h(p)}w}\ge\frac{n}{2z_p}.
\]
If \(n\ge2\), the entire orbit would contribute to \(N_{z_p}\)
if its strict group had an element of order \(p^{z_p}\).
The strict inequality in the preceding corollary excludes
this. The case \(n=1\) has already vanished. Every element of
order larger than \(p^{z_p-1}\) would have a multiple of order
exactly \(p^{z_p}\). Integral character decomposition at
\(p\nmid[U:\QQ]\) therefore bounds the exponent on \(U\)
by \(p^{z_p-1}\), using Theorem~\ref{thm:levelstrict345}.

The dyadic degree allowed in every \(E_x\) grows as a positive
power of \(V\), so \(U\) is a subfield throughout this same
interval. Field averaging identifies its contribution upstairs
by prime-to-\(p\) descent. A half-residue weight in the
given range has \(p\le 2t(x)+1<P(x)\) for sufficiently large
\(x\). Also \(z_p-1\le v_p(r)\le\log_2 t(x)<b_0(x)\).
The quotient removing the \(p^{b_0(x)}\)-torsion therefore
kills this entire summand, including after the other commuting
field projections.
\end{proof}

For example, at \(h(p)=1\) and \(a_{\rm o}=1\), divisibility
of a half-residue weight by \(p\) suffices on this mixed field.
At weights with fewer available digits the relative class group
still determines their initial torsion. A uniform bound on the
length beyond this precision and the estimate for the full
fields in Section~\ref{sec:missing} remain open.


\subsection{A single class number for the half-residue weights of the interval}
\label{sec:jointweights644}

Fix one of the intervals of Theorem~\ref{thm:conductor232}, with
largest field \(E=E_*\), cyclotomic modulus \(N=N_*\), degree
\(n_*=[E:\QQ]\), and \(R=t_*=\lfloor V^{1203/2}\rfloor\).
Fix \(p\ge5\) with \(p\nmid n_*N\), and put \(s=(p-1)/2\).
Only \(p\le2R+1\) can have a weight \(1\le r<R\) congruent to
\(s\bmod(p-1)\). Let \(\eta_p\) be the quadratic character of
conductor \(p\), let \(D_p\) be its quadratic field, and put
\[
                H=(E D_p)^+,\qquad H_0=E^+ .
\]
The superscript denotes the maximal real subfield. Write \(c\)
for complex conjugation. We consider the individual characters
\(\alpha\) of \(E\) with \(\alpha(c)=(-1)^s\). If this list is
empty, all the sums below are empty.

\begin{lemma}[The class group for all the characters of the given parity]
\leavevmode\label{lem:onehalfclass644}
For a nonempty list, \(H/H_0\) is quadratic. Let \(\iota\)
be its nonidentity automorphism. Put
\[
 C=\tfrac12(1-\iota)\Cl(H)\{p\},\qquad B_p=\log_p|C|.
\]
For an orbit of \(\alpha\mapsto\alpha^p\) on the list, put
\(\tau=\alpha^{-1}\eta_p\), \(O_\tau=\ZZ_p[\tau]\),
and \(f_\tau=[O_\tau:\ZZ_p]\). Let \(h_\tau\) be the
\(O_\tau\)-length of the corresponding factor of \(C\).
Then
\[
   \sum_{\rm orbits} f_\tau h_\tau=B_p
              =v_p(h_H/h_{H_0}).
\]
The compact module for \(S(\alpha,r)\) in
Lemma~\ref{lem:primitiveclass502} has constant specialization
of length \(h_\tau\).
\end{lemma}
\begin{proof}
The prime \(p\) is unramified in \(E\) and ramified in \(D_p\).
Thus \(E\cap D_p=\QQ\), and the Galois group of \(E D_p\) is
the direct product. Its element acting only on \(D_p\) descends
to \(\iota\) in \(H\). The characters on the list give exactly
the characters of the product trivial on \(c\) and nontrivial
on this element, by writing them as \(\alpha^{-1}\eta_p\).
Nonemptiness shows that \(\iota\) is nontrivial. Its fixed field
in \(H\) is \(E\cap H=H_0\).

The order of \(\operatorname{Gal}(H/\QQ)\) is prime to \(p\).
Its integral group algebra over \(\ZZ_p\) splits into unramified
integer rings indexed by Frobenius orbits of characters.
Indeed the exponent is prime to \(p\), so the roots of unity
in a splitting ring remain distinct after reduction. The
character projectors over that ring are integral, and their
Frobenius orbits descend to the indicated factors. A factor
of length \(h_\tau\) therefore has order \(p^{f_\tau h_\tau}\).
Extension and norm for \(H/H_0\) identify the plus part of
\(\Cl(H)\{p\}\) with \(\Cl(H_0)\{p\}\), since their composites
are \(2\) and \(1+\iota\) and \(2\) is a unit at \(p\).
This gives the displayed equality.

For the constant specialization of a character let \(L=K_\tau\).
The calculation at \(T=0\) in Lemma~\ref{lem:classdet466}
identifies its length with the \(\tau\)-factor of
\(\Cl(L)\{p\}\) modulo the classes at \(p\). Here inertia
contains the image of \(\iota\), which fixes every prime above
\(p\) and acts by \(-1\) on the factor. The prime classes
consequently make no contribution. For \(H/L\), extension
and norm identify \(\Cl(L)\{p\}\) with the invariants upstairs:
the degree is prime to \(p\), and the group sum divided by the
degree is the projector onto invariants. The \(\tau\)-factor
upstairs already belongs to these invariants. This proves the
constant specialization assertion.
\end{proof}

For the following elementary estimate, normalize each extension
of \(v_p\) by \(v_p(p)=1\). Put
\[
 W=1+\lfloor\log_p R\rfloor,\qquad
 d_r=1+v_p(r),\qquad
 H_\tau(r)=h_\tau+(W-d_r)\lfloor h_\tau/d_r\rfloor .
\]
Write \(k_\alpha(r)=\operatorname{length}_{O_\tau}S(\alpha,r)\).
In an inequality this length can be read as \(+\infty\) if
the specialized module has positive rank.

\begin{theorem}[A weighted bound over characters and weights]
\leavevmode\label{thm:allhalfweights644}
Count every character label, including the distinct conjugates
over \(\ZZ_p\). For every \(R\ge2\) as above,
\[
 \sum_{\alpha}\ 
 \sum_{\substack{1\le r<R,\ r\equiv s\ ({\rm mod}\ p-1)\\
                   k_\alpha(r)>H_\tau(r)}}
                   (1+v_p(r))\ \le\ B_p .
\]
The same inequality holds for any subset of the characters and
weights. In particular one need only count characters of \(E\)
once when passing through the nested fields \(E_x\).
\end{theorem}
\begin{proof}
Fix one orbit. Its compact module \(B\) is finite type over
\(O_\tau[[T]]\); its specialization at \(t_r=(1+p)^{-r}-1\)
is dual to the strict group by Lemma~\ref{lem:primitiveclass502}.
For \(h=h_\tau=0\) the module vanishes by compact Nakayama.
For \(h>0\) there is a relation determinant \(A(T)\) with
\[
 v_p(A(0))=h,\qquad k_\alpha(r)\le v_p(A(t_r)).
\]
Indeed take a finite free cover of \(B\). At \(T=0\) the
images of its relations form a full rank lattice of colength
\(h\). Lifting a basis of this lattice to relations gives a
square matrix. The cokernel of its specialization surjects
onto the specialized module, which proves both claims.

Weierstrass preparation writes \(A\) as \(p^\mu\) times a
unit times a distinguished polynomial. In a splitting field
write its roots, with multiplicity, as \(\beta\).
Each has \(v_p(\beta)>0\), none is zero, and
\(\mu+\sum_\beta v_p(\beta)=h\). For \(u,v\) distinct positive
integers the binomial theorem gives
\[
 v_p(t_u)=1+v_p(u),\qquad
 v_p(t_u-t_v)=1+v_p(u-v).
\]
In fact taking a power prime to \(p\) preserves the valuation
of a principal unit minus one, and taking its \(p\)-th power
increases that valuation by one. This also applies to inverses.
In particular the distinct \(t_r\) under consideration have
mutual differences of valuation at most \(W\).

Only roots of valuation \(d_r\) can make
\(v_p(t_r-\beta)>v_p(\beta)\). There are at most
\(\lfloor h/d_r\rfloor\) such roots. If every distance valuation
is at most \(W\), their total increase above their contribution
at zero is at most \((W-d_r)\lfloor h/d_r\rfloor\); all other
roots have contribution at most their value at zero. It follows
that a weight with \(k_\alpha(r)>H_\tau(r)\) has a root with
\(v_p(t_r-\beta)>W\), necessarily of valuation \(d_r\).
One root cannot serve two such weights, by the bound for
\(v_p(t_r-t_{r'})\) and the ultrametric inequality. Assigning
one root to each weight gives \(\sum d_r\le\sum_\beta v_p(\beta)
\le h\) within this orbit. Conjugation preserves the lengths.
Multiplication by the orbit size \(f_\tau\) and summation give
the theorem by Lemma~\ref{lem:onehalfclass644}.

Finally a character inflated from \(E_x\) defines the same
scalar Galois module over \(\QQ\) and the same local conditions
as the original character. Thus its exceptional weights are
already counted at the largest field.
\end{proof}

\begin{corollary}[Application on the full period fields]
\leavevmode\label{cor:periodfieldweights644}
On the polynomial-conductor intervals, for all good primes
\(5\le p\le 2R+1\),
\[
 B_p\ll n_*\bigl(1+\log N_*/\log p\bigr).
\]
For \(p\ge V\) the weighted number of exceptional pairs in
Theorem~\ref{thm:allhalfweights644} is
\(O_{\psi,\varepsilon}(n_*)\) per prime, including characters
involving either moving quadratic field. For every
\(x\in[V,V^{3/2}]\), the total logarithmic order of the
nonexceptional half-residue components of these good primes
at \(1\le r<t(x)\) is \(O_{\psi,\varepsilon}(V^A)\) for a
fixed \(A\). Each such component vanishes after quotienting
by the subgroup killed by \(p^{b_0(x)}\), for sufficiently
large \(V\).
\end{corollary}
\begin{proof}
The absolute discriminant of \(H\) is at most
\((pN_*)^{[H:\QQ]}\) by the conductor-discriminant formula.
Minkowski's bound supplies an integral ideal of norm at most
\(\sqrt{|\operatorname{disc}(H)|}\) in each class, since \(H\)
is totally real. The number of ideals of norm at most \(Y\ge1\)
is at most \(Y^2\zeta_H(2)\), and the Euler factors give
\(\zeta_H(2)\le\zeta(2)^{[H:\QQ]}\).
Use \(B_p\le\log_p h_H\) and \([H:\QQ]\le2n_*\).
The bound for \(p\ge V\) follows since \(N_*\) is polynomial
in \(V\) on these intervals.

At a nonexceptional weight the length in a single orbit is
at most \(Wh_\tau\). Even if different orbits are exceptional
at different weights, the sum of \(f_\tau\) times the
nonexceptional lengths for a fixed \(p,r\) is at most \(WB_p\).
The logarithmic order further multiplies this by \(\log p\).
There are at most \(R\) weights and \(2R+1\) possible primes,
and \(W\le1+\log_5 R\). All these quantities, \(n_*\), and
\(\log N_*\) are polynomially bounded in \(V\). Prime-to-\(p\)
descent to the scalar factors on \(E_x\) uses the same exact
group algebra decomposition as in Lemma~\ref{lem:onehalfclass644},
or the descent in Theorem~\ref{thm:halfclassdigits643}.
Finally the exponent of a finite module of length \(k\) over
an unramified integer ring divides \(p^k\). Uniformly in \(p\),
\(WB_p\) is polynomial in \(V\), whereas
\(b_0(x)=\lfloor\exp(x^\varepsilon/8)\rfloor\) exceeds every
fixed power of \(V\) eventually. This proves the vanishing.
\end{proof}

For the odd powers of one \(\alpha\) on \(K C_j\), the same root
argument uses \(v_p(h_L/h_F)\) in Theorem~\ref{thm:halfclassdigits643}.
Thus under the hypotheses of Corollary~\ref{cor:allhalfprimes643}
with \(n\ge2\) the sum of \(1+v_p(r)\) over all its exceptional
label and weight pairs is strictly less than \(n/2\), even when
all half-residue weights in the interval are included at once.
For \(n=1\) the same corollary gives vanishing. On the full
period fields the bound above controls the number and the
weight divisibility of the exceptional pairs. Condition
\textup{(L)} also requires an upper bound for the lengths
at the exceptional pairs themselves.


\subsection{Quantitative unit indices at split coefficient primes}
\label{sec:gaudronsplit646}

The zero-residue weights involve an index of local logarithms
when the coefficient prime splits completely in the character
field. We give a bound for this index in terms of the conductor,
the degree and the number of character components being tested.
The bound also controls the invariant ordinary ideal classes
at every level of the cyclotomic tower of that field.

Let \(L\) be a totally real cyclic field of degree \(d\ge2\)
and conductor \(M\). Let \(p\ge5\) be prime to \(dM\) and split
completely in \(L\). For a faithful character \(\tau\) of
\(\Delta=\operatorname{Gal}(L/\QQ)\), write
\(O_\tau=\ZZ_p[\tau]\). Recall from
Theorem~\ref{thm:splitamb506} that \(p^{c_\tau}O_\tau\) is
the image of the completed ordinary units in the
\(\tau\)-component of the local cyclotomic exponent lattice.
This lattice is identified with \(\ZZ_p[\Delta]\) by numbering
the places over \(p\). List all \(\varphi(d)\) faithful
characters at one embedding of the character values into
\(\overline{\QQ}_p\), including individual Frobenius conjugates,
and order their \(c_\tau\) in decreasing order, denoting the
result by \(c^{(1)},c^{(2)},\ldots\). Set
\[
             A(d,p,M)=d^2\log(2d)+d\log(2pM).
\]

\begin{theorem}[Simultaneous bounds for the split unit indices]
\leavevmode\label{thm:quantlog646}
There is an absolute constant \(C\) such that for
\(1\le m\le\varphi(d)\),
\[
 \log(2+c^{(m)})\le C\left(
       \frac{A(d,p,M)}{m}+\log(2d)+\log\log(3pM)\right).
\]
In particular for a character of coefficient residue degree
\(f_\tau\) the bound with \(m=f_\tau\) holds for \(c_\tau\).
\end{theorem}
\begin{proof}
Take a generator \(g\) of \(\Delta\), fix a place \(v_0\) over
\(p\), and label the places by \(v_j=g^jv_0\), \(0\le j<d\).
Put
\[
 w=N_{\QQ(\zeta_M)/L}(1-\zeta_M),\quad
 u=g(w)/w,\qquad
 z_j=\log_p\big((g^{-j}u)^{p(p-1)}\big)
\]
in the embedding at \(v_0\). The unit calculation in
Lemma~\ref{lem:explicitunits523} gives
\(h(u)\le2M\log(2M)\) for the absolute logarithmic Weil
height, and the logarithmic archimedean vector of \(u\)
has nonzero coefficient at every faithful character. Indeed
the norm on \(w\) unfolds the coefficient to the Gauss sum
times the nonzero Dirichlet value \(L(1,\chi)\) of the
primitive even character. The factor \(g-1\) multiplies it
by a nonzero scalar. At \(v_0\) the conjugates of \(u\)
are in \(\ZZ_p^\times\). Their \(p(p-1)\)-th powers belong
to \(1+p^2\ZZ_p\), so the \(z_j\) have absolute value at
most \(p^{-2}\), and their exponentials are the indicated
algebraic units.

Let \(T_z\subset\QQ^d\) be the smallest rational subspace
whose scalar extension contains the vector \(z=(z_j)\).
We need the rank, on this subspace, of the Fourier forms
\(\ell_\tau(Z)=\sum_{j=0}^{d-1}\tau(g)^jZ_j\).
For an integral vector \((a_j)\), the relation
\(\sum a_jz_j=0\) is equivalent to
\(\prod_j(g^{-j}u)^{p(p-1)a_j}=1\). This follows by the
logarithm and exponential on \(1+p^2\ZZ_p\) and injectivity
of the field embedding. The relations are therefore stable
under cyclic shifts. They all vanish also on the vector of
real numbers \((\log|g^{-j}u|)_j\), using any fixed real
embedding. Clearing denominators treats rational relations.
Thus \(T_z\otimes\mathbb C\) contains that real vector and
its cyclic shifts. Projection in the complex group algebra
onto each faithful eigenline gives a nonzero vector by the
archimedean calculation. Every faithful eigenline is therefore
contained in the complexification. The same inclusion holds
over \(k=L(\zeta_d)\): these eigenlines are defined over \(k\)
and membership in \(T_z\) is tested by rational linear
equations. Distinct Fourier forms pair with distinct such
lines. In particular any \(m\) faithful forms restrict to
rank \(m\) on \(T_z\otimes k\).

Apply Gaudron's simultaneous \(p\)-adic logarithm bound over
\(k\) to these \(m\) forms and the \(d\) logarithms \(z_j\).
We spell out the parameter dependence because both the field
and \(p\) vary. The degree \(D=[k:\QQ]\) is at most \(d^2\),
the coefficients of the forms have height zero, and the
heights of the exponentials are at most
\(H=2p^2M\log(2M)\). The rank parameter in the theorem is
the codimension on \(T_z\), hence is \(m\). Its real
radius parameter may be taken equal to \(\exp(1)\), since
\(p^{-2}\le \exp(-1)p^{-2/(p-1)}\) for \(p\ge5\).
Take each logarithmic height parameter equal to \(H\) and
the logarithmic coefficient parameter equal to \(D\).
The auxiliary integer in the bound is then at most
\(a=2+\lceil D\log(\exp(1)+D+dH)\rceil\).
For any set \(\mathcal T\) of \(m\) distinct faithful labels
the theorem gives, with an absolute constant,
\[
 -\log\max_{\tau\in\mathcal T}|\ell_\tau(z)|_p
 \ \ll\ (10d)^{59d^2/m}\,a^{1/m}
       (D+a+D\log(2D))(1+DH)^{d/m}.
\]
Here the product in the logarithm theorem runs over a
maximal rationally independent subfamily, of size at most
\(d\). This application does not require multiplicative
independence of all the conjugates. The published estimate
and its parameter conventions are in
\href{https://smf.emath.fr/system/files/2017-08/smf_bull_142_1-62.pdf}{%
Gaudron, \emph{Bulletin de la Soci\'et\'e Math\'ematique de France}
142 (2014), Theorem~2.1, p.~5}.

For \(u\) itself the image in the local exponent lattice at
the \(\tau\)-factor is the evaluation of its logarithm vector
at \(\tau(g)\), namely
\(\ell_\tau(z)/(p(p-1)\log_p(1+p))\).
The denominator has valuation two under \(v_p(p)=1\).
The image ideal of all ordinary units contains this element,
so \(v_p(\ell_\tau(z))\ge2+c_\tau\). For a singleton the
rank computation and Gaudron's theorem also give nonvanishing.
Taking the \(m\) indices of largest \(c_\tau\) lower bounds the
left side of the preceding display by \((2+c^{(m)})\log p\).
Take ordinary logarithms. The power of \(10d\) contributes
\(O(d^2\log(2d)/m)\); the product contributes
\(O(d\log(2dpM)/m)\). The remaining factors contribute
\(O(\log(2d)+\log\log(3pM))\) by the bound on \(a\).
Absorbing \(d\log(2d)\) into \(d^2\log(2d)\) proves the
claim. Conjugating coefficients preserves the index and the
orbit has \(f_\tau\) members, which gives the last assertion.
\end{proof}

\begin{corollary}[Uniformity in the cyclotomic layer at the split primes]
\leavevmode\label{cor:splitlayers646}
Use the intervals and full fields \(E_x\subseteq E_*\) of
Theorem~\ref{thm:conductor232}. For each cyclic \(L\subset E_*^+\)
as above and each good split prime \(p\le t_*+1\), put
\[
 L_i=L\QQ_i,\quad G_i=\operatorname{Gal}(L_i/L),\qquad
 h_\tau^0=\operatorname{length}_{O_\tau}\Cl(L)\{p\}_\tau ,
\]
where \(\QQ_i\) is the degree \(p^i\) cyclotomic layer.
Except for at most
\[
         O_{\psi,\varepsilon}\!
         \left(\frac{d^2\log(2d)+d\log(2V)}{V^\varepsilon}\right)
\]
faithful character labels for this \(L,p\), both
\[
 \sup_{i\ge1}\operatorname{length}_{O_\tau}
                  \big((\Cl(L_i)\{p\}_\tau)^{G_i}\big)
 \quad\text{and}\quad
 \operatorname{length}_{O_\tau}(B/TB)
\]
are at most \(h_\tau^0+\exp(V^\varepsilon/16)<b_0(x)\)
throughout the interval for sufficiently large \(V\).
Here \(B\) is the compact module for the zero-residue weights
of \(\alpha=\tau^{-1}\). For every fixed \(0<\delta<1\) the
exceptional fraction tends uniformly to zero for
\(2\le d\le V^{(1-\delta)\varepsilon}\).
There are no exceptional labels for \(d\le V^{\varepsilon/3}\)
once \(V\) is sufficiently large.
\end{corollary}
\begin{proof}
The hypotheses give \(p,M,n_*\le V^{C_{\psi,\varepsilon}}\).
If there are \(m>0\) indices larger than
\(\exp(V^\varepsilon/16)\), the theorem bounds \(m\) as
stated: its additive logarithmic terms are \(O_{\psi,\varepsilon}
(\log V)=o(V^\varepsilon)\). This argument includes primes
dividing class numbers or indices of circular units.
Theorem~\ref{thm:splitamb506} gives the first length at level
\(i\) as \(h_\tau^0+\min(i,c_\tau)\).
Theorem~\ref{thm:sunitbase507} gives the second as an
\(S\)-class length plus the local index for \(S\)-units,
where \(S\) consists of the primes over \(p\). Ordinary units
belong to the \(S\)-units and the \(S\)-class group is a
quotient of the ordinary class group. Both displayed
quantities are consequently bounded by \(h_\tau^0+c_\tau\).
Minkowski's ideal bound and the conductor-discriminant
formula, as in Corollary~\ref{cor:periodfieldweights644}, give
\(h_\tau^0\le\log_p h_L\ll d\log(2M)/\log p\), a polynomial
bound in \(V\). This proves the comparison with \(b_0(x)\).

The degree \(d\) has odd part dividing the fixed order of
\(\psi\), since \(E_*\) is obtained from \(K_\psi\) by
extensions of degree a power of two. Thus
\(\varphi(d)\gg_\psi d\). Dividing the bound on the number
of exceptions by \(\varphi(d)\) gives a quantity tending
uniformly to zero in the stated degree range. In the
smaller range the bound for the number itself tends to
zero. Subfields occurring at any \(x\) are already
included at the top field.
\end{proof}

At a positive zero-residue weight the actual strict group is
dual to \(B/(T-((1+p)^{-r}-1))B\). The corollary bounds the
constant specialization and the invariant ordinary classes at
all levels. Controlling the exceptional positive-weight
specializations in Section~\ref{sec:missing} still requires
information on the cyclotomic action.


\subsection{One class group quotient for the sum at the half-residue weights}
\label{sec:productweights647}

We retain the lengths at the exceptional weights of
Theorem~\ref{thm:allhalfweights644}. Their sum, including
every digit through an arbitrarily prescribed precision,
can be measured by imposing one polynomial relation on one
finite ordinary class group. The additive error below is
independent of that precision.

Use \(E=E_*,D_p,H=(ED_p)^+\), \(H_0=E^+\), \(\iota\), and
\(B_p\) as in Lemma~\ref{lem:onehalfclass644}. We consider
only primes with a nonempty list in that lemma. Let
\(\QQ_j\) be the cyclotomic layer of degree \(p^j\), set
\(H_j=H\QQ_j\) for \(j\ge1\), and write \(H\) itself for
the base layer. In particular \(H_0=E^+\) continues to
mean the fixed field of \(\iota\).
Let \(\Delta=\operatorname{Gal}(H/\QQ)\), and choose
\(\gamma\) in the cyclotomic direction with principal
cyclotomic value \(1+p\). Put
\[
 C(j)=\tfrac12(1-\iota)\Cl(H\QQ_j)\{p\}\quad(j\ge0),
 \qquad \QQ_0=\QQ,\qquad u_r=(1+p)^{-r}.
\]
Then \(\log_p|C(0)|=B_p\). Fix a union of coefficient
Frobenius orbits of the allowed characters \(\alpha\).
Let \(e\in\ZZ_p[\Delta]\) be the sum of the idempotents
for their characters \(\tau=\alpha^{-1}\eta_p\).

\begin{lemma}[Finite precision in the ordinary class group]
\leavevmode\label{lem:normmeasure647}
There is a norm isomorphism
\[
           C(j)/(\gamma-1)C(j)\simeq C(0).
\]
Let \(r\equiv(p-1)/2\bmod(p-1)\) be positive, \(k\ge1\),
and \(j\ge\max(0,k-1-v_p(r))\). Then
\[
 \sum_{\alpha\ {\rm selected,\ one\ per\ orbit}}
      \log_p |S(\alpha,r)[p^k]|
   =\log_p\left|eC(j)/(p^k,\gamma-u_r)eC(j)\right|.
\]
\end{lemma}
\begin{proof}
The finite Galois group of \(H\QQ_j/\QQ\) is
\(\Delta\times\langle\gamma\rangle\), the second factor
of order \(p^j\), because \(p\nmid[H:\QQ]\).
The indicated lower bound on \(j\) makes the principal part of
the coefficient character trivial modulo \(p^k\) on the
absolute Galois group of \(H\QQ_j\), by
\(v_p(u_r^{p^j}-1)=1+v_p(r)+j\).
On the finite quotient the actions are \(\tau\) and \(u_r\).
The central element \(\iota\) acts by \(-1\) and lies in the
inertia image at every prime over \(p\): its tame image was
already present in \(H\), and a subgroup of a product of groups
of coprime orders contains the pure factors of each of its
elements. Thus global invariants of the divisible coefficient
and invariants at \(p\) vanish as in
Theorem~\ref{thm:halfclassdigits643}. One can use the lift
of tame inertia whose principal cyclotomic part is one, or
just use its action modulo \(p\), whose difference from one is
a unit. Multiplication by \(p^k\) identifies \(p^k\)-torsion
in the global and \(p\)-local first cohomology with first
cohomology on the finite coefficients.

The finite quotient and the local decomposition quotient at
\(p\) have vanishing positive-degree cohomology on these
coefficients. Indeed first take cohomology of the central
subgroup \(\langle\iota\rangle\): averaging kills positive
degrees and its invariants vanish. Inflation--restriction
therefore gives the equivariant homomorphisms on the upper
field and detects the zero local class. Away from \(p\) the
inertia quotient is of order prime to \(p\). Over the upper
field, inertia away from \(p\) acts trivially also on the
divisible coefficients, since the finite character there
is trivial on inertia and the cyclotomic part is unramified.
Consequently the unramified condition descends and ascends
by restriction, just as in Theorem~\ref{thm:halfclassdigits643}.
The real places contribute no odd-primary inertia.
Reciprocity over \(H\QQ_j\) now computes the homomorphisms
from its ordinary class group that kill the classes above
\(p\). These prime classes are fixed by \(\iota\) and already
map to zero in every selected factor. In the \(\tau\)-factor,
equivariance imposes \(\gamma=u_r\) and the coefficient
precision imposes \(p^k=0\). Finite duality preserves the
order of each factor over its unramified integer ring.
Multiplying over the selected orbits proves the second formula.

For the norm assertion take the base of cohomology to be
\(E^+\) and the divisible coefficient \(\QQ_p/\ZZ_p\) twisted
only by the character of \(H/E^+\). On
\(\operatorname{Gal}(H\QQ_j/E^+)=
\langle\iota\rangle\times\langle\gamma\rangle\), the two
actions are now \(-1\) and \(1\). The very same global and
local averaging arguments identify the dual of the strict
group with \(C(j)/(\gamma-1)C(j)\). Computing over \(H\)
instead gives \(C(0)\). Restriction of the homomorphisms
corresponds to ideal norm by functoriality of reciprocity.
Dualizing therefore gives the indicated norm isomorphism,
which is compatible with the \(\Delta\)-action.
\end{proof}

\begin{theorem}[The product of the weight relations]
\leavevmode\label{thm:productmeasure647}
Let \(\mathcal R\subset[1,R)\) be any finite set of positive
half-residue weights and put
\[
 P_{\mathcal R}(Y)=\prod_{r\in\mathcal R}(Y-u_r),\qquad
 D(\mathcal R)=\sum_{\substack{r,s\in\mathcal R\\r<s}}
                                      (1+v_p(s-r)).
\]
For \(1\le k\le m\) use one layer \(j=m-1\) and define
\[
 q_k=\log_p\left|eC(j)/(p^k,P_{\mathcal R}(\gamma))eC(j)\right|,
 \qquad
 a_k=\sum_{r\in\mathcal R}\ 
       \sum_{\alpha\ {\rm selected,\ one\ per\ orbit}}
                        \log_p |S(\alpha,r)[p^k]|.
\]
Then
\[
          0\le a_k-q_k\le \log_p|eC(0)|\,D(\mathcal R).
\]
In particular the two differences \(a_m-a_k\) and
\(q_m-q_k\) differ in absolute value by at most this same
bound, independent of \(m\). For the empty set the quotient
by the empty product is zero and the logarithms of orders
are zero.
\end{theorem}
\begin{proof}
Set \(A=eC(j)/p^k eC(j)\) and enumerate the weights as
\(r_1,\ldots,r_n\). Write \(D_i=\gamma-u_{r_i}\) acting
on \(A\), \(Q_0=1\), and \(Q_i=D_iQ_{i-1}\).
By the lemma, \(a_k=\sum_i\log_p|\ker D_i|\), since an
endomorphism of a finite group has equal kernel and cokernel
orders. The order of \(Q_{i-1}A/Q_i A\) is
\(|\ker D_i\cap Q_{i-1}A|\); this is the kernel-cokernel
equality for the restriction of \(D_i\) to \(Q_{i-1}A\).
These latter logarithms sum to \(q_k\). Their termwise
deficits are nonnegative.

On \(\ker D_i\) the operator \(Q_{i-1}\) acts by the
nonzero scalar \(\prod_{h<i}(u_{r_i}-u_{r_h})\), of
valuation \(b_i=\sum_{h<i}(1+v_p(r_i-r_h))\). Therefore
\(\ker D_i\cap Q_{i-1}A\) contains \(p^{b_i}\ker D_i\).
The deficit at \(i\) is at most \(b_i\) times the number
of cyclic generators of the finite abelian \(p\)-group
\(\ker D_i\). All \(u_r\) are one modulo \(p\); the subgroup
killed by \(p\) in \(\ker D_i\) is precisely
\(A[p]^{\gamma=1}\). Its dimension is at most
\(\log_p|\ker(\gamma-1:A\to A)|\), which equals
\(\log_p|A/(\gamma-1)A|\). The latter is at most
\(\log_p|eC(0)|\) by the norm assertion and right exactness
of quotienting. Summing over \(i\) gives the inequality.
Subtracting two numbers in the same interval
\([0,\log_p|eC(0)|D(\mathcal R)]\) proves the assertion
about differences.
\end{proof}

\begin{corollary}[A finite test for the entire half-residue tail]
\leavevmode\label{cor:halftail647}
At \(x\in[V,V^{3/2}]\) let \(H(x)\) be the contribution to
\(\mathcal R(x)\) of Proposition~\ref{prop:res23} from good
\(p\ge5\) and half-residue weights. Let \(X>0\) be any
proposed bound in natural logarithmic order, \(b=b_0(x)\),
and, for each contributing prime, take
\[
              m_p=b+\lfloor X/\log p\rfloor+1 .
\]
At fixed \(p\), split the half-residue weights \(1\le r<t(x)\)
into those with \(r\le x^\eta\) and those with \(r>x^\eta\).
On each set take \(e\) to select precisely the characters
inflated from \(E_x\) surviving \(1-J_{p,r}\) in
Section~\ref{sec:missing}. Use the theorem on each such set,
at \(m=m_p\), and add \((q_{m_p}-q_b)\log p\) over the sets
and primes, obtaining \(Q(x,X)\). Define
\[
 \mathcal D(x)=\sum_{p,\,\mathcal R}
       (\log p)\log_p|eC(0)|\,D(\mathcal R).
\]
There is \(A=A_{\psi,\varepsilon}\) with
\(\mathcal D(x)\ll_{\psi,\varepsilon}V^A\) uniformly on
the interval and independently of \(X\). One has both
implications
\[
 \begin{split}
    Q(x,X)+\mathcal D(x)\le X &\ \Longrightarrow\ H(x)\le X,\\
    H(x)\le X &\ \Longrightarrow\ Q(x,X)\le X+\mathcal D(x).
 \end{split}
\]
\end{corollary}
\begin{proof}
The idempotent is constant on each indicated set by the
definition of \(J_{p,r}\). Prime-to-\(p\) character descent
identifies the contribution on \(E_x\) with the selected
scalar groups: averaging is exact globally and locally, and
the inertia quotients away from \(p\) have order prime to
\(p\). All these primes satisfy \(p\le2R+1<P(x)\) for large
\(V\), so their truncation is at \(b\). Capping the
elementary divisors at \(m_p\) gives a truncated total
\(H_{\rm cap}=\sum_{p,\mathcal R}(a_{m_p}-a_b)\log p\),
with \(0\le H_{\rm cap}\le H(x)\). Its difference from
\(Q(x,X)\) is bounded by \(\mathcal D(x)\) by the theorem.
For each set \(D(\mathcal R)\le R^2(1+\log_p R)\).
Corollary~\ref{cor:periodfieldweights644} bounds
\(\log_p|eC(0)|\) by \(O(n_*(1+\log N_*/\log p))\).
There are at most \(2(2R+1)\) sets, and the interval
parameters here are polynomial in \(V\). This proves the
uniform error bound.

If any elementary divisor in the sum were longer than
\(m_p\), its contribution alone to \(H_{\rm cap}\) would
be at least \((m_p-b)\log p>X\). Thus \(H_{\rm cap}\le X\)
forces \(H_{\rm cap}=H(x)\). This proves the first
implication, and the second follows from \(H_{\rm cap}\le H(x)\).
\end{proof}

The action of \(\gamma\) is essential to this test. For
example, as a calculation with finite modules, fix a weight
\(r_0\) prime to \(p\) and an integer \(M\ge2\). On
\(\ZZ/p^{j+1}\ZZ\) let \(\gamma\) act by \(u_{r_0}+p^M\).
Its difference from one has valuation one, so this action
factors through order \(p^j\); reduction gives the maps
between the layers. The orders \(p^{j+1}\) and the
coinvariant orders \(p\) are independent of \(M\), whereas
the length at \(u_{r_0}\) is \(\min(M,j+1)\). The arithmetic
quotients in the corollary retain the actual action of the
period fields. Their required size on arbitrarily distant
intervals has yet to be bounded.


\subsection{All weights at a good prime in one ordinary class group}
\label{sec:allgoodmeasure648}

We include the zero residue at the completely split primes
and retain the original averaging operator at each weight.
Write \(E=E_*\), \(N=N_*\), \(R=t_*\), and take
\(5\le p\le2R+1\), \(p\nmid n_*N\), on a common interval.
Put
\[
 F=(E\QQ(\mu_p))^+,\quad F_j=F\QQ_j,\quad
 C_j=\Cl(\mathcal O_{F_j}[1/p])\{p\},\quad
 \Delta=\operatorname{Gal}(F/\QQ).
\]
Thus \(C_j\) is an ordinary ideal class group modulo the
classes of the primes above \(p\). Its degree-zero logarithmic
analogue will not be needed. The cyclotomic layers \(\QQ_j\)
and generator \(\gamma\) are normalized as in
Section~\ref{sec:productweights647}. The degree of \(F\) is
prime to \(p\). In particular \(F\) is disjoint from
\(\QQ_\infty\), and locally at \(p\) its intersection with
the totally ramified cyclotomic \(\ZZ_p\)-extension is trivial.

Given a character \(\alpha\) of \(E\) and \(0\le s<p-1\)
with \(\alpha(c)=(-1)^s\), the finite character
\(\tau=\alpha^{-1}\omega^{-s}\) is an even character of
\(E\QQ(\mu_p)\), where \(\omega\) is the Teichm\"uller
cyclotomic character. Every character of \(F\) occurs once
this way: \(E\cap\QQ(\mu_p)=\QQ\) by unramifiedness and
total ramification at \(p\). Frobenius conjugation of the
coefficients preserves \(s\). Write \(e_\tau\) for the
integral orbit idempotent in \(\ZZ_p[\Delta]\).

\begin{lemma}[Ordinary class groups at the split places]
\leavevmode\label{lem:allordinary648}
Suppose \(\tau\ne1\), \(m\ge1\) and \(r>0\) is congruent
to \(s\bmod(p-1)\). Put \(u_r=(1+p)^{-r}\). For \(j\ge m\),
the finite \(O_\tau\)-dual of \(S(\alpha,r)[p^m]\) is
\[
              e_\tau C_j/(p^m,\gamma-u_r)e_\tau C_j.
\]
This holds also when \(\tau\) is trivial on the entire
decomposition group at \(p\).
\end{lemma}
\begin{proof}
Write \(D=(\operatorname{Frac}(O_\tau)/O_\tau)
(\tau\langle\kappa_p\rangle^{-r})\) and \(d_r=1+v_p(r)\).
The finite quotient over \(\QQ\) is
\(\Delta\times\langle\gamma\rangle\), with the second
factor of order \(p^j\). Over \(F_j\) the coefficient on
\(D[p^m]\) is trivial because \(j+d_r\ge m\).
The global invariants of \(D\) vanish: a nonidentity value
of the prime-to-\(p\) character is different from one modulo
\(p\). Hence \(H^1(\QQ,D)[p^m]=H^1(\QQ,D[p^m])\).
Averaging over \(\Delta\) gives zero invariants and zero
higher cohomology on the finite coefficient, so restriction
identifies this first cohomology with equivariant
homomorphisms on \(G_{F_j}\).
Away from \(p\) the inertia quotient has order prime to \(p\)
and inertia over \(F_j\) acts trivially even on \(D\).
Consequently unramifiedness of the divisible class is
equivalent to vanishing of the homomorphism on that inertia.

At \(p\) the decomposition quotient is
\(\Delta_w\times\langle\gamma\rangle\). Its projections are
onto and their orders are coprime. If \(\tau\) is nontrivial
on \(\Delta_w\), averaging there proves local injectivity,
and local invariants of \(D\) are zero. Suppose instead that
\(\tau|_{\Delta_w}=1\). The principal cyclotomic image on
\(G_{\QQ_p}\) is all of \(1+p\ZZ_p\). Its image on
\(G_{F_{j,w}}\) is the subgroup of \(p^j\)-th powers:
the local degree of \(F\) is prime to \(p\).
A finite cocycle whose divisible local class is zero has
the form \(g\mapsto(\theta(g)-1)y\) locally, with
\(\theta=\tau\langle\kappa_p\rangle^{-r}\).
Then \(p^m y\) lies in \(D^{G_{\QQ_p}}=D[p^{d_r}]\).
On \(G_{F_{j,w}}\) the difference \(\theta(g)-1\) is
divisible by \(p^{j+d_r}\). The assumption \(j\ge m\)
therefore makes the finite cocycle itself zero there.

Conversely a local cocycle that vanishes on this upper
group has divisible class inflated from the local quotient
on \(D^{G_{F_{j,w}}}=D[p^{j+d_r}]\). Average out
\(\Delta_w\), which is acting trivially in this case.
For the remaining cyclic group, first cohomology is the
kernel of \(1+u_r+\cdots+u_r^{p^j-1}\) modulo the image
of \(u_r-1\). Their valuations are \(j\) and \(d_r\);
the first kernel and the second image on \(D[p^{j+d_r}]\)
are both \(p^{d_r}D[p^{j+d_r}]\). Thus the divisible local
class is zero. The real places impose no odd-primary
condition. Class field theory now gives \(C_j\), with
equivariance imposing \(\tau\) and \(\gamma=u_r\).
Taking finite \(O_\tau\)-duals proves the formula.
\end{proof}

Fix \(x\). For each nontrivial orbit \(\tau\), let
\(\mathcal R_\tau\) consist of the weights \(1\le r<t(x)\)
in its residue for which \(\alpha\) is inflated from \(E_x\)
and survives \(1-J_{p,r}\). Put \(\mathcal R_1=\varnothing\):
the only trivial pair is \(\alpha=1,s=0\), whose scalar group
is zero by Lemma~\ref{lem:primitiveclass502}.
Define one polynomial over the integral group algebra and
one finite group at the chosen precision:
\[
 \begin{aligned}
   P(Y)&=\sum_\tau e_\tau\prod_{r\in\mathcal R_\tau}(Y-u_r),\\
   Z_m&=C_m/(p^m,P(\gamma))C_m .
 \end{aligned}
\]
Empty products are one. Let \(Y_m\) be the direct sum of
the groups on the right of Lemma~\ref{lem:allordinary648}
at \(j=m\) over all the indicated orbits and weights. Set
\[
 a_p=\max_{\tau,\ r\in\mathcal R_\tau}
           \sum_{\substack{r'\in\mathcal R_\tau\\r'\ne r}}
                                (1+v_p(r-r')),
\]
taking zero when there are no weights.

\begin{theorem}[Two-sided bounds for all the surviving elementary divisors]
\leavevmode\label{thm:alltails648}
There are homomorphisms between \(Z_m\) and \(Y_m\) in both
directions whose two composites are multiplication by
\(p^{a_p}\). Thus for every integer \(b\ge a_p\),
\[
        |p^{b+a_p}Z_m|\ \le\ |p^bY_m|\
                                   \le\ |p^{b-a_p}Z_m|.
\]
If \(g_p\) is the number of primes of \(F\) above \(p\),
the base-\(p\) logarithms of the two outer orders differ by
at most
\[
                 2a_pR\,\big(v_p(h_F)+g_p-1\big).
\]
This error is uniform in the cyclotomic level \(m\), including
at primes that divide the indices of global units.
\end{theorem}
\begin{proof}
Work on one nonempty orbit factor \(A=e_\tau C_m/p^m e_\tau C_m\).
List its distinct roots \(u_1,\ldots,u_n\), let
\(P_0(Y)=\prod_i(Y-u_i)\), and let \(P_i(Y)=P_0(Y)/(Y-u_i)\).
The diagonal quotient map from \(A/P_0(\gamma)A\) goes to
\(\bigoplus_i A/(\gamma-u_i)A\). For the reverse map, lift
the entries \(x_i\) and send them to the class of
\[
       \sum_i \frac{p^{a_p}}{P_i(u_i)}P_i(\gamma)x_i .
\]
All the coefficients lie in \(\ZZ_p\) by the difference
valuation in Theorem~\ref{thm:allhalfweights644}; that
valuation holds for any two distinct positive integers.
Changing a lift changes the sum by a multiple of
\(P_0(\gamma)\). Evaluation in the \(i\)-th quotient gives
\(p^{a_p}x_i\), and the interpolation identity
\(\sum_i p^{a_p}P_i(Y)/P_i(u_i)=p^{a_p}\) gives the other
composite. Summing the maps over the orbit factors proves
the assertion for \(Y_m,Z_m\). Factoring multiplication by
\(p^{b+a_p}\) through \(p^bY_m\), and multiplication by
\(p^b\) on \(Y_m\) through \(p^{b-a_p}Z_m\), proves the
two inequalities.

It suffices to bound the number of abelian \(p\)-group
generators of \(Z_m\), since lowering a power of \(p\) by
one increases the logarithmic order by at most that number.
Modulo \(p\), on each factor the polynomial is
\((Y-1)^{|\mathcal R_\tau|}\).
Consequently \(Z_m/pZ_m\) is a quotient of
\(C_m/(p,(\gamma-1)^R)C_m\). The filtration by powers
of \(\gamma-1\) has at most \(R\) quotients, each an image
of \(C_m/(p,\gamma-1)C_m\).

For completeness the latter dimension is bounded at the
bottom of the tower even when its full coinvariant length
is large. All the primes over \(p\) in \(F\) are totally
ramified in \(F_m/F\). In the abelian \(p\)-class field of
\(F_m\) split above \(p\), a decomposition group at one
such prime lifts \(\langle\gamma\rangle\) isomorphically.
The Galois group over \(F\) is thus the semidirect product
of \(C_m\) and this lift. Its abelianization replaces \(C_m\)
by \(C_m/(\gamma-1)C_m\). To make every prime over \(p\)
split, kill the chosen lift and the differences with the
other \(g_p-1\) decomposition generators. This gives
exactly the split class group of \(F\): its class field
base-changes into the corresponding split class field of
\(F_m\), and the constructed quotient has the required
ramification and splitting. Hence at most \(g_p-1\) more
generators are needed than for \(\Cl(\mathcal O_F[1/p])\{p\}\),
which itself needs at most \(v_p(h_F)\). This proves the
stated width without an upper bound for any unit index.
\end{proof}

\begin{corollary}[A simultaneous ordinary class group test for the good primes]
\leavevmode\label{cor:goodtest648}
Let \(G(x)\) be the part of \(\mathcal R(x)\) in
Proposition~\ref{prop:res23} from all primes \(p\ge5\)
prime to \(n_*N_*\), including all weights. For \(X>0\)
put \(b=b_0(x)\) and \(m_p=b+\lfloor X/\log p\rfloor+1\).
Over the good primes through \(2R+1\) form the two quantities
\[
 T_{\rm low}=\sum_p\log|p^{b+a_p}Z_{m_p}|,\qquad
 T_{\rm up}=\sum_p\log|p^{b-a_p}Z_{m_p}|.
\]
For all sufficiently large \(V\) these are defined, and
\(0\le T_{\rm up}-T_{\rm low}\ll_{\psi,\varepsilon}V^A\)
for some fixed \(A\), independently of \(X\).
On the same fields and intervals one has
\[
 \begin{split}
 T_{\rm up}+\log J(V)\le X &\ \Longrightarrow\ G(x)\le X,\\
 G(x)\le X &\ \Longrightarrow\ T_{\rm low}\le X ,
 \end{split}
\]
where \(J(V)\) is the integer of Lemma~\ref{lem:fullinteger517},
with \(\log J(V)\ll_{\psi,\varepsilon}V^C\).
\end{corollary}
\begin{proof}
For \(p\le2R+1\) we have \(p<P(x)\) eventually, and
\(a_p\le R(1+\log_p R)<b\). The theorem and prime-to-\(p\)
descent as in Corollary~\ref{cor:halftail647} bracket the
actual capped total on these primes between the two sums.
The corresponding \(F\) has degree at most \(n_*(p-1)\)
and conductor dividing \(pN_*\). Hence both \(g_p\) and
\(v_p(h_F)\) are polynomially bounded in \(V\), by the
Minkowski estimate of Corollary~\ref{cor:periodfieldweights644}.
Summing the displayed width of the theorem with the factors
\(\log p\) gives the asserted polynomial width.

Every prime beyond \(2R+1\) satisfies \(p-1\nmid2r\) for
every weight in question. For sufficiently large \(V\) it
also lies outside the fixed finite exceptional set of
Lemma~\ref{lem:fullinteger517}. That lemma bounds the total
logarithm of the full orders there by \(\log J(V)\).
If the capped total is at most \(X\), an elementary divisor
longer than \(m_p\) at a prime in the test would contribute
at least \((m_p-b)\log p>X\) to that total. Thus the cap
has lost nothing at those primes. Combining these facts
gives the first implication. For the second, \(T_{\rm low}\)
is at most the capped total, itself at most \(G(x)\).
\end{proof}

The two numbers now measure the contribution at all good
primes within polynomial uncertainty, using the original
weight-dependent field averages. Their absolute size still
depends on the action on ideal classes in the common
cyclotomic layer. An upper bound for that size on arbitrarily
distant intervals, the other coefficient primes, and the
integral signed period comparison remain necessary.


\subsection{Prime ideal generators at exact cyclotomic depth}
\label{sec:simultaneousdepth652}

The splitting restriction in Corollary~\ref{cor:detectingsplit650}
has a converse for the scalar groups whose residual character
is nontrivial at the coefficient prime. We construct rational
primes at the least possible common depth for an entire list.
Several weights on the same character and several coefficient
primes may occur together.

Fix one interval with largest field \(E_*\), degree \(n_*\)
and modulus \(N_*\). Use the scalar characters and weights
of Section~\ref{sec:allgoodmeasure648}. An entry consists of
a good prime \(p\ge5\), \(p\nmid n_*N_*\), a positive weight
\(r<t(x)\), and a character \(\alpha\) of \(E_x\) with
\(\alpha(c)=(-1)^r\). Inflate to \(E_*\), and let
\(\tau=\alpha^{-1}\omega_p^{-r}\) be nontrivial, with
\(\omega_p\) the Teichm\"uller character. One can restrict
the list to the characters selected by the original field
averages. Write \(F=(E_*\QQ(\mu_p))^+\),
\(F_j=F\QQ_j\), where \([\QQ_j:\QQ]=p^j\) in the cyclotomic
\(\ZZ_p\)-extension. Put \(d_r=1+v_p(r)\) and \(u_r=(1+p)^{-r}\).
For an integer \(m\ge1\) let \(Q\) be the finite dual of
\(S(\alpha,r)[p^m]\) over its unramified coefficient ring.
For an abelian \(p\)-group we use the additive convention.

\begin{lemma}[A class field at the precision of the twist]
\leavevmode\label{lem:twistdepth652}
Choose the layer index
\[
 j=
 \begin{cases}
  \max(0,m-d_r),&\tau\ \text{nontrivial on a decomposition group at }p,\\
  m,&\tau\ \text{trivial there}.
 \end{cases}
\]
There is a finite Galois extension \(M/\QQ\), abelian over
\(F_j\), with \(\operatorname{Gal}(M/F_j)=Q\). The class
map on \(F_j\) is the quotient
\[
 e_\tau C_j/(p^m,\gamma-u_r)e_\tau C_j,\qquad
                C_j=\Cl(\mathcal O_{F_j}[1/p])\{p\}.
\]
The largest subfield of \(M\) abelian over \(\QQ\) is \(F_j\).
\end{lemma}
\begin{proof}
Lemma~\ref{lem:allordinary648} gives this class quotient for
\(j=m\). We justify the smaller layer when \(\tau\) is locally
nontrivial. Over \(F_j\) the principal cyclotomic character
raised to \(-r\) is trivial modulo \(p^m\), since
\(v_p(u_r^{p^j}-1)=j+d_r\). The remaining finite character
already factors through \(F\). Globally and at \(p\) the
finite quotient is the product of its prime-to-\(p\) part
and its cyclic \(p\)-part. At \(p\) this follows also from
the fact that a subgroup of a finite abelian group contains
the prime-to-\(p\) and \(p\)-parts of each of its elements.
In each quotient a prime-to-\(p\) element acts on the finite
coefficient by a root of unity differing from one by a unit.
Averaging over its cyclic subgroup kills the invariants and
all higher cohomology. Inflation--restriction now gives
equivariant homomorphisms over \(F_j\) and detects zero
localization at \(p\). The invariants of the divisible
coefficient vanish in the same two places, so the exact
sequence for multiplication by \(p^m\) identifies the
divisible and finite conditions through that precision.
Away from \(p\) the inertia quotient in \(F_j/\QQ\) has
order prime to \(p\). On the upper inertia the finite and
divisible coefficients have trivial action, since the
cyclotomic extension is unramified there and the finite
character is trivial on that inertia. Thus the condition
is unramifiedness, by restriction and averaging. Class
field theory gives the asserted quotient, with the scalar
relation on \(\gamma\) coming from equivariance.

The map from the ordinary class group to \(Q\) has
Galois-stable kernel. Its abelian class field \(M/F_j\)
is consequently Galois over \(\QQ\). Choose an element of
the prime-to-\(p\) factor whose \(\tau\)-value is nontrivial.
It acts on \(Q\) with its difference from the identity a
unit on the whole orbit component. The commutators with
any lift of this element therefore generate \(Q\).
Conversely all commutators in \(\operatorname{Gal}(M/\QQ)\)
belong to \(Q\), since \(F_j/\QQ\) is abelian. Its
commutator subgroup is exactly \(Q\), proving the last assertion.
\end{proof}

\begin{theorem}[Simultaneous generators with prescribed Frobenius depth]
\leavevmode\label{thm:exactclassprimes652}
Let \(i=1,\ldots,s\), \(s\ge1\), be a finite list as above,
with entries \(p_i,r_i,\tau_i,m_i,j_i,Q_i\). Assume that
\(Q_i\) has exponent \(p_i^{m_i}\). This can always be
arranged for a nonzero truncated group by reducing its
precision to its exponent. For each occurring \(p\) define
\[
              N_p=1+\max_{i:p_i=p}j_i .
\]
Let \(d=\max_i\dim_{\mathbf F_{p_i}} Q_i/p_iQ_i\).
There are \(h=d+\lceil\log_2(2s)\rceil\) distinct rational
primes \(q_a\), outside any prescribed finite set, such that
all split in \(E_*\) and
\[
 v_p(q_a-1)=N_p\quad\hbox{for every }a,p .
\]
For each \(i\) one can choose primes over these \(q_a\) in
\(F_{i,j_i}\) whose images generate \(Q_i\) as an abelian
group. The choices over different fields can be taken from
one prolongation per rational prime.

If every entry at \(p\) has \(\tau_i\) nontrivial locally
at \(p\), then
\[
           N_p=\max\bigl(\{1\}\cup
                    \{m_i-v_p(r_i):p_i=p\}\bigr)
\]
is the least possible common depth for generating these
groups by primes split in \(E_*\) and congruent to one
modulo \(p\). The same necessity holds when one requires the
multiplied prime classes only to generate \(p^{b_i}Q_i\),
provided \(0\le b_i<m_i\).
\end{theorem}
\begin{proof}
For the finite set of coefficient primes take
\[
 A=\QQ\bigl(\mu_{N_*\prod_p p^{N_p}}\bigr),\qquad
 C=\QQ\bigl(\mu_{N_*\prod_p p^{N_p+1}}\bigr).
\]
Every \(F_{i,j_i}\) lies in \(A\) by the choice of \(N_p\).
The Chinese remainder theorem prescribes an automorphism
\(\sigma\) of \(C\) which is one modulo \(N_*\) and
\(1+p^{N_p}\) modulo \(p^{N_p+1}\) at every \(p\).
It fixes \(A\). Let \(K=C\prod_i M_i\), using the class
fields of the lemma, and \(G=\operatorname{Gal}(K/\QQ)\).
Since \(C/\QQ\) is abelian,
\(M_i\cap C=F_{i,j_i}\). Consequently restriction maps
\(\operatorname{Gal}(K/C)\) onto each \(Q_i\). This follows
by extending an automorphism first to \(M_i C\) over \(C\)
and then to the Galois extension \(K/C\).

The lifts of \(\sigma\) to \(G\) form a coset of this
subgroup. Each fixes all \(F_{i,j_i}\). A uniform choice
in the coset therefore maps uniformly to each \(Q_i\).
Take \(h\) independent choices. At a given \(i\), failure
to span the reduction modulo \(p_i\) puts all \(h\) entries
in one hyperplane. For dimension \(d_i\) there are
\((p_i^{d_i}-1)/(p_i-1)\) hyperplanes, each of relative
size \(1/p_i\). The failure probability is at most
\(p_i^{d_i-h}\le2^{d-h}\). The sum of these bounds over
the list is at most \(1/2\). Thus a tuple spanning every
reduction exists, including when some weight quotients are
correlated. Spanning modulo \(p_i\) implies generation of
the full finite group, by iteration of multiplication by
\(p_i\) on its quotient by the span.

Apply Chebotarev's theorem in \(K/\QQ\) to each chosen
element. There are infinitely many unramified rational
primes in its conjugacy class, and conjugating the
prolongation realizes the element itself. Choose the
primes successively with the stated exclusions.
Their actions on \(C\) give the exact valuations and
complete splitting in \(A\). Thus their residue degrees
over each \(F_{i,j_i}\) above \(\QQ\) are one. Their
restrictions to \(M_i\) are precisely the relative
arithmetic Frobenius elements, hence the ideal class
images under reciprocity. These are the generating
elements just chosen.~\cite{XI-VIII-I-tz50}

For necessity, take any candidate prime \(q\) unramified
in the fields in question, split in \(E_*\), and congruent
to one modulo \(p\). Its Frobenius on \(F\) is trivial.
In the cyclotomic direction it acts on \(Q_i\) by \(q^{-r_i}\).
Indeed \(q\) is a principal \(p\)-adic unit, and the
relation on the chosen generator is
\(\gamma=(1+p)^{-r_i}\). Frobenius fixes the class of
each prime above \(q\) in \(F_{i,j_i}\). Thus every image
is killed by \(q^{r_i}-1\), also when the prime does not
split completely at level \(j_i\). Its order has
base-\(p\) exponent at most \(v_p(q-1)+v_p(r_i)\), by
binomial expansion. To generate a group of exponent
\(p^{m_i}\) at least one generator must have that order.
To generate \(p^{b_i}Q_i\) with \(b_i<m_i\), at least one
multiplied class must have order \(p^{m_i-b_i}\), with
the same consequence for the original class. Taking the
maximum gives the lower bound for any common depth. In
the locally nontrivial case our chosen \(1+j_i\) is
exactly \(\max(1,m_i-v_p(r_i))\), proving attainment.
\end{proof}

This also gives the cost of a simultaneous choice of such
primes for the independent symplectic inertia characters of
Theorem~\ref{thm:pairedtrace651}. For a fixed \(p\), the \(h\)
primes just constructed have local separating ideal
\(p^{hN_p}\mathcal O\). Any set of probes of depth at most
\(D\) at \(p\) misses generators for a component with
\(m_i>D+v_p(r_i)\). This statement concerns the prime ideal
images in the indicated class groups. The construction lets
the probes depend on the actual exponents. An upper bound for
those exponents, or for their weighted sum on the interval,
is still required in Section~\ref{sec:missing}.

\clearpage
\section{References}\label{nav:bibliography-XI}
\begingroup\def\section*#1{}\begin{thebibliography}{99}\small
\bibitem{XI-VIII-liu375} T. Liu, \emph{Torsion \(p\)-adic Galois representations and a conjecture of Fontaine}, Ann. Sci. \'Ecole Norm. Sup. (4) \textbf{40} (2007), 633--674, Sections 1 and 7.3. \href{https://www.math.purdue.edu/~tongliu/pub/fontainel.pdf}{Original paper}.
\bibitem{XI-VIII-VI-kisin121} M. Kisin, \href{https://people.math.harvard.edu/~kisin/dvifiles/licm.pdf}{\emph{The structure of potentially semi-stable deformation rings}}, Theorem~1.2.2. See also \emph{Potentially semi-stable deformation rings}, J. Amer. Math. Soc. \textbf{21} (2008), 513--546.
\bibitem{XI-VIII-VI-kato120} K. Kato, \href{https://archive.numdam.org/item/AST_2004__295__117_0.pdf}{\emph{\(p\)-adic Hodge theory and values of zeta functions of modular forms}}, Ast\'erisque \textbf{295} (2004), 117--290, Theorem~14.2(1) and Section~14.10.
\bibitem{XI-VIII-VI-nt119} J. Newton and J. A. Thorne, \href{https://arxiv.org/pdf/1912.11265}{\emph{Adjoint Selmer groups of automorphic Galois representations of unitary type}}, J. Eur. Math. Soc. \textbf{25} (2023), 1919--1967, Theorem~A and Section~2.5.
\bibitem{XI-VIII-geraghty363} D. Geraghty, \emph{Modularity lifting theorems for ordinary Galois representations}, Math. Ann. \textbf{373} (2019), 1341--1427, Section 2. \href{https://doi.org/10.1007/s00208-018-1742-4}{Published paper}. Proposition and lemma numbers refer to the \href{https://citeseerx.ist.psu.edu/document?doi=8faa2d675c94ae78d0ac8ebb3f9d4f69ede42d44\&repid=rep1\&type=pdf}{original manuscript}.
\bibitem{XI-VIII-VI-bk117} S. Bloch and K. Kato, \href{https://virtualmath1.stanford.edu/~conrad/BSDseminar/refs/BKTamagawa.pdf}{\emph{L-functions and Tamagawa numbers of motives}}, in \emph{The Grothendieck Festschrift}, Vol.~I, Progr. Math. \textbf{86} (1990), 333--400, Section~3.
\bibitem{XI-VIII-VI-pan40} L. Pan, \emph{The Fontaine--Mazur conjecture in the residually reducible case}, J. Amer. Math. Soc. \textbf{35} (2022), 1031--1169. \href{https://arxiv.org/abs/1901.07166}{Author's paper}.
\bibitem{XI-VIII-I-ac123} J. Arthur and L. Clozel, \href{https://www.claymath.org/library/cw/arthur/pdf/30.pdf}{\emph{Simple Algebras, Base Change, and the Advanced Theory of the Trace Formula}}, Annals of Mathematics Studies 120, Princeton University Press, 1989, Chapter III.
\bibitem{XI-VIII-lgp361} T. Barnet-Lamb, T. Gee, D. Geraghty and R. Taylor, \emph{Local-global compatibility for \(l=p\), II}, Ann. Sci. \'Ec. Norm. Sup\'er. (4) \textbf{47} (2014), 165--179, Theorem A. \href{https://www.numdam.org/articles/10.24033/asens.2212/}{Published paper}.
\bibitem{XI-VIII-V-shin} S. W. Shin, \emph{Galois representations arising from some compact Shimura varieties}, Ann. of Math. \textbf{173} (2011), 1645--1741. \href{https://annals.math.princeton.edu/2011/173-3/p09}{Theorem 1.2}.
\bibitem{XI-VIII-VI-saito40} T. Saito, \emph{Modular forms and \(p\)-adic Hodge theory}, Invent. Math. \textbf{129} (1997), 607--620. See also \href{https://www.ms.u-tokyo.ac.jp/~t-saito/pp/h.pdf}{\emph{Hilbert modular forms and \(p\)-adic Hodge theory}}, Compos. Math. \textbf{145} (2009), 1081--1113.
\bibitem{XI-VIII-chenevier364} G. Chenevier, \emph{The \(p\)-adic analytic space of pseudocharacters of a profinite group and pseudorepresentations over arbitrary rings}, in \emph{Automorphic Forms and Galois Representations}, Vol.~1, Cambridge University Press, 2014, 221--285, Introduction and Theorem A. \href{https://gaetan.chenevier.perso.math.cnrs.fr/articles/determinants.pdf}{Original paper}.
\bibitem{XI-VIII-katz361} N. Katz, \emph{\(p\)-adic properties of modular schemes and modular forms}, in \emph{Modular Functions of One Variable III}, Lecture Notes in Mathematics 350, Springer, 1973. \href{https://web.math.princeton.edu/~nmk/old/padicpropMFMS.pdf}{Original paper}.
\bibitem{XI-VIII-billerey361} N. Billerey and R. Menares, \emph{On the modularity of reducible mod \(\ell\) Galois representations}, Math. Res. Lett. \textbf{23} (2016), 15--41, Section 1. \href{https://lmbp.uca.fr/~billerey/Research/Published/Menares_Modularity.pdf}{Published paper}.
\bibitem{XI-VIII-rapoport361} P. Deligne and M. Rapoport, \emph{Les sch\'emas de modules de courbes elliptiques}, in \emph{Modular Functions of One Variable II}, Lecture Notes in Mathematics 349, Springer, 1973, 143--316. \href{https://publications.ias.edu/book/export/html/367}{Original paper}.
\bibitem{XI-VIII-V-newforms} F. Diamond and J. Im, \emph{Modular forms and modular curves}, in \emph{Seminar on Fermat's Last Theorem}, CMS Conf. Proc. \textbf{17} (1995), 39--133. \href{https://wstein.org/edu/Fall2003/252/references/diamond-im/Diamond-Im-Modular_forms_and_modular_curves.pdf}{Authors' text}.
\bibitem{XI-VIII-V-deligne} P. Deligne, \emph{Formes modulaires et repr\'esentations \(\ell\)-adiques}, S\'eminaire Bourbaki, Exp. 355 (1969). \href{https://www.numdam.org/item/SB_1968-1969__11__139_0.pdf}{Paper}.
\bibitem{XI-VIII-I-compat123} H. Carayol, \href{https://www.numdam.org/articles/10.24033/asens.1512/}{\emph{Sur les repr\'esentations \(\ell\)-adiques associ\'ees aux formes modulaires de Hilbert}}, Ann. Sci. \'Ecole Norm. Sup. (4) \textbf{19} (1986), 409--468. T. Saito, \emph{Modular forms and \(p\)-adic Hodge theory}, Invent. Math. \textbf{129} (1997), 607--620.
\bibitem{XI-VIII-ordinary361} C. M. Skinner and A. J. Wiles, \emph{Ordinary representations and modular forms}, Proc. Natl. Acad. Sci. USA \textbf{94} (1997), 10520--10527. \href{https://doi.org/10.1073/pnas.94.20.10520}{Published paper}.
\bibitem{XI-VIII-I-gj123} S. Gelbart and H. Jacquet, \href{https://www.numdam.org/item/ASENS_1978_4_11_4_471_0/}{\emph{A relation between automorphic representations of \(\mathrm{GL}(2)\) and \(\mathrm{GL}(3)\)}}, Ann. Sci. \'Ecole Norm. Sup. (4) \textbf{11} (1978), 471--542.
\bibitem{XI-VIII-finiteunr443} P. B. Allen and F. Calegari, \emph{Finiteness of unramified deformation rings}, Algebra Number Theory \textbf{8} (2014), 2263--2272, Theorem~1. \href{https://arxiv.org/abs/1411.3759}{Original paper}.
\bibitem{XI-VIII-rayclass444} J. S. Milne, \emph{Class Field Theory}, version 4.03 (2020), Chapter V, Sections 1 and 3, and the reciprocity law. \href{https://www.jmilne.org/math/CourseNotes/CFT.pdf}{Author's notes}.
\bibitem{XI-VIII-zhang445} X. Zhang, \emph{Modular points and dimensions of Eisenstein deformation spaces}, arXiv:2512.21249v4 (2026), Appendix~A, Proposition~A.1.4 and Theorem~A.3.1. \href{https://arxiv.org/pdf/2512.21249v4}{Author's preprint, version~4}.
\bibitem{XI-VIII-I-tz50} J. Thorner and A. Zaman, \emph{A unified and improved Chebotarev density theorem}, Algebra Number Theory \textbf{13} (2019). \href{https://msp.org/ant/2019/13-5/ant-v13-n5-p02-p.pdf}{Theorem 1.4 and Theorem 3.1}.
\bibitem{XI-VIII-ambiguous505} F. Lemmermeyer, \emph{The ambiguous class number formula revisited} (2013), Theorem~1 and Proposition~1. \href{https://arxiv.org/abs/1309.1071}{Author's paper}.
\bibitem{XI-VIII-brumer330} A. Brumer, \emph{On the units of algebraic number fields}, Mathematika \textbf{14} (1967), 121--124. \href{https://doi.org/10.1112/S0025579300003703}{Original paper}.
\bibitem{XI-VIII-weberFK524} T. Fukuda and K. Komatsu, \emph{Weber's class number problem in the cyclotomic $\mathbb Z_2$-extension of $\mathbb Q$, II}, J. Th\'eor. Nombres Bordeaux \textbf{22} (2010), 359--368, Corollary 1.1. \href{https://www.numdam.org/article/JTNB_2010__22_2_359_0.pdf}{Published paper}.
\bibitem{XI-VIII-grasdyadic525} G. Gras, \emph{Weber's class number problem and $p$-rationality in the cyclotomic $\widehat{\mathbb Z}$-extension of $\mathbb Q$}, arXiv:2009.05278, v3 (2020), \S 7. \href{https://arxiv.org/pdf/2009.05278}{Author's paper}.
\bibitem{XI-VIII-sutherland529} A. V. Sutherland, \emph{The analytic class number formula}, MIT 18.785, Lecture 19 (2018), Theorem~19.12. \href{https://math.mit.edu/classes/18.785/2018fa/LectureNotes19.pdf}{Lecture notes}.
\end{thebibliography}\endgroup
\end{document}
